% ===================================================================
%  Premiere partie — MORFOLOGIO E SINTAXO
%  Feuillet 15 = folio 11 (le folio n'est pas imprime sur cette page).
%
%  Mesures (outils/page.py, feuillet 15) :
%    titre ALFABETO.      y = 459 px ; cap 2,540 mm ; largeur 18,54 mm
%                         -> 10,63 pt, interlettrage -19 (gras)
%    titre PRONUNCO...    cap 2,455 mm ; largeur 41,99 mm
%                         -> 10,27 pt, interlettrage +8 (gras)
%    1re ligne de corps   y = 548 px = 46,40 mm
%    titre PRONUNCO...    y = 1005 px = 85,09 mm
%    filet de note        y = 1550 px = 131,2 mm
% ===================================================================

% --- Feuillet 15 (folio 11) ----------------------------------------
\begin{VUpage}[15]{}
% Les notes sont placees a leur ordonnee mesuree ; elles occupent une
% hauteur nulle et se declarent donc en tete de page (cf. preambule).
\VUnotes{131.37mm}{%
(1) Por obtenar plu granda fixeso en la pronunco di ta du vokali,\nl
on darfas sequar la yena indiko, qua esas ne regulo ma konsilo\nl
utila por preventar hezito :

\VUblancAlinea
« On \textit{klozas} la vokali \VUgras{e}, \VUgras{o} en silabo \textit{apertita} (finanta per vokalo).\nl
Ex. : \textit{mé, bôné, dômô.}

\VUblancAlinea
On \textit{apertas} la vokali \VUgras{e}, \VUgras{o} en silabo \textit{klozita} (finanta per konso\cc
nanto). Ex. : \textit{sèmpré, fénèstrô, pòrdô, òmna, mèntô.} »

Tale la pronunco esas plu fixa ed eufonioza. Ma co esas nur\nl
konsilo e nule regulo obliganta. Nur importas, ma \textit{tre importas}, la\nl
klara pronunco di ta vokali e lia rispektiva dicernebleso de la ceteri\nl
(Segun \textit{Progreso}, VII, 400).%
}

\VUsaut{14.7mm}
\VUtitre{10.50pt}{-81}{ALFABETO.}
\VUsaut{3.0mm}

\VUblancAlinea
1. --- L'\textit{alfabeto} di la linguo kontenas 26 literi : kin vokali,\nl
duadek-e-un konsonanti.

\VUblancAlinea
La kin vokali esas : \VUgras{a, e, i, o, u.} La duadek-e-un konso\cc
nanti esas : \VUgras{b, c, d, f, g, h, j, k, l, m, n, p, q, r, s, t, v,}\nl
\VUgras{w, x, y, z.} Esas tre rekomendata desegnar la literi en maniero\nl
qua impedas konfundar, inter la mayuskuli, F a T o a C, e\nl
I a J. En l'imprimo on atencez, ke Ido ne divenez Jdo. Gene\cc
rale, e precipue pri la adresi, on sorgez formacar la literi\nl
maxim bone e lekteble.

\VUblancAlinea
% Blancs mesures au feuillet 15 : 114 px au-dessus du titre, 89 px
% au-dessous (le compose donnait 92 et 85). Comme au feuillet 28, les
% blancs qui encadrent un titre se mesurent, ils ne se deduisent pas.
\VUsaut{5.26mm}
\VUtitre{10.15pt}{-24}{PRONUNCO DIL VOKALI.}
\VUsaut{3.54mm}

\VUblancAlinea
2. --- En Ido, la vokali havas meza sonoduro, quale en\nl
l'Italiana. On pronuncas li :

\VUblancAlinea
\VUgras{a} ne tro apertita, nam ol havus graveso desagreabla; ne\nl
tro klozita, nam ol esus ne sat bone dicernebla. On atencez\nl
precipue, ke la dezinenco \VUgras{a} dil adjektivi esez nultempe pro\cc
nuncata quale \textit{ä} Germana, \textit{è} o \textit{é} Franca. Ol sonez \VUgras{a} tre pure,\nl
quale en l'Italiana.

\VUblancAlinea
\VUgras{e}, \VUgras{o} apertita o klozita, se nur on audigas li en maniero\nl
donanta dicernebleso senduba. --- Notez bone, ke la vokalo \VUgras{e}\nl
nultempe esas muta, quale en Franca linguo (1).

\end{VUpage}

% --- Feuillet 16 (folio 12) ----------------------------------------
%  Mesures : titre PRONUNCO DIL KONSONANTI E DIGRAMI. a y = 1305 px,
%  cap 2,455 mm, largeur 72,81 mm -> 10,27 pt, interlettrage +9 (gras).
%  Filet de note a y ~ 1830 px = 154,9 mm.
\begin{VUpage}[16]{12}
\VUnotes{156.38mm}{%
(1) Videz plu fore, ye w, « Pronunco dil konsonanti e digrami ».%
}

\VUgras{u} sempre devas esar pronuncata quale \textit{u} Italiana o Ger\cc
mana; nultempe quale \textit{u} Franca o \textit{ü} Germana. Ol sonas\nl
quale la Franca grupo \textit{ou}.

%% aucun blanc au fac-simile
Kande \VUgras{u} sequas nemediate \VUgras{a} o \VUgras{e}, ol formacas kun ici dif\cc
tongo, qua esas pronuncata \textit{unsilabe}. La sono \VUgras{u} devas au\cc
desar quik pos \VUgras{a}, \VUgras{e} quaze sen intertempo, ma \VUgras{a} e recevas\nl
la chefa esforco di la voco : \VUgras{Australia, Europa, laute, neutra.}

%% aucun blanc au fac-simile
Ma, se la renkontro di \VUgras{a}, \VUgras{e} kun \VUgras{u} rezultas de prefixo\nl
o sufixo adjuntita a radiko, quale en \VUgras{neutila} (\textit{ne-utila}),\nl
\VUgras{kreuro} (\textit{kre-uro}), lore \VUgras{au}, \VUgras{eu} ne plus facas diftongo, e singla\nl
vokalo apartenas a silabo partikulara. Konseque singla esas\nl
aparte pronuncata, quale indikesas per (\textit{ne-utila} e \textit{kre-uro})\nl
supere.

\VUblancAlinea
Sequanta \VUgras{q} o \VUgras{g}, la litero \VUgras{u} = \VUgras{w} (1) avan vokalo. Ol do\nl
sonas quale \VUgras{u} en la duesma silabo di \textit{aquatic} Angla,\nl
\textit{aquatique} Franca, \textit{acquatico} Italiana, \textit{acuatico} Hispana.\nl
Exemple : \VUgras{qua, quar, quo, qui, guidar, linguo, lingue, lin}\cc
\VUgras{guala} = \textit{qwa, qwar, qwo, qwi, gwidar, lingwo, lingwe,}\nl
\textit{lingwala.}

%% aucun blanc au fac-simile
Du sama vokali sucedanta (\textit{ii, ee, oo}) devas ne ligesar\nl
konfuze, ma pronuncesar separite : \VUgras{alopatiisto, antee, heroo.}\nl
To esas explicite fixigita dal decidi akademiala 815 e 816\nl
(\textit{Progreso}, V, 723).

\VUblancAlinea
% Blancs mesures au feuillet 16 : 124 px au-dessus du titre, 90 px
% au-dessous (le compose donnait 78 et 87).
\VUsaut{6.49mm}
\VUtitre{10.15pt}{-13}{PRONUNCO DIL KONSONANTI E DIGRAMI.}
\VUsaut{3.65mm}

\VUblancAlinea
3. --- \VUgras{B} = \textit{b} en l'Italiana e la Franca.

\VUblanc{2.82pt}% mesure au fac-simile
\VUgras{c} = \textit{c} Germana en \textit{Ceres} o \textit{z} Italiana en \textit{zio}, o \textit{ts}. Konseque\nl
\VUgras{ca, ce, ci, co, cu} = \textit{tsa, tse, tsi, tso, tsu.}

\VUblancAlinea
\VUgras{d} = \textit{d} en l'Italiana e la Franca. Ol sempre devas esar\nl
tre klare dicernebla de \textit{t}.

\VUblancAlinea
\VUgras{f} = \textit{f} en l'Italiana, la Franca, l'Angla. Ol sempre devas\nl
esar tre klare dicernebla de \textit{v}.

\VUblancAlinea

\VUblancAlinea
\VUgras{g} = \textit{g} sempre guturala, quale en la Germana \textit{geben}, en la\nl
Angla \textit{give}, od en la Franca \textit{gant;} nultempe quale \textit{g} Franca
\parplein
\end{VUpage}

% --- Feuillet 17 (folio 13) ----------------------------------------
%  Filet de note a y ~ 1325 px = 112,2 mm.
%  Regle d'enrichissement : gras = la lettre idiste en vedette ;
%  italique = la lettre, le son ou le mot cite d'une autre langue.
\begin{VUpage}[17]{13}
\VUnotes{113.23mm}{%
(1) Segun la decido 646 : « On fixigas kom normala sono di\nl
la litero \VUgras{j}, la \textit{j} Franca; nur per tolero (por ti qui ne povas pronun\cc
car ol tale) on admisas la pronunco di la \textit{j} Angla (quan ni repre\cc
zentos per \textit{dj}). »

\VUblancAlinea
Yen la motivi di ca decido : 1\textsuperscript{e} un litero devas havar un sono;\nl
or \textit{tolero} di pronuncado ne violacas ta principo; 2\textsuperscript{e} per ta metodo on\nl
povas transskribar exakte propra nomi, quon on ne povus facar\nl
kun \VUgras{j} havanta du soni, nek kun \VUgras{j} havanta duopla sono (\textit{dj}). --- Ek\nl
\textit{Progreso}, IV, 689.

On remarkez, ke \VUgras{j} esas necese pronuncata quale \textit{j} Franca en plura\nl
Ido-vorti kontenanta \textit{dj : adjuntar, adjutanto, adjudikar, budjeto} ec.\nl
--- Kande \textit{j} sequas \textit{n}, on preske koakte sonigas \textit{d} inter la du konso\cc
nanti (ex. : en \textit{manjar : man-}(d)\textit{-jar}).

(2) Altre pronuncar \VUgras{gn} esus violacar la principo generala dil Ido\cc
pronunco. Pluse, kontre ke omni povas pronuncar ta grupo \textit{g-n},\nl
milioni de homi ne sucesus pronuncar ol unsone, quale la Franca en\nl
\textit{agneau} o l'Italiana en \textit{ogni}.%
}

\noindent en \textit{gens}. Konseque \VUgras{ga, ge, gi, go, gu} = \textit{ga, ghe, ghi, go, gu}\nl
dil Italiana, o \textit{ga, gué, gui, go, gou} di la Franca.

\VUblancAlinea
\VUgras{h} = \textit{h} sempre \textit{vere} aspirata, quale en la Germana e l'Angla\nl
\textit{hand}. Per olu tre diferas en Ido la vorti \textit{horo} e \textit{oro}, \textit{hosto}\nl
e \textit{osto} e. c. La Franci, l'Italiani e ti omna qui ne havas\nl
la \textit{h} aspirata en sua lingui devas tre atencar ta litero. Altre\nl
li ne nur pronuncus ne juste, ma en kazi pasable frequa\nl
li ne komprenesus bone.

\VUblancAlinea
\VUgras{j} = \textit{j} en la Franca, o \textit{s} en la vorto Angla \textit{vision}. Ma, se on\nl
ne povas pronuncar ol tale, on darfas donar a ta konsonanto\nl
la sono di \textit{g} en la Angla \textit{gin}, o di \textit{g} en l'Italiana \textit{giardino} (1).

\VUblancAlinea
\VUgras{k} = \textit{k} en la Germana, Franca, Angla (\textit{keck, képi, keep}, e. c.)\nl
e generale en la lingui uzanta ica konsonanto.

\VUblancAlinea
\VUgras{l} = \textit{l} en la Angla, Franca, Germana, Italiana lingui e. c.

\VUblancAlinea
\VUgras{m}, \VUgras{n} = \textit{m}, \textit{n} en l'Italiana, to esas sempre artikulata, nul\cc
tempe kun nazal sono. Konseque \VUgras{am, an; em, en; im, in;}\nl
\VUgras{om, on; um, un} klare audigas \VUgras{m}, \VUgras{n}, quale se li esus skribita\nl
\textit{amm, ann; emm, enn; imm, inn; omm, onn; umm, unn.} To\nl
esas la rezulto naturala di la principo, ke nula litero esas\nl
muta en Ido. Singla sempre sonas, sive en la komenco, sive\nl
en la mezo o fino dil silabi, kun sua sono alfabetala. --- \VUgras{Gn},\nl
konseque ta principo, audigas sua du literi, \VUgras{g} quale indikesis\nl
supere, e \VUgras{n} quale ni jus dicis. Ex. : \VUgras{regno} = \textit{reg-no}; \VUgras{digna}\nl
= \textit{dig-na} (2).
\end{VUpage}

% --- Feuillet 18 (folio 14) ----------------------------------------
%  Filet de note a y ~ 1450 px = 122,8 mm.
\begin{VUpage}[18]{14}
\VUnotes{124.52mm}{%
(1) Videz la prelasta alineo koncernanta \VUgras{u}, p. 12.--- La adopto di la\nl
litero \VUgras{q} liberigas de multa \VUgras{k} e lasas al vorti la formo internaciona :\nl
\textit{equilibro, equinoxo, equatoro, equaciono} ec. certe plu bona kam :\nl
\textit{ekvilibro, ekvinoxo, ekvatoro, ekvacio.}

(2) Autori qui admisas du soni por \VUgras{s} en la Linguo internaciona,\nl
ne volante chanjetar l'etimologial ortografio, obligas per la fakto\nl
ipsa la adepti lernar ne nur duopla regulo di pronunco, ma anke\nl
general ecepto, por imitar ula lingui nacionala, e mem ecepti dil\nl
ecepto. Tale \VUgras{s} duras sisar forte en \textit{asymptoto} (asimptoto), \textit{prae}\cc
\textit{sentir} (presentar), exemple, quankam l'ecepto generala dicas, ke \VUgras{s}\nl
sonas \VUgras{z} inter du vokali. Konkluze, por pronuncar \VUgras{s} konvene, en\nl
simila kazi, on mustas konocar l'etimologio dil vorto, o konsultar\nl
specala repertorio.%
}

\VUgras{p} = \textit{p} en l'Italiana, la Franca, la Hispana e. c. Ol devas\nl
esar tre klare dicernebla de \VUgras{b}.

\VUblancAlinea
\VUgras{q} sempre sequesas da \VUgras{u}, quale en la Latina ed en la lingui\nl
uzanta ta letro, qua sonas \textit{k} (1).

\VUblancAlinea
\VUgras{r} = \textit{r} Italiana. Se uli pronuncos lu kun kartavo, co ne\nl
havos en Ido efekto plu mala kam en la Franca linguo, che\nl
qua preske omna nordani kartavas pronuncante ta litero.

\VUblancAlinea
\VUgras{s} sempre sisas forte, sive en la komenco, sive en la mezo,\nl
sive en la fino dil vorti, e mem inter du vokali. Konseque ol\nl
nultempe recevas la sono di \VUgras{z} Franca od Angla. Per to la\nl
vorti \VUgras{roso, friso} e \VUgras{sono}, exemple, diferas de \VUgras{rozo, frizo, zono.}\nl
Do tre atencez pronuncar en Ido nultempe \VUgras{s} kun la sono\nl
di \VUgras{z} (2).

\VUblancAlinea
\VUgras{t} = \textit{t} dil Italiana, dil Angla, dil Franca e. c. en la si\cc
labo \textit{ta}. Ol konservas ica sono mem en la silabi \textit{tio, tia,}\nl
quin ula lingui pronuncas \textit{sio, sia} o \textit{tsio, tsia,} se li sonas pos\nl
altra silabo.

\VUblancAlinea
\VUgras{v} = \textit{v} en l'Italiana, l'Angla, la Franca, e \textit{w} en la Ger\cc
mana. La Hispani devas atencar aparte ta litero e ne pro\cc
nuncar ol quale \VUgras{b}, quo konfundigus \textit{volo} a \textit{bolo}, \textit{valo} a \textit{balo},\nl
\textit{voko} a \textit{boko} e. c.

\VUblancAlinea
\VUgras{w} = \textit{w} en la Angla vorti \textit{west, well, wist}. Ol adoptesis por\nl
konservar kun plu justa formo ula vorti diveninta interna\cc
ciona : \textit{westo, wato} (elektr.), \textit{warfo, wiskio} e. c.

\VUblancAlinea
\VUgras{x} esas pronuncata \textit{ks} o \textit{gz} segunvole, sen ula detrimento.\nl
Ca litero diminutas grandege la nombro dil \textit{k} quin on be\cc
\end{VUpage}

% --- Feuillet 19 (folio 15) ----------------------------------------
%  Filet de note : 1412 px au scan, 188 px pour la 1re ligne de corps,
%  soit 24,30 + (1412-188) x 25,4/300 = 127,97 mm. Mesurer l ordonnee
%  RELATIVEMENT a la 1re ligne, jamais sur le y brut du scan.
%  Titre ACENTO TONIKA. : cap 2,540 mm, largeur 29,13 mm
%                         -> 10,63 pt, interlettrage -34 (gras).
%  Petites capitales dans les notes (noms propres).
\begin{VUpage}[19]{15}
\VUnotes{130.00mm}{%
(1) En Esperanto « Doktoro \textsc{Zamenhof} skribas lore \textit{ks}, lore \textit{kz}\nl
(eksplodi, ekzameni). To esas un del kazi en qui, dicas la granda\nl
linguisto \textsc{Otto Jespersen}, il regardas nia lingui ocidentala tra sua\nl
Rusa binokli, l'ortografio Rusa ne posedante \VUgras{x} e skribante lore \textit{ks},\nl
lore \textit{kz} (quin la Rusi pronuncas \textit{gz}). Similajo eventas por \textit{oi} F., quan\nl
la Doktoro reprezentas per \textit{ua} omnafoye, kande la Rusa adoptis la\nl
vorto ed uzas ica ortografio (\textit{trotuaro, tualeto, vualo}), e per \textit{oi}\nl
o \textit{oj} en omna altra kazo (\textit{foiro, fojo, fojno}) ». --- Prefaco F. dil\nl
« Dictionnaire international français », p. IX (1908).

(2) To demonstresis en « Phonétique de l'Esperanto », p. 299,\nl
\textit{Bulletin Français-Ido}, n\textsuperscript{os} 77-78.%
}

\VUcontinue zonus sen olu; pluse, ol lasas al vorti lia aspekto ed orto\cc
grafio internaciona (1).

\VUblancAlinea
\VUgras{y} = \textit{y} en \textit{yeux} F., \textit{yes} E. En Ido ta litero nultempe divenas\nl
vokalo, quale en altra lingui. Pro fonetikal motivi tre grava,\nl
nultempe ol sequas \textit{a}, \textit{e}, \textit{o}, \textit{u} en la sama silabo, quale agas\nl
erore pri \textit{j}, olua korespondanto, la linguo Esperanto (2).

\VUblanc{2.77pt}% mesure au fac-simile
\VUgras{z} = \textit{z} en la Franca, l'Angla, la Portugalana lingui, o \VUgras{s} en\nl
\textit{Rose} Germana, \textit{rosa} Italiana.

\VUblancAlinea
\VUgras{ch} = \textit{ch} en la Angla \textit{church}, en la Hispana \textit{macho}, o \textit{tch}\nl
en la Franca : \textit{tchèque}, o \textit{c} en l'Italiana \textit{cibo}.

\VUblancAlinea
\VUgras{sh} = \textit{sh} en la Angla \textit{fish}, o \textit{sch} en la Germana \textit{Fisch}, o\nl
\textit{ch} en la Franca \textit{chat}, o \textit{sc} en l'Italiana \textit{asceta}.

\VUblancAlinea
Ma, se la du literi \VUgras{s-h} apartenas a du radiki diversa, en\nl
kompozita vorto, on separas li per streketo en la skribo ed\nl
on pronuncas konseque : \VUgras{chas-hundo}, ne \textit{cha-shundo}. On\nl
darfus anke uzar la formo : \textit{chaso-hundo}.

\VUblancAlinea
La nomi di la literi esas, por la vokali, lia sono ipsa :\nl
\VUgras{a, e, i, o, u}, e por la konsonanti, lia \textit{sono} sequata da e (\textit{é} F.) :\nl
\VUgras{bé, ce, de, fe, ge, he, ke... que... we, xe, ye, ze.}

% Blancs MESURES sur les lignes de base (outils/lignesbase.py) :
%   pas courant 41,19 px ; avant le titre 123,7 px, apres 95,0 px,
%   soit 6,98 mm et 4,56 mm a ajouter au pas.
\VUsaut{6.98mm}
\VUtitre{10.63pt}{-34}{ACENTO TONIKA.}
\VUsaut{4.56mm}

\VUblancAlinea
4. --- En Ido, la \textit{acento tonika} ne esas plualtigo o plu\cc
longigo dil sono, ma plufortigo di la voco sur un silabo dil\nl
vorti. On pozas lu « sur la lasta silabo dil infinitivi\nl
e sur la prelasta silabo di la cetera vorti » : \VUgras{amâr, venîr,}
\parplein
\end{VUpage}

% --- Feuillet 20 (folio 16) ----------------------------------------
%  Filet de note : ordonnee 107,95 mm (outils/filet.py).
\begin{VUpage}[20]{16}
\VUnotes{108.07mm}{%
(1) La cirkonflexo indikas nur la \textit{tonika acento}, ne la longeso\nl
di la silabo. Yen la texto ipsa dil decido (57), koncernanta ica\nl
acento : « L'acento esas sur la lasta silabo dil infinitivi, e sur la\nl
prelasta silabo di la cetera vorti. Ma en plursilaba radiki, \VUgras{i} e \VUgras{u}\nl
nemediate avan vokalo ne povas recevar l'acento. » --- \textit{Progreso},\nl
III, 322.

En la Hispana, Italiana, Portugalana lingui, l'infinitivo \VUgras{en-ar}\nl
recevas la acento. Ido imitas ta lingui e, pro analogeso, extensas ta\nl
acentizo a sua cetera infinitivi : \VUgras{-ir}, \VUgras{-or}.

(2) Kad on esas obligata pronuncar \VUgras{i} e \VUgras{u} quale \VUgras{y} e \VUgras{w} rispektive\nl
avan vokalo ? Nule : on darfas pronuncar indiferente \textit{naci-ono} o\nl
\textit{nacyono, mu-elar} o \textit{mwelar. La sola kozo importanta, do obligala,}\nl
\textit{esas la justa pozo di l'acento; on pronuncez do quiete :} \VUgras{âlio} o \VUgras{âlyo},\nl
\VUgras{mânuo} o \VUgras{mânwo} (\textit{Progreso}, IV, 142).

(3) Evidente \textit{mardio} ne esas \textit{mar}(o)\textit{dio, saturdio satur}(o)\textit{dio},\nl
e per la 5 \textit{sundio, lundio. merkurdio, jovdio, venerdio}, ni ne intencas\nl
dicar : la dio dil suno, dil luno, di Merkurius, di Jupiter, di Venus.\nl
(Segun \textit{Progreso}, VI, 135.)

(4) Videz p. 12, ye la 2-ma alineo,%
}

\VUcontinue \VUgras{skribôr; fenêstro, bovîno, malâda, trôno, mêa, puêri} opi\cc
\VUgras{niôno}, e. c.

%% aucun blanc au fac-simile
« Ma en plursilaba radiki, \VUgras{i} e \VUgras{u} nemediate avan vokalo\nl
ne darfas recevar l'acento » : \VUgras{Kâspia, râdie, fîlii, mistêrio;}\nl
\VUgras{tênua; îlua, êlua, sêxuo, stâtui, lînguo} (1).

%% aucun blanc au fac-simile
Pro ke la finali \textit{ia, ie, ii, io}, e \textit{ua, ue, ui, uo} di ca lasta\nl
vorti formacas nur un silabo, la \textit{prelasta} esas reale ta,\nl
qua preiras li, ed ol portas la acento tote konforme al\nl
regulo (2).

%% aucun blanc au fac-simile
En \VUgras{dio, pia, duo, gluo} ed altri simila, \VUgras{i} e \VUgras{u} esas vere la\nl
prelasta silabo. Li konseque recevas la acento, segun la\nl
regulo. Nam li kontenas ne plursilaba, ma \textit{unsilaba} radiko :\nl
\textit{di, pi, du, glu.}

%% aucun blanc au fac-simile
Se \textit{monosilaba} (unsilaba) radiko divenas lasta parto di\nl
kompozajo o derivajo, la vokali \VUgras{i}, \VUgras{u} portas la acento, quale\nl
kande li trovesas en simpla vorto. Ex. : \VUgras{butontrûo, arbor}\cc
\VUgras{glûo, despîa, cadîe.}

%% aucun blanc au fac-simile
Ma, pro ke la nomi di la sep dii semanala ne esas kompozaji\nl
(quale \textit{festo-dio, ca-die}) li recevas la acento sur la silabo qua\nl
preiras \textit{di} e qua esas reale la prelasta. On do pronuncas :\nl
\VUgras{sûndio, lûndio, mârdio, jôvdio, venêrdio, satûrdio} (3).

%% aucun blanc au fac-simile
Kande \VUgras{au}, \VUgras{eu} formacanta diftongo (4) renkontresas kom\nl
silabo prelasta, la acento restas sur ta diftongo, kun la chefa
\parplein
\end{VUpage}

% --- Feuillet 21 (folio 17) ----------------------------------------
%  Filet de note : ordonnee 125,14 mm (outils/filet.py).
\begin{VUpage}[21]{17}
\VUnotes{124.79mm}{%
(1) On bone remarkez, ke la regulo fakte inkluzas mem la mono\cc
silaba adjektivi, t. e. l'adjektivi elizionita, havanta monosilabo kom\nl
radiko : \textit{tal', nul', bel', pur', irg', grand'}, e. c.

(2) Okazione di « Profilaktol », \textit{metan, formol, amin}, qui recevas\nl
internacione la acento sur la lasta silabo, on trovas en \textit{Progreso}, II,\nl
679, la konsidero sequanta :

« On devas distingar du kazi o klasi : la \textit{propra nomi}, quale \textit{Profi}\cc
\textit{laktol}, qui esas stranjera vorti en nia linguo e konseque havas specala\nl
ortografio ed anke acentizo; e la \textit{komuna nomi}, qui apartenas a nia\nl
linguo e devas sequar la generala reguli, sive pri la gramatikal\nl
finali, sive pri l'acentizo. »

Ta remarkigo esas guido pri ta nomi stranjera relate Ido. Sam\cc
tempe ol furnisas respondo por la distraktita kritikeri qui objecionas%
}

\VUcontinue esforco dil voco sur \textit{a}, \textit{e}. Ex. : \VUgras{lâubo, kâuzo, lâuro, nêutra,}\nl
\VUgras{psêudo}. Evitez sorgoze facar ek \textit{au}, \textit{eu} du silabi en ta vorti\nl
ed altri simila.

%% aucun blanc au fac-simile
En la nombri, inter 20 e 100, la acento restas sur la\nl
radiko \textit{du}, \textit{tri}, \textit{quar} e. c. Ex. : \VUgras{dûa-dek} (o \VUgras{dûadek}), \VUgras{trîa-dek}\nl
(o \VUgras{trîadek}), \VUgras{quâra-dek} (o \VUgras{quâradek}), e. c. --- Komprenebla en\nl
\VUgras{dek-e-un, dek-e-du, dek-e-tri}, e. c. ne la konjunciono \textit{e}, ma\nl
la nombro \textit{dek} esas acentizata.

%% aucun blanc au fac-simile
La monosilabi darfas esar acentizata o ne (tam en prozo\nl
kam en poezio) segun la kuntexto e l'intenco dil autoro (1).

%% aucun blanc au fac-simile
Por eufonio e mem por bona interkompreno, l'acento\nl
tonika havas en Ido tre granda importo. Konseque omna\nl
Idisti devas tre atencar olu, e la populi qui pozas la acento\nl
sur la lasta silabo devas merkar e memorar bone, ke en Ido\nl
nur l'infinitivi \textit{-ar}, \textit{-ir}, \textit{-or} (qui esas lasta silabi) darfas kom\nl
tala recevar olu.

%% aucun blanc au fac-simile
Ma, kontraste, altri devas atencar, ke l'aplikado dil acento\nl
ne incitez li ad engluto dil silabi finala ed a neklara pro\cc
nunco di olia vokali. Nam to genitus grava erori pri la tempi\nl
en la verbo, o pri la speco gramatikala dil vorti e lia rolo.\nl
Kun atenco e sorgo (or mem plu kam linguo naciona la\nl
helpolinguo postulas la du), on tre bone povas konciliar\nl
justa acentizo e klara pronunco.

%% aucun blanc au fac-simile
Sentiginte la acento sur la silabo prelasta, evitez sentigar\nl
quaze duesma acento sur la dezinenco : \textit{Esperântô, harmôniô}.\nl
Certe la dezinenco devas audesar klare, ma por ico ne esas\nl
necesa, ke on acentizez lu anke. L'Italiani pruvas lo kon\cc
stante en sua linguo; ni imitez li pri co en Ido (2).
\end{VUpage}

% --- Feuillet 22 (folio 18) ----------------------------------------
%  Filet de note a y = 896 px. Ordonnee calculee DEPUIS LE FOLIO (y=116),
%  dont la position composee est fixe, et non depuis \VUmargeSup : cette
%  page ouvre par un titre, sa 1re ligne de corps n'est donc pas a la
%  marge habituelle. 18,71 + (896-116) x 25,4/300 = 84,75 mm.
%  Calculee a tort depuis \VUmargeSup, elle donnait 71,63 mm et le bloc
%  de notes remontait dans le texte courant.
%  Titre ARTIKLO. : cap 2,455 mm, largeur 15,66 mm
%  -> 10,27 pt, interlettrage -17 (gras).
%  La 1re note continue celle du folio 17.
\begin{VUpage}[22]{18}
\VUnotes{84.75mm}{%
\VUcontinue a ni ta vorti, quale se oli apartenus a la linguo e devus havar la\nl
dezinenci e la acento dil nomi komuna. Per quo \textit{Profilaktol}, \textit{etil}\nl
ed altra vorti kemiala apartenas a Ido plu multe kam \textit{Kashmir},\nl
\textit{Kamerun, Bengal, Portugal, Tonkin} e. c. qui havas nek la acentizo,\nl
nek la dezinenci dil Ido vorti komuna ?

Pri la decido koncerne la tonika acento en Ido e la studiado qua\nl
preparis lu, videz l'apendico unesma, ye la fino dil gramatiko.\nl
L'acentizo en Espo esas pura imito dil acentizo Polona.

(1) Pro kustumo generala, ni duras uzar en ica gramatiko l'adjek\cc
tivi-participa « definita », « nedefinita », pri l'artikli ed ula pronomi,\nl
quankam certe plu justa epiteto esas dezirinda.

(2) L'exempli kun \VUgras{le}, tale kam mult altri analoga quin on povus\nl
donar, montras per su, ke tro granda simpleso gramatikala povas\nl
meritar forjeto. Nam, exemple, l'unikeso dil artiklo definita en\nl
Esperanto impedas tradukar la supera frazi, tamen ne desfacila e\nl
tre ordinara. Ma vere, kad esas laudinda e quon valoras grama\cc
tikal simpleso sakrifikanta, od alteranta l'expresado dil pensi?\nl
Cetere uli mondolinguana kelke blinde alegas la facileso e simpleso\nl
gramatikala. Preske sempre li oblivias, ke la facileso por lerno ne\nl
koincidas necese kun la facileso por apliko. Mem eventas, ke l'unesma\nl
ofras reale nur pura trompilo, se ol nocas o jenas la duesma. Kun\nl
la duimo di sua reguli e vorti Espo e Ido certe esus plu facile\nl
lernebla. Ma kad oli esus pro ico plu bona ? On lernas helpo-linguo\nl
dum kelka hori, dum kelka dii, ma on aplikas lu dum yari. Do ques\cc
tionesas : quo meritas prefero? lerno kelkete plu kurta, ma aplikado\nl
entravata e defektoza; o studio kelkete plu longa, ma aplikado sen\cc
manka e skopokonforma? En l'unesma kazo, \textit{un} artiklo definita\nl
suficas; en la duesma kazo, certe \textit{du} artikli esas necesa.%
}

% Ordonnees mesurees RELATIVEMENT AU FOLIO, dont la position composee est
% fixe (\VUfolioY = 18,71 mm). Sur le scan : folio y=116, titre y=243,
% 1re ligne de corps y=337. Donc titre a 18,71 + 127x25,4/300 = 29,46 mm,
% corps a 18,71 + 221x25,4/300 = 37,42 mm.
% Mesurer depuis le bord de l'IMAGE donnait 20,57 mm pour le titre, qui
% venait alors chevaucher le folio : chaque image du scan a son propre
% decalage vertical.
\VUsaut{5.16mm}
\VUtitre{10.27pt}{-17}{ARTIKLO.}
\VUsaut{4.71mm}

\VUblancAlinea
5. --- L'\textit{artiklo definita} (1) esas \VUgras{la} por la du nombri; do\nl
lu esas nevariebla : \VUgras{la domo} (singularo), \VUgras{la domi} (pluralo).

%% aucun blanc au fac-simile
Kande nul altra vorto indikas pluralo, sive per sua formo\nl
(finalo \VUgras{-i}), sive per sua senco (nombro-nomo o nedefinita\nl
pronomo), on uzas \VUgras{le}; nam altre on ne savus (per \textit{la}) kad\nl
parolesas pri un individuo o pri pluri. Ex. : \VUgras{le} \textit{Gracchus},\nl
\VUgras{le} \textit{Cato} (Kato); \VUgras{le} \textit{x}, \VUgras{le} \textit{y}, \VUgras{le} \textit{z}; \textit{la cifri di ca konto esas tanta}\nl
\textit{male formacita, ke} \VUgras{le} \textit{3 e} \VUgras{le} \textit{5 esas konfundebla a} \VUgras{le} \textit{8} (2).

%% aucun blanc au fac-simile
6. --- On darfas elizionar la \VUgras{a} final dil artiklo, remplasi\cc
gante lu per apostrofo, tam egale avan konsonanto kam avan\nl
% Ligne presque entierement grasse : elle depasse de peu la justification.
% On la resserre, comme le fait le typographe, plutot que de laisser TeX
% la recouper (cf. \VUserre au preambule).
\VUserre{0.62}{vokalo : \VUgras{la charmo di la infanto} o \VUgras{di l'infanto}; \VUgras{la plumi di}}\nl
\VUgras{la ucelo} o \VUgras{di l'ucelo.}
\end{VUpage}

% --- Feuillet 23 (folio 19) ----------------------------------------
%  Filet de note : ordonnee 144,36 mm (outils/filet.py).
%  Premier jet OCR utilise comme brouillon, puis verifie et enrichi a l'oeil.
\begin{VUpage}[23]{19}
\VUnotes{144.36mm}{%
(1) Dec. 588 : On admisas l'eliziono di l'artiklo avan konsonanto\nl
dop la vorti \textit{da}, \textit{de}, \textit{di}. Ex. : \textit{da l'regulo, de l'regulo, di l'regulo.}

Dec. 713 : On adoptas la formi : \textit{dal regulo e da l'regulo}, kun o\nl
sen apostrofo, pos la tri prepozicioni \textit{da}, \textit{de}, \textit{di}.

Dec. 949 : On admisas \VUgras{al} (un vorto) kom abreviuro di \textit{a la}, apud\nl
l'abreviuro \textit{a l'} ja existanta. Pro analogeso a \textit{dal}, \textit{del}, \textit{dil} ja admisita.%
}

Ma on atencez ne elizionar la artiklo, se ol destruktas la\nl
aspiro di la litero \textit{h}. Do ne uzez \VUgras{l'homo}, \VUgras{l'hosti}, ma \VUgras{la homo,}\nl
\VUgras{la hosti.}

%% aucun blanc au fac-simile
Atencez anke evitar la miskompreno posibla. Do ne uzez :\nl
\VUgras{la duro di l'afero}, nam aude on povus komprenar : \VUgras{la duro}\nl
\VUgras{di la fero}. Dicez do : \VUgras{la duro di la afero.}

%% aucun blanc au fac-simile
Cetere l'eliziono en la artiklo esas \textit{darfo}, nule obligo.

%% aucun blanc au fac-simile
On darfas uzar la formo \VUgras{a l'}, \VUgras{da l'}, \VUgras{de l'}, \VUgras{di l'} e la kon\cc
traktaji \VUgras{al, dal, del, dil}, vice \VUgras{a la, da la, de la, di la}, sempre\nl
reguloza, do senhezite uzebla (1).

\VUblanc{5.00pt}% mesure au fac-simile
7. --- On uzas l'artiklo definita en la du sequanta\nl
kazi :

%% aucun blanc au fac-simile
1\textsuperscript{e} Kande la substantivo (singulara o plurala) indikas la\nl
tota speco, od omna individui di la speco : \VUgras{la leono ne esas}\nl
\VUgras{tam kruela kam la tigro; la uceli flugas en la aero, quale}\nl
\VUgras{la fishi natas en l'aquo.}

%% aucun blanc au fac-simile
2\textsuperscript{e} Kande ol indikas un o plura individui \textit{determinita} di\nl
la speco : \VUgras{la libri dil profesoro} (la komplemento : \textit{dil pro}\cc
\textit{fesoro} determinas \textit{libri}); \VUgras{querez la mediko} (la mediko kustu\cc
mata, o qua ja venis; \textit{mediko} signifikus : \textit{ula} od \textit{irga}\nl
\textit{mediko}).

%% aucun blanc au fac-simile
Ecepte ta du kazi, on devas ne uzar l'artiklo, e mem esas\nl
konsilata omisar olu, kande la substantivo havas senco\nl
generala ne determinita, exemple en la proverbi : \VUgras{Kontenteso}\nl
\VUgras{valoras plu multe kam richeso; povreso ne esas vicio.}

%% aucun blanc au fac-simile
Konseque ico, on nultempe bezonas uzar l'artiklo kun la\nl
nomi di abstraktita enti, qualesi, vertui, e. c., nam ta kon\cc
cepti ne esas koncepti di \textit{speci}, e ne korespondas ad individua\nl
objekti : \textit{fido, espero, karitato; kurajo, energio, esprito}, e. c.,\nl
e. c. (On remarkis ke en l'anciena Franca, tala vorti esis uzata\nl
sen artiklo). Same pri la nomi di cienci, qui esas quaze\nl
propra nomi : \textit{filologio, geometrio, fiziko}, e. c. Ma segun la\nl
regulo memorigita supere, ta nomi devas prenar l'artiklo, se
\parplein
\end{VUpage}

% --- Feuillet 24 (folio 20) ----------------------------------------
%  Filet de note a y = 1813 px, folio a y = 146 px
%  -> 18,71 + (1813-146) x 25,4/300 = 159,85 mm.
%  Page qui continue le folio 19 : le 1er alinea reprend au fer a gauche.
\begin{VUpage}[24]{20}
\VUnotes{159.85mm}{%
(1) Ica lasta alineo esas prenita ek \textit{Progreso}, I, 491.%
}

\VUcontinue li indikas \textit{un} aparta kozo inter plura kozi : \textit{la espero di Petro}\nl
(Petrus), \textit{la kurajo di Alexandro} (Alexander), \textit{la filozofio di}\nl
\textit{Epikuro} (Epikurus), \textit{la esprito di Voltaire}, e. c. Se on sequos\nl
ta regulo tre logikal, on sparos multa artikli, ed on igos la\nl
diskurso plu frapanta e plu « nervoza » (1).

%% aucun blanc au fac-simile
On anke ne uzas l'artiklo definita kun la propra nomi\nl
omnaspeca (mem di fluvii, monti, e. c.) nek kun la nomi\nl
komuna qui esas reale propra nomi, quale ti di la astri,\nl
sezoni, monati, dii.

%% aucun blanc au fac-simile
Kande la propra nomi esas preirata da titulo, on ne uzas\nl
l'artiklo : \VUgras{rejo Henrikus IV\textsuperscript{a}}, \VUgras{Papo Pius X\textsuperscript{a}}. Ma on uzas lu,\nl
kande la propra nomo esas nur apoziciono a nomo komuna :\nl
\VUgras{la genioza poeto Dante}, ed anke se la propra nomo akom\cc
panesas (preirata o sequata) da ula adjektivo : \VUgras{la bela}\nl
\VUgras{Helena, Alexander la Granda.}

%% aucun blanc au fac-simile
Rezume la artiklo definita uzesas nur kun substantivo\nl
expresata o tacata, tale ke en ica lasta kazo lu semblas rem\cc
plasar ta substantivo : \VUgras{Yen rozi; prenez la maxim bela} = \textit{la}\nl
\textit{maxim bela rozo} --- \VUgras{prenez le maxim bela} = \textit{la maxim bela}\nl
\textit{rozi.}

%% aucun blanc au fac-simile
On ne pensez ke « \VUgras{la} » esas sempre necesa avan \VUgras{maxim};\nl
nam ol esas nedependanta de ca adverbo, quale on quik vidos.\nl
Advere on dicas : \VUgras{ta homi sentas su la maxim felica} (\textit{de la}\nl
\textit{homi}, neexpresata); ma on dicas anke : \VUgras{la homi sentas su}\nl
\VUgras{maxim felica}, \textit{kande}... Konseque on devas nultempe repetar\nl
\VUgras{la} dop substantivo : \VUgras{la homi maxim felica} (e ne : \VUgras{la homi}\nl
\VUgras{la maxim felica}). Same on devas ne uzar \VUgras{la} avan \VUgras{maxim},\nl
sequata da adverbo : \VUgras{venez maxim frue} (e ne : \VUgras{la maxim}\nl
\VUgras{frue}). Fine on devas uzar \VUgras{la} kun la posedal pronomi nur\nl
kande to esas postulata dal senco (videz § 33).

%% aucun blanc au fac-simile
8. --- L'\textit{artiklo nedefinita} (F. \textit{un}, A. \textit{a}, I. \textit{un, uno, una})\nl
ne existas en Ido. La senco nedefinita indikesas dal fakto,\nl
ke l'artiklo \VUgras{la} ne preiras la substantivo. Kande on volas\nl
insistar pri la nedetermineso, on uzas \VUgras{ula}, e por nedeter\cc
mineso kompleta, \VUgras{irga}. Ex. : \VUgras{querez ula mediko, mem irga}\nl
\VUgras{mediko en la urbo, ma ne retrovenez sen mediko, nam sola}\nl
\VUgras{ni ne salvos l'infanto.}
\end{VUpage}

% --- Feuillet 25 (folio 21) ----------------------------------------
%  Filet de note a y = 1086 px, folio a y = 147 px
%  -> 18,71 + (1086-147) x 25,4/300 = 98,21 mm.
%  Titre SUBSTANTIVO. : cap 2,455 mm, largeur 24,05 mm
%  -> 10,27 pt, interlettrage -27 (gras).
%  La note (1) court sur le folio 22 : sa derniere ligne est pleine.
\begin{VUpage}[25]{21}
\VUnotes{98.21mm}{%
(1) Pri la pluralo per \VUgras{i} (e ne per \textit{s}) on trovos la motivi en la\nl
duesma apendico. Pri la \textit{genro} videz la 3-ma apendico : « Genro e\nl
maskulismo ». --- Pri la \textit{substituco} (e ne adjunto dil vokalo \VUgras{i})\nl
ni dicas :

\VUblancAlinea
Se \VUgras{o} sola suficas por reprezentar dezinence la ideo dil substantivo\nl
e dil singularo, pro quo \VUgras{i}, kun la sama radiko, ne suficus por repre\cc
zentar la ideo dil substantivo e ta dil pluralo, quale omna instante\nl
en l'Italiana ed en Slava lingui ? Se on opinionas nelogikala la procedo\nl
por pluralo, ol esas tala por singularo. Kad l'Esperantisti blamanta\nl
ne remarkis, ke la \VUgras{i} di lia infinitivo desaparas en la cetera formi dil\nl
verbo ? Nu, se \VUgras{o} devus konservesar en pluralo, kom signo dil substan\cc
tivo, totsame \VUgras{i} devus konservesar en la tota Espo-konjugo. Do\nl
l'Esperantisti devus havar : \textit{am-i-as, am-i-is, am-i-os, am-i-us, am-i-u,}\nl
\textit{am-i-anta} ec. Ma vere, kad la verbo esas min rikonocebla, quankam\nl
ol ne konservas en omna tempi e modi un karakterizivo komuna ?\nl
Pro quo agar tale pri la verbo, ma totaltre pri la substantivo ?

Fine, ka li povus justifikar quale la \VUgras{o}, signo di lia substantivo e\nl
samtempe di \textit{singularo}, darfas trovesar en lia \textit{pluralo} apud \textit{j}, signo di\nl
ca nombro ? Kad la duesma ne kontredicas l'unesma, od inverse ? Kad\nl
oli ne destruktas l'una l'altra ? Qua linguo esas justa hike, Espo o\nl
Ido ? On remarkez ke la \textit{j} ne esas substitucebla ad \textit{o} en Espo, nam\nl
\textit{patrij, lingvj}, e. c. esus ne pronuncebla; kontre ke \textit{patri, lingui} di
\parplein
}

\VUblancAlinea
% Alinea renfonce (x0 = 56 px releve) et non suite du folio precedent :
% troisieme \VUcontinue de trop, celui-ci trouve par le controle 10.
Kande on volas precize indikar la nombro 1, on uzas \VUgras{un.}\nl
Ex. : \VUgras{Un franko suficos.}

\VUblanc{5.75pt}% mesure au fac-simile
9. --- \textit{Artiklo partitiva} ne existas en Ido : \VUgras{donez a me}\nl
\VUgras{pano} = \textit{donez a me la kozo nomizata pano}. Se on volas in\cc
dikar parto o quanto nedeterminita, on uzas la prepoziciono\nl
\VUgras{de} : \VUgras{donez a me de vua pano, de vua pomi} (parto de vua\nl
pano, de vua pomi). Se on dicus : \VUgras{vua pano, vua pomi}, la\nl
senco esus : \textit{vua tota pano, vua omna pomi.}

\VUblancAlinea
Same on uzas \VUgras{de} kun pronomo (por ta ideo partitiva) :\nl
\VUgras{Yen kremo, prenez de olu} (poke o multe? pri co on ne pre\cc
cizigas; la quanto restas tote nedeterminita). Ma on povus\nl
precizigar, \textit{se on volus} : \VUgras{prenez kelke, multe de olu.}

\VUblancAlinea
\VUsaut{6.47mm}
\VUtitre{10.27pt}{-27}{SUBSTANTIVO.}
\VUsaut{4.10mm}

\VUblancAlinea
10. --- La \textit{substantivo} finas en singularo per \VUgras{o}, en plu\cc
ralo per \VUgras{i} : \VUgras{tablo, infanto; tabli, infanti} (1).
\end{VUpage}

% --- Feuillet 26 (folio 22) ----------------------------------------
%  Le folio (y = 165) echappe au detecteur sur cette page : il est pale.
%  Ordonnees recalculees depuis sa position relevee a la main :
%    corps  : 18,71 + (246-165) x 25,4/300 = 25,57 mm -> saut 1,27 mm
%    filet  : 18,71 + (1521-165) x 25,4/300 = 133,52 mm
%  Le bloc de notes s'ouvre par la fin de la note (1) du folio 21.
\begin{VUpage}[26]{22}
\VUnotes{133.90mm}{%
\VUcontinue Ido esas tre facile pronuncebla, mem plu facile kam \textit{patroj}, \textit{lingvoj}\nl
di Esperanto.

La logikisti, a qui ni jus respondis, pri la Ido-pluralo per substi\cc
tuco, devus prefere explikar, quale lia « \textit{tiu} », demonstrativo por\nl
foreso, misterioze divenas demonstrativo por proximeso, per la\nl
simpla adjunto di partikulo : \textit{chi}, qua tale recevas rolo stranjega,\nl
nam \textit{ulo} ne povas esar samtempe fora e proxima.

\VUblancAlinea
(1) Darfas uzesar metafore : l'ociereso esas la patro di omna vicii.\nl
\VUgras{Vere me ne povus dicar kad il od el esis plu vere mea patro spiri}\cc
\VUgras{tala pri la helpolinguo.}%
}

\VUsaut{1.27mm}
11. --- La genro (gramatikala) ne existas en Ido. La\nl
substantivo reprezentanta ento anmoza indikas l'animala\nl
speco, ma nule la sexuo. Altravorte, la sexuo di ta ento restas\nl
nedeterminita : \VUgras{infanto} ne indikas maskulo plu multe kam\nl
\VUgras{femino}; e tale pri omna substantivi di ento anmoza. Kom\cc
preneble omna substantivi di kozi esas neutra.

% Enumeration : chaque entree est un ALINEA a part entiere. Elle porte
% donc le renfoncement (mesure : +4,5 mm de la marge, soit \VUalinea) et
% sa derniere ligne reste courte. Chainees par une coupure explicite,
% ces lignes perdaient leur renfoncement ET se faisaient justifier de
% force : le fac-simile
% les donne a 29, 32, 45, 48, 56, 59, 63, 68 et 75 % de la justification.
% +36,7 px releves devant ce paragraphe (outils/cardage.py). L'entree
% suivante etant un alinea a elle seule, le poseur automatique ne
% trouve pas de frontiere ordinaire ou l'inscrire : pose a la main.
\VUblanc{8.84pt}
12. --- La sexuo povas esar determinata nur \textit{per sufixi} :

%% aucun blanc au fac-simile
\VUgras{-ul} por la maskuli;

%% aucun blanc au fac-simile
\VUgras{-in} por la femini.

%% aucun blanc au fac-simile
Exemple : \VUgras{infanto, kato, finko} (sexuo nedeterminita);

%% aucun blanc au fac-simile
\VUgras{infantulo, katulo, finkulo} (maskuli);

%% aucun blanc au fac-simile
\VUgras{infantini, katini, finkini} (femini);

%% aucun blanc au fac-simile
\VUgras{la homo} (sexuo nedeterminita : \textit{homa ento});

%% aucun blanc au fac-simile
\VUgras{la homulo} (homo \textit{maskula});

%% aucun blanc au fac-simile
\VUgras{la homino} (homo \textit{femina}).

%% aucun blanc au fac-simile
\VUgras{patro} (un del du genitanti) (1);

%% aucun blanc au fac-simile
\VUgras{patrulo} (la genitanto o patro maskula);

%% aucun blanc au fac-simile
\VUgras{patrino} (la genitanto o patro femina) qua tote darfas\nl
nomizesar \VUgras{matro};

%% aucun blanc au fac-simile
\VUgras{la patri} (la genitanti qui esas o maskula o femina). Ex. :\nl
\VUgras{On opinionas sempre plu generale, ke la patri dil Amerikani}\nl
\VUgras{aborigena esis Aziani.}

\VUblanc{8.81pt}% mesure au fac-simile
13. --- Kande la sexuo ja esas indikita da vorto o sufixo,\nl
esas tote neutila indikar lu itere : \VUgras{Andreas, mea kuzo, volas}\nl
\VUgras{esar advokato} (nek \VUgras{kuzulo}, nek \VUgras{advokatulo}). --- \VUgras{Mea fra}\cc
\VUgras{tino esas sekretario} (ne : \VUgras{sekretariino}), \VUgras{e mea fratulo esas}\nl
\VUgras{profesoro} (ne : \VUgras{profesorulo}). --- \VUgras{El esas helpanto di medi}\cc
\VUgras{kino} (ne : \VUgras{helpantino}).
\end{VUpage}

% --- Feuillet 27 (folio 23) ----------------------------------------
%  Filet de note a y = 1566 px, folio a y = 144 px
%  -> 18,71 + (1566-144) x 25,4/300 = 139,10 mm.
\begin{VUpage}[27]{23}
\VUnotes{139.10mm}{%
(1) Se parolesas pri pagana muliero sacerdota, e ne pri sacer\cc
doto homula (Kristana o ne), kompreneble on uzas \textit{sacerdotino}.

\VUblancAlinea
(2) Pri \VUgras{-o}, \VUgras{-a} e. c. kom sexual karakterizivi on trovos kompleta\nl
studiuro en la triesma apendico.

(3) On tote darfas dicar anke : \VUgras{siniorulo rejo, princo; siniorino}\nl
\VUgras{rejo, princo} o mem : \VUgras{siniorulo reja, princa; siniorino reja, princa.}\nl
Ma, en la du manieri, on evitez repetar \textit{-ul} o \textit{-in}. On uzez li \textit{nur}\nl
\textit{unfoye}. La sexuo ya ne bezonas indikesar \textit{dufoye}.%
}

14. --- Kompreneble on ne adjuntas \VUgras{-in} a radiko dicanta\nl
per su la sexuo feminala : \VUgras{muliero, amazono, subreto}, e. c.

%% aucun blanc au fac-simile
Same on ne adjuntas \VUgras{-ul} a radiko dicanta per su la sexuo\nl
maskulala : \VUgras{viro, oficiro, generalo, notario, sacerdoto,}\nl
e. c. (1).

%% aucun blanc au fac-simile
Fine on uzas nek \VUgras{-ul}, nek \VUgras{-in}, kande nulo postulas, ke on\nl
konocigez la sexuo. Exemple, se on skribos a homulo o\nl
homino prezidera, samideana, kunlaboranta, e. c., on ne\nl
dicos : \textit{kara preziderulo, samideanino, kunlaborantulo}, ma\nl
nur : \VUgras{kara prezidero, samideano, kunlaboranto.} Ilu ed elu\nl
ya perfekte konocas lia sexuo. Precize pro ke \textit{prezidero,}\nl
\textit{samideano, kunlaboranto} esas per su sensexua, li tote kon\cc
venas en tal okaziono, sive por viro, sive por muliero (2).

%% aucun blanc au fac-simile
15. --- Kom honor-titulo (parolante a persono altaranga,\nl
o pri olu), on uzas la vorto \VUgras{sinioro} (e \VUgras{siniorulo, siniorino},\nl
se to esas necesa por evitar konfundo). On adjuntas, segun\cc
bezone, la nomo di lua alta ofico sociala : \VUgras{sinioro rej(ul)o,}\nl
\VUgras{rej(in)o; sinioro princ(ul)o, princ(in)o} (3); \VUgras{sinioro epis}\cc
\VUgras{kopo.} Kompreneble on uzas nek \VUgras{-ul}, nek \VUgras{-in} kande, quale pri\nl
\VUgras{episkopo}, to esas neutila. Same on procedas pri \VUgras{sioro}, polita\nl
titulizo uzata por omna personi a qui, o pri qui on parolas :\nl
\VUgras{Sioro Ludovikus R.} Pro ke esas konocata la sexuo dil viro a\nl
qua on parolas, on ne dicas a lu \VUgras{siorulo}, ma nur \VUgras{sioro.}\nl
Ma, pro ke la letro destinata a spozulo povus donesar a\nl
spozino, od inverse, on enuncas \VUgras{-ul} od \VUgras{-in} en la adreso, se\nl
on skribas a spozi vivanta kune, o se on timas, ke altre\nl
agante, la letro ne atingos juste la destinario.

%% aucun blanc au fac-simile
16. --- \VUgras{Siorino} uzesas por muliero mariajita e por mu\cc
liero ne mariajita, kande on skribas o parolas \textit{ad eli}. Ma,\nl
kande on parolas \textit{pri eli}, on uzas \VUgras{damo} por l'unesma e \textit{dam}\cc
\parplein
\end{VUpage}

% --- Feuillet 28 (folio 24) ----------------------------------------
%  Filet de note a y = 1587 px, folio a y = 146 px
%  -> 18,71 + (1587-146) x 25,4/300 = 140,73 mm.
%  Titre PROPRA NOMI. : cap 2,455 mm, largeur 25,06 mm
%  -> 10,27 pt, interlettrage -18 (gras).
\begin{VUpage}[28]{24}
\VUnotes{138.92mm}{%
(1) \VUgras{Siorino} dicata a \textit{damo} ed a \textit{damzelo} dispensas de selekto ofte\nl
ne facila inter la du vorti. Tale \textit{damzelo} ne esos ofensata konfundesar\nl
a \textit{damo}, od inverse.

(2) Remplasigante \textit{c} harda per \textit{k} e \textit{ph} per \textit{f}.

(3) Nur la transkribo « Iesu » (vice \textit{Jesu} quan ni uzis erore,\nl
dum tempo pasable longa) esas justa. Ol meritas omnarelate prefer\cc
esar, e pro la formo, e pro la sono. On memorez ke \textit{J} evas nur de la\nl
\textsc{xvii}\textsuperscript{a} yarcento e ke ta litero nultempe uzesis da la Latini.%
}

\VUcontinue \textit{zelo} por la duesma, se on volas precizigar, ke ita mariajis\nl
su, e ke ica esas ankore celiba (1).

%% aucun blanc au fac-simile
En asemblajo kompozita ek homuli e homini, on dicas a\nl
li : \VUgras{siori.} Same, parolante pri viro e pri muliero kune, on\nl
dicas \VUgras{siori.} Exemple : \VUgras{vua vicini, siori B.}

%% aucun blanc au fac-simile
Pro ke l'expresuro \VUgras{gesiori B.} povas indikar la viri e la\nl
mulieri dil tota familio B. (komparez \VUgras{gefrati}), la sola\nl
maniero certa dicar \VUgras{siorulo} e \VUgras{siorino B.} intermariajita esas :\nl
\VUgras{la gespozi B.}

% Un titre courant n'est pas une frontiere d'alinea ordinaire : les
% blancs qui l'encadrent se mesurent, ils ne se deduisent pas. Ecarts
% releves au feuillet 28 : 52 px au-dessus du titre, 81 px au-dessous,
% pour un pas regulier de 41,43 px.
\VUsaut{1.00mm}
\VUtitre{10.27pt}{-18}{PROPRA NOMI.}
\VUsaut{3.49mm}

17. --- La \textit{propra nomi} omnaspeca devas principe konsi\cc
deresar kom « vorti stranjera » a la linguo. La nomi per\cc
sonal precipue, pro ke li esas la \textit{proprietajo} dil personi qui\nl
nomesas per oli, devas restar netushebla. Konseque on trans\cc
skribas li segunlitere, kande li esas skribita per l'alfabeto\nl
Romana, mem la Greka nomi, di qui la transskribo Latina\nl
esas klasika. On riproduktas, se on povas, la diakritika\nl
signi ed on indikas, \textit{segun quante on povas}, la pronunco\nl
inter parentezi. Se li apartenas a linguo ne uzanta l'alfabeto\nl
Romana, on transskribas li fonetike (\textit{maxim bone posible}).\nl
Por ico on uzas specal alfabeto plu richa kam la Romana\nl
alfabeto e posedanta diakritika signi (\VUgras{ä, ö, ü}, e. c.) e digrami\nl
(\VUgras{dh, th, kh}, e. c.) Exemple : \VUgras{Caesar, Cicero, Scipio, Grac}\cc
\VUgras{chus, Sokrates, Demosthenes, Pythagoras, Phryne; Goethe,}\nl
\VUgras{Shakespeare, Corneille, Boileau; Mickievicz, Leszczynski,}\nl
\VUgras{Przemysl; Pushkin, Pashich, Tolstoy, Shchavinskiy.}

%% aucun blanc au fac-simile
Por la propra nomi trovata en la olda ed en la nova Testa\cc
mento on konsilas anke la transkribo Latina, ja konocata;\nl
ma nur la vera qua uzas \VUgras{i}, ne \VUgras{j}. Ex. : \VUgras{Ierusalem, Iudas, Iob,}\nl
\VUgras{Beniamin, Isaias, Ioannes, Iozef, Zakarias} (2), \VUgras{Nazareth,}\nl
\VUgras{Bethlehem, Elizabeth}, e. c. (3).
\end{VUpage}

% --- Feuillet 29 (folio 25) ----------------------------------------
%  Filet de note a y = 1681 px, folio a y = 122 px
%  -> 18,71 + (1681-122) x 25,4/300 = 150,71 mm.
\begin{VUpage}[29]{25}
\VUnotes{149.49mm}{%
(1) Remplasigante \textit{c} harda per \textit{k} e \textit{ph} per \textit{f}.

(2) \VUgras{Avan-Azia} por la Latina « Asia Minor ».

(3) Por simpligar e generaligar, Ido uzas ta formo (\textit{-ia}) anke en\nl
\VUgras{Chinia, Japonia, Brazilia, Mexikia} (Chef-urbo : \textit{Mexiko}).%
}

Por la \textit{baptonomi} on konsilas anke transskribar li segun\nl
la formo Latina (1). To esas l'unika moyeno eskapar divergi\nl
sennombra. Cetere quale selektar inter \textit{Giovanni, Jean, Johann,}\nl
\textit{Jan, Hans, John, Ivan}, e. c.? Ka ni kreos vortaro specala por\nl
la baptonomi? Ni cherpez de la Latina lia nominativo e dicez :\nl
\VUgras{Ioannes, Ioanna; Iulius, Iulia; Iakobus; Andreas; Lukas;}\nl
\VUgras{Antonius, Antonia; Gabriel, Rafael; Ludovikus, Petrus,}\nl
\VUgras{Paulus, Elizabeth, Anna, Sofia, Maria}, e. c.

\VUblanc{5.75pt}% mesure au fac-simile
18. --- Devas konsideresar kom \textit{vorti stranjera} ed esar\nl
konservata en lia stranjera formo, \textit{tam singulare kam plurale},\nl
omna vorti exkluzive nacionala o lokala, ed anke la nomi\nl
di moneti, ponderili e mezurili ne apartenanta al metral\nl
sistemo : \VUgras{pasha, lama, ulema, geisha, kimono; uryadnik,}\nl
\VUgras{nagayka, troïka; cicerone} (pl. \VUgras{ciceroni}), \VUgras{lazzarone} (pl. \VUgras{laz}\cc
\VUgras{zaroni}); \VUgras{Targui} (pl. \VUgras{Touareg}); \VUgras{mehari} (pl. \VUgras{mehara}); \VUgras{cent}\nl
(pl. \VUgras{cents}), \VUgras{para, duro, peseta} (pl. \VUgras{pesetas}); \VUgras{pound, pud,}\nl
\VUgras{klaft, shtof, verst}, e. c.

%% aucun blanc au fac-simile
Tamen ulkaze la pluralo darfas indikesar per \VUgras{-i} adjuntita\nl
al vorto stranjera e separita de olu per streketo. Exemple, se\nl
on ne konocas ta pluralo, o kande la vorto stranjera havas\nl
singulare la sama formo kam plurale. Per \VUgras{-i} on lore evitas\nl
miskompreno. Ex. : \VUgras{Ni povas pagar vua cheko; kad vu deziras}\nl
\VUgras{franki o mark-i?}

\VUblanc{6.00pt}% mesure au fac-simile
19. --- La maxim multa nomi di landi e ti di la kin\nl
parti dil mondo konservas sua Latina formo historiala, kom\nl
fakte konocata internacione :

%% aucun blanc au fac-simile
\VUgras{Europa, Afrika, Amerika, Azia} (2), \VUgras{Oceania, Anglia,}\nl
\VUgras{Belgia, Bolivia, Dania, Francia, Germania, Grekia, Hispania,}\nl
\VUgras{Italia, Rusia, Skotia, Suedia; Arabia, Armenia, Australia,}\nl
\VUgras{Laponia}, e. c. (Videz la lexiko) (3).

\VUblanc{5.98pt}% mesure au fac-simile
20. --- La nomi di la habitanti derivas de oli per l'ad\cc
junto dil sufixo \VUgras{-an}, pos supreso dil \VUgras{a} finala : \VUgras{Europ-ano,}\nl
\VUgras{Itali-ano, Bolivi-ano}, e. c.
\end{VUpage}

% --- Feuillet 30 (folio 26) ----------------------------------------
%  Filet de note a y = 1624 px, folio a y = 149 px -> 143,59 mm.
%  C'est la page qui a servi a la calibration initiale du corps (§ 4.2
%  du LISEZ-MOI) : ses lignes 1 a 6 sont en partie GRASSES, ce qui avait
%  fausse cette calibration. La voici enfin relevee pour de bon.
\begin{VUpage}[30]{26}
\VUnotes{141.87mm}{%
(1) Nur la formo per \VUgras{-ano} esas generala e generaligebla; la nomi\nl
di populi quale \textit{Anglo}, \textit{Ruso} esas nur abreviuri di la reguloza \textit{Angli}\cc
\textit{ano}, ec. e devas konsideresar kom aparta kazi. (\textit{Progreso}, VI, 95.)

(2) La \VUgras{c} transskribesas per \VUgras{k} en \VUgras{Nikaragua, Kanada, Maroko}, pro\nl
ke en Ido ni havus (kun \textit{c}) \textit{Nitsaragua, Tsanada, Marotso}, quo certe\nl
ne konkordus kun l'internaciona sono di ta vorti. Ma qua homo,\nl
vidante li kun \textit{k}, vice \textit{c}, havas mem ombro di hezito por rikonocar li?%
}

\VUblancAlinea
La formi \VUgras{Angliano, Franciano, Germaniano, Rusiano,}\nl
\VUgras{Suediano, Daniano, Poloniano, Grekiano, Arabiano, His}\cc
\VUgras{paniano, Skotiano, Laponiano, Suisiano} esas tote reguloza,\nl
do permisata. Ma per imito dil uzado internaciona, qua por\nl
ta nomi di rasi o populi havas formi plu kurta, Ido preferas\nl
\VUgras{Anglo, Franco, Germano, Ruso, Suedo, Dano, Polono, Greko,}\nl
\VUgras{Arabo, Hispano, Skoto, Lapono, Suiso} (1).

%% aucun blanc au fac-simile
De li venas reguloze l'adjektivi \VUgras{Angla, Franca, Germana,}\nl
\VUgras{Rusa, Sueda, Dana, Polona, Greka, Araba, Hispana, Skota,}\nl
\VUgras{Lapona, Suisa.}

\VUblanc{4.79pt}% mesure au fac-simile
21. --- La nomi landala qui ne havas ta susteno inter\cc
naciona dil formo \textit{ia} (Latina origine) konservas sua formo\nl
nacionala : \VUgras{Honduras, San Salvador, Nikaragua, Venezuela,}\nl
\VUgras{Uruguay, Paraguay, Maroko, Kanada} (2), \VUgras{Chili, Peru, Por}\cc
\VUgras{tugal} (Videz la lexiko).

%% aucun blanc au fac-simile
La nomi finanta per \textit{land} kompreneble recevas \VUgras{o} finala,\nl
por kompletigar la vorto segun la fizionomio Idala di\nl
« lando ». Ex. : \VUgras{Finlando, Holando, Irlando, Islando, Neder}\cc
\VUgras{lando, Nova-Zelando, Zulu-lando}, e. c.

%% aucun blanc au fac-simile
Segun la principo generala, la nomi di la habitanti esas\nl
formacata per la sufixo \VUgras{-an}. Ex. : \VUgras{Honduras-ano, Venezue}\cc
\VUgras{lano, Uruguayano, Marokano, Kanadano, Chiliano, Peruano,}\nl
\VUgras{Portugalano, Finlandano, Holandano, Irlandano}, e. c.

%% aucun blanc au fac-simile
\VUgras{Usa}, to esas la tri literi U. S. A., abreviuro di \textit{Unionita}\nl
\textit{Stati} (di) \textit{Amerika} (nordala), uzesas en Ido, vice ta nomo\nl
tro longa, por nomar la granda republiko. Per to ni simple\nl
imitas la postal uzado internaciona, qua skribas U. S. A.\nl
(Usa) sur omna sendaji ad ica lando.

\VUblanc{4.67pt}% mesure au fac-simile
22. --- Pro praktikal kurteso on darfas uzar : \VUgras{la Angla,}\nl
\VUgras{la Franca, la Germana, la Araba}, e. c., por indikar la lingui\nl
Angla, Franca, Germana, Araba (tacante la vorto \textit{linguo},\nl
quan on tote darfas expresar). Ta dicomaniero povas pro
\parplein
\end{VUpage}

% --- Feuillet 31 (folio 27) ----------------------------------------
%  Filet de note a y = 1491 px, folio a y = 117 px -> 135,04 mm.
%  Titre ADJEKTIVO QUALIFIKANTA. : cap 2,455 mm, largeur 47,58 mm
%  -> 10,27 pt, interlettrage +2 (gras).
\begin{VUpage}[31]{27}
\VUnotes{132.87mm}{%
(1) Omna adjektivi di rasi o di populi sempre komencas en Ido\nl
per mayuskulo. Tale on evitas omna eroro o hezito inter du reguli,\nl
quale en la Franca, ube on skribas : \textit{le peuple anglais}, la Angla\nl
populo; \textit{l'anglais} (la linguo), la Angla; \textit{l'Anglais}, la Anglo.

(2) Habitanto dil urbo, e \VUgras{Luxemburgiano}, habitanto dil Dukio\nl
« Luxemburgia »; quale \VUgras{Mexikano}, habitanto di \textit{Mexiko}, urbo, e\nl
\VUgras{Mexikiano}, habitanto di \textit{Mexikia}, lando.

(3) Plu exakte : a radiko \textit{nomala} on adjuntas \VUgras{-a} por igar lu\nl
gramatikale adjektivo qualifikanta.

Radiko esas \textit{nomala}, se, pro sua signifiko e natural valoro, ol
\parplein
}

\VUcontinue duktar nula ambigueso, pro ke la \textit{homo} Angla, Franca, Ger\cc
mana, Araba, nomizesas : \VUgras{la Anglo, la Franco, la Germano,}\nl
\VUgras{la Arabo} (1).

\VUblanc{4.99pt}% mesure au fac-simile
23. --- La nomi geografiala di urbi, riveri, monti, pro\cc
vinci, distrikti, e. c. esas propra nomi; li do sequas logike\nl
la regulo qua koncernas ici (§ 17) : \VUgras{Paris, London, Roma,}\nl
\VUgras{Firenze, Lisboa, Torino, Napoli, Madrid, München, New-}\nl
\VUgras{York, Cambridge, Genève, Boulogne, Kharkov, Voronef,}\nl
\VUgras{Dniepr, Kiakhta, Khabarovsk; Shang-haï, Kiao-Cheu, Che-}\nl
\VUgras{Fu, Tsu-shima, Hoang-ho}, e. c.

%% aucun blanc au fac-simile
Nur la nomi di kelka monti, mari o fluvii internaciona\nl
(quin nia lingui ne skribas same) recevas l'Idal ortografio,\nl
por evitar pri oli divergi jenanta di skribo e di pronunco :\nl
\VUgras{Alpi, Blanka Monto, Blanka Maro, Reda Maro, Nigra Maro,}\nl
\VUgras{Kaspia Maro, Kaspio; Maro Mediteranea, Mediteraneo; Atlan}\cc
\VUgras{tiko; Pacifika Oceano, Pacifiko; Glacial Maro, Oceano;}\nl
\VUgras{Rheno, Danubio.}

%% aucun blanc au fac-simile
Remarquez bone, ke omna nomi di la lasta alineo kom\nl
propra nomi ne recevas la artiklo, malgre la uzo kontrea\nl
di ula lingui.

%% aucun blanc au fac-simile
Por la habitanti dil urbi on adjuntas \VUgras{-ano} (adjektivo\nl
\textit{-ana}) : \VUgras{Paris-ano} (pl. \textit{Paris-ani}), \VUgras{Paris-ana; München-ano,}\nl
\VUgras{Genève-ano, Luxemburg-ano} (2).

\VUblancAlinea
\VUsaut{4.0mm}
\VUtitre{10.27pt}{2}{ADJEKTIVO QUALIFIKANTA.}
\VUsaut{4.6mm}

\VUblancAlinea
24. --- L'\textit{adjektivo qualifikanta} finas per \VUgras{-a} e ne varias (3).

%% aucun blanc au fac-simile
Por chanjar ol a substantivo samsenca suficas substitu\cc
car al dezinenco \VUgras{a} la vokalo \VUgras{o} (singulare), o la vokalo \VUgras{i}
\parplein
\end{VUpage}

% --- Feuillet 32 (folio 28) ----------------------------------------
%  Filet de note a y = 1414 px, folio a y = 135 px -> 127,00 mm.
%  Le bloc de notes s'ouvre par la fin de la note (3) du folio 27.
\begin{VUpage}[32]{28}
\VUnotes{127.00mm}{%
\VUcontinue povas formacar substantivi per l'adjunto di \VUgras{o}, ed adjektivi per\nl
l'adjunto di \VUgras{a}.

Radiko esas \textit{verbala}, se, pro sua signifiko e natural valoro, ol\nl
povas formacar verbi per l'adjunto dil dezinenci verbala : \VUgras{-as, -is,}\nl
\VUgras{-os} e. c.

(1) Pri la substantivigo e homigo dil adjektivo, videz la apendico\nl
quaresma.

(2) En la komenco, Ido rezervis, pro prudenteso, la yuro indikar\nl
(se to esus necesa) la pluralo en l'adjektivi per adjuntar \VUgras{-i} a lia\nl
dezinenco \VUgras{a} (do \VUgras{ai}). Ma la praktiko montrinte, ke ta special pluralo\nl
povas abandonesar, decido unanima (1218) dil akademio Idista\nl
supresis lu.%
}

\VUcontinue (plurale) : \VUgras{richa} = \textit{qua esas richa} : homo richa; \VUgras{richo} =\nl
(strikte) \textit{ento, individuo richa} e (konvencione, praktike) \textit{homo}\nl
\textit{richa}; --- \VUgras{povra} = \textit{qua esas povra} : homo povra; \VUgras{povro} =\nl
(strikte \textit{ento, individuo povra} e (konvencione, praktike) \textit{homo}\nl
\textit{povra}; --- \VUgras{blinda} = \textit{qua esas blinda} : homo blinda; \VUgras{blindo}\nl
= (strikte) \textit{ento, individuo blinda} e (konvencione, praktike)\nl
\textit{homo blinda}; --- \VUgras{dezerta} = \textit{qua esas dezerta}; \VUgras{dezerto} = \textit{ulo,}\nl
\textit{loko, spaco dezerta}; --- \VUgras{bona} = \textit{qua esas bona}; \VUgras{bono} = (kon\cc
vencione, praktike) \textit{homo bona}; --- \VUgras{mala} = \textit{qua esas mala};\nl
\VUgras{malo} = (konvencione, praktike) \textit{homo mala}; --- \VUgras{la boni} =\nl
\textit{la homi bona}; --- \VUgras{la mali} = \textit{la homi mala} (1).

\VUblanc{7.75pt}% mesure au fac-simile
25. --- Kande adjektivo aplikesas a substantivo tacita, on\nl
indikas la pluralo per la artiklo \VUgras{le}, o per nedefinita pro\cc
nomo konvenanta, o mem simple per \VUgras{-i}. Ex. : \VUgras{Yen pomi,}\nl
\VUgras{prenez le bona e lasez le mala.} Kun ideo partitiva on dicus :\nl
\VUgras{il ofris a me blanka e reda rozi, me prenis uli reda o kelki}\nl
\VUgras{reda}; o plu simple : \VUgras{me prenis redi}, quale on dicus : \VUgras{me}\nl
\VUgras{prenis kelki.} Tre certe ya \VUgras{redi} e \VUgras{kelki} povas relatar nur rozi,\nl
sola kozo pri qua on parolas. La lasta exemplo montras\nl
bone, ke en l'unesma ni povabus dicar : \VUgras{prenez la boni}\nl
\VUgras{e lasez la mali.} Nam vere quon altra kam \textit{pomi} relatus « \VUgras{la}\nl
\VUgras{boni} », « \VUgras{la mali} », en ica frazo? Kad on komprenus \textit{homi},\nl
kande traktesas nur \textit{pomi}? (2).

\VUblanc{8.01pt}% mesure au fac-simile
26. --- On darfas elizionar la \VUgras{a} final dil adjektivi, ma\nl
kun la kondiciono ke to ne produktos akumulo de konso\cc
nanti. On konsilas uzar ta eliziono precipue kun l'adjektivi
\parplein
\end{VUpage}

% --- Feuillet 33 (folio 29) ----------------------------------------
%  Page a notes tres longues : 10 lignes de corps, 38 de notes.
%  Le detecteur de filet echoue ici (le filet est haut dans la page et
%  la colonne de notes le suit de tres pres). Releve a la main : y = 610.
%  -> 18,71 + (610-138) x 25,4/300 = 58,67 mm.
\begin{VUpage}[33]{29}
\VUnotes{56.71mm}{%
(1) Decido 504 : On konservas l'eliziono di \VUgras{-a} en l'adjektivi, ma\nl
restriktante ol a la kazi, en qui ne rezultos akumulo de konsonanti;\nl
per 7 voci ek 9. (\textit{Progreso}, IV, 434.)

Decido 586 : On repulsas la supreso di l'eliziono di l'artiklo;\nl
per 6 voci ek 9. (\textit{Progreso}, IV, 561.)

Decido 587 : On repulsas admisar sempre l'eliziono di l'artiklo\nl
dop vokalo ed avan konsonanto; unanime per 8 voci. Exemple :\nl
\textit{ke l'regulo, se l'regulo}. (\textit{Progreso}, IV, 561.)

(2) « La finalo \VUgras{-a} povas (darfas) supresesar per eliziono, ed\nl
esas konsilata uzar ta moyeno por igar la parolado min monotona\nl
ed anke plu fluanta; vice dicar : \VUgras{la alerta e joyoza infanti,} on dicos\nl
tre bone : \VUgras{l'alert e joyoz infanti,} e to esas tante plu naturala, ke ta\nl
senacenta finali (dezinenci) spontane malaparas (desaparas) en\nl
parolado fluanta e poke (kelke) rapida. To esas tute (tote)\nl
konforma a la linguala sentimento di la lingui qui posedas acento\nl
tonika (di intenseso). » --- L. \textsc{Couturat}, \textit{Progreso}, I, 492.

Pro ke nur la nomala radiko darfas recevar l'eliziono, en sua\nl
formo adjektivala, nula detrimento povas naskar, por la kompreno,\nl
de ta eliziono \textit{moderata}. Ol esas permisata, konsilata, en la kondi\cc
cioni dicita supere. Ma la homi a qui ol ne plezus, havas la yuro\nl
absoluta ne praktikar olu.

L'adjektivo elizionesas en nia linguo, e ne la substantivo (kontraste\nl
kun l'uzado Esperantista), pro ke l'adjektivo ne varias en Ido,\nl
kontre ke la substantivo varias. Esperanto, pro la variebleso di lua\nl
substantivo, povas elizionar lu nur en singularo e nur en nominativo,\nl
do ecepte. Ico anke indikis klare ad Ido elizionar, ne sua substan\cc
tivo variebla, ma sua adjektivo nevariebla. Ma ico anke interdiktas\nl
elizionar lua substantivo, sempre en prozo, e quaze sempre en\nl
poezio. Cetere l'eliziono en substantivo povus konfundigar ica ad\nl
adjektivo, ed efektigar erori, miskompreni. Pro la praktiko esus\nl
kontrelogika permisar l'eliziono en du vorti tante parenta. Ni nul\cc
tempe obliviez, ke la versi praktikesas por la linguo, ma ne la\nl
linguo por la versi. Ni do nultempe sakrifikez, quale tro ofte\nl
l'Espisti, ita ad ici. L'eliziono dil substantivo restez en nia versi\nl
ecepta licenco tre rara, ed ol nultempe sakrifikez l'exakteso dil\nl
penso a rimo nule necesa. La versi povas influar la prozo; or nia\nl
linguo povus vivar eterne sen rimi, ma se ol perdus la \textit{just} expre\cc
siveso dil pensi, la morto kaptus ol tre balde.%
}

\VUcontinue derivita e partikulare kande li finas per \VUgras{-al-(a)}. Ex. :\nl
\VUgras{infantal anmo, amikal ago, kordial saluto}, e. c. (1).

%% aucun blanc au fac-simile
On devas ne uzar l'eliziono tro freque. Ol ne diplasas la\nl
acento; konseque ica restas sur \VUgras{a} en \VUgras{infantal, amikal, kor}\cc
\VUgras{dial}, quale en \VUgras{infantala, amikala, kordiala.}

%% aucun blanc au fac-simile
Diveninte monosilaba per l'eliziono, l'adjektivi dusilaba\nl
sequas, pri l'acento, la regulo di omna monosilabi, pri qui\nl
parolesas en § 4.

%% aucun blanc au fac-simile
On ne obliviez ke l'eliziono nultempe esas obligala (2).
\end{VUpage}

% --- Feuillet 34 (folio 30) ----------------------------------------
%  Le detecteur de filet echoue encore (les notes le serrent de pres) :
%  releve a la main, premiere rangee sombre y = 1419, folio y = 118.
%  -> 18,71 + (1419-118) x 25,4/300 = 128,86 mm ; corrige a 129,94 mm,
%  le releve a la main du filet (seuil et fermeture differents de ceux
%  d outils/filet.py) tombant 1,08 mm trop haut. Ecart constate par le
%  controle 11 : le bloc de notes etait 12,7 px au-dessus du modele.
%  Cette page a revele un defaut de l'analyse du fac-simile : les eclats
%  de la tranche du feuillet voisin, a 40 px du dernier caractere, se
%  faisaient souder aux lignes. Corrige dans outils/page.py (bornage de
%  la colonne par densite AVANT la soudure).
\begin{VUpage}[34]{30}
\VUnotes{129.92mm}{%
(1) « La reguli di nia linguo lasas tote libera la loko (plaso) di\nl
l'adjektivo, e permisas pozar ol avan o dop la substantivo, segun\nl
volo o gusto. Ta libereso esas nule « galicismo ». Tote kontree, la\nl
« regulo » pozar ol sempre avan la substantivo esus « germanajo »,\nl
ofte nekomoda e mem tre jenanta; adminime, omnasupoze, restrikto\nl
tote neutila di la libereso di la vortordino. » (\textit{Progreso}, II, 703.)

Ni remarkigez ica stranjajo : Esperanto, qua mem per la plumo\nl
di Zamenhof imitas la Germanajo supere aludita, ne ordinas :\nl
\textit{internacia lingvo}, ma : \textit{lingvo internacia}; tale ke ol uzas l'abre\cc
viuro : \textit{L. I.} (e nule \textit{I. L.}) quale ni agas en Ido. Kun \textit{o}, konseque,\nl
la rezultajo esus \textit{Lio}, ne \textit{Ilo}.%
}
% Page cardee. Pas mesure sur les 16 pas reguliers du corps
% (outils/lignesbase.py) : 43,10 px au lieu des 41,45 habituels, soit
% 43,10 x 72,27/300 = 10,38 pt. Les \VUblanc ci-dessous sont l'EXCES du
% pas releve sur ces 43,10 px.
\VUinterlignePage{10.38pt}

27. --- Pri la \textit{plaso dil adjektivo} on juas en Ido la libe\cc
reso maxim kompleta. Nula logikal motivo impozas a ni la\nl
kustumo Angla e Germana pozar l'adjektivo sempre e ma\cc
shinatre \textit{avan} la substantivo, dum ke altra lingui pozas lu\nl
\textit{dope} (1).

\VUblancAlinea
On apene povas donar la konsili sequanta :

\VUblanc{5.61pt}% mesure au fac-simile
1\textsuperscript{e} Se l'adjektivo esas plu longa per du o tri silabi kam\nl
la substantivo, pozez lu dop ica : \VUgras{linguo internaciona}; \VUgras{rolo}\nl
\VUgras{desfacila, raporto nekomunikebla.}

\VUblancAlinea
2\textsuperscript{e} Se uzesas plura adjektivi, pozez li dop la substantivo :\nl
\VUgras{la homi instruktita, senpartisa e prudenta...}

\VUblancAlinea
3\textsuperscript{e} Se l'adjektivo havas komplemento, e precipue komple\cc
menti, pozez unesme la substantivo : \VUgras{tablo-tapiso makulizita}\nl
\VUgras{da inko e sauci diversa.}

\VUblancAlinea
4\textsuperscript{e} Se l'adjektivo esas tam longa kam la substantivo, o\nl
preske, konsultez l'eufonio, e nome evitez la hiati, se to\nl
esas posibla : \VUgras{aquo pura, alno alta, alta querko, inteligenta}\nl
\VUgras{pueri, argumenti konvinkiva,} plu plezas al orelo kam \VUgras{pura}\nl
\VUgras{aquo, alta alno, querko alta, pueri inteligenta, konvinkiva}\nl
\VUgras{argumenti.}

\VUblanc{5.80pt}% mesure au fac-simile
Che omna linguo, en la parolado sorgata, on selektas la\nl
maxim plezanta ordino dil vorti. Ma komprenebla, en fami\cc
liara diskurso o konverso, on agos tote senjene pri la plaso\nl
dil adjektivo (o participo) qualifikanta.
\end{VUpage}

% --- Feuillet 35 (folio 31) ----------------------------------------
%  PREMIER TABLEAU DU VOLUME. Releve : outils/page.py + mesures a la
%  main des abscisses (voir LISEZ-MOI, § 7 ter).
%  Lignes de base du tableau : 413, 454, 495 px pour le groupe
%  « Komparativo », 561, 602,5, 643,5 px pour le groupe « Superlativo ».
%  Le pas y vaut 41 px : le tableau est sur la grille du volume.
%  « Superlativo » et « de o ek » tombent a MI-HAUTEUR entre deux
%  lignes ; ils sont poses avec \VUdecale.
%  Graisse : la colonne des formes citees (plu, tam, min, kam, maxim,
%  minim, tre) est en gras ; les noms des degres (Supereso, egaleso,
%  infreso, relatanta, absoluta) et « de o ek » sont en romain. Verifie
%  a la loupe sur un agrandissement x3 (scan/T35c.png, T35d.png).
\begin{VUpage}[35]{31}
% Ordonnee du bloc de notes. Le filet releve donne 102,95 mm ; la
% comparaison des lignes de note composees et relevees (le controle 11 ne
% peut pas la faire ici, le decoupage en lignes du tableau ne concordant
% pas) demande 1,27 mm de plus.
\VUnotes{104.22mm}{%
(1) Pro quo Ido ne havas sintezal komparativo (\textit{di supereso})\nl
per -er, e superlativo (\textit{relatanta}) per -est, quale en D. E.? Yen la\nl
respondo trovata en \textit{Progreso}, VI, 241 :

1. La formaco di komparado per partikuli esas plu klara e plu\nl
konforma a la tendenco di l'evoluciono di omna lingui. --- 2. La\nl
lingui D. E. konocas la formaco per partikuli (\textit{more, mehr; most,}\nl
\textit{meist}), ma la cetera lingui ne konocas la formaco per sufixi (ecepte\nl
\textit{-issim} e kelka arkaika formi). 3. --- La sufixo -er esas ja okupata\nl
kun altra senco tre internaciona e tre utila. --- 4. \textit{Maxim} (di nia\nl
superlativo relatanta) esas bone konocata e tote internaciona per\nl
\textit{maximum} (same : \textit{minim} per \textit{minimum}). --- 5. On ne propozas\nl
supresar \textit{plu}, \textit{maxim} por la verbi (plu amar, maxim amar), do\nl
l'uzo di la partikuli kun l'adjektivi esas plu simpla e plu reguloza ».

Anglo igis remarkar, ke « La sufixe -\textit{er} ed -\textit{est} por komparo en\nl
l'Angla rare uzesas kun vorti di la Latina, qui komparesas quale\nl
en Ido. La plumulto dil vorti en Ido venas de la Latina. Do la uzi\nl
di ta sufixi esus tre nefamiliara por Angli ».

Se, vice -\textit{er} por la komparativo, on uzus -\textit{ior} (renkontrata en\nl
\textit{injenioro, juniora, inferioro} e. c.) e \textit{est} por la superlativo, ni havus\nl
hibridajo shokanta e qua produktus : \textit{glorioziora, ciencoziora, ocie}\cc
\textit{riora, laboreriora, mentieriora} kom komparativi di supereso qui,
% La note se poursuit au folio suivant : sa derniere ligne court donc
% jusqu'a la marge (1081 px releves pour une justification de 1083).
\parplein
}
\VUsaut{5.2mm}
\VUtitre{10.15pt}{-19}{GRADI KOMPARALA.}
\VUsaut{4.8mm}

% Accolades du tableau. Etendues relevees au fac-simile (seuil 180,
% fenetres etroites autour de chaque trait) :
%   1. « Komparativo » : y 382..498, centre 440, hauteur 116 px = 27,9 pt.
%      Elle va de Supereso a infreso en passant par egaleso ; sa pointe
%      est sur la rangee du milieu.
%   2. « relatanta / absoluta » : rangees B et C du second groupe,
%      centre 607,5, hauteur 72 px = 17,4 pt. Tres pale au scan.
%   3. « Supereso / infreso » : y 531..596, centre 566, hauteur 65 px.
%   4. « maxim / minim » : y 543..612, centre 577,5, hauteur 69 px.
%      C'est une accolade FERMANTE : elle designe « de o ek », a droite.
%      Son centre tombe sur la ligne de base de « de o ek » (581).
% Lignes de base : 413, 454, 495 (groupe 1) ; 561, 602,5, 643,5 (groupe 2).
\VUrang{%
\VUcase{4.91mm}{28.}%
\VUcase{41.15mm}{Supereso}%
\VUcase{65.87mm}{\VUgras{plu}}%
\VUcase{72.39mm}{.}\VUcase{74.34mm}{.}\VUcase{76.12mm}{.}%
\VUcase{77.98mm}{.}\VUcase{79.84mm}{.}%
\VUcase{81.03mm}{\VUgras{kam}}}
\VUrang{%
\VUcase{16.17mm}{Komparativo}%
\VUcase{37.68mm}{\VUaccolade{27.2pt}{1.23pt}}%
\VUcase{41.15mm}{egaleso}%
\VUcase{65.53mm}{\VUgras{tam}}%
\VUcase{72.56mm}{.}\VUcase{74.42mm}{.}\VUcase{76.28mm}{.}%
\VUcase{78.15mm}{.}\VUcase{80.01mm}{.}%
\VUcase{81.20mm}{\VUgras{kam}}}
\VUrang{%
\VUcase{41.15mm}{infreso}%
\VUcase{65.79mm}{\VUgras{min}}%
\VUcase{72.31mm}{.}\VUcase{74.17mm}{.}\VUcase{76.03mm}{.}%
\VUcase{77.89mm}{.}\VUcase{79.76mm}{.}%
\VUcase{80.94mm}{\VUgras{kam}}}
\VUblanc{6.02pt}% +25 px entre les deux groupes du tableau
\VUrang{%
\VUcase{36.83mm}{\VUaccolade{15.7pt}{-3.71pt}}%
\VUcase{39.79mm}{Supereso}%
\VUcase{61.64mm}{\VUgras{maxim}}%
\VUcase{77.30mm}{\VUaccoladeD{16.6pt}{-6.48pt}}%
\VUcase{81.28mm}{\VUdecale{-4.82pt}{de o ek}}}
% « Superlativo » tombe a mi-hauteur entre ce rang et le suivant. Il est
% rattache a CELUI-CI, abaisse de la valeur relevee, et non au suivant
% releve d'autant : la position composee est la meme, mais c'est ainsi
% que pdftotext le groupe, et le controle 2 peut donc le comparer.
\VUrang{%
\VUcase{0.85mm}{\VUdecale{-4.70pt}{Superlativo}}%
\VUcase{19.39mm}{\VUaccolade{17.4pt}{-3.71pt}}%
\VUcase{21.84mm}{relatanta}%
\VUcase{39.71mm}{infreso}%
\VUcase{61.55mm}{\VUgras{minim}}}
\VUrang{%
\VUcase{21.93mm}{absoluta}%
\VUcase{64.09mm}{\VUgras{tre}}}
\VUblanc{9.88pt}% une ligne blanche entre le tableau et le texte

Kande parolesas nur pri du objekti, on darfas remplasigar\nl
la superlativo relatanta per la komparativo : \VUgras{la plu yuna}\nl
\VUgras{de mea du fratuli,} vice : \VUgras{la maxim yuna de (o ek) mea}\nl
\VUgras{du fratuli.}

La Franca vorto \textit{encore}, avan komparativo, tradukesas per\nl
\textit{mem}. Ex. : \VUgras{Petrus esas plu bona kam Paulus, ma Ioannes}\nl
\VUgras{e Maria esas mem plu bona.}

\VUgras{Kam} uzesas kun omna vorti implikanta komparo, quale\nl
\VUgras{sama, tala, altra, preferar,} e. c. Ex. : \VUgras{egala o plu granda}\nl
\VUgras{kam... Me preferas autuno kam printempo} (1).
\end{VUpage}

% --- Feuillet 36 (folio 32) ----------------------------------------
%  Filet des notes releve a la main (le detecteur echoue) : y = 833,
%  folio y = 118 -> 18,71 + (833-118) x 25,4/300 = 79,25 mm.
%  Second tableau du volume, plus simple que celui du folio 31 : pas
%  d'accolade, quatre rangs sur la grille (lignes de base 434, 477,
%  518,6 et 560 px, pas regulier 41,6). Abscisses relevees une a une.
%  Cardage : exces de 22,4 px au-dessus du tableau et de 22,7 px
%  au-dessous, mesures par outils/cardage.py.
\begin{VUpage}[36]{32}
\VUnotes{77.68mm}{%
% Cette note continue celle du folio 31 : sa premiere ligne part de la
% marge (x0 = 1 px releve, contre 43 pour un debut d'alinea de note).
\VUcontinue emulesus kun la superlativi di supereso : \textit{grosieresta, prudentesta,}\nl
\textit{modestesta, pacientesta} e. c.; ed en adverbo : \textit{grosieriore, glorioziore,}\nl
\textit{ciencoziore, ocieriore; prudenteste, modesteste, diligenteste} e. c. Ne\nl
pen-valorus havar \textit{du} sistemi di komparado por tala rezulto.

Ultre co, \textit{un} plusa sufixo esus necesa anke por la komparativo di\nl
\textit{infreso, un} por la komparativo di \textit{egaleso}, omna du tam meritanta\nl
e \textit{tam frequa} kam la komparativo di supereso. Fine \textit{un} sufixo por\nl
la superlativo di infreso (la nuna \textit{minim}) tam meritanta e \textit{tam}\nl
\textit{frequa} kam la superlativo di supereso. Entote \textit{kin} plusa sufixi\nl
(komparala) esus necesa en Ido por pagar ta luxo di \textit{duesma} sistemo\nl
di komparado, kande ja 17 yari pruvas konvinkive, ke la nuna\nl
sufikas komplete.

(1) L'esperantista pronomo \textit{mi} havas la praktikal desavantajo\nl
tre grava, ke la orelo povas facile konfundar lu a \textit{ni}, aparte en\nl
telefono. Do lu esis forjetenda. Pluse, la pronomo \textit{me} di Ido, qua\nl
remplasas lu, donas la posedali \textit{mea}, \textit{mei}, certe plu komprenebla\nl
sen lerno da multa homi, nome dal latinisti, kam \textit{mia}, \textit{miaj} di Espo.

(2) Ke \VUgras{lu} servas la tri genri, segunvole, en la singularo, ne esas\nl
en su plu astonanta kam vidar \VUgras{li} servar la tri genri, segunvole, en\nl
la pluralo. Ma, kontraste, ni ulkaze bezonas distingar la pronomo\nl
« li » per \textit{ili, eli, oli}, totsame kam ni distingas la pronomo « lu »\nl
per \textit{ilu, elu, olu}. La Angla e la Germana lingui pri co pruvas nur\nl
un kozo : ke li ulkaze jenesas por expresar distingo pronomala\nl
quan altra lingui expresas facile en frazi tala kam ici : \textit{Ili} arivis\nl
plu frue kam \textit{eli}. \textit{Eli} meritis la sama puniso kam \textit{ili}. Vu departos\nl
unesme kun \textit{ili}, e me rajuntos vu kun \textit{eli}. L'Esperantisti eroras\nl
grande, kande li laudas pri sua linguo, e relate ca punto, simpleso\nl
\textit{tro granda}.

Per decido (1558) l'akademio Idista repulsis omna restrikto pri\nl
la uzo di \textit{lu}. On darfas do uzar ta pronomo totsame por kozo e por
% La note se poursuit au folio suivant : derniere ligne justifiee.
\parplein
}
\VUsaut{5.2mm}
\VUtitre{10.15pt}{-14}{PERSONAL PRONOMI.}
\VUsaut{4.9mm}

29. --- La personal pronomi esas :

\VUblanc{3.54pt}% +22,4 px au-dessus du tableau
\VUrang{%
\VUcase{47.67mm}{Singularo}%
\VUcase{70.53mm}{Pluralo}}
\VUrang{%
\VUcase{11.77mm}{1-ma persono :}%
\VUcase{49.87mm}{\VUgras{me}\ (1)}%
\VUcase{74.34mm}{\VUgras{ni}}}
\VUrang{%
\VUcase{11.60mm}{2-ma persono :}%
\VUcase{49.62mm}{\VUgras{vu}}%
\VUcase{74.51mm}{\VUgras{vi}}}
\VUrang{%
\VUcase{11.68mm}{3-ma persono :}%
\VUcase{49.87mm}{\VUgras{lu}}%
\VUcase{74.51mm}{\VUgras{li}}}
\VUblanc{3.61pt}% +22,7 px au-dessous du tableau

Por la duesma persono singulara existas anke formo fami\cc
liara : \VUgras{tu}, quan on darfas uzar nur ad amiki tre intima, a\nl
frati o parenti kun qui on uzas, en sua linguo matrala, formo\nl
familiara korespondanta.

\VUgras{Lu} (quale \VUgras{li}) esas uzata \textit{por la 3 genri} (2). Ma, se esas
\parplein
\end{VUpage}

% --- Feuillet 37 (folio 33) ----------------------------------------
%  Filet releve a la main : y = 1287, folio y = 129
%  -> 18,71 + (1287-129) x 25,4/300 = 116,75 mm.
%  ATTENTION : le decoupage en lignes du fac-simile SOUDE ici deux
%  lignes de note, « ... vice la simpla » et « lua, lia. » (bande
%  1650-1714) : la seconde est courte, entierement en italique, et ses
%  jambages se melent aux descendantes de la premiere. Le fac-simile
%  compte donc 41 lignes et non 40. Le releve ci-dessous en tient
%  compte ; les controles 9 et 11 sauteront la page en le signalant.
\begin{VUpage}[37]{33}
\VUsignature{168.74mm}{2}% signature de cahier
\VUnotes{116.75mm}{%
% Cette note continue celle du folio 32 : premiere ligne a la marge.
\VUcontinue persono di sexuo evidenta, kam por animali di sexuo nekonocata e\nl
persono qua havas nomo sensexua, quale \textit{infanto}, \textit{puero}, \textit{homo} e. c.,\nl
tam vere maskula kam femina.

La motivi di la decido donesis en « Mondo », XI, 68 : « Lu » por\nl
la singularo esas exakte lo sama kam « li » por la pluralo. Logiko,\nl
simetreso e facileso postulas ico. Konseque, same kam « li » darfas\nl
uzesar por personi, animali, kozi, omnafoye kande nulo obligas\nl
expresar la sexuo, tale « lu » darfas uzesar por personi, animali,\nl
kozi en la sama kondiciono. La distingo propozita esus subtilajo\nl
jenanta. Simila subtilajo rezultus pri la posedala pronomi, qui\nl
freque devus esar : \textit{ilua, elua, olua, ilia, elia, olia}, vice la simpla\nl
\textit{lua, lia.}

\VUblanc{3.97pt}% +24,2 px avant la note (1)
(1) Kande \VUgras{il, el, ol} trovesas ye la fino di frazo-membro (quale\nl
\textit{ilu}, hike) esas plu eufonioza uzar \VUgras{ilu, elu, olu,} kompleta (sen eliziono)\nl
kam \textit{il}, \textit{el}, \textit{ol} qui lore, sen la susteno di vokalo, restas quaze suspen\cc
data en la aero.%
}

\VUcontinue necesa (o se on volas) distingar la genri, on uzas singulare :\nl
\VUgras{ilu} (mask.), \VUgras{elu} (fem.), \VUgras{olu} (neut.), quin maxim ofte on\nl
kurtigas per \VUgras{il, el, ol}; e plurale : \VUgras{ili, eli, oli.}

En Ido, la distingo dil genri esas \textit{naturala} : on uzas la\nl
maskulal genro por la maskuli, la feminal genro por la\nl
femini, la neutra por la kozi, ed anke por la enti di qui\nl
la sexuo esas nedeterminita : \textit{infanti, pueri, homi, yuni, oldi,}\nl
\textit{personi, individui, kavali, bovi, kati,} e. c.

La triesma persono havas \textit{pronomo reflektiva} : \VUgras{su} por la\nl
singularo e la pluralo. Nultempe ol esas subjekto, e sempre\nl
ol referas la subjekto dil propoziciono en qua ol trovesas,\nl
kande ta subjekto reprezentas la triesma persono. Ex. :\nl
\VUgras{Il parolis pri su} (sama persono kam \VUgras{il}, subjekto); \VUgras{eli parolas}\nl
\VUgras{pri su} (sama personi kam \VUgras{eli}, subjekto dil propoziciono).\nl
Ma : \VUgras{Il parolas pri ilu} (1) (altra persono kam \VUgras{il}, subjekto);\nl
\VUgras{li parolas pri li} (altra personi kam \VUgras{li}, subjekto).

Infinitivo, o participo havanta komplemento, konstitucas\nl
propoziciono : \VUgras{me vidis li lavar su} o : \VUgras{me vidis li lavanta}\nl
\VUgras{su} = \VUgras{me vidis li qui lavis su.}

\VUblanc{5.85pt}% +32,0 px avant le paragraphe 30
30. --- Existas en Ido personal pronomo nedefinita : \VUgras{onu,}\nl
qua preske sempre trovesas en la formo kurtigita \VUgras{on} (segun\nl
\textit{ilu, elu, olu,} divenanta \textit{il, el, ol}). Ma nulo interdiktas uzar\nl
anke por olu la formo kompleta.

La « me » filozofiala tradukesas exakte per la substan\cc
\end{VUpage}

% --- Feuillet 38 (folio 34) ----------------------------------------
%  Filet releve a la main : y = 1264, folio y = 135
%  -> 18,71 + (1264-135) x 25,4/300 = 114,30 mm.
%  La liste des adjectifs possessifs (lignes 5 a 18) est faite d'alineas
%  courts, tous au renfoncement (x0 = 51 px releve), aucun justifie :
%  chacun est donc un alinea a lui seul, comme l'enumeration du folio 22.
\begin{VUpage}[38]{34}
\VUnotes{112.70mm}{%
(1) Pro ke « la ego » esas vera substantivo, nulo impedas donar\nl
a lu la \VUgras{-i} dil pluralo, se to esas necesa. Forme e sone « egi » ne esas\nl
plu stranja kam « eki », pluralo di \textit{eko}.

(2) On asertis nejuste, ke \VUgras{mea, tua, vua} e. c. devus esar \textit{meala,}\nl
\textit{tuala, vuala} e. c.

1\textsuperscript{e} To supozus la \textit{meo}, quale \textit{rejala} supozas la \textit{rejo}, fakte existanta.\nl
Or existas ne \textit{la meo, la tuo} e. c., ma \textit{lo mea, lo tua, lo sua} e. c.

2\textsuperscript{e} \VUgras{-ala}, quale omna sufixi o prefixi, aplikesas nur a nomala\nl
radiki od a verbala radiki. Or \VUgras{me, tu} e. c., esas nek nomal, nek\nl
verbal radiki. Kad on postulas \textit{-ala} por l'adjektivi pronomi nedefi\cc
nita? Kad on volus havar \textit{kelkala, ulala, omnala}, exemple, vice\nl
\textit{kelka, ula, omna?}

3\textsuperscript{e} Kad on adjuntas \textit{-ala} a vorti tote kompleta per su, quale\nl
\textit{avan, pos, quik} e. c.? Or \VUgras{me, tu, vu} e. c. esas ne radiki por derivaji,\nl
ma vorti kompleta apartenanta a specal klaso (quale la nombro\cc
vorti). Pro quo do on dicus \textit{meala, vuala} e. c., kande on ne dicas\nl
\textit{avanala, posala, quikala}, ma \textit{avana, posa, quika?}%
}

\VUcontinue tivo : la \VUgras{ego}, qua ritrovesas en la vorti internaciona : \textit{ego}\cc
\textit{ism(o)}, \textit{egoist(a)}, \textit{egoist(o)} (1).

\VUsaut{6.9mm}
\VUtitre{10.15pt}{-11}{POSEDAL ADJEKTIVI E PRONOMI.}
\VUsaut{4.6mm}

31. --- La posedal adjektivi esas :

\VUblanc{2.19pt}% +16,8 px avant la liste
\VUgras{Mea}, qua apartenas a \textit{me};

\VUgras{tua}, qua apartenas a \textit{tu};

\VUgras{vua}, qua apartenas a \textit{vu};

\VUgras{lua}, qua apartenas a \textit{lu};

\VUgras{nia}, qua apartenas a \textit{ni};

\VUgras{via}, qua apartenas a \textit{vi};

\VUgras{lia}, qua apartenas a \textit{li};

\VUgras{ilua}, qua apartenas ad \textit{ilu};

\VUgras{elua}, qua apartenas ad \textit{elu};

\VUgras{olua}, qua apartenas ad \textit{olu};

\VUgras{ilia}, qua apartenas ad \textit{ili};

\VUgras{elia}, qua apartenas ad \textit{eli};

\VUgras{olia}, qua apartenas ad \textit{oli};

\VUgras{sua}, qua apartenas a \textit{su}.

\VUblanc{4.04pt}% +24,5 px apres la liste
Quale on vidas, ta vorti esas nulo altra kam la pronomi\nl
personala, a qui on adjuntas la dezinenco \VUgras{a} dil adjektivi (2).
\parplein
\end{VUpage}

% --- Feuillet 39 (folio 35) ----------------------------------------
%  Filet releve a la main : y = 1127, folio y = 137
%  -> 18,71 + (1127-137) x 25,4/300 = 102,53 mm.
\begin{VUpage}[39]{35}
\VUnotes{102.09mm}{%
(1) Texte aparinta (\textit{Progreso}, V, 627) en longa expliko quan ni\nl
donis ibe, e qua duras per ico :

« Exemple, on devas dicar : « Il anuncas, ke \textit{lua} amiko departis »,\nl
e ne : \textit{sua}, nam \textit{lua amiko} esas ipsa subjekto di la propoziciono\nl
subordinita, do ne povus referar per \textit{sua} la subjekto di la chefa\nl
propoziciono (\textit{il}).

On darfas dicar : « Il venis kun \textit{sua} amiki », ma on devas dicar :\nl
« Il e \textit{lua} amiki venis » (ne \textit{sua}), nam tala propoziciono kontenas\nl
reale du propozicioni : « il venis », e « \textit{lua} amiki venis, do la sub\cc
jekto di la duesma ne povas referar per \textit{sua} la subjekto di l'unesma.

Altra kazo, en qua l'uzo di \textit{sua} esas tentiva ed erorigiva. On\nl
darfas dicar : Tala esas la metodo, quan Prof. X... uzis en \textit{sua}\nl
explori », ma on devas dicar : Tala esas la metodo \textit{uzita da} Prof. X...\nl
en \textit{lua} explori », nam \textit{sua} povus referar nur la metodo.

Rezume, on darfas uzar preske sempre \textit{lua}, \textit{lia}, mem en la kazi,\nl
en qui \textit{sua} esus korekta; do esas plu sekura uzar prefere \textit{lua}, \textit{lia},\nl
ecepte en la kazi di reala dusenceso, por qui \textit{sua} esas adoptita e\nl
vere utila. »

Ta expliko detaloza komprenigas quante Ido erorus, se ol imitus\nl
la Franca uzado pri \VUgras{sua}. E tamen ta uzado imitesas blinde da uli\nl
de nia kritikeri, qui reprochas a nia linguo imitar la Franca!%
}

32. --- \VUgras{Sua}, quale la personal pronomo \VUgras{su}, sempre re\cc
feras la subjekto di la propoziciono en qua ol trovesas,\nl
kande ta subjekto reprezentas la triesma persono. Ex. : \VUgras{Il}\nl
\VUgras{promenas kun sua amiko} (la amiko di \textit{il}, subjekto); \VUgras{il pro}\cc
\VUgras{menas kun sua amiko e sua filii} (la amiko e la filii di \textit{il},\nl
subjekto); \VUgras{il promenas kun sua amiko e lua filii} (la amiko\nl
di \textit{il}, subjekto; ma la filii dil amiko, \textit{ne} subjekto); \VUgras{amar}\nl
\VUgras{sua filii esas naturala} (sua propra filii, ne la filii di altra).

« On \textit{darfas} uzar \VUgras{sua} \textit{nur} kande ol referas la subjekto\nl
di la propoziciono, en qua ol trovesas. On \textit{devas} uzar lua o\nl
lia en la cetera kazi. Ed on \textit{darfas} uzar li mem kande li\nl
referas la subjekto di la propoziciono, se nur to genitas nula\nl
dusenceso. Konseque, on esas \textit{obligata} uzar \VUgras{sua} nur en kazo\nl
di posibla dusenceso e por evitar olu. » (1).

\VUblanc{6.02pt}% +32,7 px avant le paragraphe 33
33. --- La \textit{posedal pronomi} esas identa al adjektivi :\nl
\VUgras{mea, tua, vua,} e. c.; ma li recevas la pluralo per la chanjo\nl
del dezinenco \VUgras{-a} al dezinenco \VUgras{-i}, signo generala dil pluralo :\nl
\VUgras{mei, tui, vui, lui, sui, ilui, elui,} e. c.

% Ces deux lignes paraissent grasses sur une vue d'ensemble — c'est
% l'encre qui s'etale. Le controle 9 donnait -24 % et -17 % (le compose
% plus charge que le modele), et l'agrandissement x 1,9 le confirme :
% elles sont en romain.
De ta regulo konsequas ke :

Li darfas esar preirata dal artiklo, se l'objekto esas deter\cc
\end{VUpage}

% --- Feuillet 40 (folio 36) ----------------------------------------
%  Filet releve a la main : y = 1344, folio y = 120
%  -> 18,71 + (1344-120) x 25,4/300 = 122,34 mm.
%  Comme au feuillet 37, le detecteur manque la derniere ligne de note
%  (« XI, 296.) », courte et pale : sa projection culmine a 68 quand le
%  seuil de la page est plus haut). Le fac-simile compte 38 lignes, le
%  releve en donne 38 ; les controles 9 et 11 sauteront la page.
%  Blancs du titre : 132,0 px au-dessus, 99,2 px au-dessous (mesures sur
%  les lignes de base, outils/lignesbase.py).
\begin{VUpage}[40]{36}
\VUnotes{121.53mm}{%
(1) La decido 950 pri la posedal pronomi judikesis unanime dal\nl
Akademio en la maniero sequanta :

« Pri la decido 950 (VI, 161) l'Akademio konstatas 1\textsuperscript{e} ke, ante\nl
la decido e pos olu, uni uzis \textit{mea}, \textit{tua}, e. c. kom pronomi segun la\nl
maniero D. F. I. S. to esas : kun l'artiklo \textit{la}; altri segun la maniero\nl
Angla, to esas : sen artiklo; 2\textsuperscript{e} ke, pro oblivio, la decido reale ne\nl
propozesis nek diskutesis en Progreso dum la tri monati preirinta,\nl
quale postulis la statuti; deklaras ke per la fakto la du manieri\nl
restis e restas permisata. L'Akademio judikos ipsa, \textit{pos la periodo}\nl
\textit{di stabileso}, kad ol devas indikar un del du manieri kom unike\nl
uzenda. Komprenende, \textit{kom adjektivi posedala} (to esas avan nomo :\nl
mea domo, mea matro e. c.) la sempre antea e generala uzado per\cc
manas : nul artiklo devas preirar \textit{mea}, \textit{tua} e. c. en ica kazo. (\textit{Mondo},\nl
XI, 296.)%
}

\VUcontinue minita : \VUgras{me havas mea chapelo, prenez (la) vua; me ne}\nl
\VUgras{havas kavalo; prestez ad me un de la vui.}

Vice \VUgras{la mei, la vui, la lui, la lii, la nii,} e. c., on darfas\nl
anke uzar : \VUgras{le mea, le vua, le lua, le lia, le nia,} e. c. Ex. :\nl
\VUgras{me perdis mea gepatri, ka vu havas ankore la vui o le}\nl
\VUgras{vua?} (1)

\VUblanc{2.07pt}% +8,6 px
Esas bone remarkenda, ke la posedal adjektivi implikas\nl
l'ideo di \textit{determineso} : \VUgras{mea amiko} = ta quan me amas\nl
aparte, ta quan on konocas kom mea amiko partikulara, in\cc
tima. Se existus nedetermineso, on dicus : \VUgras{amiko di me}, to\nl
esas : un (irga) de mea amiki.

\VUsaut{6.77mm}
\VUtitre{10.15pt}{-10}{DEMONSTRATIV ADJEKTIVI-PRONOMI.}
\VUsaut{5.07mm}

34. --- La \textit{demonstrativ adjektivi} (sequata da substan\cc
tivo) havas \VUgras{a} kom dezinenco ed esas nevariebla, quale la\nl
cetera adjektivi : \VUgras{ca} od \VUgras{ica} por l'objekti, di qui on volas\nl
indikar explicite la proximeso; \VUgras{ta} od \VUgras{ita} por la fora objekti,\nl
od ordinare. Ex. : \VUgras{Me amas ica puerino, ma me odias ita}\nl
\VUgras{puerulo.}

Quale en ta exemplo, on opozas l'adjektivo \VUgras{ca} al adjek\cc
tivo \VUgras{ta}, kande parolesas pri du enti o kozi, de qui un esas\nl
proxima e l'altra esas fora. Ma, en la praktiko, on indikas\nl
ordinare omna enti od objekti per \VUgras{ta, ita}, kande on ne volas\nl
tote partikulare indikar la proximeso.
\end{VUpage}

% --- Feuillet 41 (folio 37) ----------------------------------------
%  Filet releve a la main : y = 1035, folio y = 134
%  -> 18,71 + (1035-134) x 25,4/300 = 94,99 mm.
%  Deux lignes courtes — « misas. » et « ici, iti. » — passaient sous le
%  seuil de detection ; le rattrapage des lignes pales les retrouve, et
%  le fac-simile est bien releve a 43 lignes.
\begin{VUpage}[41]{37}
\VUnotes{93.78mm}{%
(1) La demonstrativi propozesis en ta formi : \VUgras{icu, ecu, ocu, oco}\nl
(inspirita da \textit{hic, haec, hoc} sen \textit{h}); \VUgras{istu, estu, ostu, osto} (inspirita\nl
da \textit{iste, ista istud} Latina) e simetra kun \textit{ilu, elu, olu}). Kom adjektivi\nl
li esis : \VUgras{ica, ista} (\textit{ca, sta}).

La Franca substantivo « os » esis « oso » e la pluralo dil pro\cc
nomi : \textit{ici, eci, oci; isti, esti, osti}. Ma, judikante, ke ta sistemo esas\nl
komplikita e ne sate dicernebla (\textit{ica, ista}), la permananta Komisitaro\nl
quan elektabis la Komitato dil \textit{Delegitaro}, preferis nur : \VUgras{(i)ca, (i)ta}\nl
per la sequanta decido (20 di febr. 1908) :

« L'indikanta pronomi esas : por proximeso : \VUgras{(i)ca}; por foreso :\nl
\VUgras{(i)ta}. Neutra formo (por objekto nedeterminita) : \VUgras{ico, ito}. La komen\cc
canta \textit{i} darfas esar supresata, se l'eufonio permisas. Kande on\nl
bezonos indikar la genro, on prefixigos \VUgras{il, el, ol} al formo sengenra\nl
\VUgras{ca, ci, ta, ti}. Tale on obtenos: \VUgras{ilca, elca, olca; ilci, elci, olci; ilta}, e. c. »

Pose, pro ke \VUgras{ta} (\textit{ita}) esas plu facile pronuncebla kam \VUgras{ca} (\textit{ica}) da\nl
multa homi, nome dal Anglalinguani, kande ol komencas frazo\cc
membro, o sequas konsonanto, decidesis, ke \VUgras{ta} (\VUgras{ita}) sempre uzesos\nl
kom demonstrativo, ecepte kande on volos tote explicite indikar la\nl
proximeso; lore, kompreneble on uzos : \VUgras{ica, ca; ici, ci; ico, co}.

(En la kinesma Apendico trovesas la justifiko dil formi selektita\nl
en Ido por ta adjektivi-pronomi.)

(2) « Kozo » havas hike la tota ampleso di lua intima e natu\cc
rala signifiko : ol nule restriktesas ad \textit{objekto} determinita e mate\cc
riala. Exemple, kande me dicas : \VUgras{quo} \textit{donas a vu ta espero?} o : \VUgras{quo}%
}

On uzas prefere la formo plu kurta, kande l'eufonio per\cc
misas.

%% aucun blanc au fac-simile
La \textit{pronomi demonstrativa} (uzata sole) havas la sama\nl
formo kam l'adjektivi; ma plurale li divenas \VUgras{ci}, \VUgras{ti}, od\nl
\VUgras{ici, iti.}

%% aucun blanc au fac-simile
Kande on bezonas distingar la genro, on prefixigas al\nl
formi plu kurta (\VUgras{ca, ta}; \VUgras{ci, ti}) la personal pronomi \VUgras{il, el, ol} :

\VUblanc{3.87pt}% mesure au fac-simile
\VUgras{ilca} por maskulo, \VUgras{elca} por femino, \VUgras{olca} por neutro;

\VUgras{ilci} por maskuli, \VUgras{elci} por femini, \VUgras{olci} por neutri;

\VUgras{ilta} por maskulo, \VUgras{elta} por femino, \VUgras{olta} por neutro;

\VUgras{ilti} por maskuli; \VUgras{elti} por femini, \VUgras{olti} por neutri.

\VUblanc{3.80pt}% mesure au fac-simile
Fine, ta du pronomi havas \textit{neutra formo nedeterminita} :\nl
\VUgras{co} od \VUgras{ico}, \VUgras{to} od \VUgras{ito}. Ex. : \VUgras{me havas du hundi; ica esas}\nl
\VUgras{malada, e co jenas me; ma ita esas nefatigebla, e to esas}\nl
\VUgras{tre utila ad me por chasar en (i)ca lando} (1).

%% aucun blanc au fac-simile
On vidas la difero inter la neutra determinita \VUgras{olca, olta,}\nl
qua relatas ula objekto definita, e la neutra \VUgras{(i)co, (i)to}\nl
nedeterminita, qua relatas irga kozo (2) o fakto. Por dicer\cc
\end{VUpage}

% --- Feuillet 42 (folio 38) ----------------------------------------
%  Filet releve a la main : y = 1208, folio y = 134
%  -> 18,71 + (1208-134) x 25,4/300 = 109,64 mm.
%  Le folio de cette page se reduit aux deux chiffres : les cadratins
%  qui l'encadrent sont plats et larges, donc rejetes des le masque de
%  glyphes, et « 38 » soude ne fait que 60 px. Il fallait le masque de
%  secours de folio_ligne pour le retrouver.
%  Blancs du titre : 123,0 px au-dessus, 102,4 px au-dessous.
\begin{VUpage}[42]{38}
\VUnotes{108.46mm}{%
% Cette note continue celle du folio 37 : premiere ligne a la marge.
\VUcontinue \textit{eventis a vu?} me vizas nula objekto determinita : multa, diversa\nl
konsideri, reflekti povas kontenesar en la « quo » donanta espero,\nl
e plura od un sola fakto povas kontenesar en la « quo » dil ques\cc
tiono : \VUgras{quo} \textit{eventis a vu?} Do en ta frazi parolesas pri kozo, kozi,\nl
fakto o fakti tote nedeterminita e reprezentata da « quo ». Ta noto\nl
utilesas anke por komprenar juste la signifiko ed uzo di \VUgras{(i)co, (i)to}\nl
en multa kazi.

(1) \VUgras{Qua} koncernas nur l'individueso. Se on parolas pri la qualeso,\nl
kompreneble on uzas \VUgras{quala}, advere adjektivo qualifikanta nature,\nl
ma \VUgras{qua} tote darfas uzesar, e tre konvene, por questionar « quala »\nl
esas persono, animalo od objekto : \VUgras{quala il esas? kad il esas olda o}\nl
\VUgras{yuna, afabla od acerba?} Me vidis lu apene, do me ne povas dicar\nl
\VUgras{quala il esas.}

(2) Se on dicus : \VUgras{qua vu vidas}, on dubus, kad \VUgras{qua} vidas \VUgras{vu} o\nl
kad \VUgras{vu} vidas \VUgras{qua}. Pro to, kande \VUgras{qua} o \VUgras{qui} esas komplementi direta\nl
e preiras la subjekto, on adjuntas \VUgras{n} ad oli. Pro ke ta \VUgras{n} quik\nl
indikas l'inversigo dil subjekto (qua sequas vice preirar la kom\cc
plemento) ni nomizas lu \VUgras{n inversigala}. Ni studios plu ample ta\nl
punto en la sintaxo.%
}

\VUcontinue nar li praktike, suficas questionar su, kad on darfas\nl
adjuntar al pronomo ula substantivo determinita : \VUgras{ica} =\nl
\textit{ica hundo} (en l'exemplo supera); ma \VUgras{co} e \VUgras{to} = \textit{ca kozo, ca}\nl
\textit{fakto, ta kozo, ta fakto.}

Kom antecedenti di la pronomi relativa on uzas la formi\nl
en \textit{t} : \VUgras{(i)ta, (i)ti, (i)to} (V. § 35).

\VUsaut{6.85mm}
\VUtitre{10.15pt}{-6}{RELATIVA E QUESTIONALA ADJEKTIVI-PRONOMI.}
\VUsaut{6.00mm}

35. --- Kom \textit{pronomi} li esas : \VUgras{qua} (singulare), \VUgras{qui} (plu\cc
rale), \VUgras{quo} (kozo, fakto).

On uzas \VUgras{quo} nur kande parolesas pri kozo nedetermi\cc
nita; on uzas \VUgras{qua} pri kozo determinita. Ex. : \VUgras{Me renkontris}\nl
\VUgras{povra oldulo, qua demandis de me almono.}

Kom \VUgras{adjektivo} questionala, \VUgras{qua} esas nevariebla, segun la\nl
regulo generala dil adjektivi : \VUgras{qua homo venis? qua homi}\nl
\VUgras{venis?} (1).

Mem kom pronomo, lu restas nevarianta, se lu esas sub\cc
jekto : \VUgras{qua venis?} Ma, se lu divenas direta komplemento \textit{e}\nl
\textit{preiras la subjekto}, lu recevas \VUgras{n} finala, por preventar omna\nl
miskompreno : \VUgras{quan vu vidis?} En ta frazo, \textit{qua}(\textit{n}) esas kom\cc
plemento direta ed ol preiras la subjekto \VUgras{vu} (2).
\end{VUpage}

% --- Feuillet 43 (folio 39) ----------------------------------------
%  Filet releve a la main : y = 1711, folio y = 147
%  -> 18,71 + (1711-147) x 25,4/300 = 151,13 mm.
%  Blancs du titre : 123,4 px au-dessus, 94,4 px au-dessous.
\begin{VUpage}[43]{39}
\VUnotes{150.03mm}{%
(1) \VUgras{Quo} reprezentas kozo ne determinita o fakto. Do lu ne povas\nl
havar pluralo, same kam \VUgras{ico, ito} qui tre ofte preiras lu kom ante\cc
cedenti. (Videz noto 2, pag. 37.)

(2) \textit{Progreso}, VI, 238.%
}

Kom \textit{pronomo relativa}, \VUgras{qua} recevas la nombro di lua\nl
antecedento e la \VUgras{n} inversigala, se ol preiras la subjekto :\nl
\VUgras{la homi qui venis; la homo quan me vidas; la homi quin}\nl
\VUgras{me vidas.}

Same \VUgras{quo} recevas la \VUgras{n} inversigala se, esante komplemento\nl
direta, ol preiras la subjekto : \VUgras{Quo falis? Quon vu vidas?}\nl
\VUgras{To, quo eventas, dezolas me. To, quon il naracis, interesis}\nl
\VUgras{ni altagrade} (1).

\VUgras{Ita, ta, iti, ti} ed \VUgras{ilta, ilti; elta, elti; olta, olti} uzesas kom\nl
antecedenti di \VUgras{qua, qui}. Ex. : \VUgras{ta qua volas, ti qui volas, ilta}\nl
\VUgras{qua volas, ilti qui volas, elta qua volas,} e. c.

On prefixigas \VUgras{il, el, ol} a \VUgras{qua, qui} kande to esas necesa\nl
por precizigar, ed aparte por impedar nejusta refero : \VUgras{Ilta}\nl
\VUgras{qua volas mariajar su} ed \VUgras{elta qua volas mariajar su} esas\nl
plu preciza kam \VUgras{(i)-ta qua volas mariajar su.} --- \VUgras{La matro}\nl
\VUgras{di mea kuzulo, ad ilqua me parolis; la matro di mea kuzulo,}\nl
\VUgras{ad elqua me parolis.} Kun \VUgras{qua} sen \textit{il} e \textit{el} on ne savus, ka\nl
me parolis al kuzulo od al matro.

\VUsaut{6.14mm}
\VUtitre{10.15pt}{-38}{PRONOMO « LO ».}
\VUsaut{4.55mm}

36. --- « \VUgras{Lo} » nule esas pronomo « neutra », aplikebla\nl
a determinata kozi o sengenra enti; ica rolo apartenas ad\nl
\textit{olu}, e nur ad olu (2). Cetere, yen la texto ipsa di la pro\cc
pozo qua determinis la decido 948 dil Akademio : « De\nl
Beaufront e Couturat propozas... adoptar \VUgras{lo} kom pronomo\nl
ed artiklo indikanta kozo nedeterminita analoge a \VUgras{co, to} ».\nl
Sequas l'expozo dil motivi :

« Semblas a ni necesa adoptar \textit{lo} por indikar kozo tote\nl
nedeterminita (abstraktita, quale on dicas ofte nejuste);\nl
ne nur por expresar \textit{la belajo, la verajo} (triviala argumento),\nl
ma por mult altra kazi simila ad ici : « \textit{Lo} grava en ica\nl
afero... Me deziras lo maxim bona. » En altra frazo-formo\nl
on dicus : « \textit{To quo} esas grava... \textit{to quo} esas maxim bona. »
\parplein
\end{VUpage}

% --- Feuillet 44 (folio 40) ----------------------------------------
%  Filet releve a la main : y = 1631, folio y = 142
%  -> 18,71 + (1631-142) x 25,4/300 = 144,78 mm.
%  Blancs du titre : 148,0 px au-dessus, 99,3 px au-dessous.
\begin{VUpage}[44]{40}
\VUnotes{144.27mm}{%
(1) \textit{Progreso}, VI, 238-239. --- « Lo » esas nature \textit{pronomo}. Mem\nl
en \textit{lo bona, lo bela, lo vera, lo yusta} e. c., ol esas reale pronomo, ne\nl
artiklo, nam ol signifikas \textit{to quo}. Ex. : \VUgras{juntar lo agreabla a lo}\nl
\VUgras{utila} = \textit{juntar to quo esas agreabla a to quo esas utila.}

(2) \VUgras{Tala, quala, sama,} qui pozesis en l'antea listo (Franca) di\nl
la « Grammaire Complète », inter l'adjektivi-pronomi nedefinita,
% La note se poursuit au folio suivant : derniere ligne justifiee
% (1082 px releves pour 1083 de justification).
\parplein
}

\VUcontinue Do \textit{lo} esas quaze abreviuro di \textit{to quo}, e lua formo esas tote\nl
analoga, do necesa por la simetreso. Pluse, \textit{lo} esus utila\nl
por tradukar precize D. \textit{es}, kande ol referas, ne determinita\nl
objekto (hazarde neutra) ma integra frazo, t. e. fakto\nl
(segun la koncepto di D\textsuperscript{ro} \textsc{Talmey} : II, 148). On povas\nl
ya uzar \textit{co} e \textit{to}, ma ca vorti implikas demonstrativa\nl
nuanco, qua esas superflua. Exemple : « Prenez to, me\nl
volas lo. » \textit{Me volas to} esus dusenca; \textit{me volas ol} semblas\nl
referar objekto, do esus anke dusenca; on ne volas l'objekto\nl
prenenda, ma « ke vu prenez to », la preno ipsa.

« On bone remarkez, ke ni ne propozas \textit{lo} vice \textit{ol}, ma\nl
apud ed exter \textit{ol}, exakte same kam ni havas \textit{ico} apud \textit{olca},\nl
\textit{quo} apud \textit{olqua}. » (1).

Ni donez plusa exempli :

\VUgras{Lo facenda postulos longa tempo e multa lukti.} --- \VUgras{Me}\nl
\VUgras{esforcis omnamaniere por evitar lo neremediebla.} --- \VUgras{Lo}\nl
\VUgras{obtenita esas quaze nulo kompare a lo obtenenda.}

\VUgras{Il esas mortinta de tri monati, e vu ne savas lo!} (ke il\nl
esas mortinta). --- \VUgras{Restez e repozez me volas lo} (ke vu\nl
restez, e. c.) --- \VUgras{La rural domo di nia vicini esis incendiata.}\nl
\VUgras{On informis me pri lo} (ke ol esis incendiata).

Se on ne uzus « lo » en la tri unesma exempli, on mustus\nl
uzar la perifrazo : \textit{to quo esas}... E se, vice « lo », on uzus\nl
\VUgras{co, to,} en la du lasta, on obtenus altra nuanco, pro la signi\cc
fiko demonstrativa di ca pronomi.

\VUsaut{8.11mm}
\VUtitre{10.15pt}{-9}{ADJEKTIVI-PRONOMI NEDEFINITA.}
\VUsaut{5.10mm}

37. --- L'\textit{adjektivi-pronomi nedefinita} esas : \VUgras{ula, nula,}\nl
\VUgras{irga, altra, kelka, singla, omna, multa, poka, plura, tanta,}\nl
\VUgras{quanta, cetera, ipsa} (2).
\end{VUpage}

% --- Feuillet 45 (folio 41) ----------------------------------------
%  Filet releve a la main : y = 1174, folio y = 141
%  -> 18,71 + (1174-141) x 25,4/300 = 106,17 mm.
%  Page a blanc d'alinea ordinaire : sept frontieres a 7,4-9,1 px.
\begin{VUpage}[45]{41}
\VUnotes{104.48mm}{%
% Cette note continue celle du folio 40 : premiere ligne a la marge.
\VUcontinue esas plu reale \textit{qualifikanta}. Ni do eliminis li ek la listo dil adjektivi\nl
nedefinita en ica edituro.

\VUgras{Sama} relatas e la qualeso e l'individueso. Ex. : \textit{La sama qualeso} =\nl
la qualeso nediferanta. \textit{La sama homo} = la homo ne altra. Se on\nl
reflektas, on komprenas, ke homo povas divenar diferanta, ma lu\nl
ne povas divenar altra, \textit{altru}, to esas : perdar sua individueso. ---\nl
\VUgras{La samo} = la sama homo; \VUgras{lo sama} = la sama kozo, fakto. Ex. :\nl
\textit{Al samo eventis lo sama} = al sama homo eventis la sama kozo,\nl
aventuro, fakto e. c.

(1) \textit{Individual pronomi} (\VUgras{-u}) havas senco multe plu ampla kam\nl
\textit{personal pronomi}. Ci lasta implikas nur la persono (gramatikala),\nl
kontre ke l'unesmi implikas l'individuo. Or l'individuo kontenas\nl
omna ento organizita, persono, animalo o planto, relate lia speco.\nl
Konseque, se la pronomi en \VUgras{-u} reprezentas ordinare persono, homo,\nl
li tote darfas, ye l'okaziono, reprezentar anke bestio o planto. Ex. :\nl
\textit{omna individui di la hundo-speco ne prezentas la sama tipo. Li}\nl
\textit{diferas inter su, plu multe kam l'individui dil kavalo-speco.}

(2) Videz la noto 2 dil pagino 40.

(3) \VUgras{Kelka} indikas nombro min granda kam \VUgras{multa} e plu granda\nl
kam \VUgras{plura}, qua povas indikar mem nur du, okazione.%
}

Kande to esas posibla, li produktas individual pronomi,\nl
quale sube, per la simpla substituco di \VUgras{-u} (pluralo \VUgras{-i}) a la\nl
dezinenco \VUgras{-a} (1).

\VUblancAlinea
\VUgras{Ula} (L. \textit{ullus}), pronomo : \VUgras{ulu} (\textit{ula individuo}), plur. \VUgras{uli,}\nl
donas indiko nepreciza, nedefinita.

\VUblancAlinea
\VUgras{Nula} (L. \textit{nullus}), pronomo : \VUgras{nulu} (\textit{ne mem un indi}\cc
\textit{viduo}), plur. \VUgras{nuli}, esas la negativo di \VUgras{ula.}

\VUblancAlinea
\VUgras{Irga} (D. \textit{irgend ein}), pronomo : \VUgras{irgu}, plur. \VUgras{irgi}, donas\nl
l'indiko maxim ne preciza, maxim nedefinita.

\VUblancAlinea
\VUgras{Altra} (E. F. I. S.), pronomo : \VUgras{altru}, plur. \VUgras{altri}, indikas\nl
chanjo pri l'individueso, kontre ke « diferanta » chanjas\nl
la qualeso (2).

\VUblancAlinea
\VUgras{Kelka} (F.), pronomo : \VUgras{kelku}, plur. \VUgras{kelki}, donas indiko\nl
ne preciza pri mikra quanteso o nombro : \VUgras{kelka pano; yen}\nl
\VUgras{fragi; prenez kelki} (de oli) (3).

\VUblancAlinea
\VUgras{Singla} (L. \textit{singuli}, E. I.), pronomo : \VUgras{singlu}, plur. \VUgras{singli,}\nl
indikas konsiderante la unaji aparte. Ol havas senco distri\cc
butiva e signifikas : \textit{un po un, unope}, sive on uzas lu sin\cc
gulare, sive on uzas lu plurale. Ex. : \VUgras{singla soldato} o soldati\nl
recevis duopla porciono. --- \VUgras{Il parolis a singlu.}

\VUblancAlinea
\VUgras{Omna} (L. \textit{omnis}; omnibus), pronomo : \VUgras{omnu}, plur. \VUgras{omni,}
\parplein
\end{VUpage}

% --- Feuillet 46 (folio 42) ----------------------------------------
%  Filet releve a la main : y = 1334, folio y = 142
%  -> 18,71 + (1334-142) x 25,4/300 = 119,63 mm.
%  ATTENTION : la ligne « en -u. » (y ~ 1003, 90 px de large) passe sous
%  le seuil et sous le critere d'encre du rattrapage (0,07 contre 0,12
%  exiges). Le fac-simile compte 41 lignes, le detecteur n'en voit que
%  40 : meme limite qu'au feuillet 37. Le releve ci-dessous les donne
%  toutes ; les controles 9 et 11 sauteront la page.
\begin{VUpage}[46]{42}
\VUnotes{119.63mm}{%
(1) En l'expresuri : ni \textit{omna}, vi \textit{omna}, li \textit{omna}, ili \textit{omna}, eli\nl
\textit{omna}, la vorto « omna » esas duesma. En : ti \textit{sama}, ti \textit{altra}, ti \textit{kelka}\nl
qui pretendas, la vorti « \textit{sama}, \textit{altra}, \textit{kelka} » esas anke duesma.\nl
Konseque on ne vidas pro quo « omna » esus unesma en : \textit{omna} ti\nl
qui pretendas o : \textit{omni} ti qui pretendas.

Simile en : to \textit{poka} quon vu povas, agez lo por me, la vorto \textit{poka}\nl
esas duesma. Pro quo do \textit{omna}, od \textit{omno} esus unesma en : \textit{omna} to\nl
quon vu povas, \textit{omno} to quon vu povas, agez lo por me? Pro quo\nl
vere ica e nura ecepto pri la plaso koncerne \textit{omna}, \textit{omno?} Esperanto\nl
esas tante min imitinda pri ca punto, ke dicinte \textit{chiuj tiuj (kiuj)}, ol\nl
kontredicas su ipsa, dicante : \textit{tiuj kelkaj (kiuj)}.

Ni do pozez \textit{omna}, \textit{omno} ye la duesma plaso, dicante : ti \textit{omna}\nl
(qui o quin), to \textit{omna} (quo o quon). Ex. : \VUgras{ti omna qui vidis lu; ti}\nl
\VUgras{omna quin lu vidis.} --- \VUgras{To omna quo plezos a vu. Il prenis to omna}\nl
\VUgras{quon il volis.} --- \VUgras{Co omna eventis ante to, quon on rakontis a vu.}%
}

\VUcontinue indikas konsiderante l'ensemblo. Ol havas senco kolektiva,\nl
sive on uzas lu singulare, sive on uzas lu plurale. Ex. :\nl
\VUgras{Omna homo} o \VUgras{omna homi esas mortiva}, signifikas egale :\nl
\VUgras{la homi senecepte esas mortiva} (1).

\VUblancAlinea
\VUgras{Multa} (L. \textit{multus}, E. F. I. S.), pronomo : \VUgras{multi}, indikas\nl
kun ideo di granda quanteso o nombro. Ex. : \VUgras{multa vino;}\nl
\VUgras{multa homi.} Pro lua signifiko, ta adjektivo ne povas genitar\nl
pronomo en \VUgras{-u}. Un individuo ya ne povas esar « multu »,\nl
kontre ke lu povas esar « kelku » e « singlu ».

\VUblancAlinea
\VUgras{Poka} (L. \textit{paucum}, I. \textit{poco}), pronomo : \VUgras{poki}. Ol esas la\nl
kontreajo di \VUgras{multa} e konseque indikas kun ideo di mikra\nl
quanteso o nombro. Ex. : \VUgras{poka vino; poka homi.} Pro lua\nl
signifiko, \VUgras{poka} ne povas genitar pronomo singulara en \VUgras{-u},\nl
nam ica signifikus : individuo en mikra quanteso, poka\nl
individuo : senco absurda.

\VUblancAlinea
\VUgras{Plura} (L. \VUgras{Plures}, D. E. F. I. S.), pronomo : \VUgras{pluri} (\textit{plura}\nl
\textit{individui}) indikas kun ideo di plureso (adminime du). Pro\nl
lua signifiko, ol kompreneble ne povas genitar pronomo\nl
en \VUgras{-u}.

\VUblancAlinea
\VUgras{Tanta} (L. \textit{tantus}, F. I. S.), pronomo : \VUgras{tanti}, indikas kun\nl
ideo ocilanta inter granda nombro e granda forteso. Ex. :\nl
\VUgras{tanta esforci; tanta laboro.} Ica valoro interdiktas « tantu ».

\VUblancAlinea
\VUgras{Quanta} (L. \textit{quantus}, D. E. F. I. S.), pronomo : \VUgras{quanti}.\nl
Ol esas korelativo di \VUgras{tanta}. Ol indikas kun ideo pri la stando\nl
grandesala o nombrala dil indikato. Ex. : \VUgras{Il havas tanta}
\parplein
\end{VUpage}

% --- Feuillet 47 (folio 43) ----------------------------------------
%  Filet releve a la main : y = 1054, folio y = 145
%  -> 18,71 + (1054-145) x 25,4/300 = 95,67 mm.
\begin{VUpage}[47]{43}
\VUnotes{94.13mm}{%
(1) « Kozo » havas hike la tota ampleso di lua intima e naturala\nl
signifiko; ol nule restriktesas ad objekto determinita e materiala.\nl
(Videz la noto 2, pag 37.)

(2) Ni povas explikar mem \VUgras{multo, poko, tanto, quanto} per: \textit{multa}\nl
(abundanta) \textit{kozo, poka kozo, tanta kozo, quanta kozo, la cetera}\nl
\textit{kozo}. Kad ni ne dicas : \VUgras{multa} (abundanta) \VUgras{pano, poka pano, quanta}\nl
\VUgras{pano, tanta pano, la cetera pano?}

Pro ke \VUgras{irgo, kelko} e. c. indikas, quale ni dicis, ulo tam ne definita\nl
o plu juste, ulo tam ne determinita (mem spece) kam on povas\nl
supozar, on uzas kun li la relativo \VUgras{quo} (ne \textit{qua}) same kam kun \VUgras{to}\nl
e \VUgras{co}. Ta omna vorti ya implikas nedefiniteso (nedetermineso) kom\cc
pleta. Kande me dicas : \VUgras{la homo quan me odias}, on savas, ke me\nl
odias homo, on konocas l'objekto di mea odio. Ma, se me dicas :\nl
\VUgras{to quon me odias} o : \VUgras{to omna quon me odias}, on ne konocas l'objekto\nl
di mea odio; ol restas tote nedefinita. Esas do logikala, ke nede\cc
finita antecedento (\textit{co, to, ulo, nulo, omno}, e. c.) postulas relativo\nl
nedefinita : \VUgras{quo} e ne la relativo definita : \VUgras{qua}. L'atencoza lektero\nl
vidas, ke en nia expliko, la vorti « nedefinita », « nederminita »\nl
koncernas, ne la stando gramatikala, ma la signifiko ipsa dil vorto.

On lektas en \textit{Progreso}, IV, 651 : « \textit{Quo} nule equivalas \textit{olqua} :\nl
\textit{ol-qua} koncernas objekto di speco determinata o konocata (ex. :\nl
\textit{domo, libro, chapelo}, ...); \textit{quo} koncernas objekto nedeterminita o di\nl
speco nedefinita. La formi \textit{co, to, quo} \VUgras{ne esas propre neutra, ma exter}\nl
\VUgras{la tri genri}; li esas vera substantivi, ed expresas altra distingo kam%
}

\VUcontinue \VUgras{enemiki quanta amiki. Quanta penon to kustis!} Same kam\nl
pri \VUgras{tanta}, singularo en \VUgras{-u} ne esas posibla por \VUgras{quanta.}

\VUblancAlinea
\VUgras{Cetera} (la) (L. \textit{caetera}), pronomo : (la) \VUgras{ceteri}. Konocata\nl
da omni, adminime per \textit{et caetera}, ol indikas la lasta parto\nl
restanta di la enti o kozi quin on mencionas. Ex. : \VUgras{de mea}\nl
\VUgras{kin amiki quar livis me; la cetera restis e vartis mea risa}\cc
\VUgras{nesko.} Quale montras ta exemplo, la \textit{ceteru}, ne uzata, esus\nl
neutila.

\VUblanc{3.56pt}% mesure au fac-simile
\VUgras{Ipsa} (L. \textit{ipse}) juntesas a nomo o pronomo por indikar,\nl
ke lu agis, agas, agos od agus sen ul mediaco. Ex. : \VUgras{mea}\nl
\VUgras{kuzino ipsa responsos, se la afero ne sucesos.}

\VUblanc{5.71pt}% mesure au fac-simile
38. --- Per substituco dil vokalo \VUgras{-o} al vokalo \VUgras{-a}, ta adjek\cc
tivi produktas (segun lia signifiko) pronomi kozala o quan\cc
tesala. Ma, pro lia naturo esence singulara, li nultempe\nl
recevas la marko (\VUgras{-i}) dil pluralo :

\VUblancAlinea
\VUgras{ulo}, ula kozo; \VUgras{nulo}, nula kozo; \VUgras{irgo}, irga kozo; \VUgras{kelko,}\nl
kelka kozo; \VUgras{singlo}, singla kozo; \VUgras{omno}, omna kozo (1);\nl
\VUgras{multo, poko, tanto, quanto} (2). Ex. : \VUgras{Ta afero ne produktis}
\parplein
\end{VUpage}

% --- Feuillet 48 (folio 44) ----------------------------------------
%  Filet releve a la main : y = 1289, folio y = 153
%  -> 18,71 + (1289-153) x 25,4/300 = 114,88 mm.
%  Le folio de cette page est traverse par un pli du papier : le trait
%  diagonal portait l'etendue de sa bande a 396 px, centree a 0,703, et
%  folio_ligne le rejetait comme trop large. Le detecteur examine
%  desormais les AMAS de la bande, non son etendue totale.
\begin{VUpage}[48]{44}
\VUnotes{114.88mm}{%
\VUcontinue la genri, plu importanta kam la distingo di la genri. Kande on dicas:\nl
« quon vu serchas? », on nule pensas : « \textit{ol-quan}, t. e. qua libron,\nl
qua chapelon », ma : « qua kozon », tote generale. »

(1) \VUgras{Pluro} (de plura) ne semblis formacenda, pos explorado e\nl
probi.

« La distingo di la personi (per \VUgras{-u}) e di la kozi (per \VUgras{-o}) esas\nl
utila nur en la adjektivi, e ne en la substantivi; nam la senco\nl
ipsa di substantivo indikas persono o kozo, ed adjektivo sempre\nl
referas substantivo expresita o tacita. Cetere, ta distingo esas nule\nl
logikal, ma pure gramatikala : la logiko ne distingas personi e kozi,\nl
omni esas logike « objekti » od « enti ». Pro to la finalo \VUgras{o} di la\nl
substantivo ne signifikas plu kozo kam persono, ma simple « objek\cc
to » od « ento »... Ma, \textit{praktike}, en la maxim multa kazi, la senco\nl
ipsa dil adjektivo indikas a la « komuna raciono », ke ta ento esas\nl
persono o kozo : \textit{avaro, blindo, klaudikero} (unesme \textit{lamo}), \textit{strabo}\nl
povas esar nur personi; \textit{vakuo, dezerto, acido}, povas esar nur kozi. »\nl
(\textit{Progreso}, I, 555.)%
}

\VUcontinue \VUgras{multo; vere la quanto dil rezultaji ne kompensas la tanto}\nl
\VUgras{dil esforci.}

\VUblancAlinea
(La) \VUgras{cetero} = la lasta parto restanta dil kozo mencionata.

\VUblancAlinea
La naturo ipsa dil adjektivo-pronomo \VUgras{ipsa} ne plu per\cc
misas \textit{ipso} kam \VUgras{ipsu} (1).

\VUblanc{5.46pt}% mesure au fac-simile
39. --- Per substituco dil vokalo \VUgras{-e} al vokalo \VUgras{-a} la nede\cc
finita adjektivi produktas adverbi manierala o quantesala,\nl
segun sua signifiko :

\VUblancAlinea
\VUgras{Ule, nule, irge, altre} = en \textit{ula, nula, irga maniero}; ---\nl
\VUgras{cetere} = \textit{egardante la cetero}; --- \VUgras{ipse} = \textit{de su ipsa, per su}\nl
\textit{ipsa}; --- \VUgras{kelke} = \textit{en kelka grado o quanteso}; --- \VUgras{multe} =\nl
\textit{en alta grado o granda quanteso} : \VUgras{Il multe laboras, ma il}\nl
\VUgras{ganas multe.}

\VUblancAlinea
\VUgras{Poke} = \textit{en basa grado o mikra quanteso} : \VUgras{Il poke laboras,}\nl
\VUgras{ma il ganas poke.}

\VUblancAlinea
\VUgras{Tante} = \textit{en tanta grado o quanteso} : \VUgras{Il tante amoras elu!}\nl
\VUgras{Ni tante produktas ke ni ne povas vendar nia omna pro}\cc
\VUgras{dukturi.}

\VUblancAlinea
\VUgras{Quante} = \textit{en quanta grado o mezuro, ye qua preco} :\nl
\VUgras{Quante il amoras elu! Vu tante pagesos quante vu laboros.}\nl
\VUgras{Po quante vu vendas ico?}

\VUblancAlinea
\VUgras{Single} = \textit{singla, singlu aparte} : \VUgras{Cigari po 25 centimi}\nl
\VUgras{single.}
\end{VUpage}

% --- Feuillet 49 (folio 45) ----------------------------------------
%  Filet releve a la main : y = 1680, folio y = 154 -> 147.91 mm.
\begin{VUpage}[49]{45}
\VUnotes{149.78mm}{%
(1) Erore la vorto « omne » trovesas kun la signifiko « omna\cc
maniere ». Nur ica lasta adverbo kompozita esas justa por expresar\nl
la ideo « en omna maniero ». Nam, per sua naturo ipsa, la radiko\nl
\textit{omn} implikas ideo di nombro, quale \textit{singl} e \textit{plur}.%
}

% Cette ligne est un alinea a elle seule : le fac-simile la renfonce
% de 49 px, soit le renfoncement ordinaire. Ce n'est pas la suite de
% l'alinea du folio precedent.
\VUgras{Plure} = \textit{pluri kune.}

\VUblancAlinea
\VUgras{Omne} = \textit{omni kune} : \VUgras{Unesme li venis single, pose plure,}\nl
\VUgras{fine omne} (1).

\VUblanc{6.06pt}% mesure au fac-simile
40. --- La vorti \VUgras{irga, irgu (irgi), irgo, irge} devas esar\nl
uzata sole, kande li trovesas en propoziciono nedependanta\nl
e kompleta : \VUgras{donez a me irgo; irgu komprenos to; venez}\nl
\VUgras{irge, ma venez.}

\VUblancAlinea
Ma li devas esar sequata dal vorti \VUgras{qua, qui, quo, quale,}\nl
\VUgras{quante,} se li esas ligata kun subordinal propoziciono. Ex. :\nl
\VUgras{Ad irgu, qua venos, vu dicos, ke me ne esas en la hemo}\nl
\VUgras{(o heme);} --- \VUgras{irgo quon vu donacos ad ilu, to ne kon}\cc
\VUgras{tentigos lu;} --- \VUgras{irge quale vu procedos, il blamos vu;} ---\nl
\VUgras{irge quante li laboros, li ne povos finar tante balde.}

\VUblancAlinea
Komprenende la adverbo \VUgras{irge} darfas ligesar a la cetera\nl
relativi, exemple a \VUgras{qua, quala, quanta,} e. c. por donar a li\nl
la senco nedeterminita. Ex. : \VUgras{irge qua mulieron o yuninon}\nl
\VUgras{il vidas, il desprizas elu;} --- \VUgras{irge quala instrumenton vi uzos,}\nl
\VUgras{vi ne sucesos facar ta laboro;} --- \VUgras{irge quanta librin il havas,}\nl
\VUgras{il deziras ankore plu multi.}

\VUblancAlinea
N. B. --- Notez bone, ke la derival regulo, segun qua on\nl
obtenas \textit{ento} (praktike \textit{homo}, maxim ofte) substitucante \VUgras{-o}\nl
al \VUgras{-a} dil adjektivo, koncernas nur \textit{l'adjektivo qualifikanta},\nl
qua venas de nomal radiko : \VUgras{avara, avaro.} Or en la 38-ma\nl
paragrafo (p. 43), ni vidis, ne l'adjektivo qualifikanta, ma\nl
\textit{l'adjektivi-pronomi nedefinita}, qui nule venas de nomal ra\cc
diko e konstitucas special kategorio gramatikala. Pro to la\nl
regulo dil renversebleso ne atingas li, quale ol ne atingas\nl
la propra nomi, tote exter ta regulo. \textit{Altra} kam l'adjektivo\nl
qualifikanta, per sua naturo e rolo, l'adjektivi-pronomi nede\cc
finita sequas \textit{altra} e propra derivo. On do ne darfas obje\cc
cionar pri li regulo koncernanta nur l'adjektivo qualifikanta\nl
e la participo adjektiva.

(En la apendico 4-ma trovesas fundamental studiuro pri\nl
ca punto).
\end{VUpage}

% --- Feuillet 50 (folio 46) ----------------------------------------
%  Filet releve a la main : y = 940, folio y = 149 -> 85.68 mm.
\begin{VUpage}[50]{46}
\VUnotes{90.64mm}{%
(1) En gramatiko e pri verbi, la \textit{voco} esas l'ensemblo dil formi\nl
indikanta, kad la subjekto facas (voco aktiva) o subisas (voco\nl
pasiva), la ago expresata dal verbo.

(2) Remarkez ke en la tota konjugo, la vokali \textit{a}, \textit{i}, \textit{o} indikas\nl
rispektive : \textit{prezento, pasinto, futuro.}

(3) L'Idal \textit{preterito} expresas omna indikatival tempi pasinta qui\nl
ne esas antea pasinti : \textit{imparfait, passé défini, passé indéfini} di la\nl
Franca. Ica lasta linguo havas, pri la pasinta tempi dil indikativo,\nl
distingo qua ne korespondas a ta dil cetera lingui. Esabus eroro\nl
imitar lu en Ido, e tante plu ke ol ne esas reale necesa. Se stranjero\nl
dicas a Franco : « je \textit{voyais} votre mère hier, o : je \textit{vis}, vice j'\textit{ai vu}\nl
votre mère hier : me vidis vua matro hiere » kad ta Franco ne\nl
komprenas perfekte, ke la stranjero dicas \textit{vidir} lua matro? To kom\cc
prenigas, ke Ido povas kunfuzar la 3 tempi di la Franca en la tempo\nl
\VUgras{-is} (preterito) sen ula detrimento por l'ideo expresenda.

Tamen, kande to esas utila, la Franca « imparfait » (imperfekto)\nl
povas tradukesar per la preterito (\textit{esis}) di la verbo \textit{esar} e la\nl
participo prezenta dil uzata verbo : \textit{me esis manjanta, kande vu}\nl
\textit{arivis.}

(4) Pri \VUgras{s} en \textit{as, is, os, us} videz la 7-ma apendico : « Verbal rolo\nl
di \VUgras{s} en la helpolinguo. »

« L'akademio repulsis chanjar la tri tempi di l'indikativo : e ton\nl
lu facis bone. En nia konjugo la regulozeso dil tri vokali \textit{amas},\nl
\textit{amis}, \textit{amos} esas preferinda kam ula naturala formi (internaciona ne\nl
existas) pro ke per ta unika regulo on lernas anke la tri tempi di la\nl
cetera modi : \textit{amar, amir, amor; amanta, aminta, amonta; amata,}%
}

\VUsaut{5.25mm}
\VUtitre{10.15pt}{-32}{VERBO.}
\VUsaut{5.85mm}

41. --- Duspeca verbi existas en Ido : la \textit{transitivi} e la\nl
\textit{netransitivi.}

%% aucun blanc au fac-simile
La transitivi (qui povas havar « objekto » o komple\cc
mento direta) posedas du voci (1) : l'aktiva e la pasiva;\nl
la netransitivi posedas nur un voco, l'aktiva.

%% aucun blanc au fac-simile
L'aktiva voco esas formacata per la verbal radiko, a qua\nl
on adjuntas la sequanta dezinenci :

\VUblanc{2.85pt}% mesure au fac-simile
\VUgras{-ar} por l'infinitivo prezenta : \VUgras{am-ar.}

\VUgras{-ir} por l'infinitivo pasinta : \VUgras{am-ir.}

\VUgras{-or} por l'infinitivo futura : \VUgras{am-or} (2).

\VUgras{-as} por l'indikativo prezenta : \VUgras{am-as.}

\VUgras{-is} por l'indikativo pasinta : \VUgras{am-is} (3).

\VUgras{-os} por l'indikativo futura : \VUgras{am-os.}

\VUgras{-us} por la kondicionalo prezenta : me \VUgras{am-us} (4).
\end{VUpage}

% --- Feuillet 51 (folio 47) ----------------------------------------
%  Filet releve a la main : y = 906, folio y = 136 -> 83.90 mm.
\begin{VUpage}[51]{47}
\VUnotes{82.01mm}{%
\VUcontinue \textit{amita, amota}... Ta formi esas komoda e mem necesa. Pri la pasinta\nl
infinitivo \textit{amir} cadie on ne plus dubas; pri la futura formi \textit{amor},\nl
\textit{amonta, amota} on semblas ankore havar dubi, ma tote neyuste. Se\nl
me dicas : \textit{Me konjektas ke me vidos il} (t. e. en futuro), lore esas\nl
postulo di la logiko, ke on povez dicar : \textit{Me konjektas vidor il}, por\nl
expresar futuro kontraste kun la prezento : \textit{vidar il.} --- Se me dicas:\nl
« \textit{En la publikigota artiklo mencionez to} », me indikas, ke on certe\nl
publikigos artiklo, ed esas mala helpilo dicar, quale plura vivanta\nl
lingui : \textit{la publikigenda artiklo}, nam ne traktesas artiklo qua \textit{devas}\nl
esar publikigata, ma ula qua \textit{esos} publikigata. Esus regretinda\nl
kulpo imitar, sub pretexto di internacioneso, caregarde la nacionala\nl
lingui, qui ne havas tala formo, e koaktar ni dicar ulo altra kam\nl
ni pensas, ed expektar ke l'audanto divinos quon ni vizas : vere plu\nl
importanta kam l'internacioneso (propre dicite hike nur la nacio\cc
naleso) esas furnisar per nia linguo a la homi qui pensas logikale\nl
\textit{ula instrumento apta expresar lia idei klare e precize.} E ton nia\nl
konjugo kun lua tri tempi facas admirinde. Ni do konservez lua\nl
sistemo ». --- (P. \textsc{de Janko}, \textit{Progreso}, \textit{IV}, aprilo, 1911.)

(1) L'\textit{imperativo} (o, plu juste, \textit{volitivo}) \VUgras{-ez} atingas quik lekte\nl
irgu qua savas kelke la Franca : \textit{venez, dormez}, e. c., kontre ke \VUgras{-es}\nl
atingus nulu : \textit{venes, dormes.} Ma, se ulu ne povas pronuncar facile\nl
\textit{-ez} finala, lu pronuncez ol quale \textit{-es}, sen hezito. La avantajo di \VUgras{-ez}\nl
ne perdesos, same kam la avantajo di \textit{j} ne perdesas, se uli pronuncas\nl
ol quale \textit{dj}, segun explicita permiso di nia gramatiko. --- Esperantisti\nl
blamis \VUgras{-ez} por la volitivo, obliviante ke \VUgras{-u} dil sama modo en lia\nl
linguo komprenesas da nulu : \textit{venu, dormu}, e ritrovesas en \textit{iu, tiu,}\nl
\textit{chiu, neniu} (mem en \textit{chu, morgau, kontrau, lau, antau} e. c.) quo\nl
forsan ne esas kolmo di logiko.%
}

\VUgras{-ez} por l'imperativo prezenta : ni \VUgras{am-ez} (1).

\VUblanc{3.54pt}% mesure au fac-simile
La distingo dil personi esas indikata dal personal pro\cc
nomi o dal substantivo subjekta. Ma ye la duesma persono\nl
singulara o plurala dil imperativo, on darfas tacar la sub\cc
jekto : \VUgras{amez}, vice \textit{vu} o \textit{vi} amez.

%% aucun blanc au fac-simile
La \textit{participi aktiva} esas formacata per la sequanta sufixi,\nl
a qui on adjuntas la dezinenco (\textit{-a}) dil adjektivi :

\VUblanc{3.72pt}% mesure au fac-simile
\VUgras{-ant} por la participo prezenta : \VUgras{am-ant-a} (qua amas).

\VUgras{-int} por la participo pasinta : \VUgras{am-int-a} (qua amis).

\VUgras{-ont} por la participo futura : \VUgras{am-ont-a} (qua amos).

\VUblanc{3.97pt}% mesure au fac-simile
La verbo \VUgras{esar} ipsa recevas la diversa dezinenci quin ni jus\nl
vidis : \VUgras{es-ar, es-ir, es-or; es-as, es-is, es-os, es-ez; es-ant-a,}\nl
\VUgras{es-int-a, es-ont-a.}

\VUblanc{5.78pt}% mesure au fac-simile
42. --- L'\textit{antea tempi} di la voco aktiva esas formacata
\parplein
\end{VUpage}

% --- Feuillet 52 (folio 48) ----------------------------------------
%  Filet releve a la main : y = 1121, folio y = 130 -> 102.61 mm.
\begin{VUpage}[52]{48}
\VUnotes{101.49mm}{%
(1) La formi kun \VUgras{-ab-} unesme refuzita, pose admisita por probo,\nl
fine sancionita unanime dal Akademio (decido 748), sempre plu\nl
praktikesas en Ido kom plu kurta, plu lejera e plu facile docebla.

Memorez bone, ke la formo \VUgras{-abas} ne darfas uzesar, nam l'Akademio\nl
ne adoptis lu, pro ke ol ne esas necesa.

Praktike on ne havas l'okaziono dicar : \textit{esez aminta o amabez.}\nl
Ma on povas tre bone dicar : \VUgras{esez fininta o finabez} (ica laboro ante\nl
mea retroveno).

(2) \VUgras{Me esis departonta} signifikas : « me (lore) preparis me, o\nl
me esis ja pronta por departar quik », granda difero kun la simpla:\nl
me departos.

(3) \VUgras{Il ne esas mortinta} dicas, ke il ne \textit{esas en la stando} di homo\nl
qua \textit{mortis.} Simile : \textit{nun ni esas arivinta} dicas altro kam ni \textit{arivis} :\nl
ol montras la stando rezultanta de \textit{arivir.} Komparante : \VUgras{il ne esas}\nl
\VUgras{mortinta, nun ni esas arivinta} a : \textit{il ne mortis, ni arivis}, on\nl
vidas la granda difero existanta inter l'unesmi e la duesmi. Per ici\nl
on quaze konstatas ed enuncas la rezulto di iti : \textit{il ne esas mortinta,}\nl
pro ke il ne mortis, e \textit{ni esas arivinta}, pro ke ni arivis.

(En la 6-ma apendico on trovos la motivi gravega e decidigiva\nl
pro qui la verbo \textit{havar} nultempe servas kom helpanto en la Idala\nl
konjugo, malgre l'exemplo di lingui romanala e di E. D.)%
}

\VUcontinue per la verbo \VUgras{esar} kombinata kun la pasinta participo aktiva,\nl
o per \VUgras{-ab-} quan on pozas inter la radiko e la dezinenco :

\VUblanc{2.96pt}% mesure au fac-simile
\textit{antea pasinto} : \VUgras{me esis aminta} o \VUgras{me amabis.}

\textit{antea futuro} : \VUgras{me esos aminta} o \VUgras{me amabos.}

\textit{antea kondicionalo} : \VUgras{me esus aminta} o \VUgras{me amabus.}

\textit{antea volitivo} : \VUgras{esez fininta} o \VUgras{finabez} (1).

\VUblanc{2.80pt}% mesure au fac-simile
On darfas formacar altra tempi sekundara, sive en l'ak\cc
tivo, sive en la pasivo (quan ni vidos sube) per la kombino\nl
dil verbo \VUgras{esar} kun la cetera participi : \VUgras{il esis lektanta; me}\nl
\VUgras{esis departonta} (2); \VUgras{il ne esas mortinta} (3).

\VUblanc{5.84pt}% mesure au fac-simile
43. --- La \textit{participi pasiva} esas formacata per la sequanta\nl
sufixi, a qui on adjuntas la dezinenco (\textit{-a}) dil adjektivi :

\VUblanc{2.89pt}% mesure au fac-simile
\VUgras{-at} por la participo prezenta : \VUgras{am-at-a} (qua esas amata).

\VUgras{-it} por la participo pasinta : \VUgras{am-it-a} (qua esis amata).

\VUgras{-ot} por la participo futura : \VUgras{am-ot-a} (qua esos amata).

\VUblanc{2.87pt}% mesure au fac-simile
La \textit{voco pasiva} formacas lua chefa tempi per la verbo \VUgras{esar}\nl
sequata nemediate dal prezenta participo pasiva. Tale on\nl
obtenas la yena tempi e modi :

\VUblanc{2.73pt}% mesure au fac-simile
\textit{Indikativo prezenta} : \VUgras{me esas amata} (on amas me).
\end{VUpage}

% --- Feuillet 53 (folio 49) ----------------------------------------
%  Filet releve a la main : y = 763, folio y = 124 -> 72.81 mm.
%  La ligne « tita. » (courte) passe sous le seuil : le fac-simile compte
%  45 lignes, le detecteur 44. Meme limite qu'aux feuillets 37 et 46.
\begin{VUpage}[53]{49}
\VUnotes{72.81mm}{%
(1) La \textit{perfekto} propre dicata (qua expresas \textit{nuna stando} rezultanta\nl
de ago tote pasinta) logike tradukesas per : \VUgras{me esas amita} (olim).---\nl
L'exempli : \VUgras{Il ne esas mortinta; nun ni esas arivinta} di la noto\nl
3-ma, pag. 48, korespondas, en la voco aktiva, a \VUgras{me esas amita} di\nl
la voco pasiva. En ta omna exempli, l'uzo dil \textit{perfekto} expresas \textit{nuna}\nl
\textit{stando} pro ago \textit{facita} dal subjekto o \textit{subisita} da olu.

La pasinta participo pasiva (\VUgras{-ita}) indikas, ke la ago subisita dal\nl
subjekto esas nun tote finita : \VUgras{l'imprimo di ta libro duris tre longe;}\nl
\VUgras{ol esis imprimata dum sis monati ed erste de 8 dii ol esas imprimita.}\nl
--- \VUgras{Tam longe kam on serchas la solvo di problemo, ta problemo} \textit{esas}\nl
\textit{solvata;} \VUgras{ma venas instanto pos qua ol} \textit{esas solvita.} --- \VUgras{La domo esis}\nl
\VUgras{konstruktata dum un yaro; ma depos du monati ol esas konstruk}\cc
\VUgras{tita.}

« Por savar, kande vu hezitas, \textit{qua tempo devas uzesar en la pasivo,}\nl
\textit{tradukez per l'aktivo.} Kad on arestas e kondamnas nun persono?\nl
Se yes, lu esas \textit{arestato, kondamnato.} Se ne, se on ja arestis o kon\cc
damnis lu, on devas uzar la pasinto : \textit{arestito, kondamnito;} ne\nl
importas kad l'efekto di la ago ankore duras o ne. Mem liberigita,\nl
l'arestito restos arestito; mem indulgita, la kondamnito restos\nl
kondamnito (t. e. ulu, quan on kondamnis ultempe). Per ta facila\nl
procedo (qua divenas balde mashinala ed instintala per kustumo)\nl
\textit{on ne povas erorar.} « La pordo esas klozata » (on klozas ol nun);\nl
« la pordo esas klozita » (ol esas nun en la stando qua rezultas de\nl
ke on \textit{klozis} olu). --- Du participi preske sempre postulas la sama\nl
tempo : \textit{konocata} devas generale esar en la prezento, nam : « Me\nl
esas konocata » signifikas : « on konocas me » (same \textit{reputata},\nl
konsiderata, e. c.). --- Kontraste, \textit{fatigita, konvinkita} devas generale\nl
esar en la pasinto. Nam, kande on dicas : me esas fatigita, konvin\cc
kita, la efekto ne esas \textit{nun} produktata, ma \textit{ja} produktita. « Me fati\cc
gesis promenante e nun me esas fatigita. » --- « Me konvinkesis da\nl
omno, quon me audis o vidis, pro to me ne plus dubas; me esas\nl
konvinkita » (segun \textit{Progreso}, III, 396.)%
}

\textit{Indikativo pasinta} : \VUgras{me esis amata} (1) (on amis me).

\textit{Indikativo futura} : \VUgras{me esos amata} (on amos me).

\textit{Kondicionalo prezenta} : \VUgras{me esus amata} (on amus me).

\textit{Volitivo prezenta} : \VUgras{esez amata} (on amez vu).

\textit{Infinitivo prezenta} : \VUgras{esar amata.}

\textit{Infinitivo pasinta} : \VUgras{esir amata.}

\textit{Infinitivo futura} : \VUgras{esor amata.}

\VUblanc{2.71pt}% mesure au fac-simile
L'\textit{antea tempi} dil voco pasiva esas formacata per la sama\nl
helpo-verbo (esar) sequata dal \textit{pasinta} participo pasiva :

\VUblanc{2.64pt}% mesure au fac-simile
\textit{Antea pasinto} : \VUgras{me esis amita} (on amabis me).

\textit{Antea futuro} : \VUgras{me esos amita} (on amabos me).

\textit{Antea kondicionalo} : \VUgras{me esus amita} (on amabus me).
\end{VUpage}

% --- Feuillet 54 (folio 50) ----------------------------------------
%  Filet releve a la main : y = 1272, folio y = 135 -> 114.98 mm.
%  Les huit formes du § 44 sont alignees sur deux renfoncements releves :
%  119-122 px pour celles sans « me », 47-48 px pour les autres.
\begin{VUpage}[54]{50}
\VUnotes{112.54mm}{%
(1) Ultre lia funciono en la konjugo, la participi (aktiva o\nl
pasiva) pleas la rolo di adjektivi, sive kom epiteto : \textit{vazo ruptita,}\nl
sive kom atributo : \textit{il esas erudita.} Li anke povas formacar substan\cc
tivi : \textit{la mortanto, la vivanti.} Fine per \VUgras{-e} li povas divenar adverbi :\nl
\textit{mortante il dicis. Sendez omno afrankite.}

Irgu qua examenis sen prejudiko nia sis participi, en lia funcioni\nl
diversa, povas nur gratular ni havar per li, e sen kompliko, richeso\nl
tante utila por la facila e justa expresado dil pensi.

(2) En Espo la korespondanti esas ica charmanta formi : \textit{estis}\nl
\textit{estinta, estos estinta, estus estinta.} La frazo : \textit{kande ni esabis tre}\nl
\textit{obediera li rekompensis ni} tradukesas ad Espo per : \textit{kiam ni estis}\nl
\textit{estinta tre obediema ili rekompencis ni.} Esperanto regretinde tro\nl
ofte oblivias la oreli sentema.

(3) La lerno dil regulo : \textit{amesar} = \textit{esar amata} ne esas jenanta;\nl
tote kontree, ta formi esas vere tre komoda. (E. \textsc{Ferrand}, \textit{Progreso},\nl
IV, 210.) Til nun (aprilo 1911) nulu plendis pri reala dusenceso di\nl
\textit{ameso} e simila vorti. (\textit{Progreso}, IV, 90.)%
}

\textit{Antea volitivo} : (to) \VUgras{esez finita} (on finabez to) (1).

\VUblanc{2.75pt}% mesure au fac-simile
Ma per la sufixo \VUgras{-ab}, quan la verbo \VUgras{esar} recevas en :\nl
\VUgras{esabis} (antea pasinto), \VUgras{esabos} (antea futuro), \VUgras{esabus} (antea\nl
kondicionalo), \VUgras{esabez} (antea volitivo) (2) on tre bone povas\nl
indikar l'anteeso en la quar supera tempi, konservante la\nl
participo \VUgras{-ata} di la chefa tempi :

\VUblancAlinea
\VUrang{\VUcase{4.06mm}{\VUgras{me esabis amata} (antea pasinto).}}
\VUrang{\VUcase{4.06mm}{\VUgras{me esabos amata} (antea futuro).}}
\VUrang{\VUcase{4.06mm}{\VUgras{me esabus amata} (antea kondicionalo).}}
\VUrang{\VUcase{4.06mm}{(to) \VUgras{esabez finita} (pos du hori; --- antea volitivo).}}

\VUblancAlinea
Tale la tota konjugo pasiva postulas nur un sola parti\cc
cipo (\textit{-ata}), quo esas plu simpla e ne min preciza.

\VUblanc{5.55pt}% mesure au fac-simile
44. --- On povas obtenar por la voco pasiva \textit{sintezal}\nl
(exakte : \textit{kunjuntal}) formi plu kurta, kunjuntante, por\nl
abreviar, la verbal radiko e la verbo \VUgras{esar} :

\VUblancAlinea
\VUrang{\VUcase{10.33mm}{\VUgras{amesar} = (\textit{am}(\textit{at})\textit{esar}) esar amata.}}
\VUrang{\VUcase{10.25mm}{\VUgras{amesir} = (\textit{am}(\textit{at})\textit{esir}) esir amata.}}
\VUrang{\VUcase{10.25mm}{\VUgras{amesor} = (\textit{am}(\textit{at})\textit{esor}) esor amata.}}
\VUrang{\VUcase{4.06mm}{me \VUgras{amesas} = (\textit{am}(\textit{at})\textit{esas}) me esas amata.}}
\VUrang{\VUcase{4.06mm}{me \VUgras{amesis} = (\textit{am}(\textit{at})\textit{esis}) me esis amata.}}
\VUrang{\VUcase{3.98mm}{me \VUgras{amesos} = (\textit{am}(\textit{at})\textit{esos}) me esos amata.}}
\VUrang{\VUcase{3.98mm}{me \VUgras{amesus} = (\textit{am}(\textit{at})\textit{esus}) me esus amata.}}
\VUrang{\VUcase{10.08mm}{\VUgras{amesez} = (\textit{am}(\textit{at})\textit{esez}) esez amata (3).}}
\end{VUpage}

% --- Feuillet 55 (folio 51) ----------------------------------------
%  Filet releve a la main : y = 1642, folio y = 145 -> 145.46 mm.
\begin{VUpage}[55]{51}
\VUnotes{143.94mm}{%
(1) Por ta pasivi ni preferas l'expresuro « kunjuntal formi » pro\nl
ke li esas la \textit{kunjunturo} di « esar » e di verbal radiko. Ta formi\nl
esus \textit{vere} sintezala, nur se li prezentus dezinenci segun la modelo\nl
Latina : \textit{am-or, am-aris} e. c. Ma en \textit{am-esar, am-esas, am-esis} e. c. ne\nl
existas dezinenci : existas radiko e kompleta verbo (\VUgras{esar}) kunjun\cc
tita, \VUgras{unionita.}%
}

On darfas unionar al verbal radiko (en pasivo) mem la\nl
tempi \VUgras{esabis, esabos, esabus, esabez.} Ma ta uniono o kun\cc
junto esas deskonsilata por radiki \textit{plursilaba}, pro ke lore la\nl
vorto divenas longa ed ofte min klara kam en analizal formo.\nl
Exemple : \VUgras{interpretesabus} o mem \VUgras{destruktesabus}, kelke\nl
longa, ne havas la klareso di : \VUgras{esabus interpretata, esabus}\nl
\VUgras{destruktata.}

%% aucun blanc au fac-simile
Do, praktike, uzez \textit{esabis, esabos, esabus, esabez} nur kun\nl
monosilaba radiko.

%% aucun blanc au fac-simile
La kunjuntal formi (1), tre komoda pro lia kurteso (kun\nl
radiko monosilaba) precipue utilesas por tradukar la falsa\nl
verbi reflektiva (§ 49).

\VUblanc{5.62pt}% mesure au fac-simile
45. --- Malgre l'exemplo di lingui nacionala (qui agas\nl
pri la transitiveso dil verbi deskonkorde ed arbitrie) verbo\nl
esas \textit{transitiva} en Ido, se ol expresas faco, ago povanta\nl
atingar \textit{direte} ula objekto : \VUgras{dankar, mokar, nocar, obediar,}\nl
\VUgras{repugnar} \textit{ulu}, nam vere e logike \textit{ulu} povas esar direte l'ob\cc
jekto dil \textit{danko}, dil \textit{moko, noco, obedio, repugno.} Ico nekon\cc
testeble pruvesas dal fakto, ke l'objekto povas divenar la\nl
subjekto en voco pasiva : \textit{Dankez Deo, Deo esez dankata} o\nl
\textit{Deo dankesez; il mokis ta povrulo, ta povrulo esis mokata}\nl
\textit{o mokesis da ilu; el nocos vu, vu esos nocata o nocesos da}\nl
\textit{elu; fine la filio obediis la matro, la matro esis obediata dal}\nl
\textit{filio}; \VUgras{ta procedo repugnas me}, \textit{me esas repugnata o repu}\cc
\textit{gnesas da ta procedo.}

%% aucun blanc au fac-simile
Konseque, esas \textit{netransitiva} en Ido nur la verbi qui logike,\nl
nature, quale \VUgras{naskar} e \VUgras{mortar} ne povas havar « objekto »\nl
(komplemento direta). \VUgras{Restar, sejornar, repozar, kreskar,}\nl
\VUgras{dekadar, degenerar, falar, irar, marchar, venar, arivar, de}\cc
\VUgras{partar}, exemple, ne esas transitiva, ma tre certe \textit{netransitiva.}\nl
% Ces trois lignes sont en romain, SAUF les deux series de verbes
% cites, qui sont en gras — la regle d'enrichissement du volume.
% La densite d'encre ne tranchait pas ici (tout entre 7,3 et 8,9) ;
% l'EPAISSEUR DES FUTS, elle, est nette : « frapar, donar, lektar »
% mesurent 4,0 px quand « quale en la transitivi » et « ma en li »
% mesurent 3,0. Voir outils/graisse.py.
En ica lasta verbi, la subjekto facas la ago totsame kam en\nl
verbo transitiva : en \VUgras{venar, arivar, kreskar} la subjekto agas,\nl
quale en la transitivi \VUgras{frapar, donar, lektar}; ma en li, netran\cc
\end{VUpage}

% --- Feuillet 56 (folio 52) ----------------------------------------
%  Filet releve a la main : y = 1611, folio y = 148 -> 142.58 mm.
\begin{VUpage}[56]{52}
\VUnotes{143.12mm}{%
(1) Reale, to esas \textit{me}, mea agiveso qua komencas, duras, finas o\nl
cesas. (Noto 1, pag. 21, \textit{Grammarie Complète}.)

(2) « Letro \textit{duras} dum skribesar, ed anke dum lektesar); e mem\nl
plu juste, nam \textit{durar} nule implikas kontinueso. Cetere, quo esas plu\nl
naturala kam : \textit{Me duras parolar, skribar?} Ma analoge, on devas\nl
dicar : \textit{Me durigas mea diskurso o letro} (\textit{Progreso}, IV, 594).

Decido 1615 : On decidas, ke \textit{durar}, \textit{durigar} suficas, repulsante%
}

\VUcontinue sitiva, nul « objekto » povas existar, nul « objekto » povas\nl
recevar la ago dil subjekto; ol restas en lu : on ne dicas\nl
\textit{venar ulu, arivar ulu, kreskar ulu}, quale on dicas : \VUgras{frapar}\nl
\VUgras{ulu, donar ulo, lektar ulo.} Konseque la \textit{netransitivi} nultempe\nl
darfas havar participo pasiva. Do nultempe dicez : \textit{venata,}\nl
\textit{venita} o \textit{venota; arivata, arivita, arivota; kreskata, kreskita}\nl
o \textit{kreskota}, ma nur, segun la kazo : \textit{venanta, veninta, venonta;}\nl
\textit{arivanta, arivinta, arivonta; kreskanta, kreskinta, kreskonta.}\nl
E tale pri omna verbo netransitiva. Cetere, merkez e memorez\nl
bone, ke \textit{nul formo pasiva darfas donesar en Ido a verbo}\nl
\textit{netransitiva.}

\VUblanc{6.48pt}% mesure au fac-simile
46. --- \textit{Mixita} on nomizas en Ido ula verbi quin la linguo\nl
traktas kom transitiva e netransitiva. Nul ambigueso povas\nl
rezultar de to, nam se li uzesas transitive, li havas « ob\cc
jekto » : \VUgras{il turnis e riturnis sua chapelo}; e, se li uzesas\nl
netransitive, li ne havas « objekto » : \VUgras{la tero turnas sur}\nl
\VUgras{sua axo e jiras cirkum la suno.}

%% aucun blanc au fac-simile
Ta procedo supresas fonto di desfacilaji quaze nevin\cc
kebla; altraparte ol esas konforma al kustumo di multa\nl
lingui, ed ici pruvas, ke ta dualeso di senco havas nula de\cc
trimento por la klareso.

%% aucun blanc au fac-simile
Exempli : \VUgras{komencar, durar, finar, cesar : mea laboro}\nl
\VUgras{komencas, duras, finas, cesas} (senco \textit{netransitiva}); \VUgras{me}\nl
\VUgras{komencas, duras, cesas, finas mea laboro} (o laborar) senco\nl
\textit{transitiva} (1).

%% aucun blanc au fac-simile
Ma, se parolesas pri ulu quan me igas laborar, me dicos :\nl
\VUgras{me komencigas, durigas, cesigas, finigas lua laboro.}\nl
Exemple : \VUgras{la pueri cesis ludar} (o : sua ludo); ma : \VUgras{cesigez}\nl
\VUgras{la pueri ludar} (o : la ludo dil pueri), o : \VUgras{igez ke la pueri}\nl
\VUgras{cesez ludar} (o : cesez sua ludo). Quale on vidas, la Franca\nl
verbo \textit{continuer} (ne cesar to quon on agas), tradukesas per\nl
\VUgras{durar} (2). Lasta exemplo : \VUgras{Ne povante durar ipsa ta laboro,}\nl
\VUgras{me durigos olu da mea fratulo.}
\end{VUpage}

% --- Feuillet 57 (folio 53) ----------------------------------------
%  Filet releve a la main : y = 1112, folio y = 141 -> 100.92 mm.
\begin{VUpage}[57]{53}
\VUnotes{99.33mm}{%
\VUcontinue \textit{kontinuacar}, qua obtenis un voco. --- Kom noto pri ca decido tro\cc
vesas : « \textit{Grammaire Complète} (sancionita dal komisitaro konstanta\nl
dil \textit{Delegitaro}) fixigas pri durar, durigar, pag. 35, regulo 32.

Kontinuigar = igar kontinua, sen ula rupto di kontinueso. Ex. :\nl
\textit{Durez ta laboro e, se vu ne povos ipsa durar lu, durigez lu da altru.}\nl
\textit{Ma ol esez kontinuigata til parfino sen la maxim kurta interrupto.}

(1) Ta omna verbi mixita, e nome : \textit{komencar, finar} havas, quale\nl
\textit{durar}, senco netransitiva e senco transitiva. Or, nur pri \textit{durar} til\nl
nun on pensis a special radiko por la senco transitiva. Tamen \textit{durar}\nl
esas nek plu frequa, nek plu importanta kam \textit{finar} o \textit{komencar.}\nl
Ma ula lingui havas vorto specala por olua senco transitiva, e\nl
kredeble pro to on judikis tala vorto kom necesa en Ido. Kad en ta\nl
lingui existas anke special radiko por la senco transitiva di \textit{komencar,}\nl
\textit{finar, cesar, turnar, movar} e. c.? Do, nur pri la senco transitiva di\nl
\textit{durar} li facas ecepto. L'akademio esis justa opinionante, ke 16 yari\nl
de praktiko montris ico sat bone : \textit{durar} ne bezonas ecepto parti\cc
kulara; ol povas havar, quale la cetera verbi mixita, la du senci\nl
transitiva e ne transitiva en un sola formo.

(2) La slavi devas do atencar por ne falar en l'idiotismo di sua\nl
lingui, qui uzas \VUgras{su} mem por l'unesma e la duesma personi; li ne\nl
dicez : \textit{me lavas su, vi lavas su}, e. c. ma bone : \VUgras{me lavas me, vi}\nl
\VUgras{lavas vi.}%
}

Lo sama esas praktikata pri la verbi \VUgras{chanjar, movar, tur}\cc
\VUgras{nar, pendar,} e. c. (Videz la lexiko).

%% aucun blanc au fac-simile
\VUgras{Me chanjis depos mea yuneso; me chanjis mea vesti;} ---\nl
\VUgras{el chanjigis} (da ulu) \VUgras{la kolumo di sua peliso} (\textit{desfidante}\nl
\textit{a sua propra habileso}).

%% aucun blanc au fac-simile
\VUgras{Mea chapelo pendis an arboro; me pendis mea chapelo}\nl
\VUgras{an arboro;} --- \VUgras{vu pendigos me.}

%% aucun blanc au fac-simile
\VUgras{La tero movas} (jiras) \VUgras{cirkum la suno} (on do parolos\nl
pri la movo (jiro) di la tero.

%% aucun blanc au fac-simile
\VUgras{Natante la fisho movas sua flosi e kaudo; ta navo movesas}\nl
\VUgras{da} (o \VUgras{per}) \VUgras{vaporo.}

%% aucun blanc au fac-simile
\VUgras{Me turnas la roto;} --- \VUgras{me turnigas la roto} (da altru).

\VUgras{Paris komunikas telefone kun Lyon;} --- \VUgras{on komunikis a}\nl
\VUgras{me interesanta informi} (1).

\VUblanc{5.66pt}% mesure au fac-simile
47. --- La \textit{verbi reflektiva} esas formacata per la pronomi\nl
dil unesma e dil duesma personi (\textit{me, ni; vu, vi; tu}) por\nl
ta personi, e per la pronomo \VUgras{su} (nevariebla) por la triesma\nl
persono : \VUgras{me lavas me, tu lavas tu, ni lavas ni, vi lavas vi,}\nl
\VUgras{vu lavas vu, il lavas su, el lavas su, lu lavas su, li lavas su,}\nl
\VUgras{ili lavas su, eli lavas su} (o \VUgras{su lavas}) (2).
\end{VUpage}

% --- Feuillet 58 (folio 54) ----------------------------------------
%  Filet releve a la main : y = 1621, folio y = 159 -> 142.49 mm.
\begin{VUpage}[58]{54}
\VUnotes{141.48mm}{%
(1) « On ne obliviez d'uzo tre komoda di la prefixo \textit{inter} por\nl
formacar la verbi reciproka. » (\textit{Progreso}, IV, 147 : Remarki, 1\textsuperscript{a}). ---\nl
Ta verbi reciproka nature genitas substantivi reciproka : \textit{inter}\cc
\textit{konsento, internoco, interhelpo, interodio, interamo} e. c.

(2) La senco esas pasiva, on do uzas pasivo, segun la principo\nl
tante ofte memorigita por Ido : tradukez, ne sklavatre segun la\nl
vorti (di vua linguo), ma segun l'ideo (expresenda en Ido).%
}

48. --- La \textit{verbi reciproka} esas formacata per adjuntar al\nl
verbo \textit{transitiva} la vorti \VUgras{l'una l'altra} (o \VUgras{una altra}), kande\nl
parolesas pri du subjekti, e per adjuntar \VUgras{l'uni l'altri} (o \VUgras{uni}\nl
\VUgras{altri}), kande parolesas pri subjekti plu multa kam du :\nl
\VUgras{li batas l'una l'altra; amez l'uni l'altri.}

%% aucun blanc au fac-simile
On anke uzas la adverbo \VUgras{reciproke} (sen \VUgras{su}) : \VUgras{li helpez e}\nl
\VUgras{sustenez reciproke.}

%% aucun blanc au fac-simile
Kande la verbo esas \textit{netransitiva} e lua komplemento be\cc
zonas prepoziciono, on pozas ica inter \VUgras{una} ed \VUgras{altra} : \VUgras{li iras}\nl
\VUgras{una kun altra; li falis una sur altra; li kuris una kontre}\nl
\VUgras{altra; la du peci fricionas una sur altra.}

%% aucun blanc au fac-simile
Fine on darfas uzar en ula kazi verbo kompozita kun\nl
\VUgras{inter-} kom prefixo, kande nul dubo povas existar, ke la sub\cc
jekti facas inter su la ago : \VUgras{li interparolas; ni interkonsentis;}\nl
\VUgras{pro que vi interdisputas ed interkombatas?} \VUgras{La homi devus}\nl
\VUgras{ne internocar ed interodiar, ma interhelpar e mem inter}\cc
\VUgras{amar} (1).

\VUblanc{4.70pt}% mesure au fac-simile
49. --- La \textit{falsa reflektivi}, to esas la verbi reflektiva per\nl
la formo, ma \textit{pasiva per la senco} quin on uzas ofte en nia\nl
lingui, kontrelogike, devas tradukesar per la pasivo (kun\cc
juntala prefere) (2) : \VUgras{to trovesas hike; to videsas ofte; ta}\nl
\VUgras{libro lektesas facile; to manjesas plezure; la vazo ruptesis}\nl
\VUgras{dal shoko; ta vorto uzesas tre rare; la fenestro apertesis dal}\nl
\VUgras{vento; la aernavo elevesas desfacile.}

%% aucun blanc au fac-simile
Ma, pro ke esas permisata omna metafori klara e natu\cc
rala, on darfas uzar la verbo reflektiva, kande on personigas\nl
l'objekto e konsideras lu kom facanta la ago sur su ipsa;\nl
exemple on dicas : \VUgras{la suno levas su majestoze}, quale on\nl
dicas : \VUgras{Petrus levas su frue.}

\VUblanc{4.73pt}% mesure au fac-simile
50. --- La \textit{verbi unpersona} ne havas aparanta subjekto :\nl
\VUgras{oportas, importas, konvenas, decas, suficas, pluvas, ventas,}\nl
\VUgras{nivas, pruinas, frostas} e. c.
\end{VUpage}

% --- Feuillet 59 (folio 55) ----------------------------------------
%  Filet releve a la main : y = 1587, folio y = 144 -> 140.88 mm.
\begin{VUpage}[59]{55}
\VUnotes{141.13mm}{%
(1) Vice \VUgras{oportas} uzez prefere la verbo \VUgras{mustar} kun subjekti : \textit{mea}\nl
\textit{patrulo mustas departar. Me, tu il, vu} e. c. \textit{mustas departar frue.}

\VUgras{Mustar} implikas neceseso absoluta, neeskartebla : \textit{ni mustas}\nl
\textit{mortar; vu mustas repozar.}

\VUgras{Devar} implikas obligo : \textit{ni devas respektar} ed \textit{obediar nia gepatri.}

\VUgras{Povar} implikas posibleso, e \VUgras{darfar}, yuro, permiso : \textit{me povus venjar}\nl
\textit{me, ma me ne darfas, nam omno interdiktas lo ad me.}%
}

% Cette ligne est un alinea renfonce (x0 = 58 px releve, contre 54 px
% pour le renfoncement ordinaire), non la suite du folio precedent.
En la realeso por frazi tal quala : \VUgras{importas departar frue,}\nl
la subjekto esas : \VUgras{departar frue.} Pro to, en la kazi analoga\nl
on devas uzar l'adjektivo kom atributo (e ne la adverbo\nl
quale en Esperanto) : \textit{necesa esas} o : \VUgras{esas necesa departar}\nl
\VUgras{frue} (1).

\VUblanc{11.09pt}% mesure au fac-simile
Quante sorgoze la konjugo-sistemo di Ido esis longatempe\nl
diskutata en \textit{Progreso} pruvas konvinkive la numeri sequanta :\nl
10 --- 11 --- 13 --- 14 --- 15 --- 18 --- 19 --- 22 --- 23\nl
--- 24 --- 26 --- 28 --- 31-32 --- 33 --- 36 --- 37 --- 38 --- 40.

\VUblanc{2.97pt}% mesure au fac-simile
Pos ta diskuti l'Akademio decidis :

\VUblancAlinea
Decido 671 : « On replusas la konjugo-sistemi propozita :\nl
I, 377 e III, 510; --- I, 705; --- II, 399; --- II, 400; ---\nl
II, 587 e 593; --- II, 722; --- II, 723; --- III, 611; ---\nl
IV, 208; en apendico IV di cirkularo 10 (e ca lasta sistemo\nl
citesas en noto 3, p. 691, dil sama numero kontenanta la\nl
supera decido : \textit{Progreso}, IV, februaro 1912).

\VUblancAlinea
La decido dil Akademio judikesas kom tre saja e prak\cc
tikala, kande on konsideras, ke la konjugo di Ido esas tante\nl
simpla ke ol konsistas nur ek ico :

\VUblancAlinea
\VUgras{a, i, o, u} por \textit{prezento, pasinto, futuro, kondicionalo} e\nl
sequata da \VUgras{s} en la personal modi (\VUgras{-as, -is, -os, -us}), da \VUgras{r} en\nl
l'infinitivi (\VUgras{-ar, -ir, -or}), dal grupo \VUgras{-nt} en la participi aktiva\nl
(\VUgras{-ant, -int, -ont}), da \VUgras{t} en la participi pasiva (\VUgras{-at, -it, -ot}),\nl
plus \VUgras{-ez} por l'imperativo (o voltivo).

\VUblancAlinea
Vere, kad on povas trovar konjugo-sistemo plu kohe\cc
ranta, plu reguloza, plu simpla? E kad ol ne valoras irga\nl
naturaleso \textit{nacionala} (ne plu ciencoza) arbitrial e hibrida?

%% aucun blanc au fac-simile
(Videz en la 6-ma apendico : « La konjugo-sistemo di\nl
Ido » e « Verbal dezinenci Idala ».)

%% aucun blanc au fac-simile
On reprochis ad \textit{-as, -is, -os, -us}, e. c., diferar inter su
\parplein
\end{VUpage}

% --- Feuillet 60 (folio 56) ----------------------------------------
%  Filet releve a la main : y = 1739, folio y = 121 -> 155.70 mm.
%  Blancs du titre mesures AVANT composition (lecon du folio 46) :
%  136,0 px au-dessus, 108,1 px au-dessous.
%  La derniere ligne de note (« elu (v. § 52). ») passe sous le seuil :
%  35 lignes au fac-simile, 34 au detecteur.
\begin{VUpage}[60]{56}
\VUnotes{155.70mm}{%
(1) Ma, se \textit{tro, plu, min, maxim, minim} uzesas izolite, li recevas\nl
elu (v. § 52).%
}

\VUcontinue nur per vokalo. Ma kad en la Angla, la Germana, la Franca,\nl
e. c., ula verbal formi diferas per altro? Ka nun ankore la\nl
diversa populi dil mondo ne perceptas e bone dicernas en\nl
la Latina : \textit{legis, legas, leges; amabas, amabis; bibo, bibi;}\nl
\textit{lego, legi; amant, ament; amantur, amentur; amabat, amabit;}\nl
\textit{amabant, amabunt; erant, erunt; sunt, sint,} e. c.? Pri ca\nl
linguo Latina, ankore nun tante uzata en la medii katolika,\nl
ed olim la maxim uzata e parolata de la lingui, kad on\nl
ultempe dicis, ke ta difero di formi verbala per un sol\nl
vokalo jenas o jenis la kompreno? Kad on reprochis lo al\nl
naturala lingui che qui ta fakto existas?

%% aucun blanc au fac-simile
Pro quo supozar, ke en la helpolingo, stranjera ad omni,\nl
on esos min atencema pri bona pronunco, kam en altra\nl
linguo stranjera? Cetere, en qua linguo esas permisata pro\cc
nuncar tante sensorge, ke l'audanto povas konfundar \textit{as} ad\nl
\textit{os, us} ad \textit{as, is} ad \textit{es, a} ad \textit{e} o \textit{e} ad \textit{i}? En qua linguo on esas\nl
dispensata pronuncar klare, por esar bone komprenata?

\VUsaut{7.19mm}
\VUtitre{10.15pt}{-40}{ADVERBI.}
\VUsaut{5.79mm}

51. --- Quale en nia lingui, l'adverbi esas nevariebla\nl
en Ido.

%% aucun blanc au fac-simile
Lia gradi komparala indikesas segun la maniero uzata\nl
por l'adjektivi (§ 28).

%% aucun blanc au fac-simile
Existas en Ido triopla kategorio de adverbi : le \textit{simpla},\nl
le \textit{derivita}, le \textit{kompozita.}

%% aucun blanc au fac-simile
Le \textit{simpla}, esante adverbi nature e signifike, ne bezonas\nl
dezinenco : \textit{tre, tro, plu, olim, nun, nur}, e. c.(1).

%% aucun blanc au fac-simile
Le \textit{derivita} e le \textit{kompozita}, ne esante adverbi nature, havas\nl
omni la dezinenco \VUgras{-e}, qua igas li adverba. On adjuntas ica\nl
dezinenco al vorti sendezinenca, ed on substitucas lu al\nl
dezinenco di la ceteri : \VUgras{Lore, pluse} (de \textit{lor, plus}), \VUgras{vere} (de\nl
\textit{vera}), \VUgras{nokte} (de \textit{nokto}), \VUgras{koncerne} (de \textit{koncernar}), \VUgras{super}\cc
\VUgras{poze} (de \textit{super-pozar}).
\end{VUpage}

% --- Feuillet 61 (folio 57) ----------------------------------------
%  Filet releve a la main : y = 1152, folio y = 154 -> 103.21 mm.
%  La ligne de note « utileso. » passe sous le seuil : 43 lignes au
%  fac-simile, 42 au detecteur.
\begin{VUpage}[61]{57}
\VUnotes{103.21mm}{%
(1) On darfas derivar direte, se la senco permisas, adverbo de\nl
substantivo, quale ni ja vidis supere : \VUgras{nature}, de naturo; \VUgras{persone}, de\nl
persono; \VUgras{individue}, de individuo; \VUgras{sexue}, de sexuo. Ma povas divenar\nl
necesa od utila uzar la mediaco di adjektivo. Exemple, en ula kazi,\nl
vice \VUgras{logike}, povas esar utila o necesa dicar \textit{logikale} o \textit{logikoze}, por\nl
pluprecizigar. Exter ta kazi, esas plu bona derivar direte l'adverbo\nl
del substantivo, se la senco permisas, e ne plulongigar l'adverbo sen\nl
utileso.

En \textit{Progreso}, III, 218, on lektas : Kad \textit{bele agar} = bela ago,\nl
\textit{forte kantar} = forta kanto, \textit{parole promisar} = parola promiso,\nl
\textit{letre respondar} = letra respondo? Se no, me ne komprenas. Se yes,\nl
on povas (darfas) same dicar : \textit{pede irar} = peda iro, \textit{dome restar} =\nl
doma resto, \textit{jorne laborar} = jorna laboro = laboro qua esas jorno,\nl
\textit{veture promenar} = vetura promeno. Sed lor (ma lore) quon signi\cc
fikas ica expresuri? (E. \textsc{Ferrand}.)

\VUgras{Respondo.} --- « Evidente, en la duesma serio de adverbi, la formo\nl
adverbala \VUgras{-e} vicas (remplasas) prepoziciono tacita : \textit{pede} = per\nl
pedo, \textit{dome} = en domo, \textit{jorne} = dum jorno. Pro to on ne povas\nl
(darfas) derivar de li nemediate adjektivi. On devas dicar ex. :\nl
« promeno \textit{en} o \textit{per} veturo ». Sed (ma) en l'unesma serio anke \textit{parolo},\nl
\textit{letro} povas (darfas) tradukesar per : « per parolo, per letro ».%
}

Quale on vidas, l'adverbi darfas derivar de prepoziciono\nl
od adverbo (\textit{lor-e, plus-e}), de adjektivo o substantivo (\textit{ver-e,}\nl
\textit{nokt-e}), de verbo simpla o kompozita (\textit{koncern-e, super-poz-e}).

%% aucun blanc au fac-simile
Pri lia komplemento l'adverbi kondutas quale la vorti\nl
de qui li derivas : \VUgras{koncerne ico}, nam on dicas : \textit{koncernar}\nl
\textit{ulo}; \VUgras{konforme al modelo}, nam on dicas \textit{konforma al modelo};\nl
\VUgras{funde dil barelo}, nam on dicas la \textit{fundo dil barelo}; e tale\nl
sempre.

\VUblancAlinea
L'adverbi qui venas de adjektivi qualifikanta indikas\nl
generale la maniero : \VUgras{vere}, en maniero vera; \VUgras{bone}, en ma\cc
niero bona.

%% aucun blanc au fac-simile
Ti qui venas de substantivi povas expresar cirkonstanco\nl
di tempo, di loko, di maniero, e. c. Ex. : \VUgras{jorne}, dum jorno;\nl
\VUgras{nokte}, dum nokto; \VUgras{dome}, en la domo; \VUgras{heme}, en la hemo, la\nl
lojeyo familiala; \VUgras{aere}, en la aero; \VUgras{veture}, per veturo; \VUgras{pede},\nl
per la pedi, e. c. (1).

%% aucun blanc au fac-simile
Relate la senco, Ido havas adverbi \textit{di quanteso, di loko,}\nl
\textit{di afirmo, di nego o dubo, di maniero.}

\VUblanc{6.00pt}% mesure au fac-simile
52. --- \textit{Adverbi di quanteso}. On ja konocas ti, qui uzesas\nl
por la gradi komparala : \VUgras{plu, min, tam... kam, maxim,}\nl
\VUgras{minim, tre.}
\end{VUpage}

% -------------------------------------------------------------------
%  feuillet 62 — folio 58 : 17 lignes de corps, 25 de note.
%  Pas regulier 41,42 px : la page n'est pas cardee, et AUCUNE de ses
%  seize frontieres d'alinea ne porte de blanc (le plus grand pas du
%  corps vaut 42,9 px, 1,5 px au-dessus du pas courant).
% -------------------------------------------------------------------
\begin{VUpage}[62]{58}
\VUnotes{90.65mm}{%
(1) \VUgras{Plu} expresas \textit{augmento} di ula quanto (qua restas o konsi\cc
deresas kom identa a su ipsa); \VUgras{plus} (quale la signo matematikal +\nl
quan ol tradukas) signifikas adiciono, do generale \textit{adjunto} di ulo\nl
ad altro.

Segun la principo di unasenceso, ta du radiki devas konservar\nl
sua propra senco en sua omna derivaji. \textit{Plue} havas do la sama senco\nl
kam \textit{plu}; ma ol uzesas izolita quale adverbo nedependanta, kontre\nl
ke \textit{plu} uzesas avan altra vorto, quan ol modifikas : \textit{plue} equivalas\nl
proxime a \textit{plu multe} (o \textit{plu grande}, e. c.); \textit{plua} a \textit{plu multa}; e \textit{pluo}\nl
a \textit{plu multo}. L'unesma vorti esas nur min preciza, pro ke li aplikesas\nl
a speco nedefinita o tacita di quanteso, qua esas indikata plu expli\cc
cite en la expresuri \textit{plu multe, plu grande, plu larje, plu alte...}, per\nl
la vorto adjuntita a \textit{plu}. Kontree, \textit{pluse} expresas ne augmento di\nl
irga quanto, ma adjunto di ula objekto, qua povas ne esar quantajo\nl
o ne konsiderar kom quantajo. Ol tradukas do exakte la vorto F.\nl
\textit{de plus, en plus, en outre}; D. \textit{weiter, zudem, sonst}; E. \textit{moreover}.\nl
Exempli : \VUgras{On devas pagar plue} (plu multe) : \VUgras{on devas pagar pluse}\nl
(on devas adjuntar altra preco). --- \VUgras{Vino e kafeo pagesas pluse}\nl
to esas : ultre la sumo pagita por la cetero. --- Ma \VUgras{vino e kafeo}\nl
\VUgras{pagesas plue} signifikus : \textit{plu kare} (chere) (kam antee, kam altra\cc
loke). --- \VUgras{Ho! ico kustas plue} (t. e. plu multe). --- \VUgras{Me donas ico pluse}\nl
(po la sama preco) dicos vendisto adjuntante ula bagatelo (\textit{Pro}\cc
\textit{greso}, II, 666.)

On komparez : \VUgras{me prenis plu multa pomi} e : \VUgras{me prenis, pluse,}\nl
\VUgras{multa pomi.}%
}

Ni adjuntez \VUgras{plus, minus}, tre uzata en matematiko, od en\nl
la senco matematikala +, ---.

%% aucun blanc au fac-simile
\VUgras{Plu, min, maxim, minim} recevas la dezinenco \VUgras{-e}, kande\nl
li esas adverbi izolita.

%% aucun blanc au fac-simile
Li sempre sequas la verbo quan li modifikas : \VUgras{Il ne povas}\nl
\VUgras{spensar plue} (1). \VUgras{Vu ne povas donar mine. Li recevis}\nl
\VUgras{maxime, minime} (de omni).

%% aucun blanc au fac-simile
Por indikar, ke la grado konstante kreskas o diminutas,\nl
on uzas l'expresuro : \VUgras{sempre plue} (o : \textit{sempre plu multe}),\nl
\VUgras{sempre mine} (o : \textit{sempre min multe}). Ex. : \VUgras{me amas lu}\nl
\VUgras{sempre plue} (o : \textit{sempre plu multe}); \VUgras{il laboras sempre mine}\nl
(o : \textit{sempre min multe}).

%% aucun blanc au fac-simile
\VUgras{Pluse, minuse} indikas, quale la signi +, ---, adjunto o\nl
sustraciono : \VUgras{pozez 1 manjilaro pluse, 1 manjilaro minuse}\nl
\VUgras{sur la tablo, nam ni havos un repastano pluse, minuse.}

%% aucun blanc au fac-simile
\VUgras{Maxime, minime} = en la \VUgras{maxim, minim} granda quanteso\nl
a grado : \VUgras{vu laboris, produktis maxime, minime} (de omni).\parplein
\end{VUpage}

% -------------------------------------------------------------------
%  feuillet 63 — folio 59 : 23 lignes de corps, 19 de note.
%  Un seul blanc d'articulation, devant « 53. --- » (exces 21,6 px).
% -------------------------------------------------------------------
\begin{VUpage}[63]{59}
\VUnotes{107.79mm}{%
(1) \textit{Posible}, en ta expresuri, equivalas « segun posibleso ».

(2) \textit{Proxime} de l'adjektivo \textit{proxima} e prepoziciono \textit{proxim}. Ta\nl
adverbo ya kontenas ideo di \textit{proximeso} relate la nombro indikata.

(3) Ne konfundez \textit{sate} a \textit{pasable}, od inverse. \textit{Pasable} (quale lua\nl
adjektivo \textit{pasabla}) indikas ne : quanteso, grado suficanta, ma\nl
grado, valoro meza, la minim granda quan on povas aceptar o\nl
tolerar : \VUgras{Il esas sat richa por donar jeneroze, ma il esas nur pasable}\nl
\VUgras{donema.}

(4) \VUgras{Plu, min, sat, tro} uzesas sola kun la adjektivi e la adverbi;\nl
ma kun la substantivi e la verbi on devas adjuntar a li rispektive :\nl
\textit{multa o multe : plu multa aquo e min multa vino; tro multa homi;}\nl
\VUgras{me drinkis sat multe; ne manjez tro multe.}

(5) \VUgras{Tam} esas sequata da \VUgras{kam} e montras egaleso; \VUgras{tante} sequesas\nl
da \VUgras{ke} e montras granda quanteso o alta grado. Ex. : \VUgras{Il esas tam}\nl
\VUgras{habila kam tu; --- il esis tante habila, ke il ruptis nulo.}

(6) Remarkez, ke la F. vorti \textit{que, où} devas tradukesar per \textit{kande},\nl
se li havas la senco di ta adverbo. Ex. : \textit{Je l'ai vu la dernière fois}\nl
\textit{que je suis venu.} \VUgras{Me vidis lu en la lasta foyo, kande me venis.} ---\nl
\textit{Le jour où il mourut}, \VUgras{en la dio kande il mortis.}%
}

\VUgras{Maxim posible, minim posible} (1) (\textit{sen artiklo}, se ne esas\nl
substantivo expresata o tacata) : \VUgras{venez maxim balde posible;}\nl
\VUgras{la maxim granda nombro posible; la maxim bel infanto}\nl
\VUgras{posible; la maxim laborera} (viro) \VUgras{posible; bruisez minime}\nl
\VUgras{posible; en omno selektez lo maxim bona posible.} --- \VUgras{Il}\nl
\VUgras{kuris tam rapide kam posible. Yen diversa klovi; prenez}\nl
\VUgras{le maxim longa posible.}

%% aucun blanc au fac-simile
\VUgras{Admaxime, adminime} : \VUgras{Ni esos kin adminime e dek ad}\cc
\VUgras{maxime.}

%% aucun blanc au fac-simile
\VUgras{Proxime}, qua tradukas F. \textit{à peu près}, environ, E. \textit{about,}\nl
\textit{nearly}, D. \textit{ungefähr, etwa}. Ex. : \VUgras{la asistanti esis kina-dek o :}\nl
\VUgras{kina-dek e kin proxime. Esis tri kloki, o tri kloki proxime} (2).

%% aucun blanc au fac-simile
\VUgras{Sat} = en sufico : \VUgras{sat richa, sat kurajoza} (3); \VUgras{pri glorio,}\nl
\VUgras{honori ed influo il havas sat multe} (4).

%% aucun blanc au fac-simile
\VUgras{Tro} (\textit{troa, troe}) implikas ideo di eceso : \VUgras{tro richa, tro}\nl
\VUgras{benigna.}

%% aucun blanc au fac-simile
Fine la adverbi ja vidita (§ 39) : \VUgras{kelke, multe, poke,}\nl
\VUgras{quante, tante} (5), \VUgras{irge quante.}

\VUblanc{5.19pt}% mesure au fac-simile
53. --- \textit{L'adverbi di tempo} esas :

%% aucun blanc au fac-simile
\VUgras{Kande} (relativa e questionala) = ye la tempo, ye l'in\cc
stanto en qua (6).

%% aucun blanc au fac-simile
\VUgras{Dume} (de la prep. \textit{dum}) = dum ta tempo.

%% aucun blanc au fac-simile
\VUgras{Fine} (de la subs. \textit{fino}) = ye la fino, kom fino.
\end{VUpage}

% -------------------------------------------------------------------
%  feuillet 64 — folio 60 : 26 lignes de corps, 15 de note.
%  Page dense : aucune de ses seize frontieres d'alinea ne porte de
%  blanc (le plus grand pas du corps vaut 43,2 px, 1,9 px de trop).
% -------------------------------------------------------------------
\begin{VUpage}[64]{60}
\VUnotes{119.18mm}{%
(1) \VUgras{Lore... lore} = ye ul tempo, instanto... ye altra... Ex : \VUgras{lore il}\nl
\VUgras{ridas, lore il ploras.} France : tantôt il rit, tantôt il pleure.

(2) Tre diferas de \VUgras{konseque}, nam milfoye ulo venas, eventas\nl
seque, ma tote ne \textit{konseque} altro. Ni evitez pri co imitar Esper\cc
antisti. \textit{Seque} la Franca revoluciono, ma ne \textit{konseque}, eventis la\nl
regno di Napoléon.

(3) \VUgras{Ja ne} tote ne equivalas \VUgras{ne ja} (quale \VUgras{tote ne} ne equivalas \VUgras{ne}\nl
\VUgras{tote}). Ex. : \VUgras{Il ja ne esis kontenta, e nun il ne ja esas kontenta,}\nl
\VUgras{malgre omno quon ni agis por kontentigar lu.}

(4) Ne konfundez \textit{ankore} (prezenta, pasinta o mem futura duro)\nl
a \textit{itere}, qua expresas \textit{itero} (de \textit{iterar}, rifacar la ago), o a \textit{pluse}.\nl
Ex. : \VUgras{Il kantas ankore} (il ne cesis). --- \VUgras{Il kantas itere} (il cesabis,\nl
ma il ri-kantas). --- \VUgras{Il kantas un ario pluse} (il adjuntas un ario\nl
a le kantita). --- La F. expresuri : \textit{encore un, encore deux} e. c.\nl
tradukesas : \textit{un pluse, du pluse} e. c.%
}

\VUgras{Lore} (de la prep. \textit{lor}) = ye ta epoko, tempo (1).

%% aucun blanc au fac-simile
\VUgras{Nun} (sen \VUgras{-e} kom primitiva) = en la tempo prezenta, \textit{pre}\cc
\textit{zente}.

%% aucun blanc au fac-simile
\VUgras{Olim} (sen \VUgras{-e} kom primitiva) = en tempo plu o min an\cc
ciena, antiqua.

%% aucun blanc au fac-simile
\VUgras{Sempre} = en omna tempo.

%% aucun blanc au fac-simile
\VUgras{Seque} (de la subs. \textit{sequo}) = ye la sequo, kom sequo (2).

%% aucun blanc au fac-simile
\VUgras{Frue} = en tempo ne tarda.

%% aucun blanc au fac-simile
\VUgras{Balde} = pos ne longa tempo.

%% aucun blanc au fac-simile
\VUgras{Tarde} (de \textit{tarda}) = ye tempo, instanto ne proxima relate\nl
altra, fixigata o kustumala.

%% aucun blanc au fac-simile
\VUgras{Erste} = ne plu balde kam... : \VUgras{Me venos erste morge.}

%% aucun blanc au fac-simile
\VUgras{Hiere} = en la dio qua preiris ica (la nuna).

%% aucun blanc au fac-simile
\VUgras{Morge} = en la dio qua sequos ica.

%% aucun blanc au fac-simile
\VUgras{Jus} (\textit{nur kun tempo pasinta}) = juste ye l'instanto antea :\nl
il jus ekiris.

%% aucun blanc au fac-simile
\VUgras{Quik} = sen la minima ajorno o tardeso : \VUgras{facez quik ta}\nl
\VUgras{laboro, nam ol tre urjas.}

%% aucun blanc au fac-simile
\VUgras{Ja} = E. \textit{already}, D. \textit{schon}, F. \textit{déjà}, S. \textit{ya}, I. \textit{già}, Por. \textit{ja}.\nl
Ex. : \VUgras{Il ja manjis omno.}

%% aucun blanc au fac-simile
\VUgras{Ne ja} pri kozo qua ne eventis ye l'instanto prezenta : \VUgras{il}\nl
\VUgras{ne ja arivis, ma il quik arivos} (3).

%% aucun blanc au fac-simile
\VUgras{Ankore} indikas duro plu longa pri ago o stando : \VUgras{il esas}\nl
\VUgras{ankore malada} (4).

%% aucun blanc au fac-simile
\VUgras{Ne plus} esas la negativo di \VUgras{ankore} : \VUgras{Kad il vivas ankore?}\nl
\VUgras{No, il ne plus vivas : il ja esis mortinta, kande me departis.}\parplein
\end{VUpage}

% -------------------------------------------------------------------
%  feuillet 65 — folio 61 : 16 lignes de corps, 27 de note.
%  Le filet de note echappe a outils/filet.py (il le trouve, mais son
%  garde-fou « rien d'autre sur la ligne » le rejette) : releve a la
%  main a y = 926 px, d'ou 24,30 + (926-210) x 25,4/300 = 84,92 mm.
%  L'enumeration « ultempe... irgekande » n'a AUCUN blanc entre ses
%  items ; le blanc d'alinea ordinaire subsiste devant les lignes 3, 4
%  et 13 seulement.
% -------------------------------------------------------------------
\begin{VUpage}[65]{61}
\VUnotes{84.92mm}{%
(1) \textit{Decido} 1626 : Unanime (minus 1 voco) on repulsas vorto\nl
nekompozita vice « nultempe ».

« Kontre la adopto di \textit{never} (VII, 287, v. anke 206) me deziras\nl
protestar tre insiste. La vorto kreus multa kolizioni pro la simileso\nl
a \textit{nevera} : vorti tam freque bezonata ne darfas esar tro simila;\nl
en telefono ed exter ol on ne povas distingar « il parolas never »\nl
e « ... nevere » (nam in \textit{ne-vere} on tre ofte pro kontrasto acentizos\nl
\textit{ne}); e se on pronuncas rapide sen pauzo « ne vere eloquenta » e\nl
« never eloquenta », la danjero di misaudo esas tre granda. On\nl
darfas naturale formacar la vorto \textit{versimila} (quo ne esas exakte la\nl
sama kozo kam \textit{probabla}) ma \textit{ne-versimila} miskomprenesos kom\nl
\textit{never-simila}, e. c.

Duesme esas « nevera », ke « on » * never uzos l'adjektivo \textit{nevera} :\nl
pro quo on ne dicus exemple : « pro sua \textit{nevera} asisto en kunsidi il\nl
havas nula influo en nia uniono? » Me ne hezitus uzar tala kombini\nl
(sua \textit{nultempa} asisto), nam tala formacuri esas precize un ek la\nl
belaji di nia linguo. Anke « la \textit{nevereso (nultempeso)} di tala eventi »\nl
devus esar posibla formacuro.

\textit{Nultempe} semblas a me, quale a S\textsuperscript{ro} Couturat, tote suficanta : ol\nl
permisas ta lenta ed emfazoza pronunco qua harmonias tante kun\nl
la expresenda ideo (\textsc{Otto Jespersen}, en \textit{Progreso}, VII, 407).

(2) On devas ne konfundar \VUgras{irgatempe} a \textit{irgekande}. L'unesma esas\nl
nur adverbo nedefinita, ma \VUgras{irgekande} esas pluse relativa, e konseque\nl
unionas du propozicioni : \VUgras{me helpos vu irgatempe; me helpos vu}\nl
\VUgras{irgekande vu volos.} En la lasta exemplo \textit{irgatempe} ne esus bona,\nl
nam, por unionar la du propozicioni : \textit{me helpos vu} e \textit{vu volos}, esas\nl
necesa adverbo relativa; sen olu la propozicioni esus sen ligo.%
}

\VUgras{Antee} (de la prep. \textit{ante}) = en antea tempo.

%% aucun blanc au fac-simile
\VUgras{Pose} (de la prep. \textit{pos}) = en posa tempo.

\VUblancAlinea
Esas adjuntenda la sequanta adverbi kompozita o derivita :

\VUblancAlinea
ultempe = en ula tempo (pasinta o futura);

%% aucun blanc au fac-simile
uldie = en ula dio (pasinta o futura);

%% aucun blanc au fac-simile
nultempe = en nula tempo (1);

%% aucun blanc au fac-simile
kelkatempe = dum kelka tempo;

%% aucun blanc au fac-simile
pokatempe = dum poka tempo;

%% aucun blanc au fac-simile
longatempe = dum longa tempo;

%% aucun blanc au fac-simile
samtempe = en la sama tempo;

%% aucun blanc au fac-simile
irgatempe = en irga tempo;

%% aucun blanc au fac-simile
irgekande = ye irga tempo, instanto en qua (2).

\VUblancAlinea
\VUgras{Unfoye, dufoye, trifoye}, e. c., \VUgras{plurfoye, kelkafoye, multa}\cc
\VUgras{foye, singlafoye, omnafoye}, pri ago facita o subisita \textit{un foyo,}\nl
\textit{plura foyi}, e. c.

%% aucun blanc au fac-simile
\VUgras{Singladie, omnadie, omnasemane, omnamonate, omnayare,}\parplein
\end{VUpage}

% -------------------------------------------------------------------
%  feuillet 66 — folio 62 : 26 lignes de corps, 13 de note.
%  La page OUVRE sur la suite de l'alinea du folio 61 : \VUcontinue.
%  Blanc d'articulation devant « 54. --- » (exces 37,0 px).
% -------------------------------------------------------------------
\begin{VUpage}[66]{62}
\VUnotes{124.37mm}{%
(1) Kom adjektivi di \VUgras{hike, ibe} l'Akademio adoptis la formi \VUgras{hik-a,}\nl
\VUgras{hik-ala} e \VUgras{ib-a, ib-ala}, pro ke la vorto \VUgras{hike} (e \VUgras{ibe}) esas adverbo, en\nl
qua la finalo \VUgras{-e} esas karakteriziva dil adverbo, do deprenebla. Lo\nl
sama valoras pri \textit{supr-e, infr-e, bald-e, fru-e, hier-e, morg-e} e. c., qui\nl
omna esas adverbi e genitas l'adjektivi : \textit{supr-a, infr-a, fru-a, bald-a,}\nl
\textit{hier-a, morg-a}. Altraparte, l'Akademio adoptis \textit{ante-a} kom adjektivo\nl
ed \textit{ante-e} kom adverbo, pro ke \textit{ante} esas prepoziciono, do primitiva\nl
vorto, en qua la finalo \textit{e} apartenas a la radiko. Lo sama valoras pri\nl
\textit{kontre}, exemple, qua genitas l'adjektivo \textit{kontre-a} e l'adverbo \textit{kontree}.\nl
(\textit{Progreso}, V, 555.)

L'adverbo \VUgras{sempre} esas evidente en la e. c. di ca noto e darfas\nl
genitar l'adjektivo \textit{sempra} : \VUgras{fine ilua sempra rezistado tedis ni.}

(2) Videz e komparez la noto 2, pag. 61.%
}

\VUcontinue
omnajorne, omnanokte, e. c. = en singla, omna dio, se\cc
mano, e. c.

%% aucun blanc au fac-simile
\VUgras{Jorne} = dum (la) jorno; \VUgras{nokte} = dum (la) nokto;\nl
\VUgras{matine} = ye o dum (la) matino; \VUgras{vespere} = ye o dum (la)\nl
\VUgras{vespero}.

%% aucun blanc au fac-simile
\VUgras{Cadie, camatine, cavespere, casemane, canokte, camonate,}\nl
\VUgras{cayare}, e. c. = en ica dio, matino, vespero, e. c.

%% aucun blanc au fac-simile
\VUgras{Prehiere} = en la dio preirinta « hiere ».

%% aucun blanc au fac-simile
\VUgras{Posmorge} = en la dio qua sequos « morge ».

%% aucun blanc au fac-simile
\VUgras{Morge matine, morge vespere} o \VUgras{matinmorge, vespermorge.}

\VUblanc{8.90pt}% mesure au fac-simile
54. --- La \textit{adverbi di loko} esas :

\VUblancAlinea
\VUgras{Ube} (questionala e relativa) = en qua loko? o : en qua\nl
(en qui). Ex. : \VUgras{ube vu esas? La lando ube il naskis.}

%% aucun blanc au fac-simile
\VUgras{Hike} = en ica loko (proxima relate la parolanto).

%% aucun blanc au fac-simile
\VUgras{Ibe} = en ita loko (fora relate la parolanto). Ex. : \VUgras{vu}\nl
\VUgras{esis hike, dum ke me serchis vu ibe} (1).

%% aucun blanc au fac-simile
\VUgras{Ulube, ul(a)loke} = en ula loko.

%% aucun blanc au fac-simile
\VUgras{Nulube, nul(a)loke} = en nula loko.

%% aucun blanc au fac-simile
\VUgras{Omnube, omnaloke} = en omna loko.

%% aucun blanc au fac-simile
\VUgras{Irgube, irgaloke} = en irga loko.

%% aucun blanc au fac-simile
\VUgras{Altrube, altraloke} = en altra loko.

%% aucun blanc au fac-simile
La formi kun \VUgras{-ube} ne esas, quale la formi kun \VUgras{-loke}, nur\nl
nedefinita; li esas pluse relativa e konseque unionas du pro\cc
pozicioni (2). Ex. : \VUgras{Me sequos vu omnaloke; irgube vu iros,}\nl
\VUgras{me sequos vu.}

%% aucun blanc au fac-simile
\VUgras{Interne} (di) = en l'internajo;
\end{VUpage}

% -------------------------------------------------------------------
%  feuillet 67 — folio 63 : 22 lignes de corps, 16 de note.
%  Blanc d'articulation devant « 55. --- » (exces 36,8 px) et devant
%  la ligne qui le suit (12,2 px).
%  COQUILLE DE L'ORIGINAL, conservee : la note (3) ouvre par « Me
%  pagis a lu cirkume dek franki » : le guillemet ouvrant manque.
% -------------------------------------------------------------------
\begin{VUpage}[67]{63}
\VUnotes{113.90mm}{%
(1) Ne konfundez ta adverbo a \VUgras{supere} = superpozite. Ex. : \VUgras{Esper}\cc
\VUgras{anto havas sis literi ornata da signo diakritika supere.}

(2) Esas naturala, ke de la proximeso lokala, on iras al proximeso\nl
metafizikala pri valoro, qualeso, evaluo, e. c. Ex. : \VUgras{To kustis de ni}\nl
\VUgras{sisa-dek franki proxime.}

(3) Me pagis a lu \textit{cirkume} dek franki » : esus nejusta dicar : « Me\nl
pagis cirkum dek franki », nam « dek franki » esas la komplemento\nl
\textit{direta} di « pagis », e « cirkume » nur modifikas ol (igante ol min\nl
preciza).

On adoptis \VUgras{cirkum} vice \textit{cirke} kom prepoziciono precize por povar\nl
distingar la prepoziciono \textit{cirkum} de l'adverbo \textit{cirkume}. « Ni drinkis\nl
\textit{cirkume} dek glasi » (t. e. : « dek glasi cirkume » : \textit{dek glasi} esas\nl
l'objekto direta). « Ni sidis \textit{cirkum} dek glasi » (o : « cirkum la\nl
tablo ») : hike \textit{cirkum} esas prepoziciono, on ne povas sidar... la glasi\nl
o la tablo). Inverse : « Ni drinkis \textit{cirkum} la tablo », ne « \textit{cirkume} »\nl
quo implikus, ke ni drinkis (preske) la tablo. (\textit{Progreso}, VII, 479.)%
}

\VUgras{extere} (di) = ye l'exterajo;

%% aucun blanc au fac-simile
\VUgras{supre} (di) = ye la suprajo (1);

%% aucun blanc au fac-simile
\VUgras{infre} (di) = ye l'infrajo;

%% aucun blanc au fac-simile
\VUgras{avane} (de la prep. \textit{avan}) = ye la avanajo;

%% aucun blanc au fac-simile
\VUgras{dope} (de la prep. \textit{dop}) = ye la dopajo;

%% aucun blanc au fac-simile
\VUgras{latere} (de la subs. \textit{latero}) = ye latero di...;

%% aucun blanc au fac-simile
\VUgras{dextre} (di) = ye la dextra latero;

%% aucun blanc au fac-simile
\VUgras{sinistre} (di) = ye la sinistra latero;

%% aucun blanc au fac-simile
\VUgras{proxime} = en la vicineso, ne fore (2);

%% aucun blanc au fac-simile
\VUgras{fore} = ye granda disto.

%% aucun blanc au fac-simile
\VUgras{cirkume} = cirkum la enti o kozi nomata; per proxima\nl
evaluo : \VUgras{La soldati kaptis la urbo ed establisis sua kam}\cc
\VUgras{peyo cirkume.} --- \VUgras{Li sideskis cirkum la ludo-tablo e riskis}\nl
\VUgras{tria-cent franki cirkume} (3).

\VUblanc{8.86pt}% mesure au fac-simile
55. --- La \textit{adverbi di maniero} esas :

\VUblanc{2.93pt}% mesure au fac-simile
\VUgras{Anke} = simile, egale, same : \VUgras{Ni konfesas erorir; ma anke}\nl
\VUgras{vu eroris.}

%% aucun blanc au fac-simile
\VUgras{Apene} = preske ne : \VUgras{li apene regardis ni.}

%% aucun blanc au fac-simile
\VUgras{Kom} = D. \textit{als} (identifizierend); F. en qualité de, en tant\nl
que : \VUgras{Kom Kristano me devas pardonar a mea enemiki,}\nl
\VUgras{quale la Kristo pardonis al sui. Ma kom homo e mem kom}\nl
\VUgras{Kristano, me ne obligesas agar a l'enemiki tale kam me}\parplein
\end{VUpage}

% -------------------------------------------------------------------
%  feuillet 68 — folio 64 : 21 lignes de corps, 20 de note.
%  La page ouvre sur la suite de l'alinea du folio 63 : \VUcontinue.
% -------------------------------------------------------------------
\begin{VUpage}[68]{64}
\VUnotes{102.79mm}{%
(1) \VUgras{Kom} implikas, ke on esas reale to quon dicas la vorto\nl
sequanta l'adverbo : \textit{Kom Kristano}, on esas Kristano : \textit{kom homo,}\nl
on esas homo; \textit{kom rejo}, on esas rejo; \textit{kom statestro}, on esas\nl
statestro; \textit{kom privato}, on esas privato. --- \VUgras{Quale} implikas nur\nl
simileso : \textit{quale} Kristano = simile a Kristano; \textit{quale homo} = simile\nl
a homo; \textit{quale rejo} = simile a rejo; \textit{quale statestro} = simile a\nl
statestro; \textit{quale privato} = simile a privato. Konseque : uzez\nl
« quale », se vu povas remplasigar per « simile a... »; uzez « kom »\nl
en la cetera kazi.

(2) Quale on vidas lo, per sua signifiko, « mem » tre diferas de\nl
\textit{ipsa} e de \textit{sama}, a qui precipue la Franci devas ne konfundar olu.

(3) « Kelka samideani astonesas, ke on uzas lore \textit{ipsa}, lore \VUgras{ipse,}\nl
e questionas, kad ica du vorti havas diversa senci. No, li ne povas\nl
havar diversa senci pro ke li esas sama vorti sub du gramatikala\nl
formi; li havas nur diversa \textit{gramatikal} rolo. Segun la sama regulo,\nl
qua enuncas, ke participo-adverbo referas la subjekto di la propo\cc
ziciono, \textit{ipse} referas sempre la subjekto, e pro to ol esas kelkafoye\nl
plu klara kam \textit{ipsa}. Ex. : \VUgras{La mastro inspektis ipse la domo} esas plu\nl
klara kam : « ipsa »; nam on povus komprenar : \textit{la domo ipsa.}\nl
(\textit{Ipse} povas bone tradukar : F. en personne). Cetere esas sempre%
}

\VUcontinue
agas al amiki. --- On aceptis ilu quale se lu venabus kom\nl
rejo, e tamen il venis inkognite, ne kom statestro, ma kom\nl
privato (1).

%% aucun blanc au fac-simile
\VUgras{Mem} = D. \textit{sogar, selbst}; E. \textit{even}; F. \textit{même} (encore, de plus,\nl
aussi); I. \textit{anche, persino}; S. \textit{aún, hasta, tambien}. Ex. : \VUgras{Mem}\nl
\VUgras{spensante la duoplo, vu ne sucesus.} --- \VUgras{Esas mem neposibla}\nl
\VUgras{dicar proxime la dato ye qua il arivos.}

%% aucun blanc au fac-simile
La signifiko di \textit{mem} justifikas, ke on uzas lu (quale \textit{etiam}\nl
Latina, \textit{encore} Franca) avan komparativo (di supereso o di\nl
infreso) por plufortigar olu : \VUgras{El esas mem plu bela kam}\nl
\VUgras{sua fratino.} --- \VUgras{Me recevis mem min multe kam omni} (2).

%% aucun blanc au fac-simile
\VUgras{Ya.} --- Ta partikulo adverba nultempe uzesas sole, ma kun\nl
irga dico quan ol konfirmas emfaze : \VUgras{Me ya facis to} (ne altru).\nl
--- \VUgras{Pro quo vu eniris? vu savis ya, ke me esas okupata.} ---\nl
\VUgras{Pro quo vu ne respondis? vu savas ya Ido.} --- \VUgras{Por stekar}\nl
\VUgras{ica klovo, ta martelo povas ya suficar; vu ne bezonas plu}\nl
\VUgras{grosa.}

%% aucun blanc au fac-simile
Ultre ta primitiv adverbi di maniero, Ido havas mult altri\nl
derivita o kompozita, inter qui :

%% aucun blanc au fac-simile
\VUgras{Ule, nule, irge, altre, cetere, ipse} (3), \VUgras{single, plure,}\nl
\VUgras{omne.}
\end{VUpage}

% -------------------------------------------------------------------
%  feuillet 69 — folio 65 : 9 lignes de corps, 34 de note.
%  Le detecteur n'en trouve que 33 : la ligne « homi. » (italique, deux
%  syllabes, tres peu d'encre) lui echappe, comme aux feuillets 37, 46,
%  53, 57, 60 et 61. Elle est ici composee a la main ; les controles 9
%  et 11 sauteront la page, les comptes divergeant d'une unite.
%  Filet releve a la main a y = 638 px : 24,30 + (638-225) x 25,4/300.
% -------------------------------------------------------------------
\begin{VUpage}[69]{65}
\VUsignature{167.98mm}{3}% signature de cahier
\VUnotes{59.27mm}{%
plu sekura pozar \textit{ipsa} (o mem \textit{ipse} apud la vorto, quan ol referas :\nl
\textit{la mastro ipsa}. --- Lasta remarko : D. \textit{von selbst,} F. \textit{de soi même}\nl
(\textit{de lui même}, e. c.) ne tradukesas per \textit{ipse}, ma per \textit{spontane} (\textit{Pro}\cc
\textit{greso}, V, 96).

En la sama loko trovesas la remarkigo sequanta : « La sama\nl
konsilo valoras pri \textit{omna}, quan on esas tentata (segun l'exemplo di\nl
F.) pozar \textit{dop} la verbo, quankam ol referas la subjekto. Ex. : \textit{li omna}\nl
\textit{admiris splendida flori}, e ne : \textit{li admiris omna splendida flori} (ube\nl
\textit{omna} referas necese \textit{flori}). \textit{Li admiris omni...} ne esus dusenca, ma\nl
esus min klara kam \textit{li omna admiris...} On ne alegez hike, ke l'akuza\cc
tivo supresus la dusenceso : nam ol ne povus helpar kun la verbo\nl
\textit{esar : li omna esis brava soldati} havas evidente altra senco kam :\nl
\textit{li esis omna brava soldati}, e tamen on ne povus pozar \textit{bona soldati}\nl
en l'akuzativo. Nova pruvo, ke la vortordino esas plu komoda e plu\nl
sekura moyeno kam l'akuzativo. » --- Pri \textit{omna} bone memorez ke ol\nl
sempre uzesas sen artiklo : \textit{omna homi} quale : \textit{kelka homi, multa}\nl
\textit{homi.}

(1) Ye \textit{kom} ni expozis la granda difero existanta inter lu e \textit{quale}.\nl
(Videz pag. 64.)

(2) \VUgras{Partikulara} e \VUgras{aparta} (per konsequo, lia adverbi) tre diferas.\nl
\textit{Partikulara} esas la kontreajo di \textit{generala} od \textit{universala}. \textit{Aparta}\nl
signifikas quaze \textit{separita, izolita, sola}. Ex. : \textit{aparta parolo} (en teatro\nl
od en la vivo); aparta konverso (di du o tri personi en plu granda\nl
asemblo); prenar ulu \textit{aparte} (por parolar a lu specale, exter e for\nl
la ceteri). Il dormas en aparta chambro (ne en la komuna dormeyo).\nl
Li manjis an aparta tablo (ne an la komuna tablo « table-d'hôte »).\nl
(\textit{Progreso}, VI, 154.)

(3) \VUgras{Precipua} = D. \textit{hauptsächlich, Haupt...;} E. \textit{principal;} I. \textit{prin}\cc
\textit{cipale} (precipuo); S. \textit{principal}.

En Ido, \textit{princip-al} havas, quale postulas logiko e praktiko, la\nl
senco : « relatanta principo ». Ka vere ne existas esencal difero\nl
inter : questiono \textit{precipua} e questiono \textit{principala?} Kad \textit{un} e sama\nl
vorto povus suficar (sen omnasorta desavantaji) por \textit{du} idei tam\nl
diferanta kam ici?%
}

\VUgras{Quale} (de \textit{quala}) = en quala maniero? o : en la ma\cc
niero di... (1).

%% aucun blanc au fac-simile
\VUgras{Tale} \ \ \ (de \textit{tala}) = en tala maniero.

%% aucun blanc au fac-simile
\VUgras{Irgequale} = en irge quala maniero.

%% aucun blanc au fac-simile
\VUgras{Aparte} (de \textit{aparta}) = quaze separite, izolite.

%% aucun blanc au fac-simile
\VUgras{Partikulare} (de \textit{partikulara}) = la kontreajo di universale,\nl
komune, generale (2).

%% aucun blanc au fac-simile
\VUgras{Private} (de \textit{privata}) = ne publike.

%% aucun blanc au fac-simile
\VUgras{Precipue} (de \textit{precipua}) = en precipua maniero (3).
\end{VUpage}

% -------------------------------------------------------------------
%  feuillet 70 — folio 66 : 27 lignes de corps, 13 de note.
%  Blanc d'articulation devant « 55. --- » (exces 24,3 px).
%  NUMERO REPETE DANS L'ORIGINAL : ce « 55. » est le second — le folio
%  63 en porte deja un. Conserve tel quel (§ 9).
% -------------------------------------------------------------------
\begin{VUpage}[70]{66}
\VUnotes{124.78mm}{%
(1) L'adverbo \VUgras{prefere} maxim ofte \textit{tradukas} : E. \textit{rather}, D. \textit{lieber}\nl
(haben, wollen), F. \textit{plutôt}, I. \textit{piuttosto}, qui povas tradukesar anke\nl
per \textit{plu vere, plu juste, plu bone, plu precize} e. c. Ex. : \VUgras{Repozez dum}\nl
\VUgras{kelka dii, prefere kam durar ta laboro qua ja tante surmenis vu.} ---\nl
\VUgras{Il pruvis, o plu vere} (\textit{plu justadice}) \VUgras{probis pruvar la yusteso di sua}\nl
\VUgras{revendiki.}

(2) \VUgras{Retroe} (komp. \textit{troe}) trovesas en « Dictionnaire Français-Ido,\nl
p. 463 : \textit{reculons} (à), \VUgras{retroe}, e 2 linei plu supre on devas lektar\nl
anke \textit{retroe} por « en reculant », vice \textit{retro} qua esas imprimeroro.

(3) Por agnoskar la certa posibleso, ma tote ne por expresar l'ideo\nl
di dubo kontenata en \textit{forsan}, e tante distinta de l'altra, ke nia\nl
lingui ipsa dicas equivale : \textit{to forsan esas posibla, ma me tre dubas}\nl
\textit{kad il sucesos.} (Videz \VUgras{forsan} en la adverbi di afirmo, nego o dubo.)%
}

\VUgras{Kune} (de \textit{kun}) = unionite : \VUgras{Li omna kune.}

%% aucun blanc au fac-simile
\VUgras{Kontree} (de \textit{kontre}) = en maniero kontrea.

%% aucun blanc au fac-simile
\VUgras{Itere} (de \textit{iterar}) = per itero, iterante.

%% aucun blanc au fac-simile
\VUgras{Volunte} (de \textit{voluntar}) = kun volunto, D. \textit{gern}; E. \textit{wil}\cc
\textit{lingly}; F. \textit{volontiers}; I. \textit{volentieri}; S. \textit{de bune gana}; Por.\nl
\textit{de boa vontade.}

%% aucun blanc au fac-simile
\VUgras{Konsente} (de \textit{konsentar}) = per konsento;

%% aucun blanc au fac-simile
\VUgras{Prefere} (de \textit{preferar}) = per prefero (1).

%% aucun blanc au fac-simile
\VUgras{Memore} (de \textit{memorar}) = per memoro.

%% aucun blanc au fac-simile
\VUgras{Konseque} (de \textit{konsequar}) = per konsequo.

%% aucun blanc au fac-simile
\VUgras{Segunvole} = segun (la) volo.

%% aucun blanc au fac-simile
\VUgras{Kontrevole} = kontre (la) volo.

%% aucun blanc au fac-simile
\VUgras{Afrankite} = (de \textit{afrankar, afrankita}) : la sendo kustos\nl
de vu tri franki (kun la afranko) afrankite.

%% aucun blanc au fac-simile
\VUgras{Retroe} (de la prefixo \textit{retro}) = per iro inversa del avanco,\nl
desavance) : \VUgras{Marchar retroe esas kontrenatura e nefacila} (2).

%% aucun blanc au fac-simile
Kompreneble omna qualifikanta adjektivi povas e darfas\nl
genitar adverbo : saja, \VUgras{saje}; forta, \VUgras{forte}; avara, \VUgras{avare};\nl
agreabla, \VUgras{agreable}; posibla, \VUgras{posible} (3); egoista, \VUgras{egoiste};\nl
kordiala, \VUgras{kordiale}; puerala, \VUgras{puerale}; instruktiva, \VUgras{instruk}\cc
\VUgras{tive}; doloroza, \VUgras{doloroze}, e. c., e. c.

%% aucun blanc au fac-simile
Anke la participi kun -e esas fakte adverbi di maniero :\nl
\textit{manjante, kantante; laborinte, dorminte; departonte, mor}\cc
\textit{tonte; batate, karezate; ekpulsite, judiciote}, e. c.

\VUblanc{5.86pt}% mesure au fac-simile
55. --- L'\textit{adverbi di afirmo, nego o dubo} esas :

%% aucun blanc au fac-simile
\VUgras{Yes} = la sama Angla vorto; D. \textit{ja}; F. \textit{oui}; I. \textit{si}.

%% aucun blanc au fac-simile
\VUgras{No} = D. \textit{nein}; E. \textit{no}; F. \textit{non}; I. \textit{no}. Ol uzesas nur kande\parplein
\end{VUpage}

% -------------------------------------------------------------------
%  feuillet 71 — folio 67 : 22 lignes de corps (dont un TITRE), 17 de
%  note. La page ouvre sur la suite de l'alinea du folio 66.
%  Blancs du titre MESURES sur les lignes de base, avant composition :
%  pas courant 41,46 px ; 123,3 px avant le titre, 90,4 px apres, soit
%  81,8 px (6,92 mm) et 48,9 px (4,14 mm) a ajouter au pas.
%  Corps et interlettrage : outils/apparat.py 71 --regle PREPOZICIONI. 13
%  --gras -> capitale 30 px = 2,54 mm -> 10,63 pt, interlettrage -45.
% -------------------------------------------------------------------
\begin{VUpage}[71]{67}
\VUnotes{113.74mm}{%
(1) Remarkez la granda difero inter « \VUgras{tote ne vera} » e « \VUgras{ne tote}\nl
\VUgras{vera}. Altr exemplo : \VUgras{ta libro esas tote ne interesanta} (nule inter\cc
esanta); \VUgras{ol esas mem ne tote korekta} (ol ne esas tote korekta).

En Progreso : « Ne » freque renkontresas en l'expresuro \textit{ne nur...}\nl
\textit{ma (anke).} » La du parti di ta expresuro devas pozesar singla neme\cc
diate \textit{avan} la parti di frazo quan li relatas, e qui esas quaza alter\cc
nativi. Ex. : « \textit{Ne nur me hungras, ma anke mea kamaradi durstas.}\nl
(Remarkez ke, kande la verbo esas komuna a la du parti, ol pozesas\nl
ye la fino. Do la komenco (en l'unesma exemplo) : \textit{Ne nur me}\nl
\textit{hungras...} igas previdar en la duesma parto altra verbo kam\nl
\textit{hungrar}). --- \textit{Me ne nur hungras, ma anke durstas}. --- \textit{Me deziras ne}\nl
\textit{nur manjar, ma anke drinkar}. --- \textit{Me deziras drinkar ne nur aquo,}\nl
\textit{ma anke vino.} » (VII, 484.)

(2) \textit{Eble ebla, ebla eble}, e mem : \textit{eble eble} esas nur nedivinebla\nl
stranjaji, per qui Esperanto tradukas nia : \textit{forsan posibla, posibla}\nl
\textit{forsan}, e \textit{forsan posible}. La elemento \textit{ebl} anke en ta linguo esas\nl
sufixo, quale en Ido : \textit{kredebla, videbla} \VUgras{e. c.}%
}

\VUcontinue
ol konstitucas per su respondo izolita e kompleta. Ex. : \VUgras{Kad}\nl
\VUgras{vu ekiros posdimeze? No; me restos heme.}

%% aucun blanc au fac-simile
\VUgras{Ne}, qua sempre devas preirar sen ul intervalo la vorto\nl
negata : \VUgras{Vu ne facis ta laboro} (quan vu devis facar). ---\nl
\VUgras{Ne vu facis ta laboro} (ma altru). --- \VUgras{Vu facis ne ta laboro,}\nl
\VUgras{ma altra tote neutila.} --- \VUgras{Il esas ne richa, ma tamen ne}\nl
\VUgras{povra.} --- \VUgras{To esas ne tote vera} (1).

%% aucun blanc au fac-simile
\VUgras{Forsan} = E. \textit{perhaps}; D. \textit{vielleicht}; F. \textit{peut-être}; I. \textit{forse};\nl
S. \textit{puede ser, quiza}.

%% aucun blanc au fac-simile
L'ideo quan ol expresas tre diferas del ideo di posibleso,\nl
e mem la sama frazo povas kontenar la du idei. Ex. : \VUgras{Forsan}\nl
\VUgras{to esas posibla} (2).

\VUsaut{6.92mm}
\VUtitre{10.63pt}{-45}{PREPOZICIONI.}
\VUsaut{4.14mm}

56. --- Kun e pro sua rolo di relato-signi e ligili, la pre\cc
pozicioni esas renkontrata omnainstante che la helpolinguo.\nl
Quale la adverbi, li esas nevariebla.

%% aucun blanc au fac-simile
Ni studios aparte singla prepoziciono primitiva, indikante\nl
lua senco propra e specala, e lumizante lu per mult exempli.\nl
Nam, por la just interkompreno, tre importas, ke nia pre\cc
pozicioni uzesez kun lia senco tre exakta, malgre l'exemplo\nl
kontrea di lingui natural, pri ta gramatikal kategorio, en\nl
qua li esas vere tro richa de nelogikaji ed idiotismi.
\end{VUpage}

% -------------------------------------------------------------------
%  feuillet 72 — folio 68 : 32 lignes de corps, 6 de note.
%  Blancs d'articulation devant « 58. --- » (24,2 px) et « 59. --- »
%  (24,6 px).
%  COQUILLE DE L'ORIGINAL, conservee : note (1), « dicernebla.Ol » —
%  le point n'est pas suivi d'une espace.
% -------------------------------------------------------------------
\begin{VUpage}[72]{68}
\VUnotes{144.44mm}{%
(1) « Kande on decidis, pro eufonio, elizionar la tri vorti \textit{ad,}\nl
\textit{ed, od}, on decidis, ke on ne elizionos li (t. e. fakte nur \textit{ad}) en la\nl
kompozado, pro ke on timis, ke en la kompozaji \VUgras{a} ne esus sat\nl
rikonocebla o sat dicernebla.Ol povus intermixesar, sive kun la prepo\cc
ziciono \VUgras{a}, sive kun la finalo \VUgras{a} di l'artiklo o di l'adjektivo preiranta. »\nl
(Progreso, II, 165.)%
}

57. --- \VUgras{Ad} (o \VUgras{a}, kande eufonio permisas, ma nultempe\nl
en kompozajo) (1), uzesas por indikar la skopo di la ago,\nl
la loko quan on volas atingar, la destinario, la persono a\nl
qua on donas od atribuas ulo :

%% aucun blanc au fac-simile
Ex. : \VUgras{me iras a la rivero; me sendas to ad amiko; Hen}\cc
\VUgras{rikus donis a me multa flori; on imputis ad ilu ta ago}\nl
\VUgras{abomininda; la patrulo imperis a sua filii sequar ilu; el}\nl
\VUgras{agis a sua matro, ne quale filio, ma quale enemiko.}

%% aucun blanc au fac-simile
Konseque, \VUgras{ad} indikas logike l'objekto di ula penso o sen\cc
timento, per opozo a la « subjekto » qua havas li : \VUgras{pensar a}\nl
\VUgras{la futuro; la amo a Deo} (komparez : \VUgras{la amo di Deo}, por la\nl
homi, exemple); \VUgras{la envidio a la richi.}

%% aucun blanc au fac-simile
Se to postulesas da la klareso, \VUgras{ad} darfas esar unionita\nl
al prepozicioni \textit{en, sur, sub}. Ex. : \VUgras{l'infanto iris aden la gar}\cc
\VUgras{deno; la kato saltis adsur la tablo; la muso kuris adsub}\nl
\VUgras{la lito. Sen ad} on ne savus kad la subjekti, « infanto, kato,\nl
muso » iris, saltis, kuris, per chanjo di loko, \textit{aden la gar}\cc
\textit{deno, adsur la tablo, adsub la lito}, o kad li iris, saltis, kuris\nl
ibe, esante ja \textit{en, sur} o \textit{sub} la kozi nomata.

\VUblanc{5.82pt}% mesure au fac-simile
58. --- \VUgras{Alonge} = D. \textit{entlang}; E. \textit{along}; F. \textit{le long de};\nl
I. \textit{lungo; lunghesso}. Ex. : \VUgras{Irez alonge la hego e vu trovos}\nl
\VUgras{ye lua pedo multa violi.}

\VUblanc{5.94pt}% mesure au fac-simile
59. --- \VUgras{An} expresas relato di kontigueso o di apogo, tale\nl
ke la kozo kontaktas o preske : \VUgras{Ne restez nur apud la tablo,}\nl
\VUgras{ma sideskez an lu. Prenez la skalo qua jacas apud la muro}\nl
\VUgras{ed apogez ol an lu. La urbo stacas an la rivero, qua humi}\cc
\VUgras{digas lua muri. Me preferas lojar en domo situita an monto}\nl
\VUgras{o mem sur monto. An la parieti pendis desegnuri e pikturi.}\nl
\VUgras{La fenestri an la korto esas kelke mikra, ma la fenestri an}\nl
\VUgras{la placo esas tre granda. Ne marchez an la maro se vu ne}\nl
\VUgras{volas humidigar vua pedi.}

%% aucun blanc au fac-simile
Remarkez, ke en la lasta exemplo, \VUgras{alonge} expresus ideo tre\parplein
\end{VUpage}

% -------------------------------------------------------------------
%  feuillet 73 — folio 69 : 14 lignes de corps, 30 de note.
%  La page ouvre sur la suite de l'alinea du folio 68.
%  Blanc d'articulation devant « 60. --- » (exces 25,3 px).
%  Filet releve a la main a y = 859 px : 24,30 + (859-226) x 25,4/300.
% -------------------------------------------------------------------
\begin{VUpage}[73]{69}
\VUnotes{77.19mm}{%
(1) L'expresuri « \textit{ante tri monati} » --- « \textit{tri monati ante nun}, od\nl
altri simila, expresas du idei tre diferanta inter su. La sequanta\nl
exempli lo komprenigos bone :

\VUgras{On videskis la danjero ante tri monati} = ante ke pasis tri monati,\nl
on videskis la danjero; ma : \VUgras{on videskis la danjero tri monati ante}\nl
\VUgras{nun} = pasis tri monati depos ke on vidis la danjero.

\VUgras{Lu skribis a me ante ok dii} = lu ne vartis ok dii por skribar a me;\nl
ma : \VUgras{lu skribis a me ok dii ante nun} = pasis ok dii depos ke lu\nl
skribis a me.

\VUgras{Li konvinkis ni pri lia yuro ante un duimo de horo} = li ne bezonis\nl
un duimo de horo por konvinkar ni; ma : \VUgras{li konvinkis ni pri lia yuro}\nl
\VUgras{un duimo de horo ante nun} = pasis un duimo de horo depos ke li\nl
konvinkis ni.

\VUgras{El naskis ante 9 monati} = el ne vartis 9 monati por naskar; ma :\nl
\VUgras{el naskis 9 monati ante nun} = pasis 9 monati depos ke el naskis.

Do, se on parolas pri tempo pasinta depos la ago o fakto aludata,\nl
on devas uzar \VUgras{ante nun}, pos expresir ta tempo, quale en singla\nl
duesma frazo supere.

En \textit{Progreso} VII, 2, trovesas : « \textit{Me naracos epizodo di mea voyajo}\nl
\textit{en Amerika}, \VUgras{quar yari ante nun} » (ye la 7-ma lineo). Sis linei plu\nl
infre : « \VUgras{Duadek yari ante nun} \textit{me esis samskolano...} » Sep linei plu\nl
infre : « \textit{Anke me okupis me pri la ideo} \VUgras{duadek yari ante nun}. »\nl
Itere sur la sama pagino. (Final diskurso da P\textsuperscript{ro} \textsc{Otto Jespersen}\nl
al Idista lernanti di sua kurso pri nia linguo, en l'Universitato di\nl
\textit{Köbenhavn}.)

La granda linguisto, qua posedas admirinde la lingui Angla e\nl
Franca, tre certe konocas : \textit{sometime ago} e \textit{il y a quelque temps}\nl
di ta du lingui, ed anke \textit{vor einiger Zeit} di la Germana. Do il agas\nl
kun plena konoco e ni esos kun ilu en bona societo por dicar :\nl
\textit{Kelka tempo ante nun, du yari ante nun} e. c.%
}

\VUcontinue
altra. On povas ya marchar \textit{alonge} la maro, ye kelka metri de\nl
olu, e konseque sen ula risko di humidigo por la pedi.

\VUblanc{6.09pt}% mesure au fac-simile
60. --- \VUgras{Ante} = en tempo preirinta. Ex. : \VUgras{To eventis ante}\nl
\VUgras{vua departo, du monati ante nun. Il departis ante me, e}\nl
\VUgras{tamen il arivis pos me. Me esas certa, ke me arivos longe}\nl
\VUgras{ante vu. To eventis ante tri monati} (1).

%% aucun blanc au fac-simile
\VUgras{El mortis, tri monati ante nun, pos longa sufri. Qua pen}\cc
\VUgras{sabus, du yarcenti ante nun, ke la homi esos konkurencanta}\nl
\VUgras{la uceli per aeroplani e direktebli?}

%% aucun blanc au fac-simile
Remarkez ankore la uzo di \VUgras{ante} en la frazi sequanta :

%% aucun blanc au fac-simile
Me ne savas precize kande me departos, forsan erste pos\nl
un monato, ma forsan ante tri semani : to dependas de la\nl
retroveno di mea spozo. --- \VUgras{La mediko dicabis, ke il duros}\nl
\VUgras{vivar adminime un yaro, ma il mortis multe plu balde, pos}\parplein
\end{VUpage}

% -------------------------------------------------------------------
%  feuillet 74 — folio 70 : 14 lignes de corps, 29 de note.
%  La page ouvre sur la suite de l'alinea du folio 69.
%  Filet releve a la main a y = 879 px : 24,30 + (879-227) x 25,4/300.
%  La note (2) porte elle-meme un appel « * » et sa sous-note.
% -------------------------------------------------------------------
\begin{VUpage}[74]{70}
\VUnotes{79.50mm}{%
(1) La plaso di ulu od ulo esas, en loko, la parto destinita,\nl
atribuita ad oli, od okupata da oli. Ex. : \VUgras{En ta cirko imensa la}\nl
\VUgras{asistanti esis tante multa, ke on ne trovabus dek plusa plasi, sive}\nl
\VUgras{por sideskar, sive mem por stacar.} --- \VUgras{La ordino konsistas en donar}\nl
\VUgras{un plaso determinita a singla kozo, ed en konstante pozar singla}\nl
\VUgras{kozo en lua plaso.}

(2) Decido 1611 : Vizante super omno la facila aplikebleso e\nl
distingi necesa, l'akademio decidas unanime, ke \VUgras{ante} e \VUgras{pos} uzesez nur\nl
pri la sucedo di eventi e fakti, ed \VUgras{avan, dop} pri omno cetera. Pri la\nl
\textit{linguala unaji} on uzez \VUgras{ante, pos}, se on vizas la linguo parolata, pro\nl
ke lore la tempo koncernesas nedubeble, ma se on vizas la linguo\nl
skribata, on uzez la prepozicioni \VUgras{avan} e \VUgras{dop} (e la adverbi \VUgras{avane} e\nl
\VUgras{dope}) pro ke lore la reprezentata linguo esas en la spaco *.

* « Yen kelka frazi motivizanta la decido 1611 : Quankam ta cifri\nl
esas \VUgras{avan} la nomi dil vari, me tre kredas, ke li skribesis \VUgras{pos} ici.\nl
Nam, se la cifri esus skribita \VUgras{ante} la nomi dil vari, e ne \VUgras{pose}, li ne\nl
esus tante mikra. Ma, pro ke restis, \VUgras{avan} la nomi, kelka spaco libera,\nl
on profitis lo por insinuar pose la cifri avan la nomi. --- Vu trans\cc
skribos ta du kolumni de cifri, ica \VUgras{ante} ed \VUgras{avan} l'altra. Kad vu\nl
komprenis bone? Yes, me devas skribar ica \VUgras{ante} skribar ita, ma\nl
tale ke ol esez \VUgras{avane}. --- Quankam sur la listo dil rekompensoti la\nl
prenomi skribesis avan la nomi familiala, vu sorgos lektar li quale\nl
se li esus dop ici. Do vu pronuncos la prenomi pos l'altri, ed ici\nl
ante iti. »

En la multa serchi quin il facis por ica gramatiko en la 6 e duima\nl
yari di \textit{Progreso}, l'autoro sempre trovis la prepozicioni \textit{ante, avan}\nl
e \textit{pos, dop} uzata quale indikesas dal decido 1611.

Decido 1624 : On refuzas unanime apertar nova diskuto pri\nl
decido 1611.%
}

\VUcontinue
sis dii, ok monati ante nun. --- \VUgras{Vu certe ruinos vu ante}\nl
\VUgras{longe, forsan mem ante un yaro, se vu duros spensar fole}\nl
\VUgras{quale vu agas.}

\VUblanc{5.55pt}% mesure au fac-simile
61. --- \VUgras{Apud} = tre proxime (ma ne tam grande kam\nl
indikas \VUgras{an}; videz ica). Ex. : \VUgras{la kirko trovesas apud nia}\nl
\VUgras{hemo. Sideskez apud me. Qua stacas apud la pordo, an la}\nl
\VUgras{muro? Qua glutinis afisho an la pordo, apud nia nomal}\nl
\VUgras{plako.}

\VUblanc{5.97pt}% mesure au fac-simile
62. --- \VUgras{Avan} relatas la loko, la plaso (1) quan okupas\nl
en la spaco la enti o kozi, kontre ke \VUgras{ante} (ja vidita)\nl
relatas la tempo (2). Ol signifikas : D. \textit{vor} (örtl.); E. \textit{be}\cc
\textit{fore} (in space); F. \textit{devant}; I. \textit{avanti, davanti, dinanzi} (di\nl
luogo).

%% aucun blanc au fac-simile
Ex. : \VUgras{vu konstruktos la paviliono avan la domo ed ante}\parplein
\end{VUpage}

% -------------------------------------------------------------------
%  feuillet 75 — folio 71 : 37 lignes de corps, AUCUNE note.
%  outils/filet.py ne trouve pas de filet, et c'est ici la bonne
%  reponse : la page n'en porte pas. Ses trente-six pas se tiennent
%  tous entre 39 et 42 px, sans le pas court propre au bloc des notes.
% -------------------------------------------------------------------
\begin{VUpage}[75]{71}
\VUcontinue
ica. --- \VUgras{Lokizez la homini unesme, do ante la homuli, e pla}\cc
\VUgras{sizez eli avan ili.} --- \VUgras{La artiklo uzesas avan la substantivo}\nl
\VUgras{min ofte en Ido kam en la Franca.}

\VUblanc{5.74pt}% mesure au fac-simile
63. --- \VUgras{Che} = en la domo, habiteyo, lando, domeno (ma\cc
teriala o spiritala) di... Ex. : \VUgras{Me lojas che mea patrulo. Me}\nl
\VUgras{esis che mea onklino, nun me iras (ad) che mea kuzi. Irez}\nl
\VUgras{quik che la mediko. Che la Angli la veturi pasas sinistre}\nl
\VUgras{e che la Franci dextre. Me kompris ta mikra poshlexiko che}\nl
\VUgras{Isaac Pitman en London.}

\VUblanc{5.42pt}% mesure au fac-simile
64. --- \VUgras{Cirkum} = D. \textit{um, herum}; E. \textit{around, about};\nl
F. \textit{autour de..., environ}; I. \textit{intorno, all'intorno, vicino,}\nl
\textit{circa}; S. \textit{alreador, en contorno, cerca de} (en omna senci :\nl
loko, tempo e quanteso). Ex. : \VUgras{la hundo kuris cirkum ilu,}\nl
\VUgras{cirkum la urbo esas granda preurbi; to eventis cirkum mea}\nl
\VUgras{duadekesma yaro.}

\VUblanc{4.72pt}% mesure au fac-simile
65. --- \VUgras{Cis} = sur ica latero (ne trans olu). Ex. : \VUgras{cis la}\nl
\VUgras{rivero la tereno esas pasable sika, ma trans olu la tereno}\nl
\VUgras{esas marshoza. Venez cis la hego, ni konversos plu facile.}\nl
(Komp. \textit{trans}.)

\VUblanc{4.90pt}% mesure au fac-simile
66. --- \VUgras{Da} indikas la facanto, facinto o faconto di la\nl
ago. Konseque 1\textsuperscript{e} la komplemento dil verbo pasiva,\nl
2\textsuperscript{e} l'autoro. Ex. : \VUgras{il esas} (\textit{esis, esos, esus}) \VUgras{amata da omni.}\nl
\VUgras{La pikturi da Murillo. La poemi da Victor Hugo.} En ta\nl
exempli, « omni, Murillo, Victor Hugo » esas la \textit{aganti}, la\nl
produktanti dil ago; pro to li indikesas kom tala dal pre\cc
poziciono rezervita a la \textit{aganti} : « \VUgras{da} ». (Komp. \textit{de, di, per}).\nl
Vice : \VUgras{la amo di Deo a la homi} (videz \textit{ad}) on tre legitime\nl
darfus dicar : \textit{la amo da Deo a la homi}.

\VUblanc{5.96pt}% mesure au fac-simile
67. --- \VUgras{De} indikas la punto di \textit{de-veno} (en la spaco e\nl
tempo), l'origino, la dependo, la departal punto. Ex. : \VUgras{Ta}\nl
\VUgras{juvelo venas de mea matro. De ube vu adportas ico? de mea}\nl
\VUgras{rur-domo. La persiko esas importacita de Persia. To ne de}\cc
\VUgras{pendas de me. La treno de Paris a Lyon. De supre e de infre,}\nl
\VUgras{de omna lateri, de omna rangi sociala venas ad ilu kurajigi.}\nl
\VUgras{Me sufras de} (o \textit{pro}) \VUgras{nevralgio dentala. Li mortis single}\nl
\VUgras{de hungro} (o \textit{pro}).

%% aucun blanc au fac-simile
\VUgras{De lua nasko il sempre montris extrema sentemeso. El}\parplein
\end{VUpage}

% -------------------------------------------------------------------
%  feuillet 76 — folio 72 : 32 lignes de corps, 7 de note.
%  La page ouvre sur la suite de l'alinea du folio 71.
%  Aucune de ses onze frontieres d'alinea ne porte de blanc.
% -------------------------------------------------------------------
\begin{VUpage}[76]{72}
\VUnotes{141.31mm}{%
(1) On darfas supresar, sen detrimento, \VUgras{de} pos la substantivo di\nl
quanteso : \VUgras{taso teo, metro drapo}, dicas la « Grammaire Complète »\nl
ma, segun mea konoco, on ne uzis ta darfo.

(2) Same kam on dicas : \VUgras{botelo de vino}, tale on havas la darfo\nl
dicar : \VUgras{botelo plena de vino (botelo de vino, plena)}. Remarkez, ke\nl
\VUgras{botelo de vino} ne equivalas \VUgras{botelo por vino, vinbotelo}; nam\nl
l'unesma povas esar okazione \VUgras{birbotelo}, quan on plenigis per vino.%
}

\VUcontinue
esas malada de tri semani. \VUgras{Me savas to de longe. El havis}\nl
\VUgras{grava morbilo, tri yari ante nun, e de lore el restis tre febla.}\nl
\VUgras{De nun vu ne plus ekiros sen me. Ni ne vidis li de un yaro.}

%% aucun blanc au fac-simile
Kande on volas parolar pri la komencal punto di ulo even\cc
tinta pos dato, pos epoko indikata, on uzas \VUgras{depos} (facita\nl
ek \textit{de} e \textit{pos}). Ex. : \VUgras{El esis ofte malada depos sua mariajo.}\nl
\VUgras{Depos sua kronizo, nia suvereno livas rare la chefurbo.}

%% aucun blanc au fac-simile
\VUgras{De} uzesas anke kun la substantivi signifikanta mezuro,\nl
quanto, kontenanto : \VUgras{un metro de drapo; turbo de civili e}\nl
\VUgras{de soldati; taso de kafeo.}

%% aucun blanc au fac-simile
Nulu povas konfundar a : \VUgras{un metro} (venanta) \VUgras{ek drapo,}\nl
\VUgras{taso} (venanta) \VUgras{ek kafeo.} (Videz \VUgras{ek} plu fore.) (1).

%% aucun blanc au fac-simile
On uzas \VUgras{de} kun l'adjektivi \textit{plena} (2), \textit{longa, larja, alta,}\nl
\textit{profunda, dika}, e. c., qui fakte relatas mezuro, dimensiono :\nl
\VUgras{plena de vino, longa de sis metri, dika de kin centimetri}, e. c.

%% aucun blanc au fac-simile
Fine on uzas \VUgras{de} kun titulo di nobeleso : \VUgras{duko de Richelieu,}\nl
\VUgras{markezo de La Fayette.}

%% aucun blanc au fac-simile
\VUgras{On bone remarkez, ke de esas neutila e devas ne uzesar}\nl
\VUgras{kun la quantesal adjektivi ed adverbi. Ex. : multa homi;}\nl
\VUgras{poka vorti; quanta invititi?} (L'expresuri \textit{multe de homi,}\nl
\textit{poke de vorti, quante de invititi?} esus galicismi tam kontre\cc
logika kam \textit{tri de homi}.)

%% aucun blanc au fac-simile
La prepoziciono \VUgras{de} darfas kombinesar kun altri por in\cc
dikar la loko de qua on venas : \VUgras{la muso saltis desub la}\nl
\VUgras{tablo adsub la lito, e desub la lito ol fugis aden la kameno.}\nl
(Komp. \VUgras{da, di.})

%% aucun blanc au fac-simile
Esas remarkenda la tre preciza distingo per \VUgras{da} e \VUgras{de}, quan\nl
on obtenas kun la verbi signifikanta \textit{recevar, komprar,}\nl
\textit{aquirar} ed altri analoga : \VUgras{Me recevis ta libro de Alexander,}\nl
\VUgras{or ta libro esas di Ioannes; do me recevis la libro di}\nl
\VUgras{Ioannes de Alexander.} --- \VUgras{Ta varo esis komprata da me} (\textit{me}\nl
\textit{kompris lu}); \VUgras{komprata de me} (\textit{me vendis lu}). --- \VUgras{Ta kavalo}\parplein
\end{VUpage}

% -------------------------------------------------------------------
%  feuillet 77 — folio 73 : 23 lignes de corps, 18 de note.
%  La page ouvre sur la suite de l'alinea du folio 72.
%  Blanc d'articulation devant « 68. --- » (exces 23,9 px).
%  La note (2) se poursuit au folio 74 (elle y occupe quatre lignes
%  avant la note (1) de cette page-la).
% -------------------------------------------------------------------
\begin{VUpage}[77]{73}
\VUnotes{110.24mm}{%
(1) Remarkez, ke se la libristo esus l'autoro di la naraci, on dicus:\nl
\VUgras{da} e ne \VUgras{de.}

(2) Sen plajiar Ido, la linguo da Zamenhof (se ni dicus : \textit{di}\nl
Zamenhof, ol esus la Polona) nultempe povus per sua unika \VUgras{de}\nl
tradukar klare e naturale ta exempli ed altri, quin ni ja donis e\nl
povus donar ankore. Tria-dek-e-ok yari de existo e la tota evoluciono\nl
di lua lexikografi e skripteri ne ja donis ad olu ica posibleso! Ol\nl
esas ankore kondamnita, pri omna analoga kazi, ad aranji plu o\nl
min stranja ed obskura, qui sakrifikas la ideo expresenda. Quon\nl
signifikas : \VUgras{la sklavo liberigita de sia mastro?} Kad : \textit{la sklavo}\nl
\textit{liberigita} \VUgras{de} \textit{sua mastro?} o : \textit{la sklavo liberigita} \VUgras{da} \textit{sua mastro?}\nl
Nul vortal ordino povas destruktar l'ambigueso esperantala en ica\nl
kazo ed en multega altri analoga. Bela produkturo di blinda e hazar\cc
dala evoluciono : 38 yari ne emendis la linguo mem pri ta punto\nl
grava e tote ne rara!

Yen ica frazo : \textit{The governement of the people by the people}. Espo\nl
necese tradukos : la regado de la popolo \VUgras{per} \textit{la popolo!} Ma hike \VUgras{per}\nl
esas tote ne justa, nam E, dicas \textit{by}, ne \textit{through}. E la tota frazo%
}

\VUcontinue
\VUgras{esis komprata da mea patrulo de mea amiko} (t. e. \textit{mea}\nl
\textit{patrulo kompris lu de mea amiko}).

%% aucun blanc au fac-simile
De co konsequas, ke on devas uzar de (nultempe \textit{ad}, quale\nl
en la Franca) por indikar la komercisto o vendinto, la\nl
persono de qua on kompris o recevis ulo : \VUgras{Me kompris de}\nl
\VUgras{ta libristo} (ne : \textit{a ta libristo}) \VUgras{amuzanta libro a} (o \textit{por}) \VUgras{mea}\nl
\VUgras{filiineto.} --- \VUgras{Me kompris a} (o \textit{por}) \VUgras{mea filii voyajal naraci}\nl
\VUgras{de la libristo} (1) \VUgras{quan vu indikis a me.}

\VUblanc{5.75pt}% mesure au fac-simile
68. --- \VUgras{Di} indikas nur la posedo, la aparteno, o la relato\nl
generala di ul objekto (quan la genitivo expresas en la lingui\nl
flexionale) : \VUgras{la libro di Petrus; di qua esas ta domo? Di}\nl
\VUgras{mea patrulo.} (\textit{La domo di me} = \textit{mea domo}.)

%% aucun blanc au fac-simile
Remarkez, en l'exempli sequanta, quale \VUgras{di} unionesas a\nl
\VUgras{da} por distingar tre klare du diversa relati, qui sen ta du\nl
prepozicioni restus konfuza e konfundebla :

%% aucun blanc au fac-simile
\VUgras{La konquesto di Anglia da la Normandi igis la duki}\nl
\VUgras{di Normandia rivala kun la reji di Francia.} --- \VUgras{La sendo}\nl
\VUgras{di ta letro da Petrus a Ioannes efektigis la deskonkordo qua}\nl
\VUgras{nun regnas inter li.} --- \VUgras{La religial libri di la Kristani omna}\cc
\VUgras{eklezia konstante parolas pri la amo di} (o \VUgras{da}) \VUgras{Deo a la}\nl
\VUgras{homi e pri la amo di} (o \VUgras{da}) \VUgras{la homi a Deo ed al proximo.}\nl
--- \VUgras{La konkordato obtenita del papo da la rejo dil Franci}\nl
\VUgras{stipulis to tre explicite} (2).
\end{VUpage}

% -------------------------------------------------------------------
%  feuillet 78 — folio 74 : 33 lignes de corps, 6 de note.
%  Le bloc de notes OUVRE sur la fin de la note (2) du folio 73 :
%  quatre lignes sans renfoncement, avant la note (1) propre.
%  Blancs d'articulation devant « 69. --- » (15,4 px) et
%  « 70. --- » (16,0 px).
% -------------------------------------------------------------------
\begin{VUpage}[78]{74}
\VUnotes{144.14mm}{%
\VUcontinue
tradukenda esas : « ...ke la guvernado \textit{di} l'populo \textit{da} l'populo e por\nl
la populo ne perisos sur la tero ». Se on turnos altre ta frazo, quo\nl
restos de l'ideo? Vere quon valoras linguo nekapabla tradukar idei\nl
tante simpla sen alterar oli?

(1) Kompreneble, se on volas dicar ke el ne cesis esar malada tam\nl
longe kam duris elua infanteso, on dicas : \textit{dum sua} \VUgras{tota} \textit{infanteso},%
}

\VUgras{Ido povas distingar tote certe l'autori, la modeli e la pro}\cc
\VUgras{prieteri di artal verko, di portreti, statui, e. c. Ex. : la por}\cc
\VUgras{treti da Rafaël} (il facis li); \VUgras{la portreti de Napoléon} (il esis\nl
la modelo, oli reprezentas lu); \VUgras{la portreti di siorulo X...} (il\nl
esas lia proprietero, il posedas li kom kolektero od altre).

%% aucun blanc au fac-simile
« Quankam la tri (prepozicioni \textit{da, de, di}) esas necesa,\nl
existas por li, quale por mult altra propozicioni, kazo-limiti,\nl
en qui on povas tre juste hezitar inter du. Por levar ta dubi,\nl
ni povas donar la praktikal regulo sequanta.

%% aucun blanc au fac-simile
« Se la senco postulas klare l'ideo di l'aganto od autoro,\nl
uzez \textit{da}; --- se ol postulas klare l'ideo di \textit{de}veno, o di kon\cc
teno o di quanto, uzez \textit{de}; --- en la cetera kazi (do en omna\nl
dubebla kazi) uzez audace \textit{di}; nam olca esas la maxim\nl
\textit{generala} de la tri, korespondanta a la genitivo; do, se vu\nl
hezitas inter la tri, vu povas uzar ol prefere; tale vu riskos\nl
minime erorar, e vu esos komprenata » (\textit{Progreso}, II, 33).

\VUblanc{3.70pt}% mesure au fac-simile
69. --- \VUgras{Dop} relatas nur la spaco, quale \VUgras{avan} (videz ica)\nl
di qua ol esas la justa kontreajo : \VUgras{Il departis avan me, me}\nl
\VUgras{sequis e balde preterpasis lu, tale ke il arivis dop me. Iosef}\nl
\VUgras{iris al tribunalo por prizentar su avan la judiciisto; ma il}\nl
\VUgras{arivis longe ante ica e mustis vartar.} --- \VUgras{En lineo de soldati}\nl
\VUgras{qui marchas ad ni, l'unesma esas avan la duesma ed ica}\nl
\VUgras{dop l'unesma; la duesma esas avan la triesma ed ica dop}\nl
\VUgras{la duesma, e. c.} --- \VUgras{Se me lektas, irante de sinistre ad dextre}\nl
\VUgras{quale en l'ocidental lingui, la vorto} \textit{l'infanto}, \VUgras{l'artiklo esas}\nl
\textit{avan} \VUgras{i ed ica} \textit{dop} \VUgras{l'artiklo; la i esas} \textit{avan} \VUgras{n ed ica} \textit{dop} \VUgras{i, e. c.}\nl
\VUgras{Kontraste, se me lektus ta vorto} (quale en la Hebrea) \VUgras{de}\nl
\VUgras{dextre ad sinistre, la o esus} \textit{avan} \VUgras{i ed ica} \textit{dop} \VUgras{o, la t esus}\nl
\textit{avan} \VUgras{n ed ica} \textit{dop} \VUgras{t, e. c.} (Videz, ye \VUgras{avan}, la decido 1611).

\VUblanc{3.85pt}% mesure au fac-simile
70. --- \VUgras{Dum} indikas la duro di tempo en qua ulo eventas\nl
od esas facata, agata : \VUgras{il dormis dum la koncerto; il esis}\nl
\VUgras{absenta dum tri yari; il sucesis fugar dum la masakro; el}\nl
\VUgras{esis malada dum sua infanteso} (1).
\end{VUpage}

% -------------------------------------------------------------------
%  feuillet 79 — folio 75 : 39 lignes de corps, AUCUNE note.
%  Blanc d'articulation devant « 71. --- » (exces 26,4 px).
% -------------------------------------------------------------------
\begin{VUpage}[79]{75}
Nultempe tacez ta prepoziciono, malgre l'exemplo di ula\nl
lingui, di la Franca exemple, qua tre logikale dicas : \textit{il}\nl
\textit{guvernis} \VUgras{dum} \textit{kin yari la kolonio} e \textit{il studiis} \VUgras{dum} \textit{tri yari}\nl
\textit{filozofio}, ma qua dicas anke, e kontrelogike : \textit{il guvernis kin}\nl
\textit{yari la kolonio}, e : \textit{il studiis tri yari filozofio}. Ma kad on\nl
guvernas yari o lando e homi? Kad on studias yari, o cienco,\nl
arto, mestiero? Por ne dicar absurdajo, kompromisar la justa\nl
kompreno, on mustus avertar (quale agas nacionala gra\cc
matiki) ke en ta expresuri e le analoga, « kin yari » e « tri\nl
yari » esas nule direta komplimenti, ma cirkonstancal kom\cc
plementi tempala, pro ke \VUgras{dum} esas tacita. Ka ne esas pre\cc
ferebla, kom plu sekura, nultempe tacar la prepoziciono e\nl
dicar kun olu : \textit{la kin yari} \VUgras{dum} \textit{qui il guvernis la kolonio}\nl
\textit{esis le maxim prosperanta por olu; me studiis filozofio} \VUgras{dum}\nl
\textit{tri yari e la general historio} \VUgras{dum} \textit{quar yari; il retromarchis}\nl
\textit{o marchis retroe} \VUgras{dum} \textit{un mi-horo pro pario stupida quan}\nl
\textit{il facabis?}

%% aucun blanc au fac-simile
Nultempe uzez \textit{durante} vice \VUgras{dum}; nam \textit{durante} ne esas\nl
prepoziciono. Ex. : \VUgras{la diskurso durante tre longe ol certe}\nl
\VUgras{tedos l'askoltanti.}

\VUblanc{6.37pt}% mesure au fac-simile
71. --- \VUgras{Ek} = de interne ad extere di... Ex. : \VUgras{la hundo}\nl
\VUgras{saltis ek la barelo} (lu esis en la barelo); \VUgras{la hundo saltis}\nl
\VUgras{de la barelo} (lu esis nur \textit{an} la barelo). Do atencez ne kon\cc
fundar ek a de; nam, quale vu vidas, \textit{ek} supozas, ke on esis\nl
\textit{en} la loko, kontre ke \textit{de} supozas, ke on esis nur \textit{an} od\nl
\textit{apud} olu.

%% aucun blanc au fac-simile
Per extenso naturala \VUgras{ek} uzesas anke por indikar la materio\nl
di ula kozo (ek qua on imaginas, ke ol esas extraktita) : \VUgras{vazo}\nl
\VUgras{ek oro}; ma esas preferinda dicar : \VUgras{ora vazo. Domo kon}\cc
\VUgras{struktita ek petro. Quon vu facas ek ico? Quon vu facas ek}\nl
\VUgras{vua pekunio?} Ne konfundez \VUgras{ek} al punto di \textit{de}veno. Dicez :\nl
\VUgras{Me recevis letro} \textit{de} \VUgras{Paris}, ne \textit{ek Paris}.

%% aucun blanc au fac-simile
Atencez bone la difero inter : \VUgras{botelo ek...} (materio),\nl
\VUgras{botelo de...} (kontenajo), \VUgras{botelo por...} (destineso). Ex. :\nl
\VUgras{botelo ek vitro, botelo de vino, botelo por oleo.}

%% aucun blanc au fac-simile
Pos la verbo \textit{konsistar}, on uzas \VUgras{en} o \VUgras{ek}. Ex. : \VUgras{La kurajo}\nl
\VUgras{ne konsistas en ne sentar ula timo, ma plu juste en ne egardar}\nl
\VUgras{la danjero. Ek quo konsistas la afero?}

%% aucun blanc au fac-simile
\VUgras{Ek darfas uzesar metafore por indikar objekto o mem}\parplein
\end{VUpage}

% -------------------------------------------------------------------
%  feuillet 80 — folio 76 : 28 lignes de corps, 9 de note.
%  UNE LIGNE ECHAPPE AU DETECTEUR : « kazo. », deux syllabes au bas
%  d'un alinea. Elle tombe DANS LE CORPS, et non dans les notes comme
%  au folio 65 : les indices de toutes les frontieres suivantes
%  glissent d'une unite, et poser_cardage ne peut pas servir ici. Les
%  quatre blancs sont donc calcules a la main sur les lignes de base,
%  la ligne manquante retablie a son pas :
%     avant « 72. » 105,9-41,5-41,5 = 23,9 px ; avant « 73. » 27,9 px ;
%     avant « 74. » 23,7 px ; avant « 75. » 28,6 px.
% -------------------------------------------------------------------
\begin{VUpage}[80]{76}
\VUnotes{134.28mm}{%
(1) Esas kazi ube plura prepozicioni povas expresar logike la\nl
relato. Do tote senyure e nejuste la helpolinguo impozus ica prefere\nl
kam ita. Ex. : \VUgras{sufrar de} (o \textit{pro}) \VUgras{dursto; en} o \VUgras{dum la jorno.}

(2) Remarkez la difero inter : \VUgras{me venos erste morge} e \textit{me venos}\nl
\textit{nur morge}. L'unesma dicas : \VUgras{me venos ne plu balde kam morge;}\nl
la duesma dicas : \textit{me venos} (unike) \textit{morge} e \textit{ne en altra dio}.

(3) Quale nia omna nocioni, la nociono di \textit{for} esas relativa, tot\cc
same kam la nociono di \textit{proxim}. Exemple, \textit{forirar} povas okazione\nl
signifikar nur \textit{eskartar su} (de) : \VUgras{forirez del tablo.}%
}

\VUcontinue
ento apartenanta ad ensemblo, kolektajo (e quin on supozas\nl
\textit{extraktita} ek olu), exemple pri relatanta superlativo : \VUgras{la}\nl
\VUgras{maxim richa ek omni.} Ma on darfas anke uzar \VUgras{de}, quale ni\nl
vidis ye : \textit{gradi komparala}. Mem \VUgras{inter} anke konvenas en ta\nl
kazo.

\VUblanc{5.75pt}% mesure a la main (ligne « kazo. » non detectee)
72. --- \VUgras{En} = ye l'internajo di (loko, o di to quon meta\cc
fore on komparas a loko). Ta prepoziciono relatas la spaco\nl
e la tempo. Ex. : \VUgras{Il marchas en la chambro. En printempo}\nl
\VUgras{l'arbori florifas. En la jorno il dormas} (1). \VUgras{Jetez ico en la}\nl
\VUgras{fosato.} Dicar « \textit{aden} la fosato » ne esus erora, ma vere tote\nl
neutila, pro ke « jetez » indikas sat bone per su, ke « ico »\nl
ne ja esis en la fosato; konseque la chanjo di loko ne bezonas\nl
indikesar duesmafoye per « aden ».

\VUblanc{6.71pt}% mesure a la main
73. --- \VUgras{Erste} = ne plu balde (kam). Ica defino mon\cc
tras, ke ta vorto esas adverbo. Pro to ol trovesas inter l'ad\cc
verbi tempala. Ma, pro ke ol semblas a kelki quaza prepo\cc
ziciono, ni repetas lu hike kun exempli : \VUgras{Me venos erste}\nl
\VUgras{morge} (2). \VUgras{Ja en la lasta yaro me suspektis ilu; ma erste}\nl
\VUgras{dum la somervakanco di ca yaro me havis la konfirmo di}\nl
\VUgras{mea suspekto. Komencez erste pos ni, nam vu laboras plu}\nl
\VUgras{rapide} (kam ni).

\VUblanc{5.70pt}% mesure a la main
74. --- \VUgras{Exter} = ye l'exterajo di... (sen chanjo di loko).\nl
Ex. : \VUgras{Pro quo vu restas exter la domo? Enirez e venez apud}\nl
\VUgras{ni. Il habitas exter hike, che sua gepatri. Exter konkurso.}\nl
\VUgras{Exter} metafore uzesas kun la senco « \textit{ecept(it)e} » : \VUgras{exter}\nl
\VUgras{ta legi ne existas altri.}

\VUblanc{6.88pt}% mesure a la main
75. --- \VUgras{For} indikas eskarto plu o min granda (3); ol\nl
esas la kontreajo di \textit{proxim}. Ex.: \VUgras{Il habitas pasable for}\parplein
\end{VUpage}

% -------------------------------------------------------------------
%  feuillet 81 — folio 77 : 31 lignes de corps, 7 de note.
%  La page ouvre sur la suite de l'alinea du folio 76.
% -------------------------------------------------------------------
\begin{VUpage}[81]{77}
\VUnotes{140.97mm}{%
(1) Videz \textit{Progreso}, V, 723, noto 2. La \textit{forpreni} di Espo esas nur\nl
imito di \textit{weg-} Germana, en la realeso.

(2) \VUgras{Inter} supozas plura personi o kozi. Ne dicez do : \VUgras{inter ilua}\nl
\VUgras{amikaro} (nomo singulara gramatike), ma : \VUgras{inter ilua amiki.} \textit{Du}\nl
\textit{personi o kozi} adminime esas necesa por ke on darfez uzar « inter ».

(3) En noto, p. 32 di « Grammaire complète », France lektesas :\nl
\VUgras{Kontre} sempre kontenas la ideo di opozo materiala od etikala.%
}

\VUcontinue
hike. \VUgras{Ne restez tale for ni, ma venez proxim ni, e sideskez}\nl
\VUgras{apud me, an la tablo.}

%% aucun blanc au fac-simile
Esas remarkenda, ke \textit{de}, qua per su indikas nur separo\nl
e departal punto, plu bone konvenas, kom prefixo, kam \textit{for},\nl
kande on ne intencas indikar \textit{foreso} o \textit{foriro} striktasence.\nl
Exemple, por indikar nur la departo-punto, \textit{deirar, deflugar,}\nl
\textit{detranar}, e. c., esas plu bona kam \textit{for}-. E mem \textit{forprenar} esas\nl
ne uzenda (1), nam on prenas \textit{de} o \textit{ek} ma ne \textit{for}, adminime\nl
en la senco quan expresas \textit{deprenar, ekprenar}. La vorto \textit{for}\cc
\textit{prenar} povas signifikar nur, ke on prenas (ulo) fore, e ne\nl
proxime.

\VUblanc{5.80pt}% mesure au fac-simile
76. --- \VUgras{Inter} = en la spaco o tempo qua separas du o\nl
plu multa personi, kozi. Ex. : \VUgras{Inter Francia e Rusia esas}\nl
\VUgras{Germania. Il pozis su inter ni du. El venis inter ok e dek}\nl
\VUgras{matine} (2).

%% aucun blanc au fac-simile
Metafore \VUgras{inter} uzesas por indikar partigo, kambio o reci\cc
prokeso : \VUgras{li dividis inter su dek e du pomi; li kambiis inter}\nl
\VUgras{su lia vesti; li luktis, parolis inter su} (prefere : \VUgras{li interluktis,}\nl
\VUgras{interparolis}, e. c.). (Videz ye la \textit{verbi reciproka}.)

%% aucun blanc au fac-simile
On devas ne konfundar \textit{inter} a \textit{ek} qua expresas ideo tre\nl
diferanta. (Vid. ica lasta.)

\VUblanc{4.76pt}% mesure au fac-simile
77. --- \VUgras{Kontre} = opoze a (3). Ex. : \VUgras{Kad vu votas por}\nl
\VUgras{o kontre la propozo? La Hispani kombatis longe kontre la}\nl
\VUgras{Mauri. La domo esas shirmata kontre la nordal vento. Lua}\nl
\VUgras{fenestri esas kontre (la) mei.}

%% aucun blanc au fac-simile
Ma on devas ne uzar \VUgras{kontre} por la senco \textit{apud, an, ad}\nl
(F. \textit{envers}), \textit{kun}. Ex. : \VUgras{La fenestri apud} (plu bone : \VUgras{an}) \VUgras{la}\nl
\VUgras{strado. Il agis ad me tre amikale.} Se on dicus \textit{kun me}, la\nl
senco esus : \textit{il e me agis}.

\VUblanc{5.66pt}% mesure au fac-simile
78. --- \VUgras{Koram} = en la asisto, prezenteso di... Ex. :\nl
\VUgras{koram ulu; koram me; koram notario; koram la judiciisti.}
\end{VUpage}

% -------------------------------------------------------------------
%  feuillet 82 — folio 78 : 26 lignes de corps, 14 de note.
%  La page ouvre sur la suite de l'alinea du folio 77.
%  La note (3) se poursuit au folio 79.
% -------------------------------------------------------------------
\begin{VUpage}[82]{78}
\VUnotes{121.24mm}{%
(1) En la netushebla \textit{Fundamento} da Zamenhof : \textit{resti kun leono}\nl
\textit{esas danghere.}

(2) \textit{Progreso}, VI, 513, noto 2.

(3) Ta prepoziciono genitas l'adjektivo lora : \VUgras{Ti qui vivis en ta}\nl
\VUgras{epoko savas quante teroriganta esis la lora eventi.} --- Ol genitas\nl
anke la adverbo \VUgras{lore, de lore, til lore. Se me esus avan la morto,}\nl
\VUgras{mem lore me atestus lo.}

En \textit{Progreso}, V, 28 trovesas ico : « Kelkafoye on uzas la prepo\cc
ziciono lor kun vorti qui indikas \textit{dato} o tempo, t. e. kun nomi di dio,\nl
monato, yaro. To esas nejusta : en omna tala kazi on devas uzar\nl
\textit{ye} (por preciza indiko), \textit{en} o \textit{dum} (por intervali), ex. : \textit{ye la unesma}\nl
\textit{di januaro, ye du kloki; en la printempo; dum la venonta yaro}. ---\nl
Lor povas (darfas) aplikesar nur a konkreta eventi, nam ol signi\cc
fikas exakte : « en la tempo di, samtempe kam ». Ex. : \textit{lor mea}%
}

\VUcontinue
\VUgras{Il dicis to ante me, avan vu e koram vua gepatri; me do nur}\nl
\VUgras{repetis to quon il dicis l'unesma e koram ni omna.}

\VUblanc{5.73pt}% mesure au fac-simile
79. --- \VUgras{Kun} = akompanate da..., juntite a... \VUgras{Il promenas}\nl
\VUgras{kun amiki; restar kun leono esas danjeroza} (1).

%% aucun blanc au fac-simile
Nultempe uzez \VUgras{kun} vice \VUgras{per}, nam la duesma indikas in\cc
strumento, kontre ke \VUgras{kun}, quale ni vidis, indikas akompano,\nl
uniono. Do ne dicez (imitante ula lingui) : \VUgras{il frapis me kun}\nl
\VUgras{bastono}, ma : \VUgras{il frapis me per bastono.}

%% aucun blanc au fac-simile
Memorez bone, ke la verbo \textit{konfundar} postulas \VUgras{kun} o \VUgras{ad}\nl
por sua nedireta komplemento, segun la decido 1205. Spe\cc
cala noto akompananta la decido fixigas la selekto tavorte :

%% aucun blanc au fac-simile
« Se ek plura kozi on facas pelmelo per mala o ne sufi\cc
canta distingo mentala, qua mixas iti kun ici, lore la pre\cc
poziciono uzenda esas \textit{kun}. Ex. : \VUgras{Nun il konfundas en sua}\nl
\VUgras{odio, la kulpozi kun l'inocenti, la boni kun la mali.} »

%% aucun blanc au fac-simile
« Se la spirito transportas erore la qualesi od individueso\nl
di ulu od ulo ad altra, lore la prepoziciono uzenda esas \textit{ad}.\nl
Ex. : \VUgras{Vu konfundis manekino a homo. Ne konfundez lam}\cc
\VUgras{piro a lanterno.} » (2).

%% aucun blanc au fac-simile
Ne dicez « kune kun » quale l'Esperantisti. Ta stranja\nl
expresuro ne uzesas en Ido. \textit{Kun}, simple, o \textit{samtempe kam}\nl
\textit{kune, unionite, solidare} expresas certe plu bone, segun la\nl
kazo, la idei quin l'Esperantisti inkluzas en « kune kun ».

\VUblanc{5.73pt}% mesure au fac-simile
80. --- \VUgras{Lor} = en la tempo di, samtempe kam : \VUgras{Lor vua}\nl
\VUgras{nasko; lor mea mariajo; lor la tertremo di...; en la realeso,}\nl
\VUgras{la homi esas egala nur lor sua nasko e lor sua morto} (3).
\end{VUpage}

% -------------------------------------------------------------------
%  feuillet 83 — folio 79 : 30 lignes de corps, 9 de note.
%  Le bloc de notes OUVRE sur la fin de la note (3) du folio 78 :
%  quatre lignes sans renfoncement.
% -------------------------------------------------------------------
\begin{VUpage}[83]{79}
\VUnotes{135.68mm}{%
\VUcontinue
\textit{mariajo; lor} la nasko di mea unesma filio; \textit{lor} la lasta aparo di la\nl
kometo di Halley, \textit{lor} l'explozo di la kurasnavo \textit{Liberté}. On darfas\nl
anke dicar : « dum la milito Ruso-Turka », nam hike on parolas pri\nl
ulo, qua duris kelka tempo. »

(1) Uli ne dicernas bone la difero di \textit{malgre} e \textit{tamen}.

Ica lasta ne esas prepoziciono e havas nula komplemento kon\cc
seque; ol pleas la rolo di adverbo kun ideo di kontrasto, di opozo :\nl
\VUgras{il esas povra e tamen} (il esas) \VUgras{jeneroza. Il esas malada, tamen il}\nl
\VUgras{volas laborar.} --- \textit{Tamen} = malgre to.%
}

81. --- \VUgras{Malgre} = sen impedesar da..., sen cedar a... Ex. :\nl
\VUgras{Il sucesis malgre omna obstakli; il departis malgre sua}\nl
\VUgras{matro o : malgre la impero di sua patrulo.} Quale on vidas,\nl
\VUgras{malgre} povas havar kom komplemento persono o kozo (1).

\VUblanc{5.62pt}% mesure au fac-simile
82. --- \VUgras{Per} indikas l'instrumento di la ago, to quo uzesas\nl
por produktar olu, la moyeno : \VUgras{skribar per plumo, per}\nl
\VUgras{krayono; sendar per posto; mortigar per hungro; li inter}\cc
\VUgras{batis per pugni; il suocidis per revolvero; persuadar per}\nl
\VUgras{dolceso; vu kovros la amasi de betravi per palio.}

%% aucun blanc au fac-simile
Precize pro ke \VUgras{per} indikas instrumento, moyeno, lu ne\nl
darfas uzesar avan la aganto en la komplemento dil verbo\nl
pasiva, sive ta aganto esas persono od animalo, sive ol esas\nl
kozo; nur \VUgras{da} uzesas takaze : \VUgras{lu esas tote kovrita da nivo.}\nl
Se vu hezitos inter \textit{da} e \textit{per}, chanjez la formo pasiva a\nl
formo aktiva : \VUgras{la nivo tote kovris la tekti.} Do la nivo esas\nl
la aganto en la frazo pasiva; konseque on devas uzar \textit{da}\nl
(ne \textit{per}) avan olu.

\VUblanc{5.50pt}% mesure au fac-simile
83. --- \VUgras{Po} preiras la kozo kambie donata; ol indikas\nl
equivalo : \VUgras{me kompris la domo po quaradek mil franki} =\nl
\textit{me pagis quaradek mil franki po la domo}; \VUgras{sigari po dek cen}\cc
\VUgras{timi.} --- \VUgras{Me kompris dek sigari po un franko} (\textit{sume}). Se on\nl
volas dicar : \textit{sigari di qui singla kustas un franko}, on dicas :\nl
\VUgras{sigari po un franko single. Silko po kin franki (singla)}\nl
\VUgras{metro; me kambiis mea biciklo po un plu nova. Il pagis}\nl
\VUgras{la glorio po sua vivo. Me defensus ta afero mem po mea}\nl
\VUgras{sango. Po quante vu vendas ta flori?} to esas : \VUgras{po quanta}\nl
\VUgras{preco.} Se on dicus : \VUgras{quante vu vendas ta flori?} la senco\nl
esus : \VUgras{kad vu vendas li grandaquante, mikraquante, grose o}\nl
\VUgras{detale?}

%% aucun blanc au fac-simile
Por indikar l'unajo kun qua relatas la preco on uzas la\parplein
\end{VUpage}

% -------------------------------------------------------------------
%  feuillet 84 — folio 80 : 32 lignes de corps, 5 de note.
%  UNE LIGNE ECHAPPE AU DETECTEUR : « yare. », dans le corps. Les
%  blancs sont donc calcules a la main sur les lignes de base, la
%  ligne manquante retablie a son pas ; la frontiere qui la suit
%  (« Esus neutila... ») n'en porte AUCUN.
%     avant « 84. » 24,2 px ; avant « 85. » 23,7 px ;
%     avant « 86. » 24,7 px ; avant « 87. » 23,1 px.
%  Le bord droit du feuillet porte l'ombre d'un pli : les largeurs
%  relevees des lignes 11, 12 et 14 sont gonflees par elle et ne
%  valent rien.
% -------------------------------------------------------------------
\begin{VUpage}[84]{80}
\VUnotes{146.47mm}{%
(1) Atencez la granda difero inter : \textit{me havas nulo por skribar} e :\nl
\textit{me havas nulo skribenda}. Forsan nun Esperantisti plajianta Ido\nl
facas la distingo dil du idei; ma ti qui ne ja adoptis nia \textit{end} esas\nl
tote nekapabla facar ta distingo : li konocas nur : \textit{me havas nulo}\nl
\textit{por skribar.}%
}

\VUcontinue
prepoziciono \VUgras{por}. Ex. : \VUgras{On abonas ta revuo po dek e du}\nl
\VUgras{franki por yaro. On darfas anke dicar : po dek e du franki}\nl
\VUgras{yare.}

% aucun blanc au fac-simile
Esus neutila uzar \VUgras{po} (od irg altra prepoziciono) avan\nl
komplemento direta : \VUgras{omna (o singla) libro kustas tri}\nl
\VUgras{franki} (e ne : \VUgras{po tri franki}).

\VUblanc{5.82pt}% mesure a la main (ligne « yare. » non detectee)
84. --- \VUgras{Por} indikas la skopo, la koncernato, la profitanto\nl
o nur destinario. Ex. : \VUgras{Por quo vu volas havar pekunio?}\nl
\VUgras{por komprar to quon me bezonas. On manjas por vivor, on}\nl
\VUgras{ne vivas por manjor. Me kompris ludili por mea infanti.}\nl
\VUgras{Ica letropapero esas por tu e ta kuverti por me. La evento}\nl
\VUgras{esas fortunoza por ilu, ma desfortunoza por vi. Me havas}\nl
\VUgras{nulo por skribar, nek plumo, nek krayono} (1).

% aucun blanc au fac-simile
On uzas \VUgras{por} avan l'unajo kun qua relatas la preco. (Videz\nl
\VUgras{po} pri ca punto.)

\VUblanc{5.70pt}% mesure a la main
85. --- \VUgras{Pos} relatas nur la tempo. Ol esas juste la kon\cc
treajo di \VUgras{ante}. (Videz ica.) Ol signifikas « plu fore kam... »\nl
(en la tempo). Ex. : \VUgras{To eventis pos mea departo. Me certe}\nl
\VUgras{arivos pos elu. Quon li agos pos mea morto?}

% aucun blanc au fac-simile
Me atestas, ke il dicis to \VUgras{ante} tua fratulo ed \VUgras{avan} tu, nule\nl
\VUgras{pos} ilu, quale il asertas; e me savas lo tre certe, nam omno\nl
eventis \VUgras{koram} me.

\VUblanc{5.94pt}% mesure a la main
86. --- \VUgras{Preter} = pasante apud (ulu od olu) ed irante\nl
plu fore kam (li). Ex. : \VUgras{Ni rajuntis li ye la halteyo, ma iris}\nl
\VUgras{preter li e preter omni quin ni renkontris, til ke ni atingis}\nl
\VUgras{la portuo.} --- \VUgras{Ni iris preter lia fenestri, malgre lia signi e}\nl
\VUgras{voki.} --- \VUgras{Li preter-vehis ni, ma salutis ni afable, pasante.}\nl
--- \VUgras{Pro quo camatine vu iris preter mea pordo, sen haltar}\nl
\VUgras{dum kelk instanti?}

\VUblanc{5.55pt}% mesure a la main
87. --- \VUgras{Pri} = koncerne, relate. Ex. : \VUgras{libri pri filozofio.}\nl
\VUgras{Il esas tre erudita pri historio. Parolez a ni pri la linguo}\nl
\VUgras{internaciona e pri vua voyaji.}
\end{VUpage}

% -------------------------------------------------------------------
%  feuillet 85 — folio 81 : 33 lignes de corps, 3 de note.
%  UNE LIGNE ECHAPPE AU DETECTEUR : « longa. », dans le corps. Blancs
%  calcules a la main sur les lignes de base :
%     avant « 88. » 23,0 px ; avant « 89. » 23,0 px ;
%     avant « 90. » 23,3 px ; avant « 91. » 23,5 px.
% -------------------------------------------------------------------
\begin{VUpage}[85]{81}
\VUnotes{150.62mm}{%
(1) Ne imitez Esperanto dicanta en « La feino » : « \textit{kiu estis la}\nl
\textit{plena portreto de sia patro laù} (vice \textit{per}) \textit{sia boneco kaj honesteco} »,\nl
nam vi alterus l'ideo tradukenda.%
}

Nultempe uzez, vice \VUgras{pri}, sive \VUgras{de}, malgre l'exemplo di la\nl
Latina, sive \VUgras{sur} malgre l'exemplo di ula lingui vivanta.

% aucun blanc au fac-simile
Ma tre reguloze vu darfas uzar \textit{koncerne}, \textit{relate} qui\nl
expresas la sama ideo kam \textit{pri}, ma en formo triople plu\nl
longa.

\VUblanc{5.53pt}% mesure a la main (ligne « longa. » non detectee)
88. --- \VUgras{Pro} = per efiko od efekto di... Ex. : \VUgras{Il mortis}\nl
\VUgras{pro hungro; me tremas, ne pro timo, ma pro koldeso; el}\nl
\VUgras{agas tale pro jaluzeso. Pro quo tu ploras? pro ke Petrus}\nl
\VUgras{batis me.}

% aucun blanc au fac-simile
Se on atencos, ke ta prepoziciono fakte konocigas \textit{de quo}\nl
venas la efekto produktita (exemple « hungro » en : \VUgras{il mortis}\nl
\VUgras{pro hungro}) on komprenos, ke ol havas kelka afineso\nl
kun \textit{de}. Co explikas, ke en ul okazioni, ol povas esar rem\cc
plasata dal prepoziciono \textit{de}. Ex. : \VUgras{il mortis de hungro ex}\cc
\VUgras{presas la kauzo dil morto} (\textit{hungro}) \VUgras{tam juste kam : il mortis}\nl
\VUgras{pro hungro. Same : el esas malada pro o de febro.}

\VUblanc{5.53pt}% mesure a la main
89. --- \VUgras{Proxim} = ye mikra disto de la punto indikata,\nl
ne for olu. Ex. : \VUgras{Me plantacigis kelka florarbusti proxim}\nl
\VUgras{la domo. Kande il sentis su proxim la morto.}

% aucun blanc au fac-simile
Ta prepoziciono esas la kontreajo di \VUgras{for}. Kun olu la vici\cc
neso restas min granda kam kun \VUgras{apud}. (Videz ica e \VUgras{an}, qua\nl
indikas ne nur vicineso, ma kontigueso o kontrakto.)

\VUblanc{5.60pt}% mesure a la main
90. --- \VUgras{Segun} = sen eskartar de..., konforme a... Ex. :\nl
\VUgras{ni vivez segun la nova precepto : « amez l'una l'altra, quale}\nl
\VUgras{me amis vi » e ne segun l'anciena : « okulo po okulo e}\nl
\VUgras{dento po dento. --- Il agis segun sua opiniono. To ne esas}\nl
\VUgras{permisata segun la lego. Arkitekturo segun la gusto di Re}\cc
\VUgras{nesanco. Pikturo (kopiuro) segun Rafael. Nultempe uzez}\nl
\VUgras{segun vice per. Do dicez : el esis la perfekta portreto di sua}\nl
\VUgras{patrulo per} (e ne : \textit{segun}) \VUgras{sua boneso e honesteso} (1).

\VUblanc{5.65pt}% mesure a la main
91. --- \VUgras{Sen} esas la kontreajo di \VUgras{kun}. (Videz ica.) \VUgras{Ol in}\cc
\VUgras{dikas l'absenteso, la manko dil persono o kozo nomata :}\nl
\VUgras{Il arivis sen sua amiko; me ne povus vivar sen tu.}
\end{VUpage}

% -------------------------------------------------------------------
%  feuillet 86 — folio 82 : 33 lignes de corps, 3 de note.
%  DERNIERE LIGNE NON DETECTEE : « pos til. », qui acheve la note (2).
%  Signalee par outils/lignes_manquantes.py (868 px d'encre sous la
%  derniere ligne detectee, contre 4 a 65 px sur les feuillets sains).
%  Elle est la DERNIERE de la page : aucune frontiere du corps ne la
%  suit, poser_cardage reste donc sur.
% -------------------------------------------------------------------
\begin{VUpage}[86]{82}
\VUnotes{151.98mm}{%
(1) O « asumar afero » egale bona, se mem ne plu bona.\nl
(2) Segun \textit{Progreso}, IV, 661. --- Videz, p. 224, \VUgras{tra} e \VUgras{trans}, omisita\nl
pos \VUgras{til.}%
}

92. --- \VUgras{Sub} indikas la situeso di ulo (od ulu) relate to\nl
quo esas supere ed en la sama vertikal direciono : \VUgras{la kato}\nl
\VUgras{dormas sub la tablo. Me refujis sub la hangaro pro la pluvo}\nl
\VUgras{qua faleskis. La muso kuris (ad) sub la armoro.} (Komp. \textit{sur}).

\VUblanc{6.82pt}% mesure au fac-simile
93. -- \VUgras{Super} (sen kontakto kun la objekto) = D. \textit{über,}\nl
\textit{oberhalb}; E. \textit{over, above}; F. \textit{au-dessus de}; I. \textit{sopra, al disopra};\nl
S. \textit{por encima de sobre}. Ex. : \VUgras{qua nombrizos la steli qui}\nl
\VUgras{brilas en la cielo super ni? La vento pulsis l'aeroplano}\nl
\VUgras{super la maro} (o, se to esas necesa por indikar translaco,\nl
\textit{adsuper} la maro). \VUgras{La muevi flugas super la maro ed ofte}\nl
\VUgras{pozas su lejere sur olua ondi.}

% aucun blanc au fac-simile
Atencez ne uzar \textit{super} o \textit{sur} vice \textit{pri}.

\VUblanc{6.66pt}% mesure au fac-simile
94. --- \VUgras{Sur} indikas la situeso di ulo (od ulu) relate to\nl
quo esas plu infre, en kontakto kun olu ed en la sama\nl
direciono : \VUgras{La navi vehigas sur l'oceano homi, bagaji e vari.}\nl
\VUgras{La kato saltis adsur la tablo por kaptar olu.}

% aucun blanc au fac-simile
Metafore : \VUgras{havar autoritato sur judiciisti; havar yuri sur}\nl
\VUgras{la krono; prenar afero sur su} (1). (Komp. \textit{sub}.)

\VUblanc{5.58pt}% mesure au fac-simile
95. --- \VUgras{Til} indikas la termino en la spaco o tempo. Ex. :\nl
\VUgras{Ni irez til la frontiero; vartez til mea retroveno.}

% aucun blanc au fac-simile
\VUgras{Til} uzesas kun \VUgras{de} por indikar spacal o tempal intervalo :\nl
\VUgras{de lundio til jovdio; de Calais til Dover.} Ol indikas anke\nl
minimo : \VUgras{il spensis de cent til duacent franki}, o simple :\nl
\VUgras{il spensis til duacent franki.}

% aucun blanc au fac-simile
\VUgras{Til} signifikas, ke on atingas la limito indikata. Kande\nl
la limito indikata esas, ne punto, ma ula intervalo, on devas\nl
kompletigar l'indiko, dicante precize : « til la komenco o\nl
fino di... » od : « til la yaro 1912 \textit{exkluzite} o \textit{inkluzite} (od\nl
\textit{exkluzita, inkluzita}) » (2).

\VUblanc{5.94pt}% mesure au fac-simile
96. --- \VUgras{Ultre} = adjunte ad : \VUgras{Ultre mea matrala linguo,}\nl
\VUgras{me savas la Germana. Me esis tre charjita, nam ultre mea}\nl
\VUgras{valizo, me portis du grosa paki.}

% aucun blanc au fac-simile
\VUgras{Ultre esas nur prepoziciono. Ni donez un plusa exemplo :}
\end{VUpage}

% -------------------------------------------------------------------
%  feuillet 87 — folio 83 : 22 lignes de corps, 18 de note.
%  La page ouvre sur la suite de l'alinea du folio 82.
%  Le detecteur soude la derniere ligne de note (« VI, 481. ») aux
%  jambages de celle qui la precede : sa largeur relevee (1081 px) ne
%  vaut rien, mais le COMPTE des lignes est juste.
% -------------------------------------------------------------------
\begin{VUpage}[87]{83}
\VUnotes{110.00mm}{%
(1) \textit{Ultre} diferas de \textit{plus}. Advere de la matematikal vidpunto la\nl
difero ne esas granda, pro ke A + B = B + A (e mem ne sempre!).\nl
Ma, en la komuna vivo, l'ordino ne esas indiferenta, sive pro la\nl
tempo, sive pro l'importo relativa e tre neegala di l'adicionendi.\nl
Exemple, la supera frazo equivalas : « Me prenis mea repasto \textit{plus}\nl
mea kafeo. » Esus ridinda dicar : « Me prenis kafeo \textit{plus} mea\nl
repasto », kontre ke on povas dicar tre bone : « Me prenis kafeo\nl
\textit{ultre} mea repasto. »

De to konsequas, ke nultempe on bezonas adverbigar \textit{ultre} : la\nl
adverbo uzenda esas sempre \textit{pluse}, qua anuncas l'adjunto di \textit{altra}\nl
kozo a l'unesma e predicita. Exemple, F. \textit{en outre} devas sempre\nl
tradukesar per \textit{pluse}. Do on ne bezonas desquietesar, kad \textit{ultre} povas\nl
uzesar kom adverbo e tale divenar kelkafoye dusenca. Exemple :\nl
« Ultre la festino eventis gaya balo. » Ula kritikemi kredas, ke la\nl
frazo esas bisenca til la streko. No, nam en l'altra senco on devus\nl
dicar : \textit{Pluse}, la festino eventis... » (e o \textit{kun gaya balo}). On havas\nl
nul ambigueso timenda, se on uzas la justa prepoziciono. (\textit{Progreso},\nl
VI, 481.)%
}

\VUcontinue
« \VUgras{Ultre mea repasto, me prenis kafeo kun glaseto de bran}\cc
\VUgras{dio} » (1).

% aucun blanc au fac-simile
Ne konfundez \textit{ultre} a \textit{exter, trans} o \textit{ecepte}.

\VUblanc{5.79pt}% mesure au fac-simile
97. --- \VUgras{Vice} = remplase...; od indikas ago kontrea, quan\nl
on opozas ad altra. Ex. : \VUgras{Il parolis vice la prezidero; il ludas}\nl
\VUgras{vice laborar.}

% aucun blanc au fac-simile
Ta prepoziciono uzesas kom prefixo, en sua senco inter\cc
naciona : \VUgras{vice-administrero, vice-prezidero, vice-sekretario,}\nl
\VUgras{vice-rejo, vice-rejeso}, e. c.

% aucun blanc au fac-simile
Ol ludas anke la rolo di adverbo. Ex. : \VUgras{Pro ke la prezi}\cc
\VUgras{dero esis absenta, me parolis vice} (lu). \VUgras{Il devus skribar}\nl
\VUgras{sua letri, vice} (to) \VUgras{il babilas.}

\VUblanc{5.95pt}% mesure au fac-simile
98. --- \VUgras{Ye} esas prepoziciono di senco nedeterminita,\nl
quan on uzas nur en la kazi ube nul altra prepoziciono\nl
postulesas da la senco. Ol indikas nome la loko o la dato\nl
exakta di evento, di fakto. Ex. : \VUgras{Ye l'angulo di la strado;}\nl
\VUgras{ye la dekesma kilometro; ye dimezo; ye la lasta foyo.}

% aucun blanc au fac-simile
Pro ke ol indikas la loko, on uzas lu por precizigar la\nl
koncernata parto di la korpo : \VUgras{me doloras ye la kapo; il}\nl
\VUgras{prenis elu ye la tayo; il kaptis la kavalo per lazo ye la kolo.}\nl
Esus tote ne justa dicar, quale la Franca : \textit{per} la tayo, \textit{per}\nl
la kolo, nam nek la tayo, nek la kolo esas l'instrumento di\parplein
\end{VUpage}

% -------------------------------------------------------------------
%  feuillet 88 — folio 84 : 30 lignes de corps, 11 de note.
%  La page ouvre sur la suite de l'alinea du folio 83.
%  Ses blancs d'articulation sont plus etroits qu'ailleurs (11 a 16 px
%  contre 23 a 28) : la page est dense, le compositeur a serre.
% -------------------------------------------------------------------
\begin{VUpage}[88]{84}
\VUnotes{134.44mm}{%
(1) Nula prepoziciono formacas verbo direte, pro ke ol ne kontenas\nl
en su l'ideo verbala. Mem \VUgras{cirkumar} esas forjetita dal Akademio\nl
pro ta motivo. Nun vice lu on uzas \VUgras{cirkondar}. Ma co ne supresis\nl
\VUgras{cirkumajo} tote analoga a \textit{exterajo, internajo} e. c.

(2) La difero inter \textit{avano, dopo} e \textit{avanajo, dopajo} (od altri\nl
simila) esas ke per \VUgras{-o} en ta vorti on indikas nur la punto avana,\nl
dopa e. c., kontre ke per \textit{-ajo} on indikas parto plu vasta o grosa.

(3) On uzas prefere \textit{suprajo, infrajo} por indikar la geometrial\nl
extremajo, surfaco o pinto. Ex. : \textit{la supro di la komodo esas hori}\cc
\textit{zontala; la suprajo di la komodo esas marmora} (o : \VUgras{ek marmoro}).\nl
To esas nur nuanco.%
}

\VUcontinue
la ago; or ni vidis ke \textit{per} indikas sempre l'instrumento, la\nl
moyeno uzata.

\VUblanc{2.78pt}% mesure au fac-simile
99. --- Kande la prepoziciono uzata ne indikas per su\nl
la chanjo di loko e nulo altra komprenigas la translaco,\nl
\textit{ma nur lore}, on montras ta chanjo adjuntante \textit{ad} a la pre\cc
poziciono : \textit{ad-en, ad-sur, ad-sub, ad-super}, e. c., quale ni\nl
ja vidis. On darfas uzar o ne la streketo, do skribar \textit{ad-en} o\nl
\textit{aden}, e. c. Ex. : \VUgras{Il kuris de la salono aden la koqueyo.} Kun\nl
\textit{ad} sola on ne savus kad il eniris o ne en la koqueyo. Ma\nl
me dicos : \VUgras{pozez la lampo sur la tablo} (sen \textit{ad}), nam videble\nl
la lampo ne esis antee sur la tablo, altre me ne dicus pozar\nl
lu sur ta moblo. Agez tale en omna kazi analoga.

\VUblanc{3.82pt}% mesure au fac-simile
100. --- La prepozicioni formacas derivaji (per \VUgras{-a, -e, -o})\nl
se la signifiko lo permisas (1). Ex. : \VUgras{avana, avane, avano;}\nl
\VUgras{dopa, dope, dopo} (2); \VUgras{apuda, apude; cisa, cise; transa,}\nl
\VUgras{transe; fora, fore; proxima, proxime; suba, sube; sura, sure;}\nl
\VUgras{supera, supere; extera, extere; dume; antea, antee; kontrea,}\nl
\VUgras{kontree; retroa, retroe}, e. c.

%% aucun blanc au fac-simile
Quale on vidis per la supera exempli, nula prepoziciono\nl
(mem ti qui finas per \textit{e}) darfas uzesar kom adverbo sen\nl
adjuntar l'adverbal dezinenco \VUgras{-e} (kontree, antee).

\VUblanc{3.58pt}% mesure au fac-simile
101. --- On ne dicas \VUgras{ene} ma \VUgras{interne} (de l'adjektivo \textit{in}\cc
\textit{terna}); nek \VUgras{eke}, ma \VUgras{extere}, nek \VUgras{pere}, ma \textit{mediace} (\textit{di}); pro\nl
ke la vorti \textit{interne, extere, mediace} expresas l'ideo adminime\nl
tam bone ed esas quik komprenata da milioni de homi, quin\nl
\textit{ene, eke, pere} astonus e trublegus.

%% aucun blanc au fac-simile
Ne konfundez \VUgras{supere} a \VUgras{supre}, nek \VUgras{sube} ad \VUgras{infre} o \VUgras{base}.

%% aucun blanc au fac-simile
\VUgras{Supre} ed \VUgras{infre} havas senco \textit{absoluta}, ed indikas loki o parti\nl
en determinita objekto : la \textit{supro} (o suprajo), la \textit{infro} (o \textit{in}\cc
\textit{frajo}) (3).
\end{VUpage}

% -------------------------------------------------------------------
%  feuillet 89 — folio 85 : 27 lignes de corps, 13 de note.
% -------------------------------------------------------------------
\begin{VUpage}[89]{85}
\VUnotes{124.12mm}{%
(1) \textit{Supere} venas \textit{de super}, quan on distingas de \textit{sur}. (Videz ta du\nl
prepozicioni e komparez li); ma \textit{infre} opozesas a \textit{supre}. (\textit{Progreso},\nl
III, 664.) --- \textit{Sube} venas \textit{de sub}, qua indikas la situeso relate to quo\nl
esas supere ed en la sama vertikal direciono.

(2) Decido 1094 : On adoptas ica regulo : « Por formacar kompo\cc
zita prepoziciono, on adjuntas a l'adverbo la prepoziciono quan\nl
l'originala vorto (verbo o adjektivo) havas : simila \textit{ad}, do : simile\nl
\textit{ad}, diversa \textit{de}, do : diverse \textit{de}, koncernar \textit{ulo}, do : koncerne \textit{ulo}\nl
(\textit{Progreso}, VI, 212).

Decido 1183 : L'Akademio aprobas ica interpreto di la decido\nl
1094 : Kande la verbo fonta postulas nula prepoziciono, t. e. esas\nl
transitiva, l'adverbo derivita postulas anke nula prepoziciono :\nl
koncerne \textit{ico}, relate \textit{ico} (\textit{Progreso}, VI, 417).%
}

\VUgras{Supere} ed \VUgras{sube} havas senco \textit{relativa}, ed indikas nur situeso\nl
o direciono relate punto determinita : t. e. nivelo plu o min\nl
alta. Exemple, konsiderante domo kun sis etaji e parolante\nl
de la vidpunto di ta, qua lojas sur la triesma, la \textit{supera} etaji\nl
esas la quaresma, la kinesma e la sisesma, ma la \textit{supra} etajo\nl
esas nur la sisesma; la \textit{suba} etaji esas la duesma, l'unesma\nl
e la ter-etajo, ma l'\textit{infra} esas nur la lasta (1). Konseque en\nl
citado, on devas dicar : \textit{videz supere} (e ne \textit{supre}), \textit{videz sube}\nl
(\textit{videz infre} signifikus : videz ye l'infro (infra parto) di la\nl
pagino, ex. en « infra noto »).

%% aucun blanc au fac-simile
Fine remarkez, ke la \textit{supro} o \textit{suprajo} ne esas sempre nek\nl
necese la \textit{somito}, nam ol povas esar larja e plana; on ne\nl
parolas pri la somito di moblo, di domo; la \textit{suprajo} di la\nl
tero tote ne identeskas kun la \textit{somiti} di la monti. On anke\nl
devas ne konfundar \textit{infra} e \textit{basa} : infra relatas la loko o\nl
situeso, \textit{basa} la dimensiono. La \textit{infra} etajo di domo povas\nl
esar \textit{alta}, e mem la maxim alta; kontraste la \textit{supra} etajo povas\nl
esar tre \textit{basa}.

\VUblanc{6.01pt}% mesure au fac-simile
102. --- Pri la \textit{komplemento di adverbo prepoziciona}\nl
observez ico :

%% aucun blanc au fac-simile
Existas adverbi derivita qui pleas la rolo di prepozicioni\nl
e konseque havas komplemento, quale la vera ed originala\nl
prepozicioni. Pri ta komplemento on agas quale pri la kom\cc
plemento dil vorto fonta : \VUgras{Koncerne, relate, ecepte} \textit{ulu} od\nl
\textit{ulo}, pro ke on dicas \VUgras{koncernar, relatar, eceptar} \textit{ulu} od\nl
\textit{olu} (2).

%% aucun blanc au fac-simile
\VUgras{Simile ad, konforme ad}, pro ke on dicas : \textit{simila ad, kon}\cc
\end{VUpage}

% -------------------------------------------------------------------
%  feuillet 90 — folio 86 : 34 lignes de corps, 2 de note.
%  LIGNE NON DETECTEE : « tismo. », dans le corps, signalee par
%  outils/lignes_manquantes.py (pas de 114,8 px pour un pas local de
%  41,4 : ni un blanc de 73,4 px, mais une ligne manquante suivie d'un
%  blanc de 32,1 px). Verifiee a l'agrandissement.
%  Les deux blancs sont donc poses a la main, la ligne retablie a son
%  pas : avant « 103. » 32,2 px ; avant « 104. » 32,4 px.
% -------------------------------------------------------------------
\begin{VUpage}[90]{86}
\VUnotes{153.25mm}{%
(1) \textit{Me ne povos facar to, ne esante sustenata} inspirus la kom\cc
preno : \textit{me ne povos... pro ke me ne esas sustenata}.%
}

\VUcontinue
\textit{forma ad}. Do juste la prepoziciono preiranta la komplemento\nl
dil vorto fonta. Same pri : \VUgras{diverse de} (o \textit{kam}), \VUgras{aparte de,}\nl
\VUgras{dextre di, sinistre di, funde di, okazione di}, pro ke on dicas :\nl
\textit{diversa de} (o \textit{kam}), \textit{aparta de, la dextro di, la sinistro di, la}\nl
\textit{fundo di, l'okaziono di}, e. c.

% aucun blanc au fac-simile
Ma preferez la propozicioni primitiva a l'adverbi prepo\cc
ziciona, kande nula specal motivo konsilas ici. Exemple, \VUgras{pri,}\nl
\VUgras{segun, inter} esas preferinda kam, \textit{koncerne, relate, kon}\cc
\textit{forme a, meze di}. Mem ica lasta darfas uzesar nur se la\nl
senco esas vere : \textit{en la mezo di}, exemple : \VUgras{meze di la cham}\cc
\VUgras{bro} = en la mezo di la chambro. Ma \textit{meze di mea amiki} nule\nl
equivalas : \textit{inter mea amiki}.

% aucun blanc au fac-simile
Sempre dicez \VUgras{pro} e ne \textit{kauze di}, plu o min Franca idio\cc
tismo.

\VUblanc{7.76pt}% mesure a la main (ligne « tismo. » non detectee)
103. --- \VUgras{La plaso dil komplemento di irga prepoziciono}\nl
(primitiva, derivita o kompozita) fixigesas da ica general\nl
regulo : ol devas sempre \textit{sequar nemediate} la prepoziciono,\nl
sen ul ecepto.

% aucun blanc au fac-simile
To esas tre importanta por evitar l'obskuraji ed ambi\cc
guaji, quin produktas en la Angla ed en la Germana, exemple,\nl
la prepozicioni qui \textit{sequas} lia komplemento, o mem \textit{preiras}\nl
e \textit{sequas} ta komplemento.

\VUblanc{7.81pt}% mesure a la main
104. --- Omna prepozicioni darfas uzesar avan infini\cc
tivo ed en la sama kazi por qui on uzus li avan equivalanta\nl
substantivo. Ex. : \VUgras{il kantas pos drinkir} (quale : \textit{pos drinko});\nl
\VUgras{il manjas ante departar} (quale : \textit{ante sua departo}); \VUgras{me ne}\nl
\VUgras{povus facar to sen esar sustenata} (1) (quale : \textit{sen susteno});\nl
\VUgras{vu facas ad ilu troa honoro per diskutar kun ilu} (quale :\nl
\textit{per vua diskuto}); \VUgras{il esas malada pro tro laborir} (quale :\nl
\textit{pro troa laboro}); \VUgras{mea laboro konsistos en montrar} (quale :\nl
\textit{en montro}.

% aucun blanc au fac-simile
Ma kande la prepoziciono indikata dal senco esas \VUgras{di} o \VUgras{ad},\nl
esas ordinare plu bona ne expresar olu, nam la kuntexto\nl
sufias por sugestar lu.
\end{VUpage}

% -------------------------------------------------------------------
%  feuillet 91 — folio 87 : 34 lignes de corps, 4 de note.
%  LIGNE NON DETECTEE : « pro li. », dans le corps, signalee par
%  outils/lignes_manquantes.py avec le verdict le plus sur qu'il rende
%  — « double net, aucun blanc » : 82,8 px pour un pas local de 41,5.
%  La frontiere qui la suit ne porte donc AUCUN blanc, et le seul de
%  la page (23,1 px, avant « 106. ») est pose a la main.
% -------------------------------------------------------------------
\begin{VUpage}[91]{87}
\VUnotes{148.84mm}{%
(1) Ol havas kom bazo ca principo logikala (quale ni vidis ye la\nl
verbi) ke malgre kontrea exempli di nia lingui, la verbo devas esar\nl
transitiva, kande lua ago povas atingar direte objekto : \textit{mokar,}\nl
\textit{nocar, atencar} e. c. ulu, ulo.%
}

105. --- Vice l'expresuro « la fakto di » sequata da in\cc
finitivo, on darfas uzar nur l'artiklo kun ta infinitivo. Ex. :\nl
« \textit{Ne la fakto (di) ebriigar su konstitucas l'alkoholismo, ma}\nl
\textit{la fakto (di) drinkar kustume alkoholo} » povas e darfas esar\nl
expresata ne min klare e plu kurte per : \VUgras{Ne la ebriigar su}\nl
\VUgras{konstitucas l'alkoholismo, ma la drinkar kustume alkoholo;}\nl
ed altramaniere : \VUgras{la alkoholismo konsistas ne en la ebriigar}\nl
\VUgras{su, ma en la drinkar kustume alkoholo.} --- \VUgras{Altr exemplo :}\nl
« \textit{Ne la fakto (di) timar nulo konstitucas kurajo, ma la fakto}\nl
(\textit{di}) \textit{afrontar danjeri malgre la pavoro, quan on sentas pro}\nl
\textit{li} »; plu kurte per : \VUgras{Ne la timar nulo konstitucas kurajo,}\nl
\VUgras{ma la afrontar danjeri malgre la pavoro quan on sentas}\nl
\VUgras{pro li.}

% aucun blanc au fac-simile
Ma ca exempli di rara kazi ne devas instigar vu ad uzar\nl
l'artiklo nediskrete e tedante avan l'infinitivo (quale se la\nl
substantivi ne plus existus), e dicar exemple : \textit{la studiar e}\nl
\textit{la muzikar esis lua delici}, vice : \VUgras{la studio e la muziko esis}\nl
\VUgras{lua delici.}

\VUblanc{5.57pt}% mesure a la main (ligne « pro li. » non detectee)
106. --- \textit{Prepozicioni kun verbi}. --- L'uzado di prepozi\cc
cioni kun verbi esas un del maxim granda desfacilaji dil\nl
vivanta lingui.

% aucun blanc au fac-simile
En nia linguo ica desfacilajo esas grandaparte vinkita\nl
per l'injenioz institucuro (1), ke la verbo esas transitiva\nl
en omna posibla kazi.

% aucun blanc au fac-simile
Yen kelka exempli ube nia linguo fixigis la verbo kom\nl
transitiva : \VUgras{aludar, bezonar, disponar, diletar, regnar...}

% aucun blanc au fac-simile
(Altra avantajo di la verbi transitiva esas, ke on povas\nl
uzar li en pasivo, Ex. : se on dicas : \VUgras{regnar lando, homi}, on\nl
povas dicar : \VUgras{lando regnata} (da ta suvereno) e parolar pri\nl
\VUgras{regnati.} Se on dicas : \VUgras{plura autori laboris ica libro}, on\nl
povas dicar : \VUgras{ica libro laboresis da plura autori.}

% aucun blanc au fac-simile
Duesma moyeno faciligar l'uzado di prepozicioni, e la\nl
sola regulo sequenda en nia linguo esas : \textit{selektar la pre}\cc
\textit{pozicioni segun lua propra senco}, tote ne segun la verbi od\parplein
\end{VUpage}

% -------------------------------------------------------------------
%  feuillet 92 — folio 88 : 31 lignes de corps (dont un TITRE), 5 de note.
%  La page ouvre sur la suite de l'alinea du folio 87.
%  LIGNE NON DETECTEE : « sur...). », dans le corps. Signalee par
%  outils/lignes_manquantes.py (pas de 83,1 px pour un pas local de
%  41,3) et verifiee au fac-simile. La frontiere qui la suit ne porte
%  aucun blanc : 83,1 - 2 x 41,3 = 0,5 px.
%  Blancs du titre MESURES sur les lignes de base AVANT composition :
%  114,9 px au-dessus (exces 73,6 = 6,23 mm), 90,5 px au-dessous
%  (exces 49,2 = 4,17 mm).
%  Corps et interlettrage : outils/apparat.py 92 --regle KONJUNCIONI. 20
%  --gras -> capitale 29 px = 2,455 mm -> 10,27 pt, interlettrage -43.
% -------------------------------------------------------------------
\begin{VUpage}[92]{88}
\VUnotes{147.24mm}{%
(1) \VUgras{Ma} indikas opozo, restrikto, \textit{difero} inter la du propozicioni\nl
o vorti quin ol unionas : \textit{Lu esas severa ma yusta. Me tre deziras}\nl
\textit{facar ta voyajo ma me ne povas lo. Fidez ad ili, ma ne ad eli.}

(2) \VUgras{Nam} anuncas expliko, motivo, pruvo relate lo dicata en la\nl
propoziciono a qua ol unionas la sua, Ex, : \textit{Tu vartos vane, nam il}%
}

\VUcontinue
altra vorti, qui akompanas lu. Ex. : \textit{absolvar ulu} de \textit{kulpo;}\nl
\textit{adicionar un nombro} \VUgras{ad} \textit{altra; konvinkar ulu} pri \textit{ulo}, e. c.

% aucun blanc au fac-simile
Fine on devas ne iterar sen neceseso la prepoziciono;\nl
generale, ne esas utila uzar lu \textit{en} e \textit{dop} la verbo. Suficas\nl
exemple dicar : \textit{irar ek la chambro, portar ulo ek la chambro}\nl
(o : \textit{forportar}, se on volas expresar l'ideo di forigo); \textit{pozar}\nl
\textit{lampo sur tablo} (vice : \textit{ekirar ek..., ekportar ek..., surpozar}\nl
\textit{sur...}).

% aucun blanc au fac-simile (ligne « sur...). » non detectee)
Pri la prepoziciono \textit{ye}, pro ke ol esas tante komoda, on\nl
ne devas tro uzar olu. Exemple, on povas dicar plu bone :\nl
\textit{abundar de} (quale \textit{plena de, richa de}, nam \textit{de} indikas la kon\cc
\textit{tenajo}); \textit{acendar de} (indikas l'origino); \textit{konvinkar pri; kom}\cc
\textit{prenar per}, o \textit{sub; reverencar ad} (\textit{ad} indikas la skopo di\nl
l'ago); e forsan mem (quale agas korekte kelka samideani) :\nl
\textit{chanjar ad, edukar ad}; nam to esas konforma a l'esencal ideo\nl
di la prepoziciono \textit{ad}. Aparte, \textit{chanjar ad}... esas certe plu\nl
logikala kam \textit{chanjar en}..., malgre l'exemplo di D. F. (cetere,\nl
D. uzas anke \textit{zu} ed E. \textit{to}). --- (Segun L. \textsc{Couturat} e P. \textsc{de}\nl
\textsc{Janko}, \textit{Progreso}, II, 458).

\VUsaut{6.23mm}
\VUtitre{10.27pt}{-43}{KONJUNCIONI.}
\VUsaut{4.17mm}

Ido havas \textit{koordinal} konjuncioni, \textit{subordinal} konjuncioni\nl
ed un \textit{questional} konjunciono.

\VUblanc{5.83pt}% mesure au fac-simile
107. --- Yen la \textit{koordinal} konjuncioni :

% aucun blanc au fac-simile
\VUgras{Do} = E. \textit{then} (therefore), consequently; D. \textit{also, demnach,}\nl
\textit{somit, folglich}; F. \textit{donc} (par conséquent); I. \textit{dunque, quindi;}\nl
S. \textit{luego, per consiguiente.}

% aucun blanc au fac-simile
\VUgras{Ed} o \VUgras{e} = D. \textit{und}; E. \textit{and}; F. \textit{et}; I. \textit{e, ed}; S. \textit{y, é.}

% aucun blanc au fac-simile
\VUgras{Ma} = D. \textit{aber, allein} (nach Verneinung), \textit{sondern}; E. \textit{but;}\nl
F. \textit{mais}; I. \textit{ma}; S. \textit{mas, pero} (1).

% aucun blanc au fac-simile
\VUgras{Nam} = D. \textit{denn}; E. \textit{for}; F. \textit{car, en effet}; I. \textit{diffati}; S. \textit{en}\nl
\textit{efecto} (2).
\end{VUpage}

% -------------------------------------------------------------------
%  feuillet 93 — folio 89 : 13 lignes de corps, 32 de note.
%  LIGNE NON DETECTEE : « S. ó, ú. », dans le corps. Verdict « double
%  net » de outils/lignes_manquantes.py (83,1 px pour un pas local de
%  41,4), verifie a l'agrandissement : la frontiere qui la suit ne
%  porte aucun blanc.
%  Le filet echappait a outils/filet.py avant l'elargissement de sa
%  fenetre de recherche (§ 8 quateretvicies) : il est a y = 783.
%  Le bloc de notes OUVRE sur la fin de la note (2) du folio 88.
% -------------------------------------------------------------------
\begin{VUpage}[93]{89}
\VUnotes{72.22mm}{%
\VUcontinue
ne venos, pro ke il esas malada. La konjunciono \VUgras{nam} unionas la\nl
duesma propoziciono : \textit{il ne venos} a l'unesma : \textit{tu vartos vane}\nl
e donas expliko, motivo pri lo dicata en ica. --- \VUgras{Me ne esperas, ke}\nl
\VUgras{il sucesos; nam, pro ke il esas tre indolenta, il ne demarshos sat}\nl
\VUgras{multe.}

(1) « On demandas ofte de ni expliko di ca konjunciono, pro ke\nl
multa lingui ne havas l'equivalanta, e supleas ol per altra konjun\cc
cioni de diversa senci. \textit{Or} uzesas \textit{nur} en la rezoni, por enduktar un\nl
plusa premiso, t. e. motivo. Ex. : « Omno quo foligas esas veneno;\nl
\textit{or} l'alkoholo foligas; \textit{do}... » On devas ne intermixar \textit{or} kun \textit{do},\nl
qua esas quaze opozita, nam ol induktas la konkluzo; nek kun \textit{ma},\nl
qua anuncas od expresas, ula opozo (ex. : opozanto povus dicar :\nl
« ma l'alkoholo fortigas, nutras... »). En la praktiko, on ne enuncas\nl
formale la rezoni, on tacas una de la premisi, od on renversas lia\nl
ordino. Kande la premiso (motivo) sequas la konkluzo, \textit{or} rempla\cc
sesas da \textit{nam;} ex. : « l'alkoholo esas veneno, nam ol foligas ». (On\nl
tacas la majora, qua konsideresas kom tro konocata od evidenta. ---\nl
Plura lingui supleas \textit{or} per la konjunciono qua signifikas \textit{nun}\nl
(D. \textit{nun} (\textit{aber}), E. \textit{now}); ma on vidas facile, ke lore ol nule indikas\nl
la tempo... Ta konjunciono esas do aquirajo di la logikala spirito,\nl
qua sempre plu dominacas la moderna e civilizita lingui, pro la\nl
progresi di la cienco. (\textit{Progreso}, VI, 141.)

(2) L'intima senco di \textit{tamen} esas : opoze ad ico, malgre lo dicita,\nl
malgre ta obstaklo. Ex. : \textit{Il esas tre richa, tamen il almonas nul}\cc
\textit{tempe. Me savas, ke me riskas la morto, tamen me iros adibe.}

(3) Quale pri \VUgras{ed, e}, on uzas prefere la formo sen \textit{d}, kande l'eufonio\nl
permisas.

(4) « Segun mea (personal) opiniono, skribis S\textsuperscript{o} \textsc{Couturat}, \textit{sive}\nl
havas esence la senco etimologiala, di \textit{od se}; e se on uzas (e darfas\nl
uzar) ol kelkafoye vice simpla \textit{od}, ico venas de elipso (di la verbo) ».\nl
--- \textit{Progreso}, VII, 411.

Segun ta expliko, sive esas plu vere subordinala kam koordinala.%
}

\VUgras{Or} = D. \textit{nun} (aber); E. \textit{now}; F. \textit{or}; I. \textit{ora}; S. \textit{ahora} (1).

% aucun blanc au fac-simile
\VUgras{Tamen} = D. \textit{dennoch, gleichwohl, nichtsdestoweniger,}\nl
\textit{trotzdem}; E. \textit{however, yet, still}; F. \textit{cependant, pourtant, néan}\cc
\textit{moins}; I. \textit{pero, tuttavia, ciò non ostante}; S. \textit{con todo, sin}\nl
\textit{embargo, a pesar de} (2).

% aucun blanc au fac-simile
\VUgras{Od, o} (3) = D. \textit{oder}; E. \textit{or}; F. \textit{ou, ou bien}; I. \textit{o, od;}\nl
S. \textit{ó, ú.}

% aucun blanc au fac-simile (ligne « S. ó, ú. » non detectee)
\VUgras{O(d)... o(d)} = D. \textit{entweder... oder}; E. \textit{either... or;}\nl
F. \textit{ou... ou.}

% aucun blanc au fac-simile
\VUgras{Sive... sive} = D. \textit{ob... ob..., sei es} (dass)... \textit{sei es} (dass)...;\nl
E. \textit{either... or, whether... or whether, be it that... be it that;}\nl
I. \textit{sia... sia...}; F. \textit{soit... soit, soit... que... soit que}; S. \textit{bien...}\nl
\textit{bien, sea que... sea que} (4).
\end{VUpage}

% -------------------------------------------------------------------
%  feuillet 94 — folio 90 : 16 lignes de corps, 25 de note.
%  Le bloc de notes OUVRE sur la suite de la note (4) du folio 89.
% -------------------------------------------------------------------
\begin{VUpage}[94]{90}
\VUnotes{89.95mm}{%
\VUcontinue
Uzez do \textit{o(d)... o(d)...} se li unionas vorti ne verba : \textit{me drinkos o}\nl
\textit{vino o biro}, qua egalesas : \textit{o vino o biro esos mea drinkajo}. Ma uzez\nl
\textit{sive... sive... avan verbo} expresata o tacata : \textit{Sive tu volos, sive tu}\nl
\textit{ne volos, ni certe vizitos li. Sive (tu esas) richa, sive (tu esas) povra,}\nl
\textit{tu devas helpar tua proximo segun tua moyeni.}

(1) Separas la neganta propozicioni, o diversa vorti di propo\cc
zicioni neganta : \textit{Me nek prizas nek estimas ta homo.} --- \textit{Il havas nek}\nl
\textit{parenti nek amiki.}

(2) Pos examenir lo dicita pri \textit{yen} S\textsuperscript{o} \textsc{Couturat}, en \textit{Progreso},\nl
VII, 208, konkluzas :

« On uzez \textit{yen} kande ol guvernas \textit{nomo}, e \textit{yen ke} kande ol guvernas\nl
propoziciono. »

Altralatere, en solida artiklo pri \VUgras{yen ke}, barono S.\textsc{de Szent}\cc
\textsc{kereszty} dicas juste : « La maxim preciza defino esas certe (VII,\nl
208) : « Interjeciono esas izolita vorto, qua equivalas integra frazo. »\nl
Ja de co rezultas, ke \textit{yen ke} esas lo justa, kande \textit{yen} guvernas propo\cc
ziciono. » (VII, 289.)

Ma nur prepozicioni recevas \textit{ke} por divenar konjuncioni; do \textit{yen}\nl
esas reale prepoziciono. Cetere, se ol ne esus prepoziciono, kad ol\nl
havus komplemento? Ka ni ne dicas : \textit{yen mea matro; yen ili}, tot\nl
same kam ni dicas : \textit{por mea matro, kun ili?} Do \textit{yen}, sola, ne esas\nl
konjunciono.

Konsiderante \textit{yen} kom prepoziciono, l'akademio decidis (aprilo\nl
1925) ke ol divenas konjunciono per l'adjunto di \textit{ke}, quale omna\nl
prepoziciono en Ido. (Vid. § 109.)%
}

\VUgras{Nek... nek} = D. \textit{weder... noch...}; E. \textit{neither... nor...}; F. \textit{ni...}\nl
\textit{ni...}; I. \textit{nè... nè...}; S. \textit{ni... ni...} (1).

% aucun blanc au fac-simile
\VUgras{Lore... lore} = D. \textit{bald... bald...}; E. \textit{now... now...}; F. \textit{tantôt...}\nl
\textit{tantôt}; I. \textit{ora... ora...}; S. \textit{ya... ya...}

% aucun blanc au fac-simile
\VUgras{Yen.} (Videz infre la noto 2.)

\VUblanc{5.64pt}% mesure au fac-simile
108. --- Yen la \textit{subordinal} konjuncioni :

\VUblancAlinea
\VUgras{Ke} = D. \textit{das}; E. \textit{that}; F. \textit{que}; I. \textit{che}; S. \textit{que.}

% aucun blanc au fac-simile
\VUgras{Se} = D. \textit{wenn, sofern, falls}; E. \textit{if}; F. \textit{si}; I. \textit{se}; S. \textit{si.}

% aucun blanc au fac-simile
\VUgras{Se ne} = D. \textit{wo nicht, sonst}; E. \textit{if not}; F. \textit{sinon}; I. \textit{se non,}\nl
\textit{se no}; S. \textit{si no.}

% aucun blanc au fac-simile
\VUgras{Se nur, nur se} = D. \textit{wenn... nur}; E. \textit{if only}; F. \textit{pourvu que;}\nl
I. \textit{purchè, a condizione che}; S. \textit{con tal que, siempre que.}

% aucun blanc au fac-simile
\VUgras{Ecepte se} = D. \textit{ausserwenn}; E. \textit{except if}; F. \textit{à moins que,}\nl
\textit{excepté si}; I. \textit{eccetto che, salvo che}; S. \textit{exceptuado si.}

% aucun blanc au fac-simile
\VUgras{Quale se} = D. \textit{als wenn, als ob}; E. \textit{as though, as if;}\nl
F. \textit{comme si}; I. \textit{come se}; S. \textit{come si.}
\end{VUpage}

% -------------------------------------------------------------------
%  feuillet 95 — folio 91 : 27 lignes de corps, 9 de note.
%  PAGE CARDEE : pas regulier 42,91 px = 10,34 pt, contre 41,45 px
%  ailleurs. Vingt-sept lignes seulement pour une justification qui en
%  porte quarante-quatre : le compositeur a reparti le surplus en
%  lames uniformes. D'ou \VUinterlignePage.
% -------------------------------------------------------------------
\begin{VUpage}[95]{91}
\VUinterlignePage{10.34pt}
\VUnotes{134.99mm}{%
(1) Kande la Franca vorto \textit{où} relatas tempo, on devas tradukar\nl
lu per \textit{kande} : \VUgras{la tempo kande me esis felica} o : \textit{la tempo en qua}\nl
\textit{me esis felica.}

(2) Generale la vorti D. \VUgras{warum}, E. \VUgras{why}, F. \VUgras{pourquoi}, I. \VUgras{perchè},\nl
S. \VUgras{porque} tradukesas per \textit{pro quo;} ma D. \VUgras{wozu} e F. \VUgras{pour quoi} devas\nl
tradukesar per la vorti : \textit{por quo.}

(3) Nam subordinal propoziciono komencanta per \VUgras{ke} equivalas\nl
substantivo od infinitivo : \VUgras{me expektas, ke il venos} = \textit{me expektas}\nl
\textit{ilua veno;} same : \VUgras{pos ke il venis} = \textit{pos ilua veno.}%
}

\VUgras{Quankam} = D. \textit{obgleich, obschon, obwohl, trotzdem dass;}\nl
E. \textit{as though, as if}; F. \textit{quoique, bien que}; I. \textit{quantunque,}\nl
\textit{benchè, sebbene}; S. \textit{aunque.}

% aucun blanc au fac-simile
Ni adjuntez la konjuncioni questional-relativa ja konocata\nl
da ni kom adverbi :

% aucun blanc au fac-simile
\VUgras{Ube, kande} (1), \VUgras{quale, quante.}

% aucun blanc au fac-simile
Fine la konjuncioni kompozita :

% aucun blanc au fac-simile
\VUgras{Pro quo} = \textit{pro qua kauzo;}

% aucun blanc au fac-simile
\VUgras{Por quo} = \textit{por qua skopo;}

% aucun blanc au fac-simile
A qui korespondas : \VUgras{pro co, pro to} = \textit{pro ica, pro ita}\nl
\textit{kauzo;} \VUgras{por co, por to} = \textit{por ica, por ita skopo} (2).

\VUblanc{8.24pt}% mesure au fac-simile
109. --- La plu multa prepozicioni genitas konjuncioni\nl
samsenca per la simpla adjunto di ke (3) :

\VUblanc{1.49pt}% mesure au fac-simile
\VUgras{Pro ke} (kauzo); \VUgras{por ke} (skopo); \VUgras{per ke} (per la fakto\nl
\VUgras{ke) : glacio diferas de aquo per ke ol esas solida} = \textit{per sua}\nl
\textit{solideso.}

% aucun blanc au fac-simile
\VUgras{De ke} (de la fakto ke, de ico ke) : \VUgras{lua febreso venas de}\nl
\VUgras{ke il esis recente malada} = \textit{de lua recenta maladeso.}

% aucun blanc au fac-simile
\VUgras{Ante ke; pos ke; depos ke o de kande : depos ke, de kande}\nl
\VUgras{il venis, me ne plus ekiras; omnafoye ke.}

% aucun blanc au fac-simile
\VUgras{Dum ke : dum ke il esis malada} = \textit{dum lua maladeso.}

% aucun blanc au fac-simile
\VUgras{Til ke; sen ke; ultre ke.}

% aucun blanc au fac-simile
\VUgras{Malgre ke} (kf. \textit{quankam}); \VUgras{ecepte ke.}

% aucun blanc au fac-simile
\VUgras{Segun ke; kontre ke} (qua expresas opozo, kontrasto).

\VUblanc{8.24pt}% mesure au fac-simile
110. --- Altra konjuncioni esas formacata per adverbi e\nl
la konjunciono ke :

\VUblanc{1.61pt}% mesure au fac-simile
\VUgras{Tale ke; tante ke; kaze ke; kondicione ke; supoze ke;}\parplein
\end{VUpage}

% -------------------------------------------------------------------
%  feuillet 96 — folio 92 : 30 lignes de corps, 7 de note.
%  La page ouvre sur la suite de l'alinea du folio 91.
% -------------------------------------------------------------------
\begin{VUpage}[96]{92}
\VUnotes{140.36mm}{%
(1) \VUgras{Time ke} (kun timo ke) diferas de \textit{timante ke}, qua relatas la\nl
subjekto di la frazo.

(2) \VUgras{Tante ofte ke} e \VUgras{tante longe ke} inkluzas nula komparo,\nl
kontre ke : \VUgras{tam ofte kam} e \VUgras{tam longe kam} esas komparanta.\nl
\VUgras{Exemple : Il venas tante ofte, ke il tedas ni; venez tam ofte kam}\nl
\VUgras{vu povos.} --- \VUgras{Il diskursis tante longe, ke fine la asistanti ne plus}\nl
\VUgras{askoltis lu; tam longe kam me vivos, me esforcos sucesigar nia ideo.}%
}

\VUcontinue
\VUgras{time ke} (1); \VUgras{tante plu ke; tante ofte ke; tam ofte ke; tante}\nl
\VUgras{longe ke; tam longe ke} (2).

\VUblancAlinea
Remarkez anke la sequanta konjuncional expresuri :

% aucun blanc au fac-simile
\VUgras{Lore kande; quik kande; same kam; tante plu... quante}\nl
\VUgras{plu = quante plu... tante plu; tante min... quante min =}\nl
\VUgras{quante min... tante min. Ex. : Quante plu me konocas li,}\nl
\VUgras{tante plu me estimas li (= me tante plu estimas li, quante}\nl
\VUgras{plu me konocas li); quante plu me vidas ilu, tante plu il}\nl
\VUgras{plezas a me; quante min il laboras, tante plu il volas re}\cc
\VUgras{pozar; quante min me dormas, tante min me somnolas.}

% aucun blanc au fac-simile
\VUgras{Segun quante : segun quante me povos, me certe helpos vu.}

% aucun blanc au fac-simile
On atencez bone dicernar la koordinal konjuncioni de le\nl
subordinala e nome \VUgras{nam} de \VUgras{pro ke}. Ex. : \VUgras{Pro ke vu esis}\nl
\VUgras{absenta, on ne vartis vu.} La konjunciono \VUgras{nam} esus ne justa\nl
hike. (Videz la noto 2 dil pag. 88.)

\VUblanc{8.60pt}% mesure au fac-simile
111. --- La \textit{questional} konjunciono esas \VUgras{kad} (\VUgras{ka}) sive\nl
por la \textit{direta} questiono, sive por la \textit{nedireta}.

% aucun blanc au fac-simile
\textit{Direta questioni} : \VUgras{Kad la mediko esas arivinta? Ka tua}\nl
\VUgras{matro konsentas?}

% aucun blanc au fac-simile
\textit{Nedireta questioni} : \VUgras{Dicez a me, kad la mediko esas ari}\cc
\VUgras{vinta. Me questionas tu, ka tua matro konsentas.}

% aucun blanc au fac-simile
Quale on vidas, la konjunciono \textit{kad, ka} sempre komencas\nl
la propoziciono (chefa o subordinala) en qua lu trovesas.

% aucun blanc au fac-simile
Ma on ne uzas lu, se la propoziciono ja havas vorto ques\cc
tionala (\VUgras{qua, quo, quala, quanta}, e. c.).

% aucun blanc au fac-simile
\VUgras{Kad ne o ka ne?} esas fakte elipso vice : \VUgras{kad ne} (tala esas\nl
la kozo)? \VUgras{Ka ne} (tale vu pensas, opinionas, e. c.)? Ol tra\cc
dukas la F. \textit{n'est-ce pas?}

% aucun blanc au fac-simile
Co explikas, ke on ne dicas : \textit{ka no}. (Videz ica lasta vorto\nl
ye la adverbo di afirmo, nego o dubo).
\end{VUpage}

% -------------------------------------------------------------------
%  feuillet 97 — folio 93 : 5 lignes de corps sous un TITRE, 38 de note.
%  Trois defauts du detecteur se cumulent sur cette page, et les trois
%  ont ete signales par outils/lignes_manquantes.py avant composition :
%    -- le FILET tombe a 0,29 H (la page est aux trois quarts en notes) :
%       il echappait a outils/filet.py avant l'elargissement de sa
%       fenetre (§ 8 quateretvicies). Il est a y = 565.
%    -- la ligne « hike. » n'est pas detectee (verdict « double net » :
%       66,6 px pour un pas local de 33,0), verifiee a l'agrandissement.
%    -- la DERNIERE BANDE porte DEUX lignes soudees : 57 px pour une
%       mediane de 27 (§ 8 sesetvicies).
%  Le compte vrai est donc : titre + 5 lignes de corps + 38 de note.
%  Blanc sous le titre MESURE sur les lignes de base : 91,5 px pour un
%  pas de 41,4, soit un exces de 50,1 px = 4,24 mm.
%  Corps et interlettrage : outils/apparat.py 97 --regle INTERJECIONI. 1
%  --gras -> capitale 29,5 px = 2,498 mm -> 10,45 pt, interlettrage -36.
% -------------------------------------------------------------------
\begin{VUpage}[97]{93}
\VUnotes{52.83mm}{%
(1) \VUgras{Ve} = la Latina e D. \textit{weh! ach!} --- E. \textit{alas!} --- F. \textit{malheur!}

(2) \VUgras{Nu} = D. \textit{nun! nu! wohlan!} --- E. \textit{well! very well!} --- F. \textit{eh}\nl
\textit{bien! allons!} --- I. \textit{ebbene! andiamo!} --- S. \textit{orsu! pues!}

(3) \VUgras{Fi} = D. \textit{pfui!} --- E. \textit{fie! fy!} --- F. \textit{fi! foin} (de)! \textit{pouah!} ---\nl
I. \textit{oibò! eh! via!} --- S. \textit{fuera! quita allá! vaya!}

(4) *\VUgras{Hola} D. E. F. I. S. uzesas por vokar : \VUgras{Hola! puero, venez}\nl
\VUgras{hike.}

(5) *\VUgras{Hem} anke \textit{hum}, \textit{hm} D. E. F. expresas dubo, hezito : \VUgras{Kredar}\nl
\VUgras{ico? hem! me esus tro naiva.}

(6) *\VUgras{Hop} D. E. F. uzesas por ecitar al kuro, por saltigar trans\nl
obstakli : \VUgras{Hop! hop! brava bestio, adavane, saltez.}

(7) *\VUgras{Krak} E. F. S. expresas subita evento, o bruiso di krako :\nl
\VUgras{Juste kande me apertis la pordo, krak! Medor fugis inter mea}\nl
\VUgras{gambi.} --- \VUgras{Dum ke il kuris e saltis sur la planko, krak! olu ruptesis}\nl
\VUgras{ed il falis en la fosato.}

(8) *\VUgras{Krik} imitas la bruiso di kozo lacerata: \VUgras{Il prenis mea naztuko}\nl
\VUgras{e krik! il laceris lu.} --- \VUgras{La vento lastafoye inflas la velo e ni audas}\nl
\VUgras{krik! olua lacereso.} --- \VUgras{Per du stroki di hakilo, krik-krak! il abatis}\nl
\VUgras{l'eshafodo.}

(9) *\VUgras{Uf} F. I. S. expresas la alejeso, la liberigeso pri ulo sufriganta,\nl
penoza, tedanta : \VUgras{Uf! me finis ta rudega laboro!} --- \VUgras{Uf! ni esas}\nl
\VUgras{liberigita de ta viro tedera!}

(10) *\VUgras{Paf} D. F. expresas frapo, stroko, subit acidento : \VUgras{Dum ke}\nl
\VUgras{me turnis la kapo, paf! me recevas pugnofrapo sur la dextra vango.}

(11) *\VUgras{Plump} D. expresas la bruiso di falo, di explozo : \VUgras{Plump!}\nl
\VUgras{il falabis en l'aquo. Ante ke me kontis dek, plump! la pulvereyo}\nl
\VUgras{explozis.}

(12) *\VUgras{Sus} F. I. S. por ecitar al defenso, al persequo : \VUgras{Sus! sus!}\nl
\VUgras{al armi, al enemiki.}

(13) *\VUgras{Psit} e \VUgras{pst} D. F. por atencigar persono qua iras avan vu : \VUgras{Me}\nl
\VUgras{audis psit! dop me e, turninte me, vidis vua patrulo.}

(14) *\VUgras{Shut} e \VUgras{sht} F. por tacigar o silencigar : \VUgras{shut! shut! yen elu,}\nl
\VUgras{tacez.} --- \VUgras{Shut, pueri, ne bruisez : via matro dormeskas.}

(15) *\VUgras{Hu, hu} D. F. por mokar ulu : \VUgras{Hu! hu! mentiero.}

(16) *\VUgras{Ba} D. E. F. indikas ne sucio pri, desprizo : \VUgras{Ba! lasez lu}\nl
\VUgras{agar quale il volas; to ne importas.}

(17) *\VUgras{Aye} F. E. S. I. klamesko di doloro : \VUgras{Aye! aye! quante me}\nl
\VUgras{sufras! Aye! vu pikas me.}%
}

\VUsaut{1.24mm}
\VUtitre{10.45pt}{-36}{INTERJECIONI.}
\VUsaut{4.24mm}

112. --- L'\textit{interjecioni} esas onomatopei, o klameski plu\nl
o min internaciona quale \VUgras{ha! he! ho! ve!} (1) \VUgras{nu!} (2) \VUgras{fi!} (3)\nl
\VUgras{hura! hola!} (4) \VUgras{hem!} (5) \VUgras{hop!} (6) \VUgras{krak} (7) \VUgras{krik} (8) \VUgras{krik-krak!}\nl
\VUgras{uf!} (9) \VUgras{paf!} (10) \VUgras{plump!} (11) \VUgras{sus!} (12) \VUgras{psit!} (13) \VUgras{shut!} (14)\nl
\VUgras{hu-hu!} (15) \VUgras{ba!} (16) \VUgras{aye!} (17)
\end{VUpage}

% -------------------------------------------------------------------
%  feuillet 98 — folio 94 : 21 lignes de corps (dont un TITRE et un
%  TABLEAU de cinq rangs), 19 de note.
%  outils/lignes_manquantes.py a signale DEUX pas doubles ici, aux
%  lignes 6 et 7. Ce sont les blancs du TITRE, non des lignes absentes :
%  un titre precede d'un blanc de 40 px donne un pas de 82 px, que le
%  test lit comme un double. C'est la limite du depistage, et c'est
%  pourquoi il rend deux lectures au lieu de trancher.
%  Blancs du titre MESURES sur les lignes de base AVANT composition :
%  82,0 px au-dessus (exces 40,2 = 3,40 mm), 89,6 px au-dessous
%  (exces 47,8 = 4,05 mm).
%  Corps et interlettrage : outils/apparat.py 98 --regle NOMBRI. 6
%  --gras -> capitale 30 px = 2,54 mm -> 10,63 pt, interlettrage -81.
%  TABLEAU : abscisses relevees a l'amas, marge du bloc a x=225 px.
%  Les points de conduite tombent sur une grille reguliere de pas
%  4,68 mm (21,61 / 26,28 / 30,96 / 35,63 / 40,30 mm) ; chaque rang
%  porte ceux qui suivent son mot.
% -------------------------------------------------------------------
\begin{VUpage}[98]{94}
\VUnotes{108.54mm}{%
(1) Nomizar \textit{nulo} la cifro \VUgras{o} esabus ridindajo; nam \textit{nulo} = nula\nl
kozo. Ma kad \VUgras{o} ne esas ulo, cifro? Kad ol ne pleas importanta rolo?\nl
Kad, kande me skribas la cifro \VUgras{o}, me skribas nulo? Se vice : skribez\nl
« zero », ni dicus : skribez « nulo », omni komprenus, ke on devas\nl
skribar \textit{nula kozo}, do ne skribar?

(2) \textit{Kin} ne povis esar \textit{quin}, pro ke ica lasta esas la plural akuza\cc
tivo di \textit{qua}. Ni do havabus okazione : \textit{la quin, quin il donis ad vu,}\nl
\textit{esas plu bela kam la quin, quin il donis ad me.} O : \textit{le quin} (5) \textit{quin}\nl
\textit{vu skribis sur ta pagino esas omni tre male desegnita.} Nia skopo ne\nl
esas konkordar maxim posible kun la Latina \textit{quin-que}, ma donar\nl
formi sat dicernebla de altri. Plu importas por ni certa dicernebleso\nl
dil diversa formi, kam sklavatra imito di vorti o formi latina efekti\cc
ganta al vera koncernati di L. I. konfundi e miskompreni.

(3) \VUgras{Milion, miliard, bilion} e. c. recevas \textit{-i} dil pluralo, se to esas\nl
necesa : \VUgras{on spensis milioni por ta edifico.} --- Ma : \VUgras{il recevis plura}\nl
\VUgras{milion,} pro ke plura sat indikas la pluralo.

Ta nombri (kande on ne adjuntas \textit{-o}, quo ne esas \textit{necesa}) havas la\nl
acento sur la final silabo, pro ke li esas reale abreviuro di \textit{miliono,}\nl
\textit{miliardo} e. c.%
}

Altra vorti apartenanta a diversa gramatikal kategorii\nl
darfas uzesar kom interjecioni, quale : \VUgras{ya, yen, vere, certe,}\nl
\VUgras{brave, bone, (ad)avane, (ad)dope, fore, abase, haltez, tacez,}\nl
\VUgras{silencez, helpo, sokurso, shamo, kurajo} (o \VUgras{kuraje}), \VUgras{adio}\nl
(o \VUgras{adie})!

\VUsaut{3.40mm}
\VUtitre{10.63pt}{-81}{NOMBRI.}
\VUsaut{4.05mm}

113. --- La \textit{nombri kardinala} esas : \VUgras{zero} (1), \VUgras{un, du,}\nl
\VUgras{tri, quar, kin} (2), \VUgras{sis, sep, ok, non, dek; cent; mil;} le\nl
maxim uzita.

% aucun blanc au fac-simile
\VUgras{Pose : milion; miliard} (\textit{mil milioni}); \VUgras{bilion} (1,000,000\textsuperscript{2});\nl
\VUgras{trilion} (1,000,000\textsuperscript{3}); \VUgras{quadrilion} (1,000,000\textsuperscript{4}); \VUgras{quintilion,}\nl
\VUgras{sextilion, septilion, oktilion, nonilion, decilion} (singla de ta\nl
nombri egalesante la preirinta multipliko per un milion (3).

% aucun blanc au fac-simile
Ni devas explikar, pro quo ni adoptis por \VUgras{biliono}, \VUgras{tri}\cc
\VUgras{liono},... la senco \textit{Germana} prefere kam la senco E. F. I. S.\nl
plu internaciona.

\VUrang{\VUcase{45.80mm}{En E. F. I. S.}\VUcase{73.41mm}{En D.}}
\VUrang{\VUcase{7.70mm}{\VUgras{Miliono}}\VUcase{21.61mm}{.}\VUcase{26.28mm}{.}%
\VUcase{30.96mm}{.}\VUcase{35.63mm}{.}\VUcase{40.30mm}{.}%
\VUcase{50.63mm}{1,000\textsuperscript{2}}\VUcase{69.34mm}{1,000,000}}
\VUrang{\VUcase{7.79mm}{\VUgras{Biliono}}\VUcase{21.61mm}{.}\VUcase{26.28mm}{.}%
\VUcase{30.96mm}{.}\VUcase{35.63mm}{.}\VUcase{40.30mm}{.}%
\VUcase{50.55mm}{1,000\textsuperscript{3}}\VUcase{69.34mm}{1,000,000\textsuperscript{2}}}
\VUrang{\VUcase{7.70mm}{\VUgras{Triliono}}\VUcase{26.28mm}{.}%
\VUcase{30.96mm}{.}\VUcase{35.63mm}{.}\VUcase{40.30mm}{.}%
\VUcase{50.55mm}{1,000\textsuperscript{4}}\VUcase{69.26mm}{1,000,000\textsuperscript{3}}}
\VUrang{\VUcase{7.87mm}{\VUgras{Quadriliono}}%
\VUcase{30.96mm}{.}\VUcase{35.63mm}{.}\VUcase{40.30mm}{.}%
\VUcase{49.78mm}{1,000\textsuperscript{5}}\VUcase{69.26mm}{1,000,000\textsuperscript{4}}}
\end{VUpage}

% -------------------------------------------------------------------
%  feuillet 99 — folio 95 : 33 lignes de corps, 6 de note.
%  La page ouvre sur la suite de l'alinea du folio 94.
%  LIGNE NON DETECTEE : « 710.) », dans le corps. Verdict « double net »
%  de outils/lignes_manquantes.py (82,4 px pour un pas local de 41,5),
%  verifie au fac-simile. La frontiere qui la suit ne porte aucun blanc.
%  Les blancs sont donc poses a la main ; tous valent le blanc ordinaire
%  (mesures 8,7 / 8,1 / 7,9 px contre 7,7 px), sauf ceux qui manquent.
% -------------------------------------------------------------------
\begin{VUpage}[99]{95}
\VUnotes{144.36mm}{%
(1) Komprenebla la acento tonika esas nek sur \textit{e}, nek sur \textit{a}, ma\nl
sur la nombral silabo (\textit{un, du, tri, dek} e. c.) dil unesma elemento.

(2) Semblas, ke la streketo utilesas por montrar plu klare l'uniono\nl
dil nombri adicionata.

L'uniformeso konsilas uzar \textit{e} sempre (e ne \textit{ed}) en ta nombri.\nl
Do \textit{dek-e-un, dek-e-ok} quale \textit{dek-e-du, dek-e-tri}.%
}

\VUcontinue
e. c. On vidas, ke en l'unesma kolumno existas nula koincido\nl
inter la \textit{nombro} implikita en la \textit{nomo} (\textit{bi, tri, quadri},...) e\nl
la exponento di la potenco (la relato esas : l'exponento =\nl
\textit{n} + \textit{1}, \textit{n} esante la nombro implikita en la nomo); dum ke\nl
en la duesma (kolumno) ta nombri esas konstante egala.\nl
Esas do plu natural e komoda, konsiderar \VUgras{biliono}, \VUgras{tri}\cc
\VUgras{liono}, e. c., quale la sucedanta potenci di \VUgras{miliono}. Plue,\nl
on ne bezonas \VUgras{biliono} por 1,000\textsuperscript{3}, pro ke on havas ja la\nl
vorto plu vulgara \VUgras{miliardo}. Fine la Germana metodo fur\cc
nisas la moyeno expresar plu granda nombri (quo povos\nl
esar tre utila en la cienci) : \VUgras{decilion} signifikas 1,000\textsuperscript{11} en\nl
l'unesma sistemo, ma 1,000\textsuperscript{20} en la duesma. (\textit{Progreso}, I,\nl
710.)

% aucun blanc au fac-simile
Omna altra nombri, quin on nomizas kompozita, esas\nl
expresata per la moyeno di la simpla nombri, segun la\nl
maniero sequanta :

\VUblancAlinea
1\textsuperscript{e} Se li konstitucas adiciono, on indikas olu per \textit{e} : \VUgras{dek}\cc
\VUgras{e-un, dek-e-du, dek-e-tri, dek-e-quar}, e. c., o \VUgras{dek e un, dek}\nl
\VUgras{e du, dek e tri}, e. c. (1).

\VUblancAlinea
2\textsuperscript{e} Se li konstitucas multipliko on indikas olu per \textit{a} ye\nl
la fino dil unesma nombro : \VUgras{dua-dek} (20); \VUgras{tria-dek} (30);\nl
\VUgras{quara-dek} (40); \VUgras{kina-dek} (50), e. c., o \VUgras{duadek, triadek,}\nl
\VUgras{quaradek}, e. c.

% aucun blanc au fac-simile
\VUgras{Duadek-e-ok} (28); \VUgras{sepadek-e-du} (72); o \VUgras{duadek e ok,}\nl
\VUgras{sepadek e du; mil e okacent e nonadek e ok} (1898) (2).

% aucun blanc au fac-simile
Praktike on ne expresas \VUgras{un} avan \VUgras{cent, mil}, e. c.

\VUblancAlinea
\textsc{Remarki} : 1\textsuperscript{a}. Super la mili, on sequos la kustumo enun\cc
car la nombri dil mili, dil milioni, dil bilioni, e. c., kom aparta\nl
nombri. Ex. : 345,000 ne enuncesos : \textit{triacentamil e quara}\cc
\textit{deka mil e kina mil}, ma : \textit{triacent e quaradek e kinamil}. To\nl
esas videbla simpligo.

% aucun blanc au fac-simile
2\textsuperscript{a}. Por enuncar granda nombri e diktar li, suficos enun\cc
car la sucedanta cifri, quo esas nur simpligo di la normala\parplein
\end{VUpage}

% -------------------------------------------------------------------
%  feuillet 100 — folio 96 : 25 lignes de corps, 17 de note.
%  La page ouvre sur la suite de l'alinea du folio 95.
%  Un seul blanc d'articulation, devant « 114. --- ».
% -------------------------------------------------------------------
\begin{VUpage}[100]{96}
\VUnotes{114.89mm}{%
(1) Ta remarki, en \textit{Progreso}, VII, 39, finas per : « E l'aranjo\nl
propozita evitos ankore ula dusencesi : on povos distingar facile\nl
(mem aude) : \textit{du deki} e \textit{duadeki}. »

Cetere sur la 38 pagino dil sama numero trovesas la \textit{motivi} qui\nl
inspiris la sistemo propozita al Akademio ed aceptita da olu :

I. Ica aranjo konservas la perfekta regulozeso di nia nombro\cc
sistemo, e la kurteso di la simpla nombri, qui restas unsilaba. Ma ol\nl
augmentas multe l'eufonio e konseque la klareso, sen longigar\nl
remarkeble la enunco (komparez \textit{okdekquar} a \textit{okadek} e \textit{quar}); ol\nl
igas la nombri plu facile pronuncebla, audebla e komprenebla.

2. Ol impedas omna (mem semblanta) dusencesi quin on objecionis\nl
a la nuna sistemo : mem ti qui esus tentata intermixar \textit{tri dek} e\nl
\textit{dektri} ne povos miskomprenar \textit{triadek} e \textit{dek e tri}.

3. Ta procedo impedas miskompreni di plura sucedanta nombri :\nl
\textit{cent ok dek tri} povas komprenesar : 108, 13; od : 100, 83; od : 183.\nl
Ma on distingos klare : \textit{cent e ok, dek e tri; (una)cent, okadek e tri;}\nl
\textit{cent e okadek e tri} (\textit{Progreso}, II, 353).%
}

\VUcontinue
enuncado. Vice dicar : \textit{triacent e quaradek e kinamil, sisacent}\nl
\textit{e sepadek e non}, on dicos : supresante le mil, cent, dek : \textit{tri,}\nl
\textit{quar, kin, sis, sep, non}. Ico esas la maxim klara, simpla e\nl
kurta procedo, qua konsistas en diktar la sucedanta cifri vice\nl
la nombro (quale se on espelus nomo). On povos ankore\nl
enuncar tale la nombri di tri cifri, enuncante nur mil,\nl
milion, bilion. Ex. : \textit{tri-quar-kin-mil, sis-sep-non}. (S\textsuperscript{o} v.\nl
\textsc{Pfaundler}, IV, 345.)

%% aucun blanc au fac-simile
La sama procedo uzesos por enuncar o diktar la decimala\nl
nombri, enuncate nur lia sucedanta cifri. Ol konvenas pre\cc
cipue por la decimala nombri longa o senfina, exemple :

%% aucun blanc au fac-simile
= 3,14159... : \textit{tri, komo, un, quar, un, kin, non}... Kom\cc
preneble, on ne omisos enuncar insiste la komo, ed anke\nl
la zero.

%% aucun blanc au fac-simile
3\textsuperscript{a}. On povos sempre, quale nun, substantivigar la nom\cc
bro-nomi, tam en pluralo kam en singularo : \textit{dekeduo, deke}\cc
\textit{dui; duadeko, duadeki} (1).

\VUblanc{4.82pt}% mesure au fac-simile
114. --- Quale on ja vidis en la 3-ma remarko, la nombri\nl
darfas esar substantivigata per \VUgras{-o} (pluralo \VUgras{-i}): \VUgras{duo; trio;}\nl
\VUgras{dekeduo; cento; milo.} En ica kazo la unajo kompozanta jun\cc
tesas a ta substantivo per la prepoziciono \VUgras{de} : \VUgras{dekeduo de}\nl
\VUgras{ovi} (ma : \VUgras{dek e du ovi}).

%% aucun blanc au fac-simile
\VUgras{Milion, miliard, bilion}, e. c., quale la cetera nombri, devas\nl
ne sequesar dal prepoziciono \VUgras{de} : on dicas \VUgras{milion homi, du}\nl
\VUgras{milion homi}, quale \VUgras{mil homi}.
\end{VUpage}

% -------------------------------------------------------------------
%  feuillet 101 — folio 97 : 29 lignes de corps, 9 de note.
%  Le feuillet est legerement en biais dans le scan : les abscisses
%  relevees glissent de +72 px en haut a +2 px en bas. Seuls comptent
%  les ECARTS a l'interieur d'une meme region — le renfoncement
%  d'alinea y vaut 44 a 48 px partout, comme ailleurs dans le volume.
% -------------------------------------------------------------------
\begin{VUpage}[101]{97}
\VUsignature{168.91mm}{4}% signature de cahier
\VUnotes{135.47mm}{%
(1) Tre diferanta de : \textit{Duopla kanto, promeno, triopla atako}, e tre\nl
diferanta anke de \textit{duesma kanto, promeno, triesma promeno}.

(2) Segun la procedo di enunco indikita p. 96 por 679 : (\textit{sis, sep,}\nl
\textit{non}), 243 darfus esar enuncata : \textit{du, quar, tri;} konseque por kurtigar\nl
l'enunco o la dikto di 243\textsuperscript{ma}, on darfus dicar : \textit{du, quar, tri-esma.}\nl
E tale por omna nombro ordinala.

(3) Dicez: \textit{la du unesma literi dil alfabeto} --- \textit{la tri unesma monati}\nl
\textit{di la yaro}, e ne : \textit{la unesma du literi..., la unesma tri monati...}\nl
(\textit{Progreso}, II, 32.)%
}

115. --- La nombri kardinala anke povas plear la rolo\nl
di adjektivi prenante la dezinenco \VUgras{a} : \VUgras{dua kanto, dua pro}\cc
\VUgras{meno} (exekutata da du personi); \VUgras{tria atako} (facata da tri\nl
individui, homa o bestia) (1).

\VUblanc{5.51pt}% mesure au fac-simile
116. --- Fine li anke povas produktar nombro-manieral\nl
adverbi, prenante la dezinenco \VUgras{e} : \VUgras{li promenas due; li pro}\cc
\VUgras{menas ne kune, ma singlu aparte; li atakis ni trie.}

%% aucun blanc au fac-simile
\VUgras{Un} genitas la pronomo \VUgras{unu} (pl. \VUgras{uni}). Ex. : \VUgras{unu postulas}\nl
\VUgras{ico, altru ito. Li nocis l'uni l'altri.}

%% aucun blanc au fac-simile
On darfas dicar : \VUgras{omni du, omni tri, omni quar}, e. c.;\nl
ma anke : \VUgras{omna du, omna tri}, e. c. (quale on dicas \textit{omna}\nl
\textit{homi}), la pluralo esante ja indikata dal nombri \textit{du}, \textit{tri}, e. c.

\VUblanc{5.65pt}% mesure au fac-simile
117. --- La \textit{nombri ordinala} obtenesas soldante al kar\cc
dinal nombri la sufixo \VUgras{-esm} e la dezinenco \VUgras{o, a, e} segun ke\nl
on bezonas substantivo, adjektivo o adverbo : \VUgras{la unesmo,}\nl
\VUgras{duesma, triesme.} Kande la nombro esas kompozita, nur la\nl
lasta elemento recevas la sufixo e la dezinenco : \VUgras{duacent}\nl
\VUgras{e quaradek e triesma} = 243\textsuperscript{ma} (2).

%% aucun blanc au fac-simile
\VUgras{Quantesma} relatas la ordino e konseque postulas nombro\nl
ordinala, quale l'adjektivo \VUgras{quanta} relatas nombro kardinala\nl
(o quanto) : \VUgras{quantesma dio di la monato esas? la sepesma,}\nl
\VUgras{la duadekesma.}

\VUblanc{5.66pt}% mesure au fac-simile
118. --- La \textit{nombri fracionala} obtenesas soldante al\nl
kardinal nombri la sufixo \VUgras{-im} e la dezinenco \VUgras{o, a, e} segun\nl
la kazi : \VUgras{la centimo} (1/100), \VUgras{la duimo} (1/2), \VUgras{la quarima}\nl
\VUgras{parto; centime, duime, quarime} (3).

%% aucun blanc au fac-simile
On uzas \VUgras{de} avan la komplemento di fraciono, quale avan\nl
la komplemento di integra nombro substantivigita : \VUgras{(la) du}\nl
\VUgras{triimi de sis esas quar} analoge a : \VUgras{dekeduo de ovi.}
\end{VUpage}

% -------------------------------------------------------------------
%  feuillet 102 — folio 98 : 32 lignes de corps, 6 de note.
%  DEUX lignes non detectees, toutes deux dans le corps : « kam 4. »
%  et « 3 foyi. ». La premiere se lisait dans le pas de 107,8 px
%  (« deux lectures » du depistage : blanc de 66,3 px OU ligne
%  manquante + blanc de 24,8), la seconde dans un double net de
%  83,9 px. Les deux verifiees au fac-simile.
%  Blancs calcules a la main, lignes retablies a leur pas :
%     avant « 119. » 20,5 px ; avant « 120. » 25,0 px ;
%     avant « Notez bone » AUCUN (1,1 px) ; avant « 121. » 23,7 px.
% -------------------------------------------------------------------
\begin{VUpage}[102]{98}
\VUnotes{144.53mm}{%
(1) Videz pag. 95 ye la 2-ma remarko.\nl
(2) Evitez sorge uzar la sufixo distributiva \VUgras{op} en senco kolektala\nl
por qua ol esas neutila, la kardinal adverbo suficante. Dicar : \VUgras{Ta du}\nl
\VUgras{amiki promenas sempre duope} esus tam kontrelogika kam dicar\nl
France ...\textit{se promènent toujours deux à deux}. Quale li povus facar\nl
ico, nam \textit{li esas nur du?}%
}

Por enuncar fraciono komplikita, on obtenas plu granda\nl
klareso uzante la prepoziciono \VUgras{sur} inter la du nombri kar\cc
dinala, qui esas la termini di la fraciono : 15/273 = \VUgras{dekekin}\nl
\VUgras{sur duacent e sepadek e tri} o : \VUgras{dekekin sur du, sep, tri} (1).\nl
\VUgras{Dek triacentimi} tote ne povas konfundesar a : \VUgras{deketri cen}\cc
\VUgras{timi.} L'unesma esas 10/300, la duesma 13/100; l'unesma lek\cc
tesos : \VUgras{dek sur triacent}, e la duesma : \VUgras{deketri sur cent.}

\VUblanc{4.94pt}% mesure a la main
119. --- La \textit{nombri multiplikera} obtenesas per la sufixo\nl
\VUgras{opl}, a qua on soldas \VUgras{o, a, e} segun la kazi : \VUgras{la duoplo} (la\nl
nombro multiplikata \textit{per du}); \VUgras{quaropla} (qua esas multi\cc
plikata \textit{per quar}); \VUgras{triople} : \textit{la nombro 12 esas triople granda}\nl
\textit{kam 4.}

\VUblanc{6.02pt}% mesure a la main (ligne « kam 4. » non detectee)
120. --- La \textit{nombro dil foyi} indikesas per la radiko \VUgras{foy}\nl
generale sequata dal adverbal dezinenco \VUgras{e} : \VUgras{dufoye}; o, se\nl
konvenas, dal dezinenco \VUgras{a} : \VUgras{trifoya voko} = \textit{voko repetita}\nl
\textit{3 foyi.}

% aucun blanc au fac-simile (ligne « 3 foyi. » non detectee)
\VUgras{Notez bone, ke la Franca vorto « fois » ne tradukesas}\nl
\VUgras{per} \textit{foy} \VUgras{ma per} \textit{opl} \VUgras{en la senco multiplikanta : triople quar}\nl
\VUgras{esas dekedu} (\textit{3 fois 4 font 12}). On darfas anke uzar la pre\cc
poziciono \VUgras{per}, precipue kande la nombri esas kelke kom\cc
plikita: \VUgras{okadek e tri per duacent a quaradek e sis}=83×246.\nl
Plu simpla enunco : \textit{ok, tri per du, quar, sis}.

\VUblanc{5.71pt}% mesure a la main
121. --- L'expresuri \VUgras{distributiva} obtenesas per la sufixo\nl
\VUgras{op}, sequata dal dezinenco konvenanta : \VUgras{duope} = per rangi\nl
o serii de \textit{du}; \VUgras{triope} = per rangi o serii de \textit{tri}, e. c. \VUgras{Unope,}\nl
un separite, aparte, \VUgras{single. La soldati marchas quarope.} ---\nl
\VUgras{Se li formacos rangi duopa e ne quaropa, lia defilo duros}\nl
\VUgras{duople plu longe.} --- \VUgras{Marchez unope.}

% aucun blanc au fac-simile
La questional adverbo korespondanta esas \VUgras{quantope :}\nl
\VUgras{quantope vu vendas la nuci} (quantanombre per un foyo?)\nl
\VUgras{centope o micentope} (dek e kinope). On darfas anke dicar :\nl
\VUgras{pokope, multope} = \textit{ye mikra quanti, ye granda quanti} (2).
\end{VUpage}

% -------------------------------------------------------------------
%  feuillet 103 — folio 99 : 30 lignes de corps, 10 de note.
%  LIGNE NON DETECTEE : « vespere. », dans les NOTES (double net,
%  66,2 px pour un pas local de 32,7). Etant dans les notes, elle ne
%  decale aucune frontiere du corps : poser_cardage reste sur.
% -------------------------------------------------------------------
\begin{VUpage}[103]{99}
\VUnotes{133.01mm}{%
(1) Pri la \textit{horo} dil nasko o morto di ulu, pri la 60 minuti dum\nl
qui lu naskis o mortis, kompreneble on ne darfas uzar \textit{kloko}. Ma\nl
on dicos tre reguloze : \VUgras{la horo di lua morto eventis ye kin kloki}\nl
\VUgras{di mea horlojeto.}

Esas dezirinda, ke on adoptez la kontado di la hori de \textit{0} a \textit{24},\nl
por evitar la distingo jenanta dil matino e dil vespero per \textit{matine,}\nl
\textit{vespere.}

(2) L'expresuro Angla « qua tempo esas », por questionar pri la\nl
horo, certe meritas preferesar. Cetere mem en F. on uzas \textit{temps} en\nl
ica senco, precipue en la tekniko : \textit{connaissance des temps; équation}%
}

122. --- Nula prepoziciono e nula sufixo esas necesa en\nl
frazi tal quala ici : \VUgras{Por mea infanti me kompris dekedu}\nl
\VUgras{pomi, ed a singla de li me donis tri pomi.} --- \VUgras{Ta libro havas}\nl
\VUgras{sisadek pagini; se do me lektos en singla jorno dek e kin}\nl
\VUgras{pagini, me finos la tota libro en quar jorni. La vorto singla}\nl
\VUgras{suficas por indikar la distributo.}

\VUblanc{5.74pt}% mesure au fac-simile
123. --- Por l'enunco di la hori, l'Akademio adoptis ex\cc
plicite \VUgras{kloko, kloki}. Ca du vorti indikas do propre la diala\nl
dividuro markizata da la cifri di horlojo. Ex. : \VUgras{il departis}\nl
\VUgras{pos du kloki}, to esas : pos la dividuro II di la horlojo. Ma\nl
\VUgras{horo} indikas (exter omna cifroplako, e per su) la duro,\nl
l'ensemblo dil 60 minuti konstitucanta un horo. Ex : \VUgras{il}\nl
\VUgras{departis pos du hori}, t. e. pos 120 minuti, pos restir dum\nl
120 minuti.

%% aucun blanc au fac-simile
Existas nur un maniero internaciona dicar la horo, to\nl
esas enuncar ol quale on skribas lu : \VUgras{esas du kloki; du kloki}\nl
\VUgras{e quarimo; du kloki e duimo; du kloki e tri quarimi; du}\nl
\VUgras{kloki dek} (minuti); \VUgras{du kloki duadek; du kloki kinadek;}\nl
\VUgras{du kloki kinadek e non} (minuti).

%% aucun blanc au fac-simile
On devas nultempe sustracionar, se on volas esar kom\cc
prenata certe ed evitar konfundi. Do sempre \textit{adjuntez} a horo\nl
la duimo, quarimi, minuti, e. c. (1).

%% aucun blanc au fac-simile
Por questionar pri la horo, fakte existas en Ido du ma\cc
nieri, quale en la Angla : \VUgras{qua kloko esas} e \textit{qua tempo esas?}\nl
Se forsan on agus tro severe interdiktante l'unesma, ne esas\nl
dubebla, ke la duesma (\textit{qua tempo esas?}) meritas la pre\cc
fero, pro ke ol esas plu justa. On ya ne questionas nur pri\nl
la \textit{horo}, ma reale anke pri la minuti e mem ulfoye sekundi\nl
qui adjuntesas a ta horo, do fakte on questionas \textit{pri la tempo,}\nl
on deziras \textit{konocar qua tempo esas} (2).
\end{VUpage}

% -------------------------------------------------------------------
%  feuillet 104 — folio 100 : 26 lignes de corps (dont un TITRE),
%  9 de note. Le bloc de notes OUVRE sur la fin de la note (2) du
%  folio 99.
%  Les deux pas doubles signales par le depistage (lignes 20 et 21)
%  sont les blancs du TITRE, comme au folio 94 — non des lignes
%  absentes. Mesures sur les lignes de base AVANT composition :
%  124,2 px au-dessus (exces 82,8 = 7,01 mm), 91,2 px au-dessous
%  (exces 49,8 = 4,22 mm).
%  Corps et interlettrage : outils/apparat.py 104 --regle SINTAXO. 21
%  --gras -> capitale 30 px = 2,54 mm -> 10,63 pt, interlettrage -44.
% -------------------------------------------------------------------
\begin{VUpage}[104]{100}
\VUnotes{136.26mm}{%
\VUcontinue
\textit{du temps; temps sidéral, temps solaire vrai} ou \textit{moyen}. (\textit{Progreso},\nl
VII, 399.)

Samloke \textit{Progreso} dicas : « On ne povas uzar ordinala nombro\nl
(ex. : \textit{quantesma horo o kloko}) pro la motivi expozita olim en la\nl
diskuto. » (III, 26-28, 218, 353, 703-4; IV, 378.)

(1) Remarkez \VUgras{tria}, ne \VUgras{triesma}. Nam ta ideo ne esas reale ordi\cc
nala : la potenco 3\textsuperscript{a} esas la potenco di qua l'exponento esas \textit{tri}.

(2) Ma \VUgras{Mezepoko} (meza epoko, ne evo) por la tempo iranta de\nl
la historio antiqua a la moderna.%
}

Quankam : \textit{la duesma, la triesma kloko}, e. c., ne esas \textit{en su}\nl
kontrelogika, ni agos plu bone, se ni sempre konservos la\nl
nombri kardinala \textit{pri la hori}. Tale ya ni evitos mixar le\nl
ordinala kun le kardinala, quale en : \VUgras{esas la duesma e dek}\nl
\VUgras{minuti.}

% aucun blanc au fac-simile
Do ni dicez sempre kun la nombri kardinala : \textit{esas}\nl
\textit{du kloki dek, du kloki e tri quarimi}, e. c.; e ni nultempe uzez\nl
la formo adverbala \textit{kloke}, quan l'Akademio explicite repulsis.

\VUblanc{2.87pt}% mesure au fac-simile
En matematiko, on tradukas reguloze la signi +, ---, ×\nl
e : \VUgras{per plus, minus, per e sur}; 2\textsuperscript{3} = \VUgras{du potenco tria;}\nl
\VUgras{3/2} = \VUgras{radiko tria di du} (1).

\VUblanc{5.90pt}% mesure au fac-simile
124. --- Por l'enunci relatanta la yari, monati, e. c., tra\cc
vivita, on uzas la radiko \VUgras{ev(ar), ev(o)}. Ex. : \VUgras{Quante vu}\nl
\VUgras{evas? me evas duadek yari}, o : \VUgras{me esas duadek-yara; quante}\nl
\VUgras{evas l'infanto di vua fratino? ol evas nur kin monati}, o :\nl
\VUgras{ol esas nur kinmonata; ol evabis nur du semani} o : \VUgras{ol esis}\nl
\VUgras{nur dusemana, kande ol mortis; ta viro semblas grandeva}\nl
o : \VUgras{tre evoza; yes il esas preske centyara. Kad vua patrulo}\nl
\VUgras{esas grandeva? No, il esas mezeva, il atingis nur mezevo} (2).

\VUsaut{7.01mm}
\VUtitre{10.63pt}{-44}{SINTAXO.}
\VUsaut{4.22mm}

125. --- Sen esar fixigita rigoroze e neflexeble, l'ordino\nl
dil vorti esas submisata en Ido ad ula reguli, quin impozas\nl
logiko e klareso. Exemple :

% aucun blanc au fac-simile
L'\textit{artiklo} sempre devas preirar \textit{nemediate} la substantivo,\nl
l'adjektivo e la pronomo quan ol akompanas.

% aucun blanc au fac-simile
L'\textit{adjektivo} devas preirar o sequar nemediate la substan\cc
\end{VUpage}

% -------------------------------------------------------------------
%  feuillet 105 — folio 101 : 31 lignes de corps, 8 de note.
%  La page ouvre sur la suite de l'alinea du folio 100.
%  La marge inferieure de ce feuillet est piquee de rousseurs : 907 px
%  d'encre sous la derniere ligne, qui faisaient crier au manque une
%  ligne inexistante. Voir outils/lignes_manquantes.py, ou le critere
%  est devenu la CONCENTRATION VERTICALE de cette encre.
% -------------------------------------------------------------------
\begin{VUpage}[105]{101}
\VUnotes{138.89mm}{%
(1) Videz ye pagino 30 to quo dicesas pri la \textit{plaso dil adjektivo}.

(2) \VUgras{Kam} sempre devas esar \VUgras{inter} la du parti o termini dil kom\cc
paro, en qua ol uzesas. Konseque ol ne darfas preirar l'unesma. Do ne\nl
dicez, exemple : \textit{to nocas} \VUgras{kam} \textit{l'editero tam l'aboninti}, ube \VUgras{kam} ne\nl
trovesas inter la du termini : \textit{editero, aboninti}, ma preiras l'unesma.\nl
Dicez : \textit{to nocas tam l'editero} \VUgras{kam} \textit{l'aboninti}, quale vu dicus : \textit{to}\nl
\textit{nocas plu l'editero} \VUgras{kam} \textit{l'aboninti} o : \textit{to nocas min l'editero} \VUgras{kam}\nl
\textit{l'aboninti.}%
}

\VUcontinue
tivo quan ol relatas. Ma ol sequas lu, se ol esas tre longa\nl
od akompanata da komplementi (1).

% aucun blanc au fac-simile
L'\textit{adverbo} devas preirar o sequar nemediate la vorto quan\nl
ol relatas. Ma la adverbi \VUgras{ne}, \VUgras{tre} devas \textit{preirar} ol sempre :\nl
\VUgras{me ne prenis vua libro} tre diferas sence de : \VUgras{ne me prenis}\nl
\VUgras{vua libro}, ed anke de : \VUgras{me prenis ne vua libro} (ma altra),\nl
e nur l'adverbo \VUgras{ne}, preiranta nemediate la vorto quan ol\nl
relatas, indikas nedubeble ta difero. --- \VUgras{Il tre deziras riches}\cc
\VUgras{kar rapide} diferas de : \VUgras{il deziras tre richeskar rapide}, od\nl
anke de : \VUgras{il deziras richeskar tre rapide}, e nur l'adverbo\nl
\VUgras{tre}, lokizita juste, indikas nedubeble ta difero (2).

% aucun blanc au fac-simile
La \textit{participo} --- en la tempi kompozita dil aktiva o pasiva\nl
voco --- sempre devas sequar la verbo helpanta, e separesar de\nl
olu nur da adverbo relatanta la verbo : \VUgras{li esas tre amata;}\nl
\VUgras{ica soldato esis grave vundata.} Ta regulo justifikesas dal\nl
fakto, ke la du vorti unionita kompozas simpla verbal formo\nl
(\VUgras{amesas, vundesis}).

% aucun blanc au fac-simile
\textit{Adjektivo} o \textit{participo} sempre devas \textit{sequesar} da sua kom\cc
plementi direta o nedireta, e lu devas \textit{sequar} nemediate sua\nl
substantivo : \VUgras{la homo estimata da omni} (e ne : \VUgras{la da omni}\nl
\VUgras{estimata homo}, o : \VUgras{la homo da omni estimata}). \VUgras{La fervoyi}\nl
\VUgras{formacas reto kovranta la mondo} (e ne : \VUgras{reto la mondo ko}\cc
\VUgras{vranta}, e mem mine : \VUgras{la mondo kovranta reto.}

% aucun blanc au fac-simile
Lasta exemplo, qua montras bone l'importo di vortordino\nl
logikala : « \VUgras{Il ne povas tolerar to quo esas segun lua kon}\cc
\VUgras{vinkeso absolute erora.} » \textit{Erora} semblas relatar \textit{konvinkeso}.\nl
Se la ideo esas : \textit{erora segun lua konvinkeso}, on devas dicar :\nl
« \VUgras{Il ne povas tolerar to quo esas absolute erora segun lua}\nl
\VUgras{konvinkeso.} »

\VUblanc{4.65pt}% mesure au fac-simile
126. --- La \textit{normal ordino} dil vorti en propoziciono esas\nl
ica : 1\textsuperscript{e} subjekto, 2\textsuperscript{e} verbo, 3\textsuperscript{e} komplemento direta; singla\parplein
\end{VUpage}

% -------------------------------------------------------------------
%  feuillet 106 — folio 102 : 31 lignes de corps, 10 de note.
%  LIGNE NON DETECTEE : « debla. », dans les NOTES (double net, 66,2 px
%  pour un pas local de 33,1) — elle acheve la note (2). Etant dans les
%  notes, elle ne decale aucune frontiere du corps.
%  Le bord droit du feuillet est distordu par le scan : les largeurs
%  relevees des lignes 28 a 31 et 37-38 depassent la justification et
%  ne valent rien.
%  Aucune des quatre frontieres d'alinea du corps ne porte de blanc.
% -------------------------------------------------------------------
\begin{VUpage}[106]{102}
\VUnotes{134.54mm}{%
(1) Uli nomizas ta \VUgras{n} « akuzativo », ma la vorto ne esas vere\nl
justa; nam Ido ne deklinas e konseque ne havas plu multe kam la\nl
Franca o l'Italiana, exemple, nominativo, akuzativo e. c., quale la\nl
Latina, la Germana ed altra lingui.

(2) Nulo certe impedas dicar : \textit{Tu rakontas a ni historio nekre}\cc
\textit{debla.}

(3) Same nulo impedas dicar : \textit{Me ja dicis ico ad vu plurfoye.}

(4) \VUgras{Bonan}, komplemento direta, preiras la subjekto \VUgras{me}. Konseque\nl
ol recevas la \VUgras{n} inversigala. Ma on dicus tre bone : \textit{Me prenis la}\nl
\textit{maxim bona.}%
}

\VUcontinue
de ta termini esante akompanata (segun la reguli donata\nl
supere) da lia omna komplementi. La komplementi nedireta\nl
darfas pozesar irgaloke, ma dop la verbo prefere : \VUgras{me}\nl
\VUgras{rakontis bela historio a vua infantino}, o : \VUgras{me rakontis a vua}\nl
\VUgras{infantino bela historio.} En ica lasta frazo, \VUgras{bela} ne povas\nl
relatar \VUgras{infantino}, nam on dicabus takaze : \VUgras{a vua bela in}\cc
\VUgras{fantino}, ed on devabus pozar \VUgras{historio} avane.

%% aucun blanc au fac-simile
Omna violaco dil normal ordino (\textit{subjekto, verbo, kom}\cc
\textit{plemento direta}) esas nomizat inversigo, sive kande la kom\cc
plemento direta preiras la verbo, sive kande la subjekto\nl
sequas lu. Exemple, se vice : \VUgras{la vintro venos balde}, me dicas :\nl
\VUgras{balde venos la vintro}, me inversigas la subjekto, nam ol\nl
sequas la verbo vice preirar olu.

%% aucun blanc au fac-simile
L'inversigi povas esar impozata dal bezono saliigar ter\cc
mino, pozante lu avane (ico remplasas l'expresuro \textit{c'est...}\nl
\textit{qui} o \textit{que} di la Franca), o (en la traduki) dal deziro sequar\nl
l'ordino dil texto originala. Se la frazo kontenas, quale\nl
supere, nur subjekto e verbo, l'inversigo dil unesma ne\nl
povas genitar ula miskompreno. Ma, se la frazo kontenas\nl
subjekto e komplemento direta od atributo, naskas miskom\cc
preno, en ula kazi, kande on inversigas lia normal ordino.\nl
Lore, por evitar ta grava desavantajo, on distingas la \textit{kom}\cc
\textit{plemento direta} o la \textit{atributo} per la final litero \VUgras{n} (1).

%% aucun blanc au fac-simile
Exempli pri l'inversigo dil komplemento direta :

%% aucun blanc au fac-simile
\VUgras{Quon vu dicas? Quan vu vidas? Quin vu prenas?} (La uzo\nl
di \textit{quon, quan, quin} esas la maxim ofta kazo dil \VUgras{n} inver\cc
sigala). --- \VUgras{Quanta homin la alkoholismo foligas !} ---\nl
\VUgras{Quanta servistin el havas ?} --- \VUgras{Nekredebla historion tu}\nl
\VUgras{rakontas ad ni} (2). --- \VUgras{Icon me ja dicis a vu plurfoye} (3).\nl
--- \VUgras{Tun, ne ilun me bezonas.} --- \VUgras{La maxim bonan me}\nl
\VUgras{prenis} (4).
\end{VUpage}

% -------------------------------------------------------------------
%  feuillet 107 — folio 103 : 30 lignes de corps, 10 de note.
%  Un seul blanc d'articulation, devant « 128. --- ».
% -------------------------------------------------------------------
\begin{VUpage}[107]{103}
\VUnotes{133.01mm}{%
(1) To esas la konsequo naturala di ca principo : \textit{en la normal}\nl
\textit{ordino}, l'unesma substantivo o pronomo di frazo sempre esas la\nl
subjekto.

(2) Kompreneble \VUgras{n} ne soldesas ad \textit{il, el, ol} (abreviuri), ma al\nl
formi kompleta dil pronomi : \textit{ilu, elu, olu}. Pro ke \textit{su} nultempe esas\nl
subjekto, ol nultempe darfas recevar la \VUgras{n} inversigala : \VUgras{su regardar} =\nl
certe \VUgras{regardar su}.

(3) La posibla kazi di tala konfundo esas \textit{rara}; ma pro ke la\nl
\VUgras{n} inversigala remedias li, ol esas uzenda sen hezito, kande li eventas.

(4) \textit{Quid sit veritas} (Latine) ? questiono di Pilatus à Iesu-Kristo.%
}

127. --- Kompreneble la \VUgras{n} inversigala esas necesa nur\nl
kande la komplemento direta preiras la subjekto. Ma, mem\nl
se ol preiras la verbo, nula \VUgras{n} devas uzesar, kande la sub\cc
jekto konservas l'\textit{unesma plaso}. Ex. : \VUgras{tua fratulo me odias}\nl
=\textit{tua fratulo odias me}. Ico precipue aplikesas al pronomi :\nl
\VUgras{el vu amas, il me vidis} = \textit{el amas vu, il vidis me}.

% aucun blanc au fac-simile
Konseque on povas donar ica regulo : ek du pronomi\nl
(sen final \textit{n}) qui trovesas en la sama propoziciono, l'unesma\nl
esas la subjekto e la duesma esas la komplemento direta (1).\nl
Ma on dicos : \VUgras{elun me amas, men el vidis}, pro ke la sub\cc
jekto ne esas l'unesma, ma la duesma (2).

% aucun blanc au fac-simile
On uzas anke \VUgras{n}, se substantivo o pronomo riskus irga\cc
grade, sen olu, konfundesar a subjekto : \VUgras{me amas vu quale}\nl
\VUgras{mea fratulon} (\textit{me amas}), certe diferanta de : \VUgras{me amas vu}\nl
\VUgras{quale mea fratulo} (\textit{amas vu}). Remarkez, ke se on komple\cc
tigas la frazo, « fratulon » preiras la subjekto \textit{me}, en\nl
l'unesma exemplo, qua fakte prizentas inversigo e qua, sen \VUgras{n},\nl
riskus komprenesar en la senco dil duesma frazo (3).

\VUblanc{5.66pt}% mesure au fac-simile
128. --- Inversigita, l'\textit{atributo} povas produktar ambi\cc
gueso e miskompreno, quale sen \VUgras{n} la direta komplemento\nl
inversigita.

% aucun blanc au fac-simile
En ica frazo : \VUgras{longa e desfacila esis ta konquesto}, la senco\nl
esas tre klara, quankam ne nur l'atributo (\textit{longa, desfacila}),\nl
ma mem la subjekto (\textit{konquesto}) esas inversigita : on povas\nl
nule konfundar l'atributi a « konquesto ».

% aucun blanc au fac-simile
Ma en : 1\textsuperscript{e} \VUgras{Quo esos o divenos tala autoritato?} 2\textsuperscript{e} \VUgras{Quo}\nl
\VUgras{divenas aquo per varmigo?} 3\textsuperscript{e} \VUgras{Quo divenis urbo Roma?}\nl
4\textsuperscript{e} \VUgras{Quo esas la verajo} (o : \textit{lo vera})? (4) Kad la senco ne povas\nl
aparar duopla?

% aucun blanc au fac-simile
L'unesma frazo dicas France : \textit{qu'est-ce qui sera ou de}\cc
\end{VUpage}

% -------------------------------------------------------------------
%  feuillet 108 — folio 104 : 25 lignes de corps, 17 de note.
%  La page ouvre sur la suite de l'alinea du folio 103.
%  Aucune de ses frontieres d'alinea ne porte de blanc.
% -------------------------------------------------------------------
\begin{VUpage}[108]{104}
\VUnotes{113.65mm}{%
(1) \textit{Progreso}, VI, 607, Remarko.

(2) En \textit{Progreso}, VII, p. 103, en \textit{Remarko} da S\textsuperscript{ro} \textsc{Couturat} :\nl
« L'ordino ne suficas por evitar la dusenceso. Yen exemplo : « \textit{Quon}\nl
divenus l'uneso di la L. I. qua esas la kondiciono di lua progreso\nl
e di lua suceso? » Esas neposibla, en tala frazo, transportar \textit{divenus}\nl
til la fino, pos omna komplementi e relativa propozicioni, qui povus\nl
esar mem plu multa e plu longa. Do l'akuzativo esas utila kun\nl
\textit{divenar}, malgre omna teoriala objecioni. Cetere kad l'atributo di\nl
\textit{divenar} ne indikas anke la objekto, la skopo? Ol darfas do indikesar\nl
per l'akuzativo same kam ordinara komplemento. »

Ed en VII, p. 351, on trovas sub la plumo di \textsc{Paul de Janko} :\nl
« Semblas a me tote indiferenta kad on nomizas kom akuzativo o\nl
kom formo di atributo la formo F. \textit{que}, qua korespondas en Ido\nl
a la formo : \textit{quon}.

L'utileso di omna du restas la sama. Pluse me remarkigas, ke la\nl
sama rezoni kam pri \textit{divenar}, valoras anke por la verbo \textit{esar},\nl
konseque ke on devos anke uzar ta formo kun ica, malgre ke to%
}

\VUcontinue
\textit{viendra une telle autorité?} ed anke : \textit{que sera ou deviendra}\nl
\textit{une telle autorité?} --- La duesma povas recevar du respondi :\nl
\textit{vaporo, glacio}. --- La triesma havas kom respondo posibla :\nl
\textit{azilo di raptisti} o \textit{la chefurbo di rejio Italia} (1). --- La qua\cc
resma signifikas France : \textit{qu'est-ce que la vérité} o : \textit{qu'est-ce}\nl
\textit{qui est la vérité?}

% aucun blanc au fac-simile
Se ni adjuntos a \VUgras{quo} di ta quar frazi la \VUgras{n} inversigala,\nl
kun olu la unesma frazo signifikos France nur : \textit{que sera}\nl
\textit{ou deviendra}..., e sen olu (quo) : \textit{qu'est-ce qui sera ou de}\cc
\textit{viendra} une telle autorité?

% aucun blanc au fac-simile
Kun la \VUgras{n} (quon) la duesma frazo havos kom respondo :\nl
\textit{vaporo}, e sen olu : \textit{glacio}.

% aucun blanc au fac-simile
Kun la \VUgras{n} (quon) la quaresma frazo signifikos France\nl
nur : \textit{qu'est-ce que la vérité}, e sen olu (quo) : \textit{qu'est-ce qui}\nl
\textit{est la vérité?}

% aucun blanc au fac-simile
Fine, kun la \VUgras{n} (quon) la triesma frazo havos kom res\cc
pondo : \textit{la chefurbo di Italia}, e sen olu : \textit{azilo di raptisti}.

% aucun blanc au fac-simile
On remarkez, ke : \VUgras{quo esos o divenos tala autoritato?}\nl
e \VUgras{quon esos o divenos tala autoritato?} esas fakte samtipa\nl
kam : \VUgras{quo donos tala autoritato?} e : \VUgras{quon donos}, e. c.

% aucun blanc au fac-simile
Advere la « Grammaire Complète » sancionita olim dal\nl
konstanta komisitaro, a qua yure sucedis l'Akademio Idista,\nl
tacas pri ca punto dil atributo. Ma tre semblas, ke esas pre\cc
ferinda uzar la \VUgras{n} inversigala kam lasar en multa kazi existar\nl
ambigueso kun miskompreni posibla (2).
\end{VUpage}

% -------------------------------------------------------------------
%  feuillet 109 — folio 105 : 31 lignes de corps, 8 de note.
%  La page ouvre sur la suite de l'alinea du folio 104, et son bloc de
%  notes OUVRE de meme sur la fin de la note (2) de ce folio.
% -------------------------------------------------------------------
\begin{VUpage}[109]{105}
\VUnotes{138.85mm}{%
\VUcontinue
semblos a multi kom mem plu stranja kam l'akuzativo kun \textit{divenar}.\nl
Fakte se me devas dicar : \textit{Quon divenas aquo per varmigo?} me devos\nl
anke dicar : \textit{Quon esos pos varmigo la nuna aquo?} Respondo : \textit{Ol}\nl
\textit{esos vaporo}. La sama formo uzesas anke en F. : \textit{Que sera l'eau après}\nl
\textit{avoir été chauffée?} e qua shokesas da la vorto « akuzativo » povas\nl
parolar anke en Ido pri formo di atributo, nam ne la nomizo\nl
importas ma l'expreso klara e preciza. »

(1) Ni vidis lo ye la prepozicioni \textit{ad, en, sur}, e. c.%
}

Esus kulpo uzar la \VUgras{n} por indikar (quale en Esperanto)\nl
la komplemento direta, se ol ne preiras la subjekto, o por\nl
indikar la dati, la disto, la preco, e. c., o por la translaco\nl
\textit{aden} loko (1). La \VUgras{n} inversigala devas ne uzesar exter la\nl
limiti fixigita supere. Tale ol restas, ne kaptilo e kompliko\nl
neutila en 8 foyi ek 10 adminime dil uzado Esperantala, ma\nl
nur moyeno impedar, mem preventar miskompreno, o sequar\nl
en traduko --- se to utilesas --- la ordino dil texto originala.

% aucun blanc au fac-simile
La rolo di la \VUgras{n} inversigala esas nur avertar pri ul inver\cc
sigo e \textit{quik} impedar miskompreno, dicante : ica vorto, modi\cc
fikita per \VUgras{n}, ne esas subjekto, ma komplemento direta od\nl
atributo.

\VUblanc{5.98pt}% mesure au fac-simile
129. --- Pronomo relativa-questionala sempre devas ko\cc
mencar la relativa propoziciono qua dependas de olu, ed\nl
ol darfas preiresar nur da prepoziciono o da \textit{irge, irgu} :\nl
\VUgras{la persono di qua vu vidas la domeno; me questionas, di qua}\nl
\VUgras{esas ta gardeno; ad irge qua povro il almonas; pri irge quo}\nl
\VUgras{vu parolas, vu sempre interesas; ad irgu qua venos, vu dicos}\nl
\VUgras{ke me esas absenta.} Konseque, kande relativa pronomo esas\nl
komplemento direta, existas necese inversigo, e ta pronomo\nl
devas recevar la \VUgras{n} inversigala : \VUgras{la soldato, quan vu vidas;}\nl
\VUgras{la soldati quin vu vidis.}

\VUblanc{6.07pt}% mesure au fac-simile
130. --- La \VUgras{-n} posibligas granda flexebleso e granda\nl
libereso; ma on devas ne trouzar l'inversigi (quale ni ja\nl
dicis) tante plu, ke mem la \VUgras{n} inversigala ne posibligas\nl
evitar sempre omna ambigueso. Ex. : \VUgras{me vidis ilun ocidan}\cc
\VUgras{tan viron} o \VUgras{me vidis ilun ocidar viron} ne esas plu klara\nl
kam : \VUgras{me vidis il(u) ocidanta} (o \VUgras{ocidar}) \VUgras{viro; la fervoyi}\nl
\VUgras{formacas reton kovrantan la mondon} ne esas plu klara\nl
kam... \VUgras{reto kovranta la mondo;} en la du kazi, nur l'ordino\nl
di la vorti povas distingar la subjekto de la komplemento.
\end{VUpage}

% -------------------------------------------------------------------
%  feuillet 110 — folio 106 : 28 lignes de corps (dont un TITRE),
%  8 de note. LIGNE NON DETECTEE : « Ido. », dans les NOTES (double
%  net, 66,5 px pour un pas local de 32,6) — elle acheve la note (3).
%  Les deux pas doubles des lignes 22 et 23 sont les blancs du TITRE,
%  non des lignes absentes (meme cas qu'aux folios 94 et 100).
%  Blancs du titre MESURES sur les lignes de base : 125,0 px au-dessus
%  (exces 83,5 = 7,07 mm), 93,7 px au-dessous (exces 52,2 = 4,42 mm).
%  Corps et interlettrage : outils/apparat.py 110 --regle
%  "TEMPI E MODI." 14 --gras -> capitale 29 px -> 10,27 pt, +3.
% -------------------------------------------------------------------
\begin{VUpage}[110]{106}
\VUnotes{136.14mm}{%
(1) Questionante pri la qualeso o stando di ulu od ulo, on devas\nl
uzar \textit{quala} (ne : \textit{quale}). \VUgras{Quala vu judikas lu} (esar)? \VUgras{Me judicas}\nl
\VUgras{lu kom benigna ma karaktere febla.} Se on uzus \textit{quale} (vu judikas)\nl
la senco esus : kad vu judikas? bone o male, yuste o neyuste, e. c.

(2) La mediko \textit{judikis}, ne trovis, nam lu ne serchis.

(3) Multi de la lasta exempli vizas aparte Esperantisti qui divenos\nl
Idisti, imitante la nombro ja tre granda dil Espisti konvertita a\nl
Ido.

(4) France : \textit{il disait qu'il écrivait.}%
}

\VUcontinue
Same nur l'ordino (kun la komuna raciono) distingas li\nl
en frazi simila ad ici : \VUgras{me igas imprimisto(n) imprimar}\nl
\VUgras{libro(n).}

% aucun blanc au fac-simile
La \VUgras{-n} esas anke neutila en kazi analoga al sequanti : \VUgras{me}\nl
\VUgras{judikis bona la vino}, o \VUgras{me judikis kom bona la vino}, ideo\nl
tre diferanta de : \VUgras{me trovis la vino bona} = (pos serchir)\nl
\VUgras{me trovis la vino bona}, o : \VUgras{la bona vino.}

% aucun blanc au fac-simile
\VUgras{Quala vu judikas} (esar) \VUgras{ta repasto} (1). \VUgras{Me judikas lu}\nl
\VUgras{kom ecelanta. La mediko judikis mea juniora fratulo kom}\nl
\VUgras{tre danjeroze malada} (2).

% aucun blanc au fac-simile
\VUgras{Anke nula -n en : Me nomizis Adolfus mea filiulo}, o :\nl
\VUgras{me nomizis mea filiulo Adolfus} (me donis ad ilu la nomo\nl
\VUgras{Adolfus}), quo tre diferas de : \VUgras{me nomis mea filiulo}\nl
\VUgras{Adolfus} (me pronuncis ilua nomo, qua esas Adolfus) e de :\nl
\VUgras{me vokis mea filiulo Adolfus.}

% aucun blanc au fac-simile
\VUgras{Nula -n} esus utila o justifikebla en tala frazi, nek en\nl
ica : \VUgras{me facas} (fabrikas) \VUgras{vitro neruptebla}, qua diferas tote\nl
de : \VUgras{me igas neruptebla la vitro.}

% aucun blanc au fac-simile
\VUgras{Mem pos kom n ne esas uzenda : Me selektis ilu kom pre}\cc
\VUgras{zidero} (por ke il esez prezidero), qua diferas de : \VUgras{me, kom}\nl
\VUgras{prezidero, selektis ilu}, o : \VUgras{kom prezidero me selektis ilu} (3).

\VUsaut{7.07mm}
\VUtitre{10.27pt}{3}{TEMPI E MODI.}
\VUsaut{4.42mm}

131. --- Pri la \textit{tempi} e \textit{modi} esas nur un regulo generala\nl
di selekto : on uzas en propoziciono subordinita (nedireta\nl
diskurso) la sama tempo e modo quan on uzus en chefa\nl
propoziciono (diskurso direta). Ex. : \VUgras{il dicis ke il skribas} (4)\nl
(il dicis : \textit{me skribas}); \VUgras{il dicis ke il skribis hiere a vua}\nl
\VUgras{matro} (il dicis : \textit{me skribis}, e. c.); \VUgras{il dicis ke il ja esis skri}\cc
\end{VUpage}

% -------------------------------------------------------------------
%  feuillet 111 — folio 107 : 31 lignes de corps, 8 de note.
%  La page ouvre sur la suite de l'alinea du folio 106.
%  Le bloc de notes est fait de six renvois d'une ligne chacun (les
%  equivalents francais) puis d'une note de deux lignes : sept alineas,
%  que le controle 12 verifie desormais.
% -------------------------------------------------------------------
\begin{VUpage}[111]{107}
\VUnotes{138.84mm}{%
(1) France : \textit{je pensais qu'il était ici.}

(2) France : \textit{je pensais qu'il serait ici.}

(3) France : \textit{je crains qu'il ne vienne.}

(4) France : \textit{je crains qu'il ne vienne pas.}

(5) France : \textit{je doute qu'il vienne.}

(6) France : \textit{je croyais qu'il viendrait.}

(7) La du lasta alinei esas absolute la vortopa tradukuro di\nl
« Grammaire Complète », p. 47, § 105.%
}

\VUcontinue
\VUgras{binta} o : \VUgras{skribabis dufoye ante recevar respondo} (il dicis :\nl
\textit{me ja esis skribinta o skribabis}, e. c.; \VUgras{dicez ad ilu ke il}\nl
\VUgras{venez} (dicez ad ilu : \textit{venez}); \VUgras{me pensis ke il esas hike} (1)\nl
(\textit{il esas hike, me pensis}); \VUgras{me pensis ke il esos hike} (2) (\textit{il}\nl
\textit{esos hike, me pensis}); \VUgras{me esperas, ke il venos} (\textit{il venos, me}\nl
\textit{esperas}); \VUgras{me timas, ke il venos} (3) (\textit{il venos, me lo timas});\nl
\VUgras{me timas, ke il ne venos} (4) (\textit{il ne venos, me lo timas}); \VUgras{me}\nl
\VUgras{dubas, kad il venos} (5) (\textit{kad il venos, me dubas pri lo});\nl
\VUgras{me kredis ke il venos} (6) (\textit{il venos, me kredis lo}); \VUgras{me kre}\cc
\VUgras{das ke il venus, se il ne impedesus} (hike la ideo esas vere\nl
kondicionala, e ne futura); \VUgras{il dicis, ke se il savabus} (o :\nl
\VUgras{esus savinta), il venabus} (o : \VUgras{esus veninta}) \VUgras{plu frue} (\textit{se}\nl
\textit{me savabus, il dicis, me venabus plu frue}); \VUgras{konvenas} (decas),\nl
\VUgras{ke vu facez ito} (\textit{facez ito, lo konvenas}).

% aucun blanc au fac-simile
Ta regulo suficas por determinar la kazi en qui on devas\nl
uzar l'imperativo o la kondicionalo en la propozicioni\nl
subordinita.

% aucun blanc au fac-simile
Partikulare, l'imperativo (o plu juste \textit{volitivo}) indikas\nl
sempre intenco o deziro, e la kondicionalo sempre supozas\nl
kondiciono explicita od implicita (ula \textit{se}...). \textit{La linguo ne}\nl
\textit{havas subjuntivo} (7).

\VUblanc{5.90pt}% mesure au fac-simile
132. --- On povas determinar l'intima naturo dil \textit{volitivo}\nl
(-ez) dicante, ke ol esas la modo di la \textit{skopo vizata} : \VUgras{por}\nl
\VUgras{ke ni vinkez ta desfacilaji, por ke ni sucesez komplete, esas}\nl
\VUgras{necesa, ke ni perseverez malgre omno ed omni.}

% aucun blanc au fac-simile
Pri la \textit{kondicionalo} bone memorez, ke ol sempre supozas\nl
kondiciono explicita od implicita. En ta modo la fakti ne\nl
prizentesas kom certa (quale en l'indikativo),ma kom plu o\nl
min dubebla, eventuala, dependanta de kondiciono quan\nl
indikas ofte (ma ne sempre) la konjunciono \textit{se}. Ex. : \VUgras{li esus}\nl
\VUgras{kontenta, se vu konsentus. Forsan ni povus sucesar} (\textit{se}...).
\end{VUpage}

% -------------------------------------------------------------------
%  feuillet 112 — folio 108 : 21 lignes de corps, 21 de note.
%  La page ouvre sur la suite de l'alinea du folio 107.
%  Le bloc de notes n'a qu'un renvoi, (1), mais quatre alineas : les
%  trois derniers exposent les motifs, sans numero.
% -------------------------------------------------------------------
\begin{VUpage}[112]{108}
\VUnotes{102.20mm}{%
(1) Ni dicas, ke en Ido la modi indikativa, volitiva o kondicionala\nl
remplasas \textit{logike} la subjuntivo di altra lingui, pro la motivi\nl
sequanta :

Se la fakto o ideo prizentesas \textit{kom certa}, l'indikativo devas uzesar,\nl
pro ke lu esas atribuita a la certeso; se la fakto, ideo prizentesas\nl
\textit{kom kondicionala}, kom dependanta de eventualajo, de supozo ne\nl
prizentata kom certa (komparez : \textit{se il venus} a \textit{se il venos}), la\nl
kondicionalo devas uzesar, pro ke lu esas atribuita al eventualaji, al\nl
idei kondicionala.

Se la fakto o ideo prizentesas kom implikanta volo, deziro, \textit{skopo}\nl
\textit{vizata}, la \textit{volitivo} esas uzenda, pro ke lu esas atribuita al indiko dil\nl
kozi vizata, imperata, irgamaniere volata.

Ma qua fakto o ideo ne apartenas (en la modi personala) al tri\nl
supera domeni? Pro quo do ni havus quaresma modo personala\nl
(subjuntivo), logike e bone expresata dal tri enuncita, quin Ido\nl
posedas? Quale ni povus logike e praktikale konstitucar lua propra\nl
domeno? Kande on vidas quante nia lingui intermixas sua indi\cc
kativo e sua subjuntivo, quante li deskonkordas pri ta punto,\nl
quante la sama linguo posedas subtilaji, kontredici e mem stultaji,\nl
por distingar od intrikar la du domeni, on benedikas Ido ne mem\nl
probir havar subjuntivo apud la tri altra modi personala.%
}

\VUcontinue
\VUgras{Me tre volus aceptar ilu} (se me ne esus impedata) \VUgras{ma me}\nl
\VUgras{esas tro malada cadie. En ta kazo, me quik forirus. Me jurus,}\nl
\VUgras{ke lu venabus sen ta obstaklo. Sen via helpo, ni perisabus}\nl
\VUgras{infalible. Kad esus posibla, ke vua gepatri ne sokursus vu.}\nl
\VUgras{Esus plu afabla, se vu skribus ipsa, ke vu aceptas lia invito.}

% aucun blanc au fac-simile
Remarkez, en la lasta exemplo, ke se tradukas la Franca\nl
konjunciono \textit{que} di : \textit{il serait plus aimable} que \textit{vous écriviez}\nl
\textit{vous-même}, e. c., de quo rezultas ke la kondicionalo esas\nl
substitucata en Ido al subjuntivo di la Franca. Vu facus\nl
kulpo, se vu tradukus o : \VUgras{esus plu afabla ke vu skribas,}\nl
o : \VUgras{esus plu afabla ke vu skribus o skribez.} Memorez bone\nl
ta exemplo ed imitez lu omnafoye, kande \textit{que} povos chan\cc
jesar a \textit{se}.

% aucun blanc au fac-simile
Tale ordinante la vorti dil exemplo : \VUgras{se vu skribus ipsa,}\nl
\VUgras{ke vu aceptas lia invito, esus plu afabla, vu komprenos bone,}\nl
\VUgras{ke nur la kondicionalo esas hike logikala e justa.}

\VUblanc{5.88pt}% mesure au fac-simile
133. --- Quale ni dicis ye la fino di § 131, Ido ne havas\nl
\textit{subjuntivo}, modo reale ne necesa, e quan logike remplasas\nl
en Ido l'\textit{indikativo}, la \textit{volitivo} o la \textit{kondicionalo}, segun la\nl
kazi (1).

% aucun blanc au fac-simile
Nun esas remarkenda fakto tre simpliganta en Ido : \textit{la}\parplein
\end{VUpage}

% -------------------------------------------------------------------
%  feuillet 113 — folio 109 : 32 lignes de corps, 6 de note.
%  La page ouvre sur la suite de l'alinea du folio 108.
% -------------------------------------------------------------------
\begin{VUpage}[113]{109}
\VUnotes{144.37mm}{%
(1) Hike valoras anke la principo fekunda, quan me donis del\nl
komenco di Ido : ne tradukez segun la vorti, ma segun l'idei\nl
(expresenda). Ol liberigas de mil heziti ed erori sugestata da nia\nl
lingui qui inspiras ico a Petrus, ma ito a Paulus, e triesma kozo\nl
a Stefanus, segun sua kustumi.

(2) Exemple : \textit{si j'étais} F., \textit{se io fossi} I. = \VUgras{se me esus}.%
}

\VUcontinue
\textit{konjuncioni havas nul influo pri la modo uzenda}. La selekto\nl
determinesas nur dal ideo (1). Ni do dicos kun l'indika\cc
tivo : \VUgras{Quankam il} \textit{esas} \VUgras{malada, il volas laborar ankore.} Pro\nl
quo ni imitus ula lingui, ne uzante l'indikativo kun « quan\cc
kam »? La fakto esas tote certa : \textit{il esas} malada. Per quo\nl
la konjunciono yurizas uzar altra modo por fakto certa e\nl
prizentata tre afirme? Nur l'indikativo devas uzesar en tala\nl
kazi; la konjunciono ne darfas impozar altra modo.

% aucun blanc au fac-simile
Same, pro quo la formo questionala o negala dil frazo\nl
dispensus ni uzar l'indikativo (prezenta, pasinta, futura)\nl
se la fakto prizentesas kom certa? L'exemplo di la Franca\nl
o di altra linguo ne valoras hike. Ni do dicos : \VUgras{kad vu kre}\cc
\VUgras{das ke} \textit{pluvas} (nun) o : \textit{pluvos} (morge)? kun l'indikativo,\nl
same kam ni dicas : \textit{pluvas}, o \textit{pluvis}, o \textit{pluvos}. Per quo \textit{kad},\nl
questionilo, yurizus ni chanjar la modo? Fine, ni ankore\nl
uzos l'indikativo por dicar : \VUgras{me ne kredas ke} \textit{pluvas} o \textit{pluvis}\nl
o \textit{pluvos}. Nam per quo la formo negala « me \textit{ne} kredas »\nl
obligus ne plus dicar \textit{pluvas, pluvis, pluvos?} Kad ni ne dicus :\nl
« mea opiniono \textit{ne} esas ke \textit{pluvas, pluvis} o \textit{pluvos} ».

% aucun blanc au fac-simile
Konkluze : sive la frazo esos afirma, nega o questiona, sive\nl
trovesos, o ne, avan la verbo la konjunciono \textit{quankam}, od\nl
altra postulanta, en irga linguo, altra modo kam l'indikativo,\nl
Ido sequos la logiko ed uzos konstante l'indikativo, se la\nl
fakto o l'ideo prizentesos kom certa.

\VUblanc{8.68pt}% mesure au fac-simile
134. --- Analoge Ido ne imitas la lingui qui uzas altra\nl
modo kam la kondicionalo pro la konjunciono \textit{se} (2). Ido\nl
sempre uzos la kondicionalo : \VUgras{li esus kontenta, se vu kon}\cc
\VUgras{sentus; li esabus kontenta, se vu konsentabus.} Ma, kompre\cc
neble, se irga linguo, la Franca exemple, kontrelogike uzas\nl
la kondicionalo vice l'indikativo futura, Ido uzas l'indi\cc
kativo. Ex. : Il espérait qu'\textit{il trouverait}, \VUgras{il esperis, ke} \textit{il}\nl
\textit{trovos} (\textit{me trovos}, il dicis a su espere). Elle nous a écrit\parplein
\end{VUpage}

% -------------------------------------------------------------------
%  feuillet 114 — folio 110 : 39 lignes de corps, AUCUNE note.
%  outils/filet.py rend « filet non trouve », et c'est ici la bonne
%  reponse : aucun des 38 pas de la page ne descend sous 38 px, donc
%  aucun bloc de notes. Meme cas qu'aux folios 71 et 75.
%  La page ouvre sur la suite de l'alinea du folio 109.
% -------------------------------------------------------------------
\begin{VUpage}[114]{110}
\VUcontinue
qu'\textit{elle scrait ici demain}, \VUgras{el skribis a ni ke} \textit{el esos} \VUgras{hike,}\nl
\VUgras{morge} (\textit{morge me esos hike}, el skribis a ni).

\VUblanc{7.57pt}% mesure au fac-simile
135. --- Same kam Ido abandonas nek l'\textit{indikativo}, nek\nl
la \textit{kondicionalo} por imitar uzi kontrelogika di ula lingui,\nl
tale ol restas sempre fidela a sua \textit{volitivo}, kande la fakto o\nl
stando dependas de volo, deziro, impero, interdikto, prego,\nl
demando, postulo, bezono, neceso, konveno, deco, merito.

% aucun blanc au fac-simile
Exempli : \VUgras{Me volas, ke vu skribez ad elu. Me totkordie}\nl
\VUgras{deziras, ke vi sucesez. Il imperis, ke me restez ed interdiktis,}\nl
\VUgras{ke mea fratino venez che ilu. Me pregas, suplikas, ke vu}\nl
\VUgras{helpez mea filiulo. Ni bezonas ed esas necesa, ke plu multa}\nl
\VUgras{homi partoprenez l'entraprezo. Vere il meritas, ke on rekom}\cc
\VUgras{pensez lu. Konvenas, ke ni vizitez li. Decas, ke la filii res}\cc
\VUgras{pektez ed obediez sua gepatri. Ni rekomendos a li, ke li tacez}\nl
\VUgras{ico. Fine il permisis, ke ni departez. Pro quo vu konsentas, ke}\nl
\VUgras{el traktez vu tale. Il venez e rakontez a me omno, quo eventis.}

% aucun blanc au fac-simile
L'expresuro por ke sempre postulas la \textit{volitivo} pro ke,\nl
kande on uzas ta expresuro konjunciona, la fakto o stando\nl
dependas tre reale de la neceseso, bezono, konveno o volo :\nl
\VUgras{Por ke vu povez pagar ta debajo, oportas ke vu kunprenez}\nl
\VUgras{sat grosa sumo de pekunio. Por ke il elektesez, vu mustas}\nl
\VUgras{luktar kun extrema energio. Por ke ni rekompensez vi, pueri,}\nl
\VUgras{konvenas ke vi meritez lo. Me volas agar omno posibla, por}\nl
\VUgras{ke vu esez kontenta pri me. Respondez ante morge, por ke}\nl
\VUgras{me savez precize to, quon me devas dicar ad ilu.}

% aucun blanc au fac-simile
Remarkez, ke la propoziciono dependanta de \textit{volitivo} pos\cc
tulas ta modo por sua propra verbo, pro ke ica fakte de\cc
pendas de impero, prego, deziro, volo : \VUgras{Imperez, ke il venez.}\nl
\VUgras{Dicez ad elu, ke el quik departez. Atencez, ke vu ne falez.}\nl
\VUgras{Ni ne tolerez, ke il tale rezistez a ni. Ni sorgez, ke il povez}\nl
\VUgras{domajar nulo.}

% aucun blanc au fac-simile
Rezume on darfus dicar : la \textit{volitivo} (modo dil skopo\nl
vizata o persequata), havas kom propra domeno to omna,\nl
quo ne apartenas al du altra modi personala : l'\textit{indikativo}\nl
e la \textit{kondicionalo}.

% aucun blanc au fac-simile
Evitez uzi tal quala ici : \textit{la questiono esas qua formon}\nl
\textit{ni selektez} o : \textit{qua linguo selektesez;} --- \textit{qui, se li efikez, pri}\cc
\textit{sentesez da profesional aktori}.

% aucun blanc au fac-simile
L'unesma semblas, pos sercho, signifikar : \textit{la questiono}\parplein
\end{VUpage}

% -------------------------------------------------------------------
%  feuillet 115 — folio 111 : 32 lignes de corps, 7 de note.
%  La page ouvre sur la suite de l'alinea du folio 110.
% -------------------------------------------------------------------
\begin{VUpage}[115]{111}
\VUnotes{140.68mm}{%
(1) Justifikar tal exempli per alegar plu o min multa vorti\nl
tacita, to esas pruvar ne direte lo jus dicita. Kad on bezonas\nl
alegar frazo-parto tacita por justifikar la uzo di la cetera modi?\nl
E kad esas certa, ke la lektanto od audanto divinos ta vorti tacita,\nl
e konseque trovos la justa senco inter pluri posibla? Ka ne esas\nl
plu sekura agar pri la \textit{volitivo} quale pri l'\textit{indikativo} e la \textit{kondi}\cc
\textit{cionalo} : ne lasar ulo divinenda ?%
}

\VUcontinue
\textit{esas nur qua formon ni selektos} o : \textit{qua formon ni devas}\nl
\textit{selektar}, do : \textit{du senci} posibla e diferanta, quo esas tote kon\cc
trea al principo dil unasenceso.

% aucun blanc au fac-simile
La duesma : « \textit{qua linguo selektesez} » semblas, anke pos\nl
sercho, signifikar : \textit{qua linguo selektesos} o \textit{esas selektenda},\nl
o : \textit{qua linguo devas esar selektata}. Do itere ni esas koram\nl
plura senci posibla e diferanta, quon interdiktas la prin\cc
cipo dil unasenceso.

% aucun blanc au fac-simile
La triesma : \textit{qui, se li efikez}, e. c., esas komplete nebuloza\nl
e, mem pos longa sercho, on esas nule certa pri lua signifiko,\nl
qua forsan esas ica : \textit{qui, por efikar, devas prizentesar da}\nl
\textit{aktori profesione}.

% aucun blanc au fac-simile
La fakto, ke ta frazi (kredeble inspirita da irga nacional\nl
idiotismo) ne esas klara e violacas la principo dil una\cc
senceso suficas por pruvar, ke la \textit{volitivo} esas en li misuzita\nl
e ne justifikebla (1).

\VUblanc{5.71pt}% mesure au fac-simile
136. --- Lernante la formi dil Ido-konjugo, ni konstatis,\nl
ke che olu existas la tri tempi \textit{prezenta}, \textit{pasinta}, \textit{futura}, ne\nl
nur en l'\textit{indikativo} : \VUgras{-as, -is, -os}, ma en l'\textit{infinitivo} : \VUgras{-ar, -ir,}\nl
\VUgras{-or} ed en la \textit{participi} aktiva : \VUgras{-ant, -int, -ont}, o pasiva : \VUgras{-at,}\nl
\VUgras{-it, -ot}. On darfus, alegante la facileso necesa, blamar ta\nl
logikal richeso, se ol postulus granda esforco de la lernanti.\nl
Ma vere, kad esas desfacilajo merkar e memorar, ke \textit{a} sim\cc
bolizas la prezento, \textit{i} la pasinto, \textit{o} la futuro en l'infinitivo\nl
e la participi, totsame kam en \VUgras{-as, -is, -os} dil indikativo?\nl
Se on judikas ta simetra sistemo sen preopinioni, on agnos\cc
kas, ke ol esas tante simpla e facila, ke merkar e memorar\nl
olu esas \textit{ne esforco ma ludo}. Altraparte on konfesos, ke lasar\nl
sen uzo ta richeso tempala pro ca motivo, ke altra lingui\nl
esas tarelate min richa, o mem tote mizeroza, to esus nesaja\nl
nek exkuzebla. La helpolinguo ya ne kondamnesas karear\nl
ta od altra formo pro ke li ne havas lu.
\end{VUpage}

% -------------------------------------------------------------------
%  feuillet 116 — folio 112 : 39 lignes de corps, AUCUNE note.
%  PAGE CARDEE : pas regulier 42,26 px = 10,18 pt.
%  LIGNE NON DETECTEE : « vere? ». Le depistage l'avait donnee en
%  « deux lectures » (blanc de 41,7 px OU ligne manquante sans blanc) ;
%  la lecture au fac-simile tranche pour la seconde, et le profil
%  d'encre le confirme : bases a 826, 868,5 et 909,5 px, deux pas
%  reguliers. La frontiere ne porte donc aucun blanc.
% -------------------------------------------------------------------
\begin{VUpage}[116]{112}
\VUinterlignePage{10.18pt}
Pro quo, exemple, Ido devus sisteme uzar nur la peri\cc
frazo : \textit{me pensas ke me pruvis}, o : \textit{me pensas esar pruvinta},\nl
kande ta linguo povas expresar la ideo plu lejere, plu kurte\nl
e totsame klare per : \VUgras{me pensas pruvir?}

% aucun blanc au fac-simile
Pro quo ni uzus sisteme nur ta longajo : \textit{me pensas ke}\nl
\textit{me pruvos} o : \textit{ke me esas pruvonta}, kande ni darfas dicar :\nl
\VUgras{me pensas pruvor?}

% aucun blanc au fac-simile
On evitas la perifrazo por dicar : \VUgras{me pensas pruvar.} Ma\nl
on mustus uzar olu por \textit{pruvir} e \textit{pruvor!} Qua saja homo\nl
povus aprobar ico?

% aucun blanc au fac-simile
Vice la komoda e kurta « \VUgras{departonte} » \VUgras{ye kin kloki, me...}\nl
ni mustus dicar : « \VUgras{pro ke me departos} » \VUgras{ye kin kloki, me...}\nl
\VUgras{o : departante...} quale se la ago ne esus \textit{futura}, ma prezenta!\nl
La motivo di ca stupidajo? --- Ta od altra linguo, e mem\nl
multa lingui ne havas participo futura. Solida motivo, ne\nl
vere?

% aucun blanc au fac-simile (ligne « vere? » non detectee)
E pro quo ni dicus « \textit{la futura uzanti} » \textit{di la linguo}, vice :\nl
\VUgras{la uzonti di la linguo?} O forsan : \textit{ti qui uzos?} Ma « la futura\nl
uzanti » esas pura absurdajo; nam, se li esas \textit{uzanti}, li ne\nl
esas \textit{futura}, e se li esas \textit{futura}, li ne esas \textit{uzanti}.

% aucun blanc au fac-simile
E quale tradukar perifraze en maniero tolerebla : \VUgras{esas}\nl
\VUgras{plu bona l'existar kam l'existir o l'existor?}

% aucun blanc au fac-simile
Qua homo serioza ne vidas, ke la supera perifrazi naskas\nl
en nia lingui de povreso, de mizero tote ne dezirinda od\nl
imitinda en Ido?

% aucun blanc au fac-simile
Ni do uzez la tri tempi (\textit{a}, \textit{i}, \textit{o}) dil modi Idala, kande\nl
l'okaziono postulas o konsilas lo, sen questionar ni, ka ni\nl
sustenesas da lingui naturala : Ido ya ne havas kom pro\cc
gramo imitar nia lingui, ma suplear li.

\VUblanc{8.67pt}% mesure au fac-simile
137. --- Kande participo esas komplemento cirkonstan\cc
cala, ol darfas recevar o l'adjektival formo, o la formo ad\cc
verbala, segun ke ol esas plu juste epiteto, od indikas plu\nl
reale la maniero (\textit{marchante} = dum mea, tua, vua, e. c.,\nl
marcho). Ma ol darfas recevar la formo adverbala, nur se\nl
ol relatas la subjekto di la propoziciono : \VUgras{il venis ne invi}\cc
\VUgras{tita,} o : \VUgras{ne invitite; il venis ne expektita,} o : \VUgras{ne expektite;}\nl
\VUgras{vidanta} o \VUgras{vidante sua amiko, il haltis; lektinta} o \VUgras{lektinte}\nl
\VUgras{la libro, il dormeskis.} En ta exempli on darfis uzar la formo\parplein
\end{VUpage}

% -------------------------------------------------------------------
%  feuillet 117 — folio 113 : 39 lignes de corps, AUCUNE note.
%  LIGNE NON DETECTEE : « testi. », en gras. Le depistage l'avait
%  donnee en « deux lectures » (blanc de 73,2 px OU ligne manquante +
%  blanc de 31,7) ; l'agrandissement tranche pour la seconde. Le blanc
%  pose est donc celui de la lecture (b) : 31,7 px = 7,64 pt.
%  Le feuillet porte un pli dans la marge de gauche, qui rend
%  outils/graisse.py inutilisable ici — sa limite propre, documentee.
% -------------------------------------------------------------------
\begin{VUpage}[117]{113}
\VUcontinue
adjektivala tam juste kam l'adverbala, pro ke la koncernato\nl
« \VUgras{il} » esas subjekto.

% aucun blanc au fac-simile
Ma en : \VUgras{la viro vidis cigno natanta sur la lago;} nur la\nl
formo adjektivala devas uzesar, pro ke la koncernato (cigno)\nl
ne esas subjekto. Se, en ica exemplo, ni dicus : \VUgras{natante} (sur\nl
la lago), la senco esus tote altra, nam lore la frazo equi\cc
valus : \textit{la viro vidis cigno, dum ke il} (la viro) \textit{natis sur la}\nl
\textit{lago}. Por ica lasta senco esus plu bona dicar : \VUgras{natanta o}\nl
\VUgras{natante sur la lago, la viro vidis cigno.}

% aucun blanc au fac-simile
Se on volas indikar nedubeble la relato dil fakto cir\cc
konstancala kun la fakto precipua, on devas uzar konvena\cc
nanta prepoziciono; komparez : \VUgras{il arivis ne avertinte me a :}\nl
\VUgras{il arivis sen avertir me.}

% aucun blanc au fac-simile
On darfas uzar la aranjo nomizata \textit{absoluta participo}, to\nl
esas propoziciono incidenta, di qua la verbo esas participo\nl
(adverba) e di qua la subjekto ne esas parto di la chefa\nl
propoziciono. Ex. : \VUgras{la enemiki fuginte, ni transiris la ponto.}\nl
Ma on devas ne tro uzar ta frazoformo, quankam ol esas\nl
utila e mem necesa en ula kazi, nome en la stilo matema\cc
tikala e judiciala : \VUgras{donite un rekto ed un punto; audite la}\nl
\VUgras{testi.}

\VUblanc{7.64pt}% mesure a la main (ligne « testi. » non detectee)
138. --- La \textit{participi substantiva} ne darfas recevar direta\nl
komplemento. Exemple, ica frazo : « \textit{la atakanti la reli}\cc
\textit{gio...} » ne esas bona logike e gramatikale. Nam, se la par\cc
ticipo adjektiva (\textit{atakanta}) e la participo adverba (\textit{ata}\cc
\textit{kante}) restas duime verba e, pro to, darfas recevar direta\nl
komplemento, la participo substantiva (\textit{atakanto}) perdas\nl
reale sua karaktero e valoro verbala; ol divenas vera sub\cc
stantivo e kom tala ne darfas havar direta komplemento,\nl
ma nur komplemento nedireta per prepoziciono, quale omna\nl
substantivi. Konseque on devas dicar : « La atakanti \textit{di} la\nl
religio », same kam on dicus : « La defensanti \textit{di} la re\cc
ligio » e : « la predikanti (o predikeri) \textit{di} la religio », « la\nl
protektanti (o protekteri) \textit{dil} febli ». Diveninte substan\cc
tivo, la participo ne plus sequas la regulo dil verbi, ma la\nl
regulo dil substantivi.

% aucun blanc au fac-simile
Se on admisus la sintaxo : \textit{la atakanti la religio}, on devus\nl
aceptar ol, pro analogeso, por omna substantivi derivita de\nl
verbo : on darfus do dicar « la atakemi la religio »; « la\parplein
\end{VUpage}

% -------------------------------------------------------------------
%  feuillet 118 — folio 114 : 26 lignes de corps, 2 de note.
%  COQUILLES DE L'ORIGINAL, conservees : ligne 8, le guillemet ouvrant
%  manque devant « la atako di la religio » ; et la citation ouverte
%  par « On (ligne 12) n'est jamais refermee.
% -------------------------------------------------------------------
\begin{VUpage}[118]{114}
% FIN DE LA PREMIERE PARTIE : un filet court et centre, pose dans le blanc
% sous le bloc de notes. Ordonnee estimee sur celle du filet de note de la
% meme page, dont il est distant de 275 px = 23,3 mm.
\VUornement{147.14mm}
\VUnotes{123.87mm}{%
(1) Segun \textit{Progreso}, VII, p. 159.

(2) « \textit{Exercaro} », lasta frazo dil exerco XXVI.%
}

\VUcontinue
laboristi la ligno », « la fumeri sigari », « la drinkeri\nl
liquori », e. c. On vidas, ke tala sintaxo (Slavatra) esus\nl
adminime stranjera, e (precipue) ne klara, pro ke ol apud\cc
pozus la vorti \textit{sen indikar lia relato}. Or logike la relato inter\nl
substantivo e lua komplemento devas esar expresata da pre\cc
poziciono, e « la atakanto la religio » o « la amanto Deo »\nl
ne esas plu permisata kam « la atako la religio » o « la\nl
amo Deo », vice la atako \textit{di} la religio » o « la amo \textit{di} Deo »,\nl
o « la amo \textit{a} Deo » (1).

\VUblanc{5.93pt}% mesure au fac-simile
139. --- Preferez, kande to esas posibla, la \textit{formo aktiva}\nl
al formo pasiva, kom plu kurta, plu vivanta e plu konforma\nl
al nuna marcho di nia lingui. Dicez do : « \VUgras{On} (o : ni)\nl
\VUgras{sequis la voyo, acensis la kolino, vizitis la kastelo, ed ad}\cc
\VUgras{miris la bela panoramo, quan on vidas del somito} (2),\nl
prefere kam per la formo pasiva : « \textit{La voyo sequesis, la}\nl
\textit{kolino acensesis, la kastelo vizitesis, e la panoramo admiresis,}\nl
\textit{qua videsas del somito.} »

\VUblanc{5.85pt}% mesure au fac-simile
140. --- Pri la \textit{sintezala formi} atencez ico : se li esas plu\nl
kurta e konciza, li riskas divenar kelkafoye obskura.\nl
Exemple : « por bone \textit{praktikeblesar} » certe prizentas a\nl
multi kelka obskureso. L'ideo esas nekontesteble plu klara\nl
en : « por \textit{esar} bone \textit{praktikebl}a ». Evitez generale soldar\nl
la verbo \textit{esar} a vorti ja derivita o kompozita. Se on abreviis\nl
\textit{amatesar} en \textit{amesar}, to ne esas por admisar tal formi kam\nl
\textit{ameblesar} o \textit{amindesar}. Same esas preferinda dicar : \textit{divenar}\nl
\textit{amoza}, \textit{divenar mizeroza}, kam \textit{amozeskar}, \textit{mizerozeskar}.
\end{VUpage}
